\documentclass[12pt]{article}
\usepackage[T1]{fontenc}
\usepackage[margin=1in]{geometry}
\usepackage{amsmath,amssymb}
\usepackage{graphicx}
\usepackage{float}
\usepackage{array}
\usepackage{longtable}
\usepackage{booktabs}
\usepackage{threeparttable}
\usepackage{caption}
\usepackage{authblk}
\usepackage[authoryear,round]{natbib}
\usepackage{xurl}
\usepackage[hidelinks]{hyperref}
\usepackage{microtype}
\emergencystretch=1.5em
\renewcommand{\thesection}{\Roman{section}}
\renewcommand{\thesubsection}{\thesection.\Alph{subsection}}
% Possessive citation: one macro so a bibliography change cannot desynchronize
% hand-rolled \citeauthor{X}'s (\citeyear{X}) constructions.
\newcommand{\citepos}[1]{\citeauthor{#1}'s (\citeyear{#1})}

\title{Mortgage Lock-In and the Federal Reserve's Quantitative Tightening Shortfall}
\author{Eugene Ong}
\affil{Carnegie Mellon University}
\date{}
\begin{document}
\maketitle
\begin{abstract}
\noindent U.S. agency mortgage-backed securities buy To-Be-Announced market liquidity with pool homogeneity built from standardized fixed-rate contracts whose par-payoff obligation, enforced through due-on-sale, locks households into below-market coupons when rates rise. This paper measures that trade-off's cost to the Federal Reserve during Quantitative Tightening, when agency MBS runoff fell \$764.7 billion short of the phased redemption caps over June 2022--November 2025. It cost less than the shortfall suggests. A loan-level survival framework built from Freddie Mac performance data recovers 91.3\% of that shortfall on a benchmark-consistent accounting basis at the headline off-window calibration, 88.2--93.7\% across that calibration's defensible range; at the in-sample calibration point on which the framework was assembled it recovers 97.9\%. Under the production hard-maximum form a null without the lock-in elasticity still recovers 85.7\% at the headline calibration (88.7\% in-sample; 35.6\% under the additive form), because scheduled amortization plus a turnover floor undershoot a \$35 billion monthly cap under any household response to rates. The Federal Reserve's own ex-ante projection anticipated the large majority of the realized shortfall. Measured against that projection rather than the never-binding cap, lock-in accounts for roughly a quarter to a half of the genuine surprise, depending on a disclosed intra-2022 allocation choice: the settlement-aware allocation matching the realized settlement pattern gives the quarter, the pre-committed uniform-spread central allocation the half. The design identifies a bounded range for the elasticity's contribution, not a pinned magnitude, and it identifies neither a level nor monthly timing: no monthly-timing claim survives its own placebo tests, and the margin's sign is fixed by a rate configuration that leaves essentially the whole book out of the money throughout the window, not by the estimated elasticity. Re-anchoring the baseline turnover floor that dominates the level outside the evaluation window and off the locked-in population puts that contribution at $+5.6$ percentage points over the rate-inelastic baseline ($+\$42.6$ billion), $+4.3$ to $+6.8$ points across the off-window floor's defensible range under the production hard-maximum form ($+3.0$ to $+8.0$ once the floor reads' cluster-bootstrap sampling error is propagated, open below $+4.3$ on an age-standardized read; wild-cluster correction of that read's under-coverage gives $+2.9$ to $+8.7$, the widest uncertainty source with a coverage property though not the widest read of that source, while the unestimated baseline seasoning ramp spans $+0.9$ to $+15.6$ at the same calibration); under the disclosed additive form the same anchors return $+11.2$ points nearly floor-invariantly, so the form-conditional identified range is $+3.5$ to $+13.1$ points, with $+5.6$ its max-form member; the in-window production calibration returns $+9.2$ points, which I now report as an in-sample calibration point, not the headline. The margin is positive at every floor and elasticity-band combination tested---a specification check the rate configuration cannot let fail, retained as such---and its magnitude is unchanged under a floor calibrated only on the window's first nineteen months, an input-stability check rather than an outcome holdout: no outcome-holdout months exist anywhere in this design. An agent-based model of the moving decision recovers 13.6\% averaged across seeds, though a cross-design variant recalibrated on real loan covariates recovers 59.3\% on the frozen seed and 60.2\% averaged over fifty seeds (76.3\% reweighted to the book's composition). The contrast cannot be cleanly attributed to modeling paradigm. The mobility cost is real. The institutional cash-flow cost is small; its sign turns on the transplant's anchor, not on any within-anchor sensitivity sweep: importing Denmark's baseline mobility along with its payoff rule yields roughly $-\$100$ billion, while a transplant changing only the payoff rule yields $+\$61$ billion at the in-sample calibration point and $+\$28$ billion at the off-window calibration the headline marginal uses. The latter's positivity follows from anchoring the Danish leg at the U.S. zero-gap hazard on a book out of the money almost everywhere; the informative quantity there is the magnitude, not the sign.
\end{abstract}

\vspace{0.5\baselineskip}
{\footnotesize\noindent\textbf{Keywords:} mortgage prepayment; lock-in effect; quantitative tightening; mortgage-backed securities; hazard models; monetary policy implementation.\\
\textbf{JEL codes:} E52, E58, G21, R31.\par}
\newpage
\section{Introduction}\label{sec:intro}

The Federal Reserve departed from its ``Treasuries-only'' framework in 2008 and 2020, purchasing trillions in agency MBS that suppressed 30-year mortgage rates to lows of roughly 2.6\% \citep{hancock2011} and triggered a refinance boom. In 2022 it reversed course. Quantitative Tightening (QT) was a passive roll-off intended to drain up to \$35 billion per month from its SOMA portfolio, under a cap phased in at \$17.5 billion and stepped up three months later. (June 2024 and April 2025 runoff slowings cut only the Treasury cap, and above-cap reinvestment never triggered in this window. So the constant post-phase-in \$35 billion MBS schedule is policy-committed throughout.) Unlike standard Treasuries, U.S. mortgages carry asymmetric, state-dependent payoff profiles. Tightening pushed the 30-year rate past 7.0\%, and under the par-payoff rule that imposed a substantial penalty on relocation: the lock-in effect. In practice, historic housing-inventory shortages and depressed consumer sentiment compounded it. The rate rise decomposes into the 10-year Treasury yield's climb plus primary-mortgage spread widening from roughly 170 to over 300~bp at its 2022--2023 peak (FRED series MORTGAGE30US less DGS10). Spread components follow \citet{boyarchenko2019}, whose framework predates and does not document the 2022--2023 magnitudes. Mortgages did not prepay. Cumulative tightening-cycle runoff fell \$764.7 billion short of the phased caps \citep[on the balance-sheet setting, see][]{lopezsalido2023,acharya2023} (Section~\ref{sec:method-benchmark}). Sections~\ref{sec:abm} and~\ref{sec:hazard} decompose that shortfall, isolating the lock-in channel's own share.

That the lock-in effect slowed the runoff of the Federal Reserve's mortgage portfolio is not a new claim; its own staff and officials have said as much (\citealp{na2024}; \citealp{perli2024}; \citealp{hammack2025}). This paper's contribution is not that mechanism but its decomposition: an accounting exercise within a fully disclosed model, not a causal identification. Switching the lock-in elasticity off partitions the realized shortfall between an elastic channel and a rate-inelastic baseline: scheduled amortization plus baseline turnover held at its observed level. Its floor is read off realized, partly behavioral turnover, so the baseline is not a non-behavioral object. The elasticity's causal content is inherited from the external estimates calibrating it, not established here (Section~\ref{sec:identification}). Most of the shortfall sits in that baseline: with the elasticity off, the loan-level model still recovers 85.7\% of the \$764.7 billion shortfall-against-cap benchmark (Section~\ref{sec:method-benchmark}) at the calibration this paper headlines and 88.7\% in-sample, on the benchmark-consistent accounting basis (Section~\ref{sec:hazard-interp}). The Federal Reserve's own ex-ante projection implies the same composition, though not independently: it may already have embedded part of the lock-in channel. The New York Fed's May 2022 staff baseline already anticipated the large majority of the realized cap-shortfall: 75.6--88.5\%, depending on a disclosed intra-2022 allocation choice. That range is wide enough that its overlap with the rate-inelastic null's recovery carries no precision.

The elasticity's identified footprint is therefore a bounded margin, not a pinned magnitude and not a sign. The floor that dominates the level is calibrated in-window, on the same locked-in cohorts the margin decomposes. I therefore re-measure it outside that window, on pre-episode 2017--2019 discount cohorts, and re-run both legs at the result. The margin runs $+4.3$ to $+6.8$ points across the off-window floor's defensible range. It runs $+2.3$ to $+9.1$ points once that floor's own sampling error is propagated with a wild-cluster small-sample correction and the pre-committed coverage rule selects the qualifying construction. That is the binding layer (Section~\ref{sec:robustness-floor}). At the mid-grid anchor inside that interval the value is $+5.6$ points, or $+\$42.6$ billion. The design delivers the range, not that value. Measured corrections to the floor and the accounting basis run in both directions: calendar-standardizing the floor read moves the value to $+4.7$ points; matching that read's housing-activity level to the window's moves it to $+10.7$. Specification variants that are not floor corrections run both ways as well, so the identified content sits in the range and not in any interior member (Section~\ref{sec:identification}). The in-window production calibration returns $+9.2$ points, reported as the in-sample calibration point, not the headline. Its sampling intervals are in Table~\ref{tab:headline}, and the operative uncertainty remains calibration (Table~\ref{tab:uncertainty}). Neither figure is the entire shortfall (Section~\ref{sec:identification}).

A pre-committed sweep of the baseline turnover floor and the elasticity band bounds the margin between $+2.1$ and $+13.2$ points, positive at every point of the calibration box and at every off-window anchor tested. That positivity checks the specification; it does not identify the mechanism (Section~\ref{sec:identification}). A functional-form test replacing the floor's hard maximum with a competing-risks combination leaves the margin inside that box and nearly floor-invariant (Section~\ref{sec:robustness-floor}). The grid, parity gates, and interpretive thresholds were committed to the repository before any run executed; I reserve ``pre-registered'' for third-party registries. That pre-commitment is a disclosure device, not a proof: the sweep's credential should be read net of the replication package's verdict appendix, which tabulates every such rule together with how its outcome was adjudicated after the run \citep{olken2015}.

I quantify this shortfall's mechanics with two structurally distinct estimators of the aggregate level. The first is a 10,000-household Agent-Based Model with dynamic macroeconomic frictions. The second is a loan-level hazard framework that abandons the household utility function altogether, introduced after the ABM's residual proved persistent and unresponsive to added behavioral realism. They disagree sharply about how much household choice explains (Sections~\ref{sec:abm} and~\ref{sec:hazard}). The ABM recovers 13.6\% on the fifty-seed mean (10--17\% across seed noise; 11.9\% on the frozen production draw), the hazard framework 91.3\% at the headline off-window floor and 97.9\% at the in-sample calibration point (remaining variants: Table~\ref{tab:headline}). The decomposition rests on the hazard framework alone: the rate-inelastic/rate-elastic partition is a hazard-framework object, because Path A is excluded from every headline range by pre-committed specification, and the ABM's null admits two readings that disagree in sign. Path A lands at 112.4\% yet does not corroborate: no useful propagated precision, no inferential object recomputed under its production specification (Section~\ref{sec:patha}). Two qualifiers travel with the two-estimator contrast wherever I state it. The $\beta_1 = 0$ null already recovers 85.7\% (88.7\% in-sample), so the hazard framework's margin over a rate-inelastic model is $+5.6$ points at the off-window floor ($+9.2$ in-sample), not the level (Section~\ref{sec:identification}). In flow terms the margin is roughly \$1 billion per month against a \$35 billion monthly redemption cap. A cross-design variant of the ABM recalibrated on real loan covariates recovers 59.3\% on the frozen seed, 60.2\% across fifty seeds, and 76.3\% at the book's composition (Section~\ref{sec:robustness-crossdesign}). So the contrast is not clean evidence about modeling paradigm alone. The hazard framework is the instrument; the decomposition it yields is the paper's central empirical result. A counterfactual on the Danish market-value repurchase regime further shows the institutional cash-flow component of systemic duration risk to be a consequence of U.S. institutional design. Under the rule-only transplant it matches the lock-in marginal in origin, and in size at the in-sample calibration point on which both were run.

This result sits inside a two-sided exchange: the pool homogeneity behind To-Be-Announced (TBA) market liquidity \citep[the market's institutional design is reviewed in][]{fuster2023} is purchased with household mobility, because the standardized contracts TBA delivery presupposes carry the ``par-payoff'' obligation. Under that obligation, a homeowner must extinguish the existing mortgage at full face value to relocate. The obligation is a property of the U.S. fixed-rate note, not of TBA eligibility, and it is enforced by the due-on-sale clause (made broadly enforceable by the Garn--St.~Germain Act of 1982). The FHA/VA loans backing Ginnie Mae pools are assumable, a statutory carve-out whose exercise is underwritten (HUD Handbook 4000.1; 38 U.S.C.~\S~3714). Those pools are 20.4\% of SOMA book face on the paper's CUSIP-level cohort parse, and this design nowhere measures their realized take-up, so that 20.4\% share bounds the carve-out from above. The observed speeds bound it from above a second way, in the same direction. An assumed loan does not prepay, so materially high take-up would show up as \emph{slower} Ginnie voluntary speeds. Over the window Ginnie's voluntary component instead runs $1.35$ points per year \emph{above} the conventional one. Section~\ref{sec:pathb} uses that differential to bracket what a Ginnie share that is not gap-inert would contribute, on a mapping it assumes rather than measures. Both bounds are ceilings, not sizes; neither prices the frictions (second-lien gaps, equity buy-outs, qualification) between the statutory right and its exercise. The decomposition disciplines this frame: household mobility is constrained, as the household-level lock-in literature documents, while the cash-flow shortfall is governed mostly by loan-level and institutional mechanics. The trade-off's binding cost is therefore denominated in mobility, not institutional cash flow (Section~\ref{sec:conclusion}). Section~\ref{sec:abm-danish}'s Danish counterfactual prices the payoff rule itself: recalibrated against separately estimated Danish elasticities, the rule-only gap is $+\$61.2$ billion (8\% of the benchmark). That figure is the zero-refinance edge of a band swept to $+\$256.8$ billion at 3\% refinance-in-place. The sign is forced rather than found: pinned at the U.S. zero-gap hazard, the Danish leg must outrun a book out of the money almost everywhere. The informative content is the magnitude. The bracketing anchor that also imports Danish baseline mobility reverses it to $-\$99.9$ billion, and the cash-haircut incidence reading reverses it to $-\$51.0$ billion at the re-derived discount. The relief is established under face accounting only, the benchmark's own denomination (Section~\ref{sec:abm-danish}; Appendix~\ref{app:params}). The mobility-versus-cash-flow ranking rests on external evidence, not a within-paper comparison: I measure the institutional leg and import the mobility leg at an elasticity I never re-measure, so the two are nowhere commensurable. The one household-side magnitude I carry is \citepos{batzer2024} roughly \$2.4 trillion of below-market financing value forfeited on relocation. It is computed on universal relocation and therefore an upper bound, and a stock of retained financing value is not a cash flow over forty-two months. The ranking's direction is what that supports; its ratio is not.

This paper's affirmative contributions are three. First, a decomposition of the QT mortgage-runoff shortfall into rate-inelastic and rate-elastic components, the elastic part delivered as a bounded margin (Sections~\ref{sec:abm}, \ref{sec:hazard} and~\ref{sec:identification}). Second, an expectations-based complement that measures the shortfall against the Federal Reserve's own ex-ante projection, not the never-binding caps (Section~\ref{sec:method-benchmark}). Third, the cap-design arithmetic. New York Fed staff projected from the month the \$35 billion ceiling took effect that the caps would sit above achievable runoff \citep{nyfed2022}; the observation is theirs, and what I add is its quantification and ex-ante construction. The ceiling was set roughly 1.7 to 1.9 times the Federal Reserve's own contemporaneous projection of achievable runoff, and both components of the rate-inelastic path are computable in advance of the cycle, as a band (Sections~\ref{sec:identification} and~\ref{sec:discussion-capdesign}). Two further design elements are not contributions. The paired cross-design test is a validation device under interpretive thresholds fixed ex ante, and it ran against this paper's own ABM reading: on all fifty seeds the real-covariate leg crossed the pre-fixed threshold at which that reading is undercut (Section~\ref{sec:robustness-crossdesign}). The order-of-magnitude reduction of the Danish market-value counterfactual, with the incidence bracket that prices it, is a robustness result on a counterfactual. It is not a measurement of the U.S. book: the \$925.5 billion figure it corrects is this paper's own earlier extrapolation, and the corrected gap's sign is anchor- and incidence-determined (Section~\ref{sec:abm-danish}; Appendix~\ref{app:params}).

Table~\ref{tab:headline} collects the core results and their operative uncertainties in one place; every value repeats a committed, gated figure.

{\footnotesize
\begin{longtable}{@{}p{4.6cm}p{4.6cm}p{5.3cm}@{}}
\caption{Core results and operative uncertainties. Each row repeats committed values derived in the referenced exhibit; recovery percentages are quoted on the benchmark-consistent shared basis. The uncertainty hierarchy is established in Table~\ref{tab:uncertainty}: simulation-seed noise is negligible at reporting precision and sampling is quoted under the stratum-cluster scheme; the calibration box is the operative uncertainty --- except the headline marginal, whose binding layer is the floor read's own propagated sampling error (Table~\ref{tab:uncertainty}). The replication package ledger's crosswalk table maps each headline quantity to its accounting basis and source run. The Path B recovery is estimated on the Freddie 2017--2021 conventional universe; Section~\ref{sec:method-benchmark} derives that universe's coverage of SOMA book face on the book's own agency $\times$ vintage joint cells, and states the external-speed overlay that bounds each exclusion (Section~\ref{sec:pathb}).}
\label{tab:headline}\\
\toprule
Object & Headline value & Uncertainty / bounds \\
\midrule
\endfirsthead
\toprule
Object & Headline value & Uncertainty / bounds \\
\midrule
\endhead
Cap-relative benchmark, June 2022--November 2025 (Section~\ref{sec:method-benchmark}) & \$764.7 billion & -- \\
Expectations-based complement (Section~\ref{sec:method-benchmark}) & \$87.8 billion, 11.5\% of the cap-relative benchmark & Fed's ex-ante projection implied 75.6--88.5\% of the cap-shortfall, by intra-2022 allocation \\
Rate-inelastic baseline, $\beta_1 = 0$ (Section~\ref{sec:identification}; Table~\ref{tab:bases}) & 85.7\% of benchmark at the headline off-window floor; 88.7\% in-sample, both under the production floor form (35.6\% under the additive form) & 83.9--86.9\% across the off-window floor range under that form; moves one-for-one with the floor level; floor-read wild-cluster $[80.6, 88.4]$\% (run \texttt{null\_floor\_interval}) \\
Lock-in marginal, off-window floor (Sections~\ref{sec:identification} and~\ref{sec:robustness-floor}; Table~\ref{tab:oosfloor}) & $+2.3$ to $+9.1$ points, basis-invariant --- the floor read's own sampling error propagated at the central elasticity (the 6.5\% mid-band value), not an interval on the elasticity, whose dimension Table~\ref{tab:lowband} sweeps; $+5.6$ points $= +\$42.6$ billion at the mid-grid anchor inside it, an anchor convention rather than a central tendency; the range carries the identified content (Section~\ref{sec:identification}) & floor-read wild-cluster bootstrap-$t$ inversion $[+2.3, +9.1]$, the binding layer (31 clusters; percentile read $[+3.0, +8.0]$, under-covering; the corrected ladder and its degrees of freedom are Table~\ref{tab:ladder})\tnote{a} \\
Lock-in marginal, convention envelope (Table~\ref{tab:assembly}) & $+0.9$ to $+15.6$ points across named convention variants (seasoning ramp 75--150 PSA; moving-share bracket; housing-activity-scaled baseline; floor form) & convention ranges without coverage properties; the calibration box spans $+2.1$ to $+13.2$ \\
\addlinespace
Lock-in marginal, in-sample calibration point (Section~\ref{sec:identification}; Table~\ref{tab:uncertainty}) & $+9.2$ points $= +\$70.3$ billion, basis-invariant & sampling: within-stratum $[+9.17, +9.23]$, stratum-cluster $[+8.27, +10.19]$ (the operative scheme); calibration box $+2.1$ to $+13.2$; Fannie replication $+8.68$ (Section~\ref{sec:pathb-fannie}) \\
Path B central level (Table~\ref{tab:estimators}) & 91.3\% at the headline off-window floor; 97.9\% in-sample (100.0\% composed as committed, 98.1\% on a single coupon convention) & 88.2--93.7\% across the off-window floor range; covariate priors 92.2--99.2\% (in-sample); Ginnie-speed overlay at the headline floor 84.0\% central against a 79.5\% null, and 89.3\% at the in-sample point (raw; of the \$66.3B raw adjustment the Ginnie-specific piece is $+\$38.5$B, within the \$20--47B bound) \\
Production ABM (Table~\ref{tab:estimators}) & 13.6\% seed mean & 10--17\% seed band; frozen draw 11.9\%; the operative uncertainty is calibration, not seed noise --- established on the real-covariate cross-design, where one calibration choice moves recovery across 20.9--59.3\% (Section~\ref{sec:robustness-crossdesign}); that span belongs to two calibration choices on that design, not to this row's 13.6\%, and does not contain it \\
Cross-design ABM, real covariates (Table~\ref{tab:crossdesign}) & 59.3\% recalibrated; fifty-seed mean 60.2\%, 95\% band 55.1--66.0\% (76.3\% at book composition; effective sample 1{,}332 of 10{,}000, re-draw 70.6\%) & 20.9\% under the frozen scale; fifty-seed mean 22.0\% (12.6\% at book composition) \\
Danish rule-only counterfactual (Table~\ref{tab:danish}) & $+\$61.2$ to $+\$256.8$ billion over 0--3\% refinance-in-place, of which $+\$61.2$ billion is the zero-refinance edge rather than a central value & positivity is forced by the zero-gap anchor, so the refinancing sweep cannot flip it; the Danish-level bracketing anchor does, at $-\$99.9$ billion. The gap is incidence-conditional: $+\$61.2$ billion under face accounting (the benchmark's own denomination), reversing to $-\$51.0$ billion under cash accounting at a re-derived discount, and product-mix-conditional besides: under face accounting itself it reverses above a Danish interest-only share of 34.3\% (runs \texttt{buyback\_credit\_bracket} and \texttt{buyback\_discount\_rederived}; Appendix~\ref{app:params}) \\
\addlinespace
Household-side foregone payoffs implied by the marginal (Section~\ref{sec:discussion-policy}) & about 179{,}500 on the book over the window (51{,}300 a year); 233{,}000--256{,}000 a year grossed to the whole market & an \emph{implied} count, not a measured one --- the marginal divided by a representative surviving balance, inheriting every conditionality the marginal carries. Upper bound at the production $\mu = 1$ convention (108{,}300 at $\mu = 0.5$, 58{,}700 at $\mu = 0.25$, on the book; the grossed upper bound is 5.7--6.3\% of the 4.08 million annual existing-home-sale run rate; Section~\ref{sec:discussion-policy}). Scaled by the household-level literature's per-move welfare estimates, the counts are the marginal's household-side image, the magnitude attributed to that literature \citep{gerardi2024} rather than measured here. \citet{batzer2024}'s \$2.4 trillion of below-market financing value prices the same households' stakes on universal relocation and is not additive with this count. \citet{hsieh2019} give the misallocation mechanism the moves would feed, not estimated here \\
\addlinespace
Portfolio duration extension (Table~\ref{tab:wal}; Appendix~\ref{sec:robustness-wal}) & empirical-path WAL 9.4 years at the window's open and 8.5 at its close --- a 6.0-year extension against the no-shock counterfactual's 2021 refinancing-boom speeds (3.4 years) & the comparator sets the result, and this row is quoted against both. Read instead against the baseline turnover floor's off-window read, the same path runs 9.4 years against 9.5 at the open and 8.5 against 8.6 at the close, so the six years measure refinancing-versus-turnover rather than lock-in-versus-normalcy. The Danish rule-only leg shortens the portfolio by 0.6 years against the U.S. leg's 9.7 (Danish-level bracketing leg 10.8). Approximate single-pool-per-cell calculation; option-adjusted duration and convexity are declined as out of scope \\
\bottomrule
\end{longtable}
}

\noindent{\footnotesize \emph{Notes to Table~\ref{tab:headline}.} a Sensitivity catalogue for the headline marginal. Each entry is a committed figure derived where it is cited, and none of them is a sampling interval: off-window floor range $+4.3$ to $+6.8$ (max form; $+3.9$ to $+7.5$ joint with the band); additive form $+11.2$, floor-invariant; form-conditional hull $+3.5$ to $+13.1$, upper end from a 61.2\%-recovery cell and not read as an equally credentialed member (the concave transform and the additive form do not compose additively: jointly $+9.2$, interaction $-1.47$); mixture curve in the floor's strictly-involuntary share (Section~\ref{sec:robustness-floor}): $+9.7$ at $\omega{=}0.25$, $+10.7$ at $\omega{=}0.4$; Ginnie overlay $+4.4$ (conventional-share scaling $0.797\times$); floor-stability check $+\$45.1$ billion.\par}

\paragraph{Definitions used throughout.} Four accounting bases attach to every recovery percentage (Section~\ref{sec:method-benchmark}): \emph{standalone-scorer} (no adjustment), \emph{benchmark-consistent shared} (nets a common \$69.6 billion curtailment flow identically from every U.S. leg; the main-text convention), \emph{full-book} (reweights to the SOMA book's coupon mix), and \emph{composed} (both). The realized SOMA CPR is a back-out and moves with the coupon convention it is backed out at: unless a site says otherwise, every empirical CPR here is the \emph{ABM-basis} back-out---window mean 5.14\%, against a hazard-basis 5.52\% and 5.79\% at the note-rate WAC, the corrected amortization basis since the note rate is what amortizes the loan (run \texttt{coupon\_convention\_amortization}; duration pricing in Table~\ref{tab:wal}). Two floor calibrations attach to every marginal and level: the \emph{in-sample calibration point} (the 4\% baseline turnover floor, in-window on the locked-in 2023--24 cohorts; returns $+9.2$ points) and the \emph{off-window calibration} (the 4.70--5.33\% floor, on pre-episode 2018 discount cohorts; returns $+5.6$ points and is the headline). The floor is a measured object: realized-turnover reads put it at 3.8--3.9\% in-window (flat across seasoning cuts); 4.70--5.33\% on the off-window 2018 leg, all of it on cohort-months aged twelve to twenty-four, so that leg's mature-seasoning test is not computable rather than passed; 5.51\% under an imputation-heavy age standardization (84\% imputed weight, an indicative bound from the same 2018 leg); and 5.52\% on an out-of-agency Fannie Mae read of the pooled off-window period---the last two sit above the clean band, as does every mature-versus-young turnover contrast the data make measurable, hence the band's soft lower edge (Section~\ref{sec:robustness-floor}). The ``baseline turnover floor'' (sometimes called an involuntary-turnover floor) names that object without decomposing it: the reads measure total deep-discount-cohort turnover---discretionary life-cycle moves and cash-out refinancings mixed with strictly involuntary events---and the strictly-involuntary share, plausibly well under half (Section~\ref{sec:pathb}), is priced directly by the mixture curve of Section~\ref{sec:robustness-floor}. ``Trapped liquidity'' and ``shortfall'' denote the net shortfall of realized roll-off against the phased redemption caps (\$764.7 billion, June 2022--November 2025). The \emph{lock-in marginal}, the paper's identified object, is the central simulation minus the \emph{rate-inelastic baseline}, its $\beta_1 = 0$ null: scheduled amortization from the book's own coupon, age, and term composition plus baseline turnover held at its observed level, the marginal being the imported elasticity's increment above it. The partition separates baseline from increment, not mechanics from behavior---the floor is read off realized, partly behavioral turnover, so the baseline is not a non-behavioral object and I do not call it one. The marginal's operative uncertainty is the calibration-and-form envelope of Table~\ref{tab:uncertainty} rather than any sampling interval, and the envelope is a grid, not a confidence region---a load-bearing distinction: the elasticity band it sweeps is \citeauthor{liebersohn2024}'s \emph{specification} range with my adopted midpoint at its centre, not a published point estimate with a published standard error, no sampling error of theirs is propagated anywhere in this paper, and so the envelope has no coverage property. Where a genuine sampling interval exists---the floor reads' cluster bootstrap, the loan- and stratum-level bootstraps of Section~\ref{sec:pathb}---I quote it as such and rank it against the envelope explicitly rather than merging the two. One global property is stated here as well as in the limitations, because every number below inherits it: \emph{no outcome-holdout months exist anywhere in this paper}, without exception; its unavoidability is scoped to the cash-flow margin. The identified object has nothing to hold out---a difference between two simulations answers to no realized outcome---while a level-fit holdout, conceivable but not attempted, would split the benchmark itself, one cumulative path against one cap schedule; the household margin does admit outcome moments this design leaves unused, importing the mobility elasticity rather than estimating it (Section~\ref{sec:limitations}). The floor that sets Path~B's level is calibrated on in-window turnover, Path~A trains on nineteen of the forty-two QT months, and Path~B's level is evaluated on the window that sets it---each disclosed again where it arises, and the nineteen-month floor variant reported later is an input-stability check rather than an outcome holdout: what it holds out is calibration months, not the benchmark. Every recovery percentage is accordingly in-sample or calibrated in the time dimension---a description of fit rather than a prediction (Section~\ref{sec:limitations}).

\section{Literature Review}\label{sec:lit}

Prepayment theory begins with ``ruthless exercise''---instantaneous refinancing as conditions dictate \citep{dunn1981}. It then admits frictions. Within a valuation framework, \citet{schwartz1989} fit an empirical prepayment function to aggregate GNMA experience by maximum-likelihood hazard, the econometric ancestor of my loan-level survival approach (Section~\ref{sec:hazard}). \citet{richard1989} formalize industry practice: prepayment as an empirical function of refinancing incentive, seasoning, seasonality, and a pool's own prepayment history. \citet{stanton1995} add heterogeneous transaction costs, which smooth the refinancing cliff into a predictable S-curve, and \citet{agarwal2013} sharpen the rational benchmark into a closed-form optimal refinancing threshold that scores observed behavior. \citet{chernov2018} recover from MBS market prices a function driven by macroeconomic turnover alongside the rate response, and my rate-inelastic null operationalizes that turnover-versus-rate-response split on the cash-flow side. \citet{quigley1987,quigley2002} established lock-in as a persistent friction on mobility in rising-rate environments, and \citet{ferreira2010} gave early large-sample evidence that rising mortgage rates and negative home equity each depress owner mobility. The latter is an equity-lock-in channel distinct from, and complementary to, the rate channel this paper models; the equity dimension is left unmodeled here, a boundary disclosed in Section~\ref{sec:pathb}.

The 2022--2025 rate shock generated a large empirical literature documenting the breakdown of standard prepayment assumptions. \citet{fonseca2024} documented the steep decline in U.S. mobility and refinancing, and \citet{fonseca2026} extend that evidence to general equilibrium, where locked-in owners who do not downsize raise net housing demand, offsetting part of the price decline a rate rise would otherwise produce. This paper neither models nor prices that channel. \citet{batzer2024} estimate roughly \$2.4 trillion of below-market financing value held by U.S. fixed-rate borrowers---retained by staying, forfeited on relocation---and roughly 1.33 million sales prevented over 2022Q2--2023Q4 (1.72 million through 2024Q2 in their revision). \citet{liebersohn2024} put it in hazard form: each percentage point of prevailing rate above a borrower's locked-in coupon reduces quarterly mobility probability by 5.5\% to 7.7\% depending on specification. My Path B microsimulation adopts that elasticity band as calibration. \citet{graybill2026}, a Federal Reserve Bank of Philadelphia working paper, independently estimate that margin from transaction records in a housing-market-equilibrium setting, finding deep rate gaps cut the sale hazard by roughly 7--11\% per percentage point. That setting is where lock-in's most-discussed housing consequence sits: a household that does not list withholds a unit, and that incidence falls on first-time buyers and renters, not the locked-in owner, a distributional consequence this design measures nowhere. The finding corroborates the band from data that never touch credit records, but it sits at or above the band, not at its centre: only the 7\% low end falls inside 5.5\% to 7.7\%, and its midpoint lies above the high edge (Section~\ref{sec:pathb}).

A related strand asks how much of the mobility decline a rational rate-gap framework explains, leaving the rest an unexplained residual. \citet{aladangady2024} hold the 2021 rate environment fixed in a counterfactual rate-gap distribution and correct for selection bias from forgone refinancing; they find rate-gap-driven lock-in explains only 44\% of the actual decline in mortgage borrower mobility between 2021 and 2022. Section~\ref{sec:identification} maps that 44\% share into this design's own units as a calibration target for the scaled-null variant (run \texttt{scaled\_null\_housing\_activity}). My decomposition in Sections~\ref{sec:abm} and~\ref{sec:hazard} is methodologically analogous, but it runs on dynamically modeled CPR and cash-flow shortfalls rather than mobility rates, with two structurally distinct estimators. \citet{gerardi2024} similarly score a structural search-and-matching lock-in model against empirical lifecycle migration, finding it accounts for most, but not all, of the observed disruption to housing market liquidity.

A separate behavioral finance literature grounds what such residuals likely represent. \citet{keys2016} find more than 20\% of U.S. borrowers fail to refinance when it is unambiguously in their financial interest, attributing this to procrastination and uneven financial sophistication. \citet{defusco2020} add a structural counterpart: employment-documentation requirements and out-of-pocket closing costs bind hardest in recessions and inhibit refinancing pass-through. On Danish administrative data, \citet{andersen2020} find inattention and inertia, not transaction costs, behind much failure to act on refinancing opportunities. \citet{andersen2022}, on Danish housing data, show reference dependence and loss aversion relative to the original purchase price, with losses weighted roughly 2.5 times more heavily than equivalent gains. Across administrative, survey, and field-experimental data, \citet{johnson2019} find that of the behavioral factors tested, suspicion of lenders' motives---not standard rational frictions such as time preference or risk aversion---is most consistently related to households' failure to accept profitable refinancing offers, which further supports a behavioral reading of the residual. \citet{bissiri2014} formalize the distinction, defining ``behavioral risk'' as the share of prepayment decisions not driven by a fully rational financial criterion. \citet{perotti2024} applies that operationalization more recently, modeling behavioral uncertainty as a stochastic overlay on an otherwise rational prepayment framework and deriving robust hedging strategies that bound, not eliminate, the exposure it creates.

A further strand treats prepayment as a loan-level survival process shaped by pool composition and servicer behavior, not by borrower optimization. Prepayment ``burnout'' is a standard feature of industry MBS valuation models and one of the four factors of \citepos{richard1989} canonical prepayment function, alongside the Public Securities Association seasoning convention. Burnout is the tendency for a pool's aggregate refinancing sensitivity to decline over time as its most rate-sensitive borrowers prepay earlier and self-select out of the remaining pool. Yet burnout is difficult to reconcile with a fixed, representative-agent population. Recent theory formalizes why: in a heterogeneous pool the observed pool-level hazard is the survival-weighted mean of the individual hazards, and its drift carries a negative selection term proportional to the cross-sectional dispersion of those hazards \citep[arXiv preprint]{lesniewski2026}, so pool speeds decline with seasoning even when no individual hazard changes. My hazard framework (Section~\ref{sec:hazard}) is designed to accommodate this loan-level survival structure.

A separate fixed-income literature highlights the systemic risks the lock-in effect generates. \citet{jiang2023} documented that duration extension in bank asset portfolios, prominently including low-coupon MBS, generated roughly \$2 trillion in unrecognized mark-to-market losses across the U.S. banking sector, contributing to self-fulfilling solvency runs such as the collapse of Silicon Valley Bank. \citet{smith2023} demonstrated that QT relies heavily on predictable MBS runoff schedules, which these duration extensions structurally break.

Closest to this paper, the Federal Reserve's own staff and officials have already described the mechanism and measured its first-order footprint. \citet{na2024} report that ``a vast majority of the agency MBS held in the SOMA are far out-of-the-money to refinance and may be subject to the `lock-in' effect.'' They further report that because ``monthly principal payments have been below the \$35 billion-per-month redemption cap \ldots the size of the SOMA agency MBS portfolio has decreased at a much slower pace than in the case of a binding cap.'' They put June 2022 to June 2024 runoff at roughly \$450 billion against the \$820 billion a binding cap would have retired, an implied shortfall near \$370 billion. The same note reports the composition of those below-cap payments: scheduled principal accounts for roughly half, with prepayments mostly due to turnover making up the remainder. That is an external, staff-measured anchor for the scheduled-plus-turnover decomposition my $\beta_1 = 0$ null quantifies. The anchor and my null have different denominators. Na et al.'s composition describes realized principal \emph{payments}: roughly half scheduled, half turnover-driven prepayment. My null is built from the same two ingredients, scheduled amortization plus a turnover floor calibrated from realized, partly behavioral turnover. My 88.7\% (in-sample; 85.7\% at the off-window floor) instead describes the null's share of the shortfall \emph{against the caps}. A payments mix of scheduled-plus-turnover that runs below \$35 billion in essentially every month is the state in which that shortfall accrues mechanically. Their composition and my null recovery are therefore the same fact read against two denominators, and read on the composition the null is built to reproduce. \citet{perli2024}, then Manager of the System Open Market Account, stated that the cap ``has not been binding \ldots resulting in very slow prepayments,'' and \citet{hammack2025} observed that some of the low-coupon holdings ``could be very long lived.'' New York Fed staff had projected the below-cap regime from its first month: a September 2022 staff post put monthly agency MBS principal payments at \$20 to \$30 billion over the following months, below the \$35 billion cap that took effect that month \citep{nyfed2022}. I treat the qualitative mechanism as established by these sources. Four features separate my exercise from theirs. First, this characterization attributes the below-cap runoff to lock-in without isolating what the elasticity contributes, as distinct from the scheduled amortization and baseline turnover that would have undershot the cap under any rate-driven prepayment response above that baseline. Sections~\ref{sec:abm} and~\ref{sec:hazard} supply that decomposition. They find the identified lock-in margin bounded between $+2.3$ and $+9.1$ points of the benchmark at a turnover floor measured off the locked-in population. That bound is a wild-cluster interval on the floor read's 31 clusters; the narrower percentile read under-covers. At the mid-grid anchor inside that interval the margin is $+5.6$ percentage points ($+9.2$ at the in-window production calibration), and Section~\ref{sec:identification} qualifies that margin. Second, I extend the same shortfall-against-cap construction from the \$450 billion-through-June-2024 quantity to the full June 2022 to November 2025 window, where it reaches \$764.7 billion (Section~\ref{sec:method-benchmark}). Third, I measure the shortfall against the Federal Reserve's own ex-ante projection as well as against the caps, separating the component the Federal Reserve itself anticipated from the genuine surprise (Section~\ref{sec:method-benchmark}). Fourth, the decomposition turns the staff observation into a statement about cap design: scheduled amortization and a turnover floor read from realized, partly behavioral turnover undershoot a \$35 billion monthly ceiling in essentially every month, so a cap at that level could not have bound under any rate response above that floor (Section~\ref{sec:identification}). A more recent \citet{eyal2026} FEDS Note also decomposes the Federal Reserve's balance-sheet reduction, but into macroeconomic drivers (inflation, real GDP growth, and active SOMA reductions relative to nominal GDP), not into the loan-level sources of the MBS runoff shortfall itself, so it does not speak to the lock-in-versus-baseline question I take up here.

The Danish mortgage securitization system offers a tested structural alternative to negative convexity. \citet{berg2018} set out its ``Balance Principle'' of strict one-to-one match funding; under those rules \citet[\S3.3.2]{berger2026}, in a July 2026 SSRN working draft, estimate household mobility close to flat in the coupon gap, because borrowers may repurchase their outstanding mortgages at market value, neutralizing extension risk for the bondholder. The mechanics are institutional, not notional, and fix what a transplant must replicate: beyond penalty-free prepayment at par, a Danish mortgagor may retire the loan by buying the corresponding bonds in the secondary market and \emph{delivering} them to the mortgage bank (an option the U.S. contract does not carry), and the Danish contract has no due-on-sale clause \citep{frankel2004}. That is an open-market purchase settled by delivery, not a lender-side discount: its cost to the borrower is the bond's market price, its availability depends on a liquid secondary market in the specific series, and the bondholder is a \emph{voluntary seller} at market price, not called at par. Transplanted onto a passthrough holder the analogue is an outright sale at a discount, so a par-denominated redemption cap is not the natural scorer for the Danish leg, hence bracketing it under face and cash accounting. \citet{svenstrup2006} set out the design's wider reform implications. The flatness is a slope result on Danish data, so Section~\ref{sec:abm-danish} is careful what it licenses: it imports the slope but anchors the level at the U.S. hazard, Berger et al.\ estimating no moving level for U.S. households under a Danish payoff rule. The strand supplies rule and slope, not a priced U.S. counterfactual; Section~\ref{sec:abm-danish} scores the rule-only transplant on the SOMA book directly, an order of magnitude below a mechanism-extrapolated read, and Appendix~\ref{app:params} prices the buyback credit's incidence, which decides whether the institutional gap survives the accounting.

A final strand situates lock-in within the theory of monetary policy transmission. \citet{berger2021} show the Federal Reserve's power to stimulate through rate cuts is path-dependent, since prepayment depends on the outstanding distribution of in-the-money refinancing incentives. I extend that to tightening: the path dependence that weakens easing after a low-rate period weakens tightening after one too, since the resulting stock of deeply in-the-money, rate-locked mortgages cannot prepay however aggressively the central bank then raises rates. \citet{eichenbaum2022} formalize this state dependence for the easing direction, where policy potency depends on the distribution of outstanding rate gaps, and \citet{beraja2019} add regional evidence that the refinancing channel's strength varies systematically with the local distribution of in-the-money mortgages. \citet{dimaggio2020} give direct evidence for the easing-direction channel this paper mirrors: the Federal Reserve's QE1 agency-MBS purchases transmitted to the real economy specifically by inducing refinancing. The QT setting also links to how balance-sheet operations transmit through the supply of long-duration assets. There \citet{krishnamurthy2011} document several distinct quantitative-easing channels, one specific to agency MBS and running through prepayment, and \citet{vayanos2021} formalize, under preferred-habitat demand, why shifts in the outstanding maturity and supply of held bonds move the term structure at all. Lock-in is thus an input to that balance-sheet-supply channel, not separate from it. \citet{du2024} assemble the cross-country QT record---seven central banks, active and passive runoff---finding announcement effects concentrated in yields; the MBS runoff shortfall decomposed here is that record's passive-runoff margin, measured at the portfolio cash-flow level. \citet{vickery2013} trace the TBA market's liquidity premium specifically to the homogeneity of deliverable pools, with TBA-eligible securities trading at significantly lower cost than otherwise similar specified pools. That homogeneity requirement is the institutional constraint that, in Section~\ref{sec:conclusion}, I argue makes the Danish market-value repurchase structure hard to reconcile with the current U.S. secondary market absent a broader redesign of agency securitization rules.

My two estimators sit at the transparent end of two active methodological literatures. On the survival side, prepayment and default hazards are increasingly fit nonparametrically or by machine learning: random survival forests \citep{ishwaran2008}, competing-risks Bayesian additive regression trees \citep{sparapani2020}, and deep neural-network mortgage hazard models \citep{sadhwani2021}. The last benchmarks what flexibility buys at scale. On origination and monthly performance records for over 120 million U.S. mortgages, the deep network of \citet{sadhwani2021} learns highly nonlinear, interaction-rich prepayment and default responses without hand-specified functional form. Those responses include the burnout-type history dependence Path~A carries as an explicit stratum-level regressor and the ABM probes by pre-committed falsification (Sections~\ref{sec:patha} and~\ref{sec:abm-results}). I instead estimate a fixed-effects Poisson hazard and a literature-calibrated microsimulation, with functional forms and calibration inputs written out in full. What I need is not maximal predictive fit but an auditable decomposition. The null of Section~\ref{sec:identification} that switches the lock-in elasticity off is interpretable only if every term it holds fixed is visible, and a black-box hazard would obscure those terms. Those machine-learning survival methods remain a natural robustness avenue for the loan-level fit, and Section~\ref{sec:patha} probes them at their tabular end. Gradient-boosted Poisson and isotonic-recalibration comparators are each given Path A's information set on the identical pre-committed holdout, and each recalibrates levels without gaining cross-cell predictive power. Deep loan-level variants remain future work on data-availability grounds. A complementary strand attacks nonstationarity, not nonlinearity. Landmarking re-anchors a simple hazard among survivors at successive prediction times \citep{vanhouwelingen2007,putter2022}. Recent credit-risk work pairs landmark targets with explicit probability recalibration to hold survival predictions stable under covariate and base-rate drift, on the same Freddie Mac performance data my hazard paths use \citep{peng2026}. My design meets that concern with fixed, visible machinery: stratum fixed effects, calendar-month dummies, the dynamic friction index of Section~\ref{sec:patha}. It prices the residual time-variation it cannot remove by the temporal moving-block bootstrap reported there. On the agent-based side my calibration is fixed by the pre-committed parameterization of Appendix~\ref{sec:method-abm}, not by an automated estimator. It sits alongside a maturing literature on data-driven ABM calibration and estimation: simulated-minimum-distance and Bayesian methods \citep{grazzini2017,platt2020}, the empirical-validation taxonomy of \citet{fagiolo2007}. Within that literature the pre-committed stress-test design of Section~\ref{sec:abm} is one validation strategy among several. The paired design-crossing of Section~\ref{sec:robustness-crossdesign} is offered as a validation design in its own right: the ABM's decision rule re-run on real loan covariates, the same survival structure run on a fully synthetic population, interpretive thresholds fixed before the cross-design runs. Its nearest relative in that literature is indirect inference, which minimizes the distance between an auxiliary model's parameters fitted on real and on simulated data \citep{platt2020}. The pair here crosses designs not to estimate parameters but to attribute a result gap to modeling paradigm versus data source. It is not, to my knowledge, a named member of the validation taxonomy of \citet{fagiolo2007} or of the estimation approaches surveyed by \citet{grazzini2017} and \citet{platt2020}.
\section{Empirical Methodology}\label{sec:method}

\subsection{Two Structurally Distinct Estimators of the Aggregate Level}\label{sec:method-estimators}

The first estimator is a 10,000-household Agent-Based Model (ABM), specified in Appendix~\ref{sec:method-abm}. It draws a synthetic population calibrated to aggregate FRED medians and tests whether a utility-maximizing household decision rule, rational or behaviorally extended, can account for the shortfall; Section~\ref{sec:abm} reports its results. The second is a hazard framework, specified in Section~\ref{sec:hazard}, which abandons the household utility function entirely and models loan-level survival directly on real Freddie Mac loan-level data. It runs along two paths. Path A is a stratum-month discrete-time hazard model estimated by Poisson pseudo-maximum likelihood on the full 2017--2021 origination universe aggregated to stratum-months. Path B is a loan-level microsimulation with competing prepayment and default risks on a stratified 75,000-loan sample drawn from that universe.

I evaluate both estimators against the same \$764.7 billion empirical benchmark (Section~\ref{sec:method-benchmark}) and the same June 2022--November 2025 QT window, so their results are directly comparable even though they draw on different data and different behavioral assumptions. They differ in modeling paradigm --- household choice versus loan-level survival --- and, on the hazard side, in the use of real loan-level data. Population size differs too: 10,000 households against Path B's 75,000 simulated loans, with Path A estimated on the full origination universe. Section~\ref{sec:robustness} finds that this size difference does not drive the gap in their explanatory power. The data-source difference, though, remains an unresolved confound that qualifies the paradigm interpretation throughout. Figure~\ref{fig:architecture} maps the pipeline.

\begin{figure}[H]
\centering
\includegraphics[width=\textwidth]{fig7_architecture}
\caption{Pipeline architecture. Data sources (FRED macro series, NY Fed SOMA holdings, Freddie Mac loan-level performance) feed three estimators (the synthetic-population ABM and hazard Paths A and B), whose simulated roll-off paths are scored through one shared macro-accounting layer against the \$764.7 billion empirical benchmark of Section~\ref{sec:method-benchmark}. The two robustness cross-links are the cross-design test (real Freddie covariates into the ABM, Section~\ref{sec:robustness-crossdesign}) and the symmetric companion test (synthetic population into the hazard framework, Section~\ref{sec:robustness-synthesis}).}
\label{fig:architecture}
\end{figure}

The agent-based model's specification, and the structural extensions whose grounding Section~\ref{sec:abm}'s results reference, are given in full in Appendix~\ref{sec:method-abm} and Appendix~\ref{sec:method-frictions}.

\subsection{Empirical Validation Benchmark}\label{sec:method-benchmark}

I construct the empirical shortfall benchmark from the New York Fed's SOMA current face value series, netting realized roll-off against the actual phased redemption caps (\$17.5 billion per month June--August 2022, \$35 billion per month from September 2022) strictly within the active QT window, June 2022--November 2025. Accumulation is net rather than shortfall-months-only, for comparability with the Danish counterfactual of Section~\ref{sec:abm-danish}, whose simulated roll-off routinely exceeds the cap. The terminal gaps in Figure~\ref{fig:benchmark} are the cap-relative benchmark and the expectations-based complement below. One feature of the series is an artifact, and no result in this paper reads it as behavior: four months come in at exactly zero. The cause is the back-out's non-negativity clip, which pins any month whose reported roll-off falls below scheduled amortization and displaces the mass into the next month --- a month-boundary allocation artifact of the reporting series (run \texttt{h1\_zero\_months\_diagnosis}). Appendix~\ref{sec:robustness-seasonalfloor} traces the mechanism and shows that repairing the affected pairs flips this paper's frozen monthly-timing rule from pass to concede. That artifact is one reason every claim made from this benchmark is a claim about levels.

\begin{figure}[H]
\centering
\includegraphics[width=0.94\textwidth]{fig9_benchmark_construction}
\caption{The empirical benchmark. Left: realized monthly SOMA agency MBS roll-off against the phased redemption cap; the shaded area is the monthly shortfall, which accrues on a net basis (the two circled opening months are net settlement inflows). Right: the same construction cumulatively. The terminal gaps are the \$764.7 billion cap-relative benchmark and, against the New York Fed's May 2022 ex-ante projection, the \$87.8 billion expectations-based complement. Monthly series from the frozen production run's metrics; the projection's monthly path is derived from the published year-end levels under the pre-committed uniform-spread rule and asserted, in the generator, to reproduce the frozen expectations-run totals.}
\label{fig:benchmark}
\end{figure}

The empirical shortfall is approximately \$764.7 billion (unrounded: \$764.748 billion; every percentage-of-benchmark figure in this paper divides by this value), and it is the benchmark for simulated trapped liquidity in Sections~\ref{sec:abm} and~\ref{sec:hazard}. The SOMA face spans the Federal Reserve's full agency MBS book, Ginnie Mae collateral included, 20.4\% of face. Of that face, 66.2\% lies in the 2017--2021 vintages of the Freddie Mac loan-level sample (Section~\ref{sec:hazard}), so the hazard framework directly samples neither Ginnie collateral nor roughly one-third of the book by vintage. Section~\ref{sec:pathb} bounds both exclusions externally; Section~\ref{sec:limitations} states what each bound leaves open. The exclusions must be netted on the book's own joint cells, not counted separately. Out-of-window vintage is 33.7\% of face (2022 cohort 23.1\%, pre-2017 10.6\%), and the 20.4\% Ginnie share cuts across vintage. The conventional-and-in-window intersection is therefore 51.0\% on the book's own agency $\times$ vintage joint cells (bounded above by 51.8\% for cell truncation; 52.7\% under an independence approximation). The Freddie 2017--2021 conventional universe estimating Path B is about \emph{half} of SOMA book face. Two-thirds is the vintage-only marginal and overstates coverage. External-speed overlays (Section~\ref{sec:pathb}) bound Ginnie at \$20--\$47 billion by published observed speeds and out-of-window vintages at \$11.7 billion by observed cohort speeds from the replication ingest, each signed toward overstating trapped liquidity. On the simulated side an income-scaled curtailment assumption enters, averaging approximately 0.84\% annualized and consistent with GSE daily prepayment report estimates; it is added to each simulated roll-off and netted as a common dollar item in the shared accounting layer of Section~\ref{sec:robustness-hybrid}, so the empirical CPR back-out's series retains that component. Two implementations of the back-out exist --- the ABM's cohort-weighted, term-aware scheduled amortization and the hazard side's single-pool approximation --- and both amortize at the pass-through WAC, though scheduled principal follows the borrowers' note rate, that WAC plus a $+0.80$-point guarantee-fee and base-servicing wedge. On the note-rate basis the hazard-side CPR level rises $+0.27$ points, a deterministic schedule difference. Every trapped-dollar figure is bit-identical: the benchmark's dollars never pass through the CPR scale. Every peak lag and $U_2$ is unchanged at reporting precision, correlations differing only at the third decimal (run \texttt{coupon\_convention\_amortization}). Table~\ref{tab:theil} is scored on the corrected basis, and Table~\ref{tab:estimators}'s note states which rows use which.

\paragraph{Accounting bases.} Every recovery percentage in this paper is quoted on one of four accounting bases, differing only in two orthogonal choices --- pure accounting layers with no modeling content --- whether the common curtailment flow is netted out, and whether the estimation sample is reweighted to the SOMA book's coupon mix. The \emph{standalone-scorer basis} makes neither adjustment; the \emph{benchmark-consistent shared basis} nets that common curtailment flow (\$69.56 billion over the window) identically from every U.S. leg, a flat 9.1 percentage points of the benchmark for every U.S. estimator; the \emph{full-book basis} instead reweights the standalone figure to the book's weighted-average coupon (WAC) without netting; and the \emph{composed basis} applies both, the same 9.1-point netting to the full-book figure. Because the netting is bit-identical across central and null legs, it cancels exactly from the lock-in marginal and every cross-estimator gap --- an invariance at every cell of the calibration box and under a $\pm 25\%$ stress on the assumed curtailment rate (derivation and stress in Section~\ref{sec:robustness-hybrid}; the four bases per hazard leg in Table~\ref{tab:bases}). The cancellation is term-by-term monthly, not merely aggregate: the netted flow is a shared monthly rate times the shared actual-holdings path, so any curtailment time profile --- calendar-seasonal, or income elasticities differing across sub-periods --- cancels identically, moving only the common wedge. The income-to-rate mapping's functional form is immaterial to the marginal for the same reason: any alternative (concave, threshold, or piecewise in income) yields another shared profile. Main-text headline figures are quoted on the shared basis throughout.

I construct a second, expectations-based benchmark because the cap-relative one carries a tension: the caps were widely expected to be non-binding for agency MBS, so a shortfall measured against them mixes an anticipated rate-inelastic component with genuine surprise. Against the Federal Reserve's own ex-ante projection, projected agency MBS runoff is \$740.6 billion and realized runoff \$652.8 billion: an expectations-based shortfall of \$87.8 billion, 11.5\% of the cap-relative \$764.7 billion. That projection is the baseline path published in the New York Fed's May 2022 Open Market Operations annual report \citep{omo2021}, starting from the February 2022 balance sheet, incorporating the FOMC's May 2022 Plans, and modeling runoff from June 2022. I spread its year-end figures uniformly within each printed year and clip them to the window under a rule fixed before the numbers were seen. The report prints year-end levels only, and no staff projection at monthly granularity is public for any vintage covering this window. Equivalently, that projection already implied 75.6--88.5\% of the realized cap-shortfall, the range spanning defensible intra-2022 allocations of its printed annual net. The pre-committed uniform-spread central construction gives \$676.9 billion, or 88.5\%. The substantive point is that a source predating every estimate in this paper anticipated the large majority of the shortfall, and it stands on its own. A sharper reading it invites does not. Under the uniform-spread construction the 88.5\% sits 0.2 points from the 88.7\% the $\beta_1 = 0$ rate-inelastic null recovers on the benchmark-consistent basis, and I presented that near-equality as convergent evidence. I withdraw the presentation: the 0.2 points require two conditions at once. They require the in-sample 4.0\% turnover floor, because at the off-window floor the paper headlines, the null recovers 85.7\% and falls 2.8 points \emph{below} the projection, not 0.2 above it. And they require the uniform-spread allocation, because the settlement-aware allocation puts the anticipated share at 75.6\%, which the null clears by 13.2 points in-sample and 10.1 points at the headline floor. A $-2.8$ to $+13.2$ point span across two disclosed choices is not a convergence, so I now report the projection range and the null's recovery separately. Both recoveries are production-floor-form reads, so the comparison is against one form's null, not a form-invariant object. Under the additive form the null recovers 35.6\% on the benchmark-consistent basis at the headline floor (44.7\% standalone; Section~\ref{sec:robustness-floor}). What survives is coarse and is what I rely on: under the production floor form both objects put the anticipated rate-inelastic component in the large majority of the shortfall. That is a fact about where the caps were set, not the size of the rate response, since the caps sat above any prepayment path this book could have delivered. Leading with it understates lock-in; the reading that speaks to mortgage markets is the marginal's share of the expectations-based shortfall, below. The central 88.5\% rides on the uniform-spread treatment of the printed 2022 annual net, which nets first-half settlement inflows against June--December runoff. That is the less faithful of the two constructions to the realized settlement record. The settlement-aware allocation assigns those inflows to the first half (realized: holdings rose \$92.3 billion over January--May 2022 in the benchmark's own SOMA series), giving roughly three-quarters (75.6\%). It leaves the large-majority reading, and every gate and threshold of the run, unchanged. Parametric within-year ramps bound the rest. Linear 2022 ramps move the share only to 88.0--89.0\%, the December-2025 window clip spans 86.1--90.6\% between all-front and linear-back extremes, and the bracket pushing all 2025 runoff to the clipped December raises it to 115.3\%. The rate-inelastic-majority threshold survives at every allocation. Every percentage-of-benchmark figure in this paper remains a cap-basis quantity. Dollar objects and estimator rankings are basis-invariant, the decomposition rescaling by the $8.7\times$ benchmark ratio on the expectations basis. The identified lock-in marginal is 49\% of the expectations-based shortfall under the central uniform-spread construction at the headline off-window floor ($+\$42.6$ billion), 80\% at the in-sample calibration point ($+\$70.3$ billion). That shortfall is itself allocation-dependent, so the ratio falls with the anticipated share (roughly 23\% and 38\% at the 75.6\% allocation) and is undefined at the extreme 115.3\% bracket, where the expectations-based shortfall turns negative. Every one of these ratios rests on one condition I cannot test: it reads the expectations-based shortfall as a residual free of lock-in, when the projection's own prepayment model may already have embedded part of the channel. That is unrecoverable from the published path of holdings. Any lock-in the projection anticipated therefore sits inside the \$740.6 billion, not the \$87.8 billion residual, and the ratios are upper bounds on lock-in's share of the genuine surprise, exact only if the projection embedded none of the channel. A second condition is likewise untestable from the published record: the projection conditions on the spring-2022 market-implied rate path, well below what was realized, so part of the residual is rate-path surprise, not lock-in. Both are properties of the projection, not of any one quotation of it, so they govern every site where these ratios appear (introduction, this section, conclusion). Each reads as an allocation-conditional upper bound on lock-in's share of a mixed lock-in-and-rate-path surprise, not a partition of it. One question of standing precedes this accounting: a shortfall against the MBS cap is not by itself evidence that a stated policy objective was missed. The Committee's operating object was the aggregate securities portfolio's decline, Treasury runoff ran against its own caps throughout, and nothing in this paper shows that reserves or the aggregate path came in off the intended course because MBS ran slow. The \$764.7 billion measures the \emph{composition} of the decline --- MBS principal that did not arrive on the book's own schedule --- and the duration consequences it carries (Section~\ref{sec:discussion-fallout}): how the decline was delivered and what that cost in portfolio structure, not a claim that the decline itself fell short.

The benchmark measures shortfall against the redemption cap, a ceiling on runoff. The cap is not a forecast. New York Fed staff projected the below-cap regime from the month the \$35 billion cap took effect \citep{nyfed2022}, so the cap-relative shortfall is an upper bound on the shortfall relative to any reasonable ex ante expectation. The complement above quantifies it: the expectations-based shortfall is \$87.8 billion, 11.5\% of the cap-relative figure. I retain the cap as the headline counterfactual for two reasons. First, it is the only policy-committed quantity available. The projection underlying the complement is a baseline forecast, not an FOMC-stated target, and adopting any forecast---the Federal Reserve's or a model's---would be circular with respect to models of expected prepayment, the object under study. Second, it is policy-relevant: it measures the gap between the maximum balance-sheet reduction the FOMC authorized and the reduction the mortgage structure delivered. ``Trapped liquidity'' and ``shortfall'' accordingly denote the shortfall relative to the cap throughout the paper.
\section{Agent-Based Model Results: A Synthetic-Population Stress Test}\label{sec:abm}

This section is a stress test of the household-choice hypothesis, not the paper's main estimate, and the first thing to report about it is what its own cross-design leg establishes about its scope. That result governs how everything below should be read. Re-running the same decision rule on real Freddie Mac structural covariates recovers 59.3\% of the benchmark on the committed draw and 60.2\% averaged over fifty seeds---above the 50\% threshold fixed before that run at which this section's paradigm reading would be undercut, and undercutting on all fifty of them. That cross-design leg runs under the recalibration rule the model already applies elsewhere. Recovery rises to 76.3\% once the sample is reweighted to the SOMA book's own coupon composition. Three things bound that reading and travel with it wherever it appears. First, the same reweight moves the frozen-calibration leg the other way, down to 12.6\%. Second, the covariates swapped in are structural only: the loan records contain no analogue for mobility desire or loss aversion. So no leg here replaces a synthetic behavioral population with a real one. Third, the verdict the run supports is that the contrast cannot, on that test alone, be cleanly attributed to modeling paradigm rather than to calibration and data source (Section~\ref{sec:robustness-crossdesign}). What this section falsifies is therefore a \emph{synthetic-population} household-choice specification: no specification tested here explains the Federal Reserve's cap shortfall. Why it does not is not cleanly separable from the population it was tested on. A second qualifier travels with the verdict and is developed where it binds: the multi-vintage stage moved against its own predicted direction and the inversion is undiagnosed (Appendix~\ref{sec:robustness-extensions}). The paradigm reading this section was written to advance is withdrawn on the strength of that crossing, and Section~\ref{sec:robustness-crossdesign} records the withdrawal where the test lives.

\subsection{The Synthetic-Population Estimate: Baseline, Production, Extensions, and the Burnout Falsification}\label{sec:abm-results}

Section~\ref{sec:identification} \emph{declines} that same cross-design recovery as evidence that the imported lock-in elasticity is too small, since a recovery level is not what that design identifies. Read as a scope limit on this section's verdict --- which is how it is used above --- it survives that objection. Neither use licenses the other. The production estimate is that synthetic-population special case, and the case in which the test bites. Every specification tested here moves further from, or fails to close, the \$764.7 billion benchmark (Figure~\ref{fig:recovery}). A purely rational baseline projects a structural cash flow deficit of \$544.0 billion (71.1\% of benchmark). The production specification adds scheduled amortization, a live CUSIP-level SOMA parse, a corrected empirical CPR back-out, a settlement-lag kernel, and the previously excluded 15-year MBS book in cohort weighting and scheduled amortization (structural terms only). Its headline quantity is the seed distribution, not any single draw. The fifty-seed mean is \$103.7 billion, 13.6\% of the benchmark. The frozen seed-42 draw (\$91.0 billion, 11.9\%, near the 34th percentile) is the reproducibility anchor the run manifests freeze. The kernel is retained in production, and both frozen run manifests record it as applied. But it is a documented timing null, not a fit improvement. A pre-fold-in ablation left the peak cross-correlation unchanged at $r = +0.192$ at lag $-3$ with the kernel on or off, and the kernel-applied production run peaks at $+0.195$ at that lag. The production specification is not kernel-free, and every figure reported here is the kernel-applied run's. Seed dispersion dominates the third digit. The seed-level standard deviation is \$24.5 billion ($\pm$1~SD share band 10.4--16.8\%). The 95\% CI of the mean, \$96.9--\$110.5 billion, prices only simulation noise in the mean itself. The 34th percentile is pinned only to within roughly $\pm$7 percentage points from 50 draws. The ABM's share is therefore better read as on the order of 10--17\% than as a three-digit point estimate (Figure~\ref{fig:mc}). The path fit is worse than the level fit and no extension repairs it (Figure~\ref{fig:abmpaths}). Simulated mean CPR is 11.68\% against an empirical 5.14\% (ABM basis). The raw path correlation is negative at lag 0 ($r = -0.318$), and nothing in a $\pm$3-month window exceeds the +0.19 at the three-month lead. The path $R^2$ is $-6.984$. Throughout, $R^2 = 1-\mathrm{SSE}/\mathrm{SST}$ against a mean-only null, so a negative value is a fit worse than a constant at the empirical mean. The claimed object is the aggregate level. The path statistics are reported as a disclosed failure, not as a fitted quantity. No estimator in this paper reproduces the monthly path once the window's shared trend is removed (Section~\ref{sec:pathb}), and no level result anywhere in the paper is scored on path fit. Three literature-grounded behavioral extensions move the model to 54.9\%, then 45.1\%, then 33.7\% of the benchmark (Figure~\ref{fig:waterfall}). They are a debt-to-income constraint with asymmetric loss aversion and a rate-shock freeze, income-scaled curtailment, and multi-vintage coupon composition. The first is a recalibration artifact, not evidence for loss aversion: a real-world turnover floor held fixed under a steeper loss-aversion penalty mechanically inflates rate sensitivity. None touches the CPR path's central failure. Each acts on the level of simulated mobility, not its timing. A fourth variant fails more informatively. Path-dependent prepayment ``burnout,'' implemented as survivor selection in a closed population of 10,000 agents, collapsed simulated CPR to approximately zero and raised trapped liquidity to 147\% of the benchmark, because that population depletes its mobile mass in the low-rate 2020--2021 window before QT begins. That is an absorbing state incompatible with a real MBS pool that continuously receives new originations, and \citepos{lesniewski2026} selection identity in its degenerate limit. It is a negative finding, not used in production, and why burnout enters Section~\ref{sec:hazard} at the cohort level. Appendices~\ref{sec:method-abm}, \ref{sec:method-frictions}, \ref{app:abmspec}, \ref{sec:robustness-extensions} and~\ref{sec:robustness-pipeline} specify the decision rule, every behavioral parameter and its source, the macroeconomic-friction and structural extensions, the three behavioral stages' mechanisms, and the 15-year fold-in. The replication package (its run ledger) reproduces them end to end. This section reports results only.

\begin{figure}[H]
\centering
\includegraphics[width=\textwidth]{fig1_recovery_by_estimator}
\caption{Benchmark recovery by estimator, on the shared accounting basis (standalone-scorer values in parentheses, both bases collected in Table~\ref{tab:bases}; ABM-family rows run through that layer natively). The production ABM row shows the fifty-seed mean (open diamond) with its $\pm 1$~SD band alongside the frozen seed-42 draw (solid dot); Path A's dot is an aggregate-context, non-corroborating exhibit excluded from headline figures (Section~\ref{sec:patha}). Raw lag-0 correlations are reference-only (figure footer; Section~\ref{sec:pathb}); level recovery alone does not rank models.}
\label{fig:recovery}
\end{figure}

\begin{figure}[H]
\centering
\includegraphics[width=\textwidth]{fig3_abm_waterfall}
\caption{The ABM stress-test bridge. Each stage applies the next correction or extension to the household-choice model, and every pre-committed behavioral extension moves the recovery away from, not toward, the benchmark (the multi-vintage stage against its own expected direction, Appendix~\ref{sec:robustness-extensions}): rational baseline 71.1\% $\to$ behavioral gates 54.9\% (a recalibration artifact, Appendix~\ref{sec:robustness-extensions}) $\to$ income-scaled curtailment 45.1\% $\to$ multi-vintage composition 33.7\% $\to$ production corrections 13.2\% (a data-and-accounting stage, not a behavioral one) $\to$ 15-year fold-in 11.9\% (the production headline of Section~\ref{sec:abm-results}; the fold-in itself is worth $-1.3$ points) $\to$ native 15-year gate 11.1\% (the Berger-run calibration-stage variant of Section~\ref{sec:abm-danish}). Stage levels from the frozen run manifests catalogued in the replication package's run ledger.}
\label{fig:waterfall}
\end{figure}

\subsection{Interpretation}\label{sec:abm-interp}

Every ABM specification tested here widens the gap to the benchmark or leaves it unchanged, and its CPR path stays uncorrelated or negatively correlated with the empirical path throughout. Section~\ref{sec:robustness}'s scale-sensitivity replication finds no size effect on the production specification. Two possibilities remain: either household-level decision-making does not govern the majority of the shortfall, or this specification is missing a mechanism. The cross-design leg above bears directly on that either-or without settling it. Real structural covariates move the same rule's recovery from 13.6\% to 60.2\% averaged over fifty seeds, which looks like a missing-population problem rather than a missing-mechanism one; but the behavioral layer stays synthetic in both legs, so the alternative is narrowed, not eliminated. The clean falsification bypasses that free parameter entirely. The simulated-minimum-distance two-moment fit (Section~\ref{sec:robustness-crossdesign}) returns joint infeasibility under a verdict vocabulary fixed before execution, its nearest admissible fit overshooting the benchmark at 107.1\%, and it does not run through the recalibration scale that qualifies the cross-design reading. Section~\ref{sec:hazard} tests the first branch directly on a different data source and modeling paradigm. Whichever explanation holds, these results weigh against household mobility choices, rational or behavioral, as the primary price the Federal Reserve paid for secondary-market liquidity. The institutional counterpart is the securitization-policy alternative examined in Section~\ref{sec:abm-danish}: whether the par-payoff rule itself, not household behavior, drives the trap. One asymmetry deserves statement. Section~\ref{sec:identification} insists that this design identifies a differential and not a level, yet I compare the ABM with Path B on levels, because the ABM admits no counterpart to the $\beta_1 = 0$ null. The reason is structural: in \eqref{eq:pathB} the elasticity is a separable term over an independently measured floor that survives the elasticity's removal, so the null stays anchored to observed baseline turnover and lands at 85.7\% of the benchmark. The ABM has no such term. Its 4--5\% turnover floor is not a component of \eqref{eq:abmmove} but a \emph{product} of it, found by searching the mobility-desire scale $\theta$ until the penalty-bearing rule returns 4--5\% CPR at an 8\% market rate; zeroing the penalty therefore removes the floor along with the channel. With the production $\theta = 43{,}882.8$ held fixed, the resulting leg returns 40.7\% CPR at the calibration point against the production rule's 4.8\%, over-retires the book, and scores $-259.4$\% of the benchmark. Re-deriving $\theta$ under the null restores the anchor to 4.7\% but scores $+116.3$\%, because the re-search fits the null to the very turnover level the differential is meant to measure. The two readings imply marginals of $+270.4$ and $-105.3$ points: they do not bracket a magnitude; they disagree in sign, and neither is interpretable. I therefore report no ABM differential, and the estimator contrast of Section~\ref{sec:hazard} stays a contrast of levels, with the qualification that entails (run \texttt{abm\_null\_feasibility}; a feasibility probe, not a pre-committed estimate, and nothing in this paper's ranges depends on it).
\section{The Hazard Framework: A Structurally Distinct Estimator}\label{sec:hazard}

This is the estimator on which the paper's central decomposition rests. The ABM of Section~\ref{sec:abm} does not approach the hazard framework's aggregate recovery. On Section~\ref{sec:robustness-hybrid}'s shared accounting basis the microsimulation form recovers 91.3\% at the headline off-window floor and 97.9\% at the in-sample calibration point. The cohort-hazard form corroborates at 112.4\% but is excluded from headline ranges (Section~\ref{sec:patha}). The ABM's production specification recovers 11.9\% on the frozen draw (fifty-seed mean 13.6\%). Section~\ref{sec:robustness-crossdesign}'s cross-design leg leaves open whether that contrast is modeling paradigm or data source. Even its highest-recovery behavioral extension, at 54.9\%, falls well short. I attribute that figure in Section~\ref{sec:abm-results} to a mechanical recalibration artifact. The uncorrected rational baseline's 71.1\% predates that section's scheduled-amortization, CUSIP-parse, and CPR-back-out corrections and is not in the production estimate family. Both hazard paths and the ABM face the same aggregate benchmark. Only Path B's recovery is accurate on the benchmark-consistent accounting basis, and no estimator's monthly co-movement with the empirical path (Path B's included) survives removal of the window's shared trend (Section~\ref{sec:pathb}). Path A's monthly path fit is no better than the ABM's. Sections~\ref{sec:patha} and~\ref{sec:hazard-interp} quantify that limitation because it bears on each recovery figure's evidential weight.

\subsection{Motivation}\label{sec:hazard-motivation}

Section~\ref{sec:abm}'s pattern is difficult to reconcile with the hypothesis that the residual reflects a missing rational or behavioral household mechanism that a sufficiently rich utility function would eventually capture. It is more consistent with the hypothesis that the binding constraints on the empirical prepayment path lie at least partly outside the household decision function: in loan-level heterogeneity, servicer pipeline dynamics, and cohort-level path dependence that a fixed, representative-household population struggles to represent.

The hazard framework introduced in Section~\ref{sec:method-estimators} is built to test this hypothesis. It abandons the household utility function---no mobility-desire parameter, no loss-aversion coefficient, no utility-maximizing decision rule---and models loan-level survival directly, using proportional hazards (proportional above a baseline turnover floor; Section~\ref{sec:pathb}) estimated on real loan-level data; the closed-population survivor-selection failure of Section~\ref{sec:abm-results} motivates a cohort-level treatment of burnout, not a per-loan one. For comparability the framework retains the same \$764.7 billion empirical benchmark and the same QT window definition as the ABM. It comprises two paths: a stratum-month discrete-time hazard model estimated on aggregated loan-level data, and a loan-level microsimulation calibrated to literature elasticities with explicit competing risks. I use competing risks in the sense of \citet{deng2000}, who estimate prepayment and default as dependent competing hazards over unobserved borrower heterogeneity.

\begin{table}[H]
\centering
\caption{Compact specification of the two hazard paths. Every entry restates
content specified in Sections~\ref{sec:patha} and~\ref{sec:pathb} or
Appendix~\ref{app:params}.}
\label{tab:specbox}
\footnotesize
\begin{threeparttable}
\begin{tabular}{@{}p{2.2cm}p{5.9cm}p{5.9cm}@{}}
\toprule
 & Path A (estimated cohort hazard) & Path B (literature-calibrated microsimulation) \\
\midrule
Unit; sample & stratum-month; full Freddie 2017--2021 origination universe (296 four-way cells) & loan-month; stratified 75{,}000-loan draw from the same universe \\
Outcome; estimation & prepaid UPB with log-exposure offset; Poisson pseudo-maximum likelihood, log link \eqref{eq:pathA} & monthly prepayment hazard \eqref{eq:pathB}; nothing estimated on the benchmark; every parameter fixed a priori, though the dominant level-setter (the turnover floor) is calibrated in-window (Sections~\ref{sec:pathb}, \ref{sec:robustness-floor}; Appendix~\ref{app:params}) \\
Age profile & seasoning spline, knots $\{12,24,36,60,84,120\}$ months (seven columns) & 100~PSA ramp to 6\% CPR at month 30 \\
Rate-gap term & estimated: $+0.65$ standardized ($+0.29$ per 100 bp), spec v4 & imported from \citet{liebersohn2024} via \eqref{eq:beta1}; 5.5--7.7\% band \\
Stratification & 295 stratum fixed effects (vintage $\times$ coupon $\times$ FICO $\times$ LTV) & stratum-level burnout state $B_{s(i),t}$; standardized FICO/LTV ex-ante priors \\
Time effects & eleven calendar-month dummies (spec v4) and the dynamic friction index & none; competing default hazard \eqref{eq:default} and the Markov delinquency pipeline; the omitted housing-activity level term is priced rather than modelled, by the scaled-null variant of Section~\ref{sec:identification} (run \texttt{scaled\_null\_housing\_activity}) \\
Turnover floor & none (stratum fixed effects set baseline levels) & 4\% CPR floor, max form; competing-risks form tested (Section~\ref{sec:robustness-floor}) \\
Inference & stratum-cluster and temporal moving-block bootstraps (Table~\ref{tab:bootstrap}); no claim survives them---the rate-gap coefficient's sign is not distinguishable from zero under any bias-respecting construction (Section~\ref{sec:patha}) & stratified loan-level bootstrap (200 replicates) for sampling; calibration box for parameters (Section~\ref{sec:robustness-floor}) \\
Validation & temporal holdout, reported as a diagnostic only (Section~\ref{sec:patha}) & Fannie Mae external replication (marginal $+8.68$ points against Freddie's $+9.20$; the pre-committed envelope it clears is 11.06 points wide) and the synthetic companion at 106.0\% (Sections~\ref{sec:pathb-fannie}, \ref{sec:robustness-synthesis}) \\
Role & aggregate context only; does not corroborate the result; excluded from headline figures & headline estimator: at the off-window floor central 91.3\%, $\beta_1 = 0$ null 85.7\%, marginal $+5.6$ points; at the in-sample calibration point 97.9\%, 88.7\%, $+9.2$ points (all shared basis) \\
\bottomrule
\end{tabular}
\end{threeparttable}
\end{table}

\subsection{Path B: Literature-Calibrated Microsimulation}\label{sec:pathb}

\subsubsection{Specification}

The first hazard path brings competing risks, literature-calibrated elasticities, and a cleaner isolation of the Danish counterfactual, without reintroducing the utility-maximizing machinery of the ABM. It simulates a stratified sample of 75,000 loans, 40,234 of which are still active when the QT window opens (Appendix~\ref{app:params}). Each loan carries static covariates---FICO, region, original LTV, coupon, and a stratum identifier keyed on the Path A specification less its LTV axis---and dynamic state: loan age, balance, and a regime-specific rate gap. Equity enters only through the static original-LTV covariate: no updated CLTV and no local house-price feed exists anywhere in the pipeline, so equity dynamics over the window are unmodeled. That is a disclosed design boundary for this turnover-dominated, deeply out-of-the-money window; the FICO/LTV covariate ablation reported below brackets what the static block contributes. I specify each loan's monthly hazard of prepayment as
\begin{equation}\label{eq:pathB}
h^{\mathrm{prep}}_{i,t} = \max\!\left\{\, \underline{h},\; h_0\!\left(a_{i,t}\right) \exp\!\left( \beta_1\, g_{i,t} + \beta_F\, z^{\mathrm{FICO}}_{i} + \beta_L\, z^{\mathrm{LTV}}_{i} + \beta_B\, B_{s(i),t} \right) \right\},
\end{equation}
with a competing hazard of default
\begin{equation}\label{eq:default}
h^{\mathrm{def}}_{i,t} = h^{\mathrm{def}}_0 \exp\!\big( \gamma_r\, \sigma_{i,t} + \gamma_F\, z^{\mathrm{FICO}}_{i} + \gamma_L\, z^{\mathrm{LTV}}_{i} \big), \qquad \sigma_{i,t} = \max\{0,\ -\tilde g_{i,t}\},
\end{equation}
where $\sigma_{i,t}$ is rate stress ($\tilde g_{i,t} = g_{i,t}/100$ is the rate gap of \eqref{eq:pathB} expressed in decimal units, so $\sigma$ is the decimal magnitude of an out-of-the-money gap), and $h^{\mathrm{def}}_0 = 3\times10^{-4}$ monthly, a prior implying $0.36\%$ per year against a realized Freddie Mac default rate of $0.31\%$ over the window on this paper's own GMAR extraction. The calibration did not use that corroboration: the prior was fixed before that series was assembled. Every Path B parameter is fixed a priori in the configuration, and none is estimated on the benchmark; Appendix~\ref{app:params} collects the full parameter vector with provenance, together with the delinquency transition matrix. Table~\ref{tab:hazard-notation} maps every symbol here to its Path A counterpart. Here $g_{i,t}$ is the rate gap: the loan's coupon minus the prevailing market rate, in 100 bp units, negative when the borrower is locked in. The positive $\beta_1$ of \eqref{eq:beta1} therefore suppresses prepayment exactly when the par-payoff penalty binds. $B_{s(i),t} \in [0,1]$ is the cohort-level burnout state of loan $i$'s stratum $s(i)$: the stratum's cumulative prepaid UPB as a share of its original UPB, capped at one. Every stratum is initialised at zero at the window open, so the state accumulates only within-window prepayment against an original-UPB denominator that still includes pre-window-terminated loans---window-cumulative, not life-cumulative. With $\beta_B = -0.5$, the burnout state scales the hazard by $e^{-0.5 B} \in [0.61, 1]$. Finally, $\underline{h}$ is the baseline turnover floor: the monthly single-month-mortality (SMM) equivalent of 4\% annual CPR, $\underline{h} = 1-(1-0.04)^{1/12} \approx 0.0034$. It floors total turnover, not a strictly involuntary component of it: deep-discount-cohort turnover at this level includes discretionary life-cycle moves and cash-out refinancings alongside strictly involuntary events. It is also held \emph{flat in loan age}, an assumption the very ladder used to calibrate it does not support. Pooled 2017--2019 out-of-the-money turnover rises monotonically with age, from 3.68\% below twelve months to 10.99\% at twenty-four to thirty-six (Appendix~\ref{app:floormech}). I flag the assumption; I do not sweep it. The reason is the same appendix's second finding: the 10.99\% cell is carried entirely by the 2019 falling-rate leg, so the age profile is substantially refinancing-contaminated, and an age-varying floor built from it would inherit that contamination. An age dimension therefore remains the one calibration axis this paper sweeps in neither level nor time, and a clean version of it needs a rate-regime-pure age ladder this design does not have. No component decomposition exists in this design either, and the strictly-involuntary share (death, divorce, and forced relocation) is plausibly well under half. The floor's semantics do institutional work in selecting the U.S.-intercept Danish anchor (Section~\ref{sec:abm-danish}) and in retaining the max floor form (Section~\ref{sec:robustness-floor}). There the operative premise is only that baseline turnover cannot fall below its observed deep-discount level when the payoff rule changes, which the total-turnover reading supports. Those choices are conventions consistent with that premise, not consequences of a strictly-involuntary decomposition. The floor's level is anchored to the same empirical observation as the ABM's mobility calibration (Appendix~\ref{sec:method-frictions}): turnover on deeply out-of-the-money discount cohorts during 2023--2024, which lies \emph{inside} the June 2022--November 2025 evaluation window. No Path B parameter touches the benchmark dollar figure, but the specification's dominant level-setter is measured from in-window prepayment behavior. The verification content below is the marginal over the $\beta_1 = 0$ null, not the aggregate level itself. The pre-committed floor sweep of Section~\ref{sec:robustness-floor} shows that even this marginal is conditional on the floor calibration: it spans $+10.8$ to $+5.5$ points at the central elasticity across the empirically plausible 3--5\% floor range ($+9.2$ at the production 4\% floor). That interval, not any point inside it, is the identified object. Its positivity at every floor and elasticity-band point swept is a consistency check the window's rate configuration cannot let fail (Section~\ref{sec:identification}). The floor's functional form matters too. Replacing the hard maximum of \eqref{eq:pathB} with a competing-risks combination, under which the elasticity is never censored, moves the production marginal to $+11.25$ points and makes it nearly floor-invariant ($+11.29$ to $+11.21$ across the same 3--5\% range; Section~\ref{sec:robustness-floor}). The max-form marginal's floor-dependence is censoring-driven. It is not evidence about the elasticity. The floor binds in 36\% of the 1,683,124 evaluated loan-months in the production U.S. run. A substantial share of the aggregate level rides on the floor, consistent with the no-lock-in null of Section~\ref{sec:hazard-interp}, which shows the rate-inelastic model recovers 97.8\% of the benchmark under the standalone scorer (88.7\% on the shared accounting basis). Those are production-floor-form levels: under the additive form the same null recovers 44.7\% at the off-window floor against the production form's 94.8\%, both on the standalone scorer (Section~\ref{sec:robustness-floor}). Proportionality in \eqref{eq:pathB} accordingly holds above the floor, not through it. The baseline hazard $h_0$ follows the 100~PSA seasoning ramp (CPR rising 0.2 percentage points per month of loan age to 6\% at month 30), applied at monthly frequency via $\mathrm{SMM} = 1-(1-\mathrm{CPR})^{1/12}$. The rate-gap coefficient is calibrated from \citet{liebersohn2024}, who estimate that a 100 bp increase in the prevailing rate above the locked-in coupon reduces quarterly mobility probability by 5.5\% to 7.7\% depending on specification. I adopt the 6.5\% midpoint of that band as the ex-ante central scenario, with the band edges as sensitivity bounds, converted to a monthly hazard via the survival-function transform
\begin{equation}\label{eq:beta1}
\beta_1 = -\ln\!\left[\frac{1-\big(1-(1-\delta)\,P_q\big)^{1/3}}{1-\big(1-P_q\big)^{1/3}}\right],
\end{equation}
where $\delta$ is the proportional quarterly mobility decline per 100 bp of rate gap, $P_q = 0.06$ a reference baseline turnover probability occupying the transform's quarterly slot, and the inner terms are the monthly single-month mortalities implied by constant-hazard compounding within the quarter, $1 - P_q = (1 - h_m)^3$. The reference probability is a conversion auxiliary, not a measured quarterly mobility rate, and it sits well above \citepos{liebersohn2024} own 1.5\% zero-gap quarterly moving level (Section~\ref{sec:robustness-crossdesign}). I retain it because the conversion is nearly insensitive to it over $P_q \in (0, 0.12]$. Evaluated at the central decline, $|\beta_1|$ is $0.0675$, $0.0679$, $0.0686$ and $0.0701$ at $P_q = 0.015$, $0.03$, $0.06$ and $0.12$: quadrupling the assumed moving level above the source's own moving level moves the imported coefficient by under four percent (Appendix~\ref{app:params}). Equation~\eqref{eq:beta1} is the transform implemented in the codebase. At the central $\delta = 0.065$ it evaluates to $\beta_1 = 0.0686$ (unrounded $0.068571$), the value carried in the production configuration and rounded to $0.069$ in Table~\ref{tab:params}. The superseded continuous-hazard variant, a distinct object the simulation does not use, is documented in Appendix~\ref{app:params}. The transform yields the dimensionless log-hazard-ratio implied by a quarterly probability decline (it replaces a divide-by-three approximation, not a hazard itself), and \eqref{eq:pathB} then applies $\beta_1$ linearly per 100 bp of gap. At QT-era gaps of roughly $-380$ bp this compounds the local quarterly estimate to a hazard ratio of about $0.77$ ($\approx e^{3.8 \times -0.069}$), a log-linear extrapolation beyond the range in which the source elasticity was estimated. I flag it as a functional-form assumption and test it directly with a pre-committed concave-transform rerun of the paired legs. The extrapolation proves immaterial to the level and material to the marginal on the relative reading that ranking rests on, and the transform dimension travels with the marginal's operative uncertainty (Table~\ref{tab:uncertainty}; the test and its materiality reading are in Appendix~\ref{app:params}).

\subsubsection{Transport assumptions}

Two transport assumptions travel with the import, and I name both here. The first is an aggregation transport that assumes response homogeneity: the source estimand is a ZIP-code-level moving hazard applied as a loan-level log-hazard coefficient. The second transport concerns survival: the estimand is estimated on a population not conditioned on mortgage survival, while \eqref{eq:pathB} applies it to a pool from which 34,734 of the 75,000 sampled loans had already prepaid before the window opened---a survival-selected pool depleted of the rate-responsive. This transport's direction is therefore signed toward attenuation, and priced. Under the gamma-frailty selection identity the attenuation is $\eta = S_{\mathrm{pre}}^{\xi}$ at the measured pre-window survival share $S_{\mathrm{pre}}$. The off-window marginal then reads $+5.1$/$+4.6$/$+3.7$ points at $\xi = 0.25/0.5/1$ (run \texttt{attenuation\_sensitivity}). The gamma-frailty parameter $\xi$ is not estimable in this design and is displayed as the free parameter. Path A's burnout coefficient, the one object that speaks to it, is not distinguishable from zero under either bootstrap scheme.

Each month both hazards are evaluated for every active loan. Appendix~\ref{app:params} holds the production step order (fractional prepayment, the default draw and delinquency pipeline, the proportional clamp that never binds), the competing-risks taxonomy of the construction, and the parameter table (Table~\ref{tab:params}).

The transform prices how the elasticity extrapolates in the gap. A second mapping assumption governs what it applies to: \citepos{liebersohn2024} estimand is a moving probability, while \eqref{eq:pathB} applies the implied $\beta_1$ to the entire voluntary hazard, moving and refinancing alike. A pre-committed mixture bracket prices the convention at both floors (runs \texttt{moving\_share\_bracket}, \texttt{moving\_share\_bracket\_offwindow}): $1-\mu$ of the voluntary hazard at its no-response level, $\mu$ at the full response, turnover floor unscaled, parity at $\mu = 1$. The marginal is $+\$37.7$ billion ($+4.9$ points) at $\mu = 0.5$ and $+\$19.2$ billion ($+2.5$ points) at $\mu = 0.25$, against the production $+\$70.3$ billion at $\mu = 1$, close to proportional in $\mu$. At the headlined 4.991\% off-window floor it is $+\$25.7$ billion ($+3.36$ points) at $\mu = 0.5$ and $+\$13.9$ billion ($+1.82$ points) at $\mu = 0.25$, against $+\$42.6$ billion at $\mu = 1$. That scaling is mildly super-proportional, since the higher floor censors the elasticity in more loan-months and censored contributions do not scale with $\mu$. For this book's rate configuration the convention is therefore signed: applying the imported elasticity to less of the hazard can only shrink the marginal, roughly by its share. No termination-reason field exists in the ingested data \citep[cf.][on the refinancing-versus-moving mix]{amromin2020}, so $\mu$ is a bracketing parameter rather than an estimate, and the production $\mu = 1$ remains the reported form. \citet{graybill2026} corroborate the imported band on exactly this margin: their transaction-record sale-hazard estimate measures the moving component that $\mu$ varies.

The floor enters \eqref{eq:pathB} as a single constant $\underline{h}$ broadcast to every loan-month, with no rate-gap, calendar-time, or covariate argument. Baseline turnover is therefore treated as acyclical, and every rate-responsive movement in prepayment is assigned to the elastic channel. The marginal-over-null design cancels the floor's level: the central and $\beta_1 = 0$ legs carry the identical constant. It does not cancel a floor co-varying with the rate cycle. To the extent that lock-in also suppresses involuntary moves, that co-movement is absorbed into the estimated elastic marginal. The in-window marginal is therefore best read as total rate-attributable prepayment suppression, an upper bound on the strictly-voluntary elasticity's contribution. At the off-window calibration the ordering reverses. The demotion's window component (roughly \$20.9 billion; Section~\ref{sec:robustness-floor}) is baseline turnover that the off-window floor classifies as rate-suppressed and the marginal no longer counts, so the off-window $+\$42.6$ billion bounds total rate-attributable suppression from below; the matched-depth $+\$63.5$ billion restores that component. The two calibrations bracket two estimands. They are not conservative and liberal readings of one estimand. The level sweep there already brackets any plausible constant-floor excursion (the 3--5\% range keeps the marginal within $+10.8$ to $+5.5$ points at the central elasticity). A within-run cyclical variant sets $\underline{h}_t = 4\% \times (1 + \kappa z_t)$, with $z_t$ the window-standardized 30-year mortgage rate and $\kappa = 0$ nesting (and in the run exactly reproducing) the production specification. It moves the marginal between $+7.1$ and $+9.4$ points across $\kappa \in [-0.5, +0.5]$, every cell inside the level sweep's bracket. Appendix~\ref{sec:robustness-seasonalfloor} decomposes that grid with a pre-committed permutation run. Most of its span survives losing the floor path's time order, so the grid measures floor dispersion before it measures cyclicality. The surviving co-movement component is real, small, and bounded there. The window-mean invariant fails, via the floor's non-negativity clip, in the stress cell that sets the bracket's low end.

\subsubsection{Aggregation and the bounds that travel with it}

Aggregation scales simulated prepaid UPB by the ratio of the Fed's MBS book to the 75,000-loan sample balance, with no reweighting toward Ginnie Mae collateral (20.4\% of book face) or vintages outside 2017--2021. The Freddie-estimated hazards extrapolate to those segments as-is. This is not behaviorally one-to-one. Ginnie Mae pools are FHA/VA collateral with a distinct borrower composition and a servicer buyout channel absent from the conventional book, one whose exercise turns on servicer balance-sheet capacity as well as borrower behaviour \citep{aiello2022}. The framework therefore bounds the composition error of treating the Ginnie share as conventional. It does not model that error. Along the coupon dimension, Section~\ref{sec:robustness-hybrid}'s full-book reweighting (sample WAC 3.9\% mapped to the book's 2.49\%) moves Path B from 107.0\% to 109.1\%, about two percentage points. That is a mixed-basis figure: 3.9\% is an unweighted mean of borrower note rates, 2.49\% a face-weighted mean of security pass-through coupons. Converted to one convention, the move is $+0.2$ points (run \texttt{coupon\_convention\_reweight}). Along the agency dimension, Ginnie Mae fixed-rate speeds ran 1.2--2.8 percentage points above Freddie Mac's across the window (5.2\% versus 4.0\%, February 2023; 10.3\% versus 7.5\%, May 2025; Ginnie Mae Global Markets Analysis Report, Recursion data). The involuntary buyout channel is a large minority of that difference, not most of it. The forty-two-month mean Ginnie--Freddie CPR gap of $+2.14$ points splits into a voluntary CRR (conditional repayment rate, prepayment net of the default-driven channel) $+1.35$ (63\%) and a buyout CDR $+0.85$ (40\%). The shares sum past one hundred because the CPR $\approx$ CRR $+$ CDR extraction carries a small residual. The most buyout-favourable snapshot (CDR 2.1\% versus 0.4\%, May 2025) is the buyout-dominated month, not the representative one. Treating the 20.4\% Ginnie face share as conventional thus mis-states book-wide CPR by at most $0.204 \times (1.2\text{--}2.8) \approx 0.24$--$0.57$ percentage points. At the null-to-central sensitivity of roughly \$83 billion per CPR point (window-average book face times one point of annual prepayment times three and a half years), that gives a bound of \$20--\$47 billion (2.6--6.2\% of benchmark), signed toward overstating trapped liquidity. The bound is one-sided because the observed differential is: published speeds run Ginnie above Freddie throughout, so the Ginnie-slower branch, signed toward understating trapped liquidity, does not arise in this data.

The vintage dimension carries no bound of either kind. The 2022 vintage alone is 23.1\% of book face, and pre-2017 vintages another 10.6\% (Appendix~\ref{sec:robustness-wal}), both outside the 2017--2021 Freddie sample. The 2022 vintage's lock-in state is not qualitatively different from the sampled vintages', though: 72.2\% of its face carries a pass-through coupon below 3.0\% (SOMA CUSIP joint coupon $\times$ vintage cells, run \texttt{composition\_shift}), predominantly late-2021/H1-2022 production four to five points below the window's market rates, with only 13.2\% at 4.0\% or above. What the extrapolation lacks is a different \emph{seasoning} state, and neither bound above constrains that. An observed-speed bound in the same template does: per-quarter cohort cells from the Fannie Mae replication ingest carry observed QT-window prepayment for the out-of-sample vintages. Over the window the 2022 vintage prepaid at 4.48\% against the sampled universe's 4.30\% ($+0.17$ unrounded), and the pre-2017 vintages at 5.26\% ($+0.96$; the 2015--2016 cohorts are 99.99\% of that tail's exposure and reproduce its speed unchanged). The share-weighted book-CPR error is $+0.14$ points, or \$11.7 billion at the same \$83-billion-per-CPR-point sensitivity. That is 1.5\% of benchmark, below the static Ginnie bound's own lower edge, and more than a quarter of the headline lock-in marginal (27.6\% of the off-window \$42.6 billion; a sixth, 16.7\%, of the in-sample \$70.3 billion), signed the same conservative direction. Out-of-sample vintages prepaid \emph{faster}, so trapped liquidity is if anything overstated. Applied as a marginal overlay, the template confines the identified marginal to the in-window-vintage share: $+3.7$ points. The overlay scores the 33.7\% out-of-window vintage share by observed Fannie speeds in \emph{both} legs, so that share contributes zero marginal. The result is pure share scaling ($0.66\times$ against a nominal retained share of $0.663$, placebo within $0.0002$), the vintage counterpart of the Ginnie overlay's $+4.4$. The two overlays cannot be composed without the agency-by-vintage cross-tabulation, since 2.4\% of book face is the 2022--Ginnie intersection. Three caveats travel with the bound: Fannie observed speeds proxy the book's; the pre-2017 leg is a seasoned acquisition tail, not a random draw of old originations; and the 2022 leg is a full-year speed applied to an H1-tilted face.

Burnout is deliberately cohort-level; Appendix~\ref{app:params} gives the construction and why a per-loan fractional-burnout measure is not well-defined in this design.

I extend the static Ginnie bound to a time-varying overlay built from the published speed series; Appendix~\ref{app:params} gives its construction, and its book-wide image enters the assembly of Table~\ref{tab:assembly}.

This framework's pipeline carries two Danish counterfactual mechanisms, and only one is live in production; Appendix~\ref{app:params} specifies both and the anchor that selects between them (results in Section~\ref{sec:abm-danish}).

\subsubsection{Simulated recovery and the calibration band}

Run on the 40,234 window-start survivors of the 75,000-loan stratified sample from the same Freddie Mac 2017--2021 origination universe (attrition accounting in Appendix~\ref{app:params}), this specification recovers \$818.5 billion, or 107.0\%, of the empirical benchmark. Raw CPR correlation against the empirical path is $r = +0.190$ at zero lag, and peak cross-correlation $r = +0.404$ at lag $-3$. Both are descriptors of the published series, not credentials: four clip-induced zero months dominate the monthly benchmark's differenced signal, so it cannot adjudicate timing (Appendix~\ref{sec:robustness-seasonalfloor}). (A CUSIP-level factor-paydown series would remove those artifacts at the source and is the institutionally superior path for future work; the benchmark retains current-face differencing because factor data were not staged in this pipeline.) The cross-correlation at lag $k$ pairs the simulated series at month $t$ with the empirical series at month $t+k$, both demeaned over the overlap and with boundary months dropped. A direction check on synthetic data confirms that a peak at $k=-3$ means the simulated path \emph{trails} the empirical path by three months. The definition was inverted in print, but the direction claims were always computed under the implemented convention (the replication package's ledger). The opposite reading survives in \texttt{extension\_risk\_results.json}'s \texttt{lag\_interpretation} field (a negative lag means the hazard path \emph{leads}); that string is superseded and retired, left unedited because the artifact is a sha-pinned replication object. At this peak the empirical SOMA-derived CPR path moves approximately three months ahead of the simulated path---a peak selected across the seven examined lags, so its nominal significance is overstated, and no inference rests on this single coefficient. The level and band follow the rate-gap coefficient's unit convention of Appendix~\ref{sec:robustness-pipeline}; the superseded figure is catalogued, with provenance, in the replication package's ledger.

Importing Path B's elasticity unmodified has two consequences. First, the rate-gap elasticity imported from \citet{liebersohn2024} never touches the \$764.7 billion benchmark during calibration, unlike Path A's rate-gap and burnout coefficients, which are estimated directly from the Freddie Mac sample. The two are mutually reinforcing, not redundant. Path A fits the same suppression channel on the same loans ($+0.29$ per 100 bp; Section~\ref{sec:patha}). The Fannie Mae replication reproduces its sign on a second agency's loan book at roughly $3.6\times$ the per-basis-point magnitude. That is directionally consistent with the imported band, but no more. A pre-committed test of $H_0\!: \beta_g \le 0$ respecting the diagnosed bootstrap bias does not reject at any conventional level (run \texttt{patha\_sign\_test}: BCa, 296-refit jackknife, $p = 0.093$ stratum, $0.412$ temporal; 200-draw month-permutation, $p = 0.241$). The in-sample estimates are therefore not statistically distinguishable from zero under bias-respecting constructions, and the elasticity's evidential basis is the external literature alone. \citet{graybill2026} independently corroborate the band's location: a transaction-record sale-hazard response of roughly 7--11\% per percentage point of rate gap, in a housing-market-equilibrium setting touching no credit file. That corroboration comes from at or above the adopted band. Only the 7\% low end falls inside 5.5\%--7.7\%, and the midpoint is above the high edge. The central calibration is therefore conservative against this anchor as well. A prepayment-hazard coefficient is also not unit-comparable to \citepos{liebersohn2024} quarterly-mobility decline, and Path A measures prepayment, not pure mobility. What remains is directional consistency, not triangulation of sign or direction, not magnitude. A second anchor corroborates from outside the band. \citepos{fonseca2024} roughly 9\% moving-rate reduction per 100 bp, propagated through the same transform ($\beta_1 \approx 0.096$) and band machinery, yields $+\$87.9$ billion ($+11.5$ points) at the production floor. That figure is above the Liebersohn--Rothstein high edge ($+\$79.5$ billion, $+10.4$ points) yet inside the pre-committed floor-and-band box, so the central calibration is conservative among the available literature estimates. That estimand is a moving-rate reduction, not a quarterly-mobility decline, so I read it as directional corroboration, not a competing point estimate.

Path B is accordingly a structural verification exercise, not a second fit to my benchmark. Does an elasticity estimated in a different empirical setting, dropped into a PSA-baseline hazard with competing risks and no tuning against my target, reproduce the Federal Reserve's actual QT runoff shortfall? The aggregate level is floor-dominated, so landing within 8 percentage points of the benchmark is not by itself evidence for the imported elasticity: the residual the elasticity must explain over the $\beta_1 = 0$ null is itself a function of the floor calibration, and the marginal spans $+2.1$ to $+13.2$ points across the calibration box of Section~\ref{sec:robustness-floor}. The defensible verification claim is the bounded interval that box implies. The marginal is positive at every floor, band, and floor-form point tested, but under this window's rate configuration that positivity verifies the specification. It does not identify the mechanism (Section~\ref{sec:identification}); I do not claim the level validates the elasticity's magnitude. Nor is timing part of that evidence: a zero-elasticity rerun reproduces the same three-month lead (Section~\ref{sec:hazard-interp}), so the lead cannot discriminate the elasticity from the rate-inelastic model.

Second, propagating the full Liebersohn--Rothstein 5.5\%--7.7\% sensitivity band through the simulation yields a recovery range of \$809.9 billion to \$827.7 billion (105.9\%--108.2\% of benchmark under the standalone scorer; 96.8\%--99.1\% on the shared accounting basis of Section~\ref{sec:robustness-hybrid}). The peak cross-correlation lag is stable at $-3$ across the entire band, with peak $r$ between $+0.371$ and $+0.423$. Given the moving-block intervals above and the seven-lag search, these differences carry no evidential weight. Stronger assumed mobility suppression therefore maps monotonically to more trapped liquidity, as expected, and the central 107.0\% estimate should be read as a point within this interval; Figure~\ref{fig:band} plots the band. That band is a calibration interval, not sampling error. The sampling dimension is quantified separately in Appendix~\ref{app:params}, where the committed within-stratum loan bootstrap and the stratum-cluster scheme that widens it are reported with their caveats. Even at the cluster scheme's wider interval the sampling layer is less than half the width of the floor-read interval of Section~\ref{sec:robustness-floor}, so the operative uncertainty statements for Path B's level and marginal remain the calibration box and floor-form results there. The shared-basis equivalents follow by the common 9.1-point netting, which cancels from the marginal entirely.

\begin{figure}[H]
\centering
\includegraphics[width=0.82\textwidth]{fig6_elasticity_band}
\caption{Path B trapped liquidity across the Liebersohn--Rothstein elasticity band. The quarterly mobility-decline parameter $\delta$ is swept across its published 5.5\%--7.7\% range (central 6.5\% marked), mapping monotonically to recovered trapped liquidity of \$809.9--\$827.7 billion on the standalone-scorer basis (shared-basis equivalents in Table~\ref{tab:bases}) against the \$764.7 billion benchmark line: stronger assumed mobility suppression yields more trapped liquidity, with the peak cross-correlation lag stable at $-3$ across the entire band. Table~\ref{tab:floorband} extends this band across the baseline-turnover-floor grid.}
\label{fig:band}
\end{figure}

\begin{table}[!t]
\centering
\caption{The calibration box: lock-in marginal over the $\beta_1 = 0$ null, in percentage points of the \$764.7 billion benchmark, across the baseline turnover floor and the Liebersohn--Rothstein band $\delta$. The 6.0\% row's level is reached only by the refi-contaminated pooled and 2019 anchors that Section~\ref{sec:robustness-floor} disqualifies; it is retained as a grid point, not a defensible anchor, and the box's lower edge should be read from the defensible off-window rows of Table~\ref{tab:oosfloor}. The bind column is the share of evaluated U.S. loan-months at which the floor lifts the hazard (central elasticity). Every marginal is basis-invariant (the underlying levels are standalone-scorer figures; both bases in Table~\ref{tab:bases}). Pre-committed runs; artifacts catalogued in the replication package's ledger.}
\label{tab:floorband}
\begin{tabular}{lcccc}
\toprule
Floor (annual CPR) & $\delta = 5.5\%$ & $\delta = 6.5\%$ & $\delta = 7.7\%$ & Floor binds \\
\midrule
2.0\% & $+9.7$ & $+11.3$ & $+13.2$ & 3.9\% \\
3.0\% & $+9.3$ & $+10.8$ & $+12.6$ & 12.1\% \\
3.5\% & $+8.9$ & $+10.3$ & $+11.8$ & 21.6\% \\
4.0\% (production) & $+8.1$ & $+9.2$ & $+10.4$ & 36.3\% \\
4.5\% & $+6.7$ & $+7.5$ & $+8.4$ & 53.5\% \\
5.0\% & $+5.0$ & $+5.5$ & $+6.0$ & 69.1\% \\
6.0\% & $+2.1$ & $+2.3$ & $+2.4$ & 88.7\% \\
\bottomrule
\end{tabular}
\end{table}

\subsubsection{Path timing}

The peak alignment does \emph{not} show the simulated path leading SOMA cash flow. The empirical path moves first, and a settlement-delay mechanism---in which loan-level prepayment decisions show up in SOMA cash two to three months later---predicts the \emph{opposite} direction of offset. I therefore attach no settlement interpretation to the timing alignment; Appendix~\ref{app:patha-detail} carries the direct input-timing diagnosis behind that reading.

\subsection{External Replication on Fannie Mae Loan-Level Data}\label{sec:pathb-fannie}

Path B's construction admits a direct external replication on Fannie Mae's single-family loan performance data (24 acquisition quarters, 2017Q1--2022Q4; 17.6 million loans, 826 million loan-months). I stage it into the identical origination/performance layout, so the panel builder, estimator, and stratified 75{,}000-loan design run unchanged on the same 2017--2021 origination universe. The spec was fixed and committed before the run. Zero-balance code 01 maps to voluntary prepayment as for Freddie. Codes 06 and 16---non-credit removals with no Freddie counterpart---are censored; were they all prepayments, the worst-case bracket is 0.38\% of removals. The pre-committed gate required a positive Fannie marginal inside Section~\ref{sec:robustness-floor}'s Freddie floor/band envelope. It replicates: central \$821.0 billion (107.4\% of benchmark), $\beta_1 = 0$ null \$754.6 billion (98.7\%), and a lock-in marginal of $+\$66.4$ billion ($+8.68$ points of benchmark) against Freddie's $+\$70.3$ billion ($+9.20$ points)---inside the envelope. But that envelope, Section~\ref{sec:robustness-floor}'s calibration box, spans 2.11 to 13.17 points: an 11.06-point window around a 9.20-point Freddie estimate, admitting a replication 77\% below it or 43\% above. A gate that wide fails only on a sign error, and I do not treat its passing as evidence. The agreement is a different matter: 17.6 million staged Fannie loans run through the same pipeline land 0.52 points from the Freddie point, far more than the gate demanded, and that closeness is the part of this exercise that is reassuring. What it certifies needs care. Every behavioral parameter is fixed a priori and identical across runs, and the two samples are compositionally indistinguishable (population-stability indices at or below 0.018, Appendix~\ref{app:composition}), so the 0.52-point agreement certifies pipeline determinism and cross-agency compositional stability, not external validity---which lives in the sign reproduction and staging inputs, not in the closeness. Path A's fixed-effects Fannie-cohort refit reproduces the rate-gap coefficient's sign, but at about $3.6\times$ the Freddie fit's per-basis-point magnitude and with Section~\ref{sec:patha}'s no-forward-precision holdout character, so it retains aggregate-corroboration status only. I record two ex ante scope notes. Whether Freddie's code 01 absorbs exits Fannie books under 06/16 is not settled by public documentation, so the claim is sign and envelope-conditional marginal consistency, not a hazard-level match. And Fannie's terms prohibit redistribution, so the package ships ingestion and replication code only---no Fannie data and, conservatively, no derived cells. The committed artifact carries headline statistics alone.

\subsection{Path A: Stratum-Month Discrete-Time Hazard Model}\label{sec:patha}

Path A is the aggregate corroboration leg: a stratum-month discrete-time hazard fit to Freddie Mac loan-level performance data. It recovers \$928.9 billion in aggregate, 121.5\% standalone---reported as corroboration only and excluded from every headline figure. Appendix~\ref{app:patha-detail} collects its estimation detail, diagnostics, and cross-correlation exhibits.

The same caution applies to the paths themselves: Path B's $+0.190$ lag-0 levels correlation is carried by the window's shared trend ($-0.23$ differenced, $-0.20$ linearly detrended). The ABM gives $-0.20$ and $-0.30$ (from $-0.32$ in levels), Path A $-0.21$ and $-0.29$ (from $-0.44$): detrending collapses them into a cluster between $-0.20$ and $-0.30$, none positive, so Path B's distinction is its level accuracy, not its timing. Against the i.i.d.\ $1.96/\sqrt{n}$ reference ($\pm 0.30$ at $n = 42$ detrended, $\pm 0.31$ at $n = 41$ differenced), Path A ($-0.29$) and Path B ($-0.20$) sit inside the band under both empirical-series conventions. The ABM's detrended value straddles its edge ($-0.301$ common-series, $-0.303$ own back-out). Indistinguishability verdicts therefore rest on Table~\ref{tab:estimators}'s wider moving-block bands, under which none of the three is distinguishable from zero. (In levels, trend-contaminated and reported only for completeness, the simulated path co-moves most strongly with the rate from four months earlier, the empirical path contemporaneously.) No single channel carries the offset: the $\beta_B = 0$ and $\beta_1 = 0$ ablations both leave the peak at $-3$; shifting the rate input one to three months either way never reaches lag 0 ($-1$ at best, under a two-to-three-month lead), moving recovered trapped liquidity under \$1.5 billion across the entire scan. The offset is therefore smooth fitted dynamics against a seasonal empirical series (mechanism not isolated); Path A's seasonal re-estimation moving its peak from $-3$ to $-2$ (Section~\ref{sec:patha}) is consistent with seasonality carrying part of it. Figure~\ref{fig:ccf} plots the full cross-correlations. At the cohort level the residual is structural, not compositional. Across the six 30-year cohorts spanning the SOMA low-coupon mass (2.0--4.5\%), every simulated CPR tracks the contemporaneous rate change near-mechanically: first-difference correlation between $-0.87$ and $-0.94$, five of six phase-locked at lag~0, the 4.5\% coupon peaking at lag $-6$ with a weaker $r=-0.29$, reproducing the aggregate $-0.318$. By contrast, no realized cohort speed responds distinguishably from zero at any lag within $\pm 6$ months. Those realized speeds are measured independently in the Freddie 2017--2021 performance panel and the Fannie replication cohorts, with Freddie lag-0 first-difference correlations $-0.05$ to $-0.16$, Fannie $-0.03$ to $-0.14$, all inside the $\pm 2/\sqrt{n}$ band, maximum magnitude $0.34$ over all cohorts and lags, and no coherent sign or lag. The aggregate empirical non-response is therefore not a composition artifact, nor the simulated lead a datable cohort-level phase shift. One honest scope limit travels with this exhibit: realized cohort-level monthly speeds exist only in the loan-level agency panels anchoring Path B, not in the SOMA holdings the empirical CPR is backed out of, so the cohort split runs in the Path B universe, and the estimator-level timing falsification above stays at the SOMA aggregate. A cumulative counterpart completes it. Monthly differences being seasonality-dominated, the cross-section is read in cumulated levels. Age-matched Freddie realized prepayment runs from 9.43\% annualized CPR (28.1\% of balance) in the shallowest surviving gap bucket ($[-1,0)$ points, mean coupon 6.02\%) to 5.24\% (16.4\%) at $-3$ points and deeper (mean coupon 3.37\%). That is a gradient of $+4.20$ CPR points, four-way stratum-cluster 95\% interval $[+3.59, +4.66]$, one-sided permutation $p$ 0.005. This is the one realized-data exhibit in this paper that moves in the lock-in direction. It is not corroboration of the imported elasticity's magnitude: the production hazard implies $+0.94$ points at the central $\beta_1$ across the same buckets, so the realized gradient is roughly $4.5\times$ the model's. The excess carries everything travelling with coupon: credit composition, within-window refinancing on high-coupon shallow buckets, seasoning past the age cut. Detection power is 0.68: the design could have resolved $\beta_1$'s implication, and what it found was larger and confounded. The composition limb is tested: standardizing the deep bucket to the shallow's vintage $\times$ FICO $\times$ LTV composition (22 cells, 18 of them carrying both endpoint buckets) leaves the gradient at $+3.44$ points (95\% CI $[+1.17, +5.14]$, imputed weight share $0.3\%$), an exact decomposition putting $+0.77$ of the raw $+4.20$ on composition and $+3.43$ within common cells (run \texttt{episode\_confrontation\_within}). Against the implied $+0.94$, the excess is not composition in dimensions the panel can hold fixed, loan age now among them: adding 12-month age bands to the standardization cell (29 cells, the same 18 in common; run \texttt{episode\_gradient\_age\_bands}) leaves the gradient at $+3.37$ points, and the same cut on age alone puts $+0.38$ of the raw $+4.20$ --- $9.1\%$ --- on seasoning composition, age in the cell moving total composition from $+0.77$ to $+0.85$. The age cut itself is held at production, so this stratifies above the cut rather than testing where it sits. The interval moves against the exhibit even as the point survives: on the age-augmented cell it runs $[+0.60, +4.35]$ (imputed weight share $0.5\%$), covering the implied $+0.94$ where the committed $[+1.17, +5.14]$ excluded it. On run \texttt{episode\_confrontation\_within}'s own pre-committed branch map, that moves this configuration from a gap that persists to one that closes. So the exhibit weakens without collapsing: seasoning composition is measured and small, and the gradient is no longer separated from the model's implication by its own sampling interval. I had pre-committed that at least $80\%$ of the primary selection's shallow exposure would sit in single-band vintages; the realized share is $60.6\%$: age is not merely a relabelling of vintage in this book, so the control is stronger than I expected. That $+0.94$ sits at the in-sample floor and production ramp, so I recomputed it across the headline calibration (run \texttt{episode\_gradient\_recompute}): the shortfall is not an artifact of the demoted anchor. At the $4.991\%$ floor the implied gradient is $+0.76$ under the production form and $+0.89$ under the additive one, so the realized-to-implied ratio is $5.5$ and $4.7$ against the $4.5$ at the in-sample anchor, and across the floor-by-form-by-ramp grid every finite ratio lies between $3.6$ and $12.7$. One cell is not finite: at the headline floor on a $75$~PSA ramp both endpoint buckets are fully censored, so the model implies \emph{no} cross-sectional gradient and the comparison degenerates. I had pre-committed two signs and one was wrong: the max-form gradient fell as predicted, but the additive form lowered it too rather than raising it. A floor in survival space compresses proportional differences even where it never censors. I therefore read the cross-section as a signed, significant level gradient whose composition component is measured and small---consistent with lock-in operating, still unusable as a measurement of its size: within-window refinancing on high-coupon shallow buckets stays unresolved, no field here separating it from moving, and a level gradient is not a dollar marginal. The monthly-timing nulls above stand unchanged.
\subsection{Identification: the Marginal over the Null}\label{sec:identification}

\subsubsection{The estimand}

The design does not identify the aggregate recovery level, which is floor-dominated. The $\beta_1 = 0$ null---elasticity switched off, baseline turnover held at its observed level---already recovers 85.7\% of the benchmark on the shared accounting basis at the headline off-window floor, and 88.7\% at the in-sample calibration point (Section~\ref{sec:hazard-interp}; standalone figures in Table~\ref{tab:bases}). What it does identify is the \emph{marginal} of the central Path B run over that null, writing $R$ for recovered trapped liquidity,
\begin{equation}\label{eq:marginal}
\text{lock-in marginal} \;=\; R_{\text{central}} - R_{\beta_1 = 0},
\end{equation}
which is basis-invariant: the shared-accounting netting is common to both legs (Section~\ref{sec:robustness-hybrid}). Its value moves with the baseline turnover floor, the design's dominant level-setter, and that floor is calibrated in-window on the very cohorts the marginal decomposes. The headline therefore measures that floor off that population: $+\$42.6$ billion, $+5.6$ points of the benchmark, at the defensible off-window floor of Section~\ref{sec:robustness-floor}, and $+4.3$ to $+6.8$ points across its range. In window the same difference is $+\$70.3$ billion, or $+9.2$ points. The recovery \emph{level} moves with the floor too, so I quote it at both calibrations, never a level from one floor beside a marginal from another. The central leg recovers 91.3\% of the benchmark on the shared basis at the off-window floor (88.2--93.7\% across its range) against 97.9\% at the in-window calibration; the null recovers 85.7\% against 88.7\%. The concave-gap variant, the only other leg re-run off window, recovers 99.8\% central against a 94.8\% null on the standalone scorer. Every other estimator level here---Path~A, the full-book and composed estimators, the covariate-prior and elasticity bands, the Ginnie overlay, the Danish legs---was run in window only, so any level without a floor label should be read at the in-window 4.0\% calibration, the configuration I call the \emph{in-sample calibration point} throughout. Differencing does not immunize the marginal against non-stationary baselines: the baseline multiplies the very rate-gap channel the null switches off, so a break in $h_0$ scales the marginal rather than cancelling from it. Under a $\pm 20\%$ multiplicative break from January 2023 onward, imposed identically on both legs and fitting nothing to the benchmark, it runs from $+5.0$ to $+12.1$ points around the committed $+9.2$ (bit-exact at the unit factor); a $\pm 20\%$ break in the turnover floor alone moves it the other way, to $+6.6$ through $+10.4$ points. Every perturbation leaves it inside the calibration box and positive: what survives a baseline break is box membership, not the point value.

Seven qualifications fix its evidential content. First, the identified content is the marginal's \emph{bounded range}, neither its point magnitude nor its sign: the box spans $+2.1$ to $+13.2$ points, and across the empirically plausible floor range the marginal spans $+5.0$ to $+12.6$ points over the full band, $+10.8$ to $+5.5$ at the central elasticity, moving to $+11.25$ under the competing-risks floor form. Censoring moves with the floor's level. At the headline 4.991\% off-window floor the central leg's hazard is pinned to the turnover floor in 68.8\% of the 1,683,124 evaluated loan-months and the $\beta_1 = 0$ leg in 35.8\%, against 36.3\% and 14.3\% at the in-sample calibration point; the central share runs 60.0\% to 77.3\% across that floor's defensible range. In most loan-months at the headline calibration, then, the elasticity no longer sets the hazard and its contribution to the marginal is truncated at the floor---though these are unweighted loan-month counts, not balance-weighted shares, and do not say what fraction of the dollar marginal the uncensored months carry. Wherever the null leg is pinned, switching the elasticity off cannot raise the hazard at all. The marginal is thus floor-level sensitive under the production max form: form-robust at roughly $+9$ to $+11$ points at the production floor (Section~\ref{sec:robustness-floor}), but with no form-robust point at the off-window anchors, where the two forms return $+5.6$ and $+11.2$. That leaves the form-conditional range $+3.5$ to $+13.1$ points as the defensible statement. Its endpoints are those of a measured mixture curve in the floor's strictly-involuntary share, rising to $+9.7$ points by $\omega = 0.25$ and plateauing near $+11.2$: the paper's own floor semantics point at the hull's upper region, and the $\omega = 0$ headline is the curve's conservative endpoint. Its upper end comes from the additive form at the 4.695\% floor and the band's 7.7\% high edge, whose central leg recovers 61.2\% of the benchmark; Section~\ref{sec:robustness-floor} states that it should not be read as an equally credentialed member of the range. The hull holds each floor read at its point estimate, so it is not the paper's widest disclosed interval: propagating the floor reads' own cluster-bootstrap sampling error puts the max-form member at $+2.3$ to $+9.1$ points after wild-cluster correction (percentile read $+3.0$ to $+8.0$; Section~\ref{sec:robustness-floor}), and leaves it open below $+4.3$ on an age-standardized read (Section~\ref{sec:robustness-floor}). The floor read and the loan/stratum resample are reported side by side, not maximized over. Convolved as independent draws the pair spans $+2.8$ to $+9.0$ points (computed on the Webb rung, predating the binding-layer adjudication), and that convolution is the interval a reader who wants one number for sampling error should use. I do not substitute it for the floor read, because independence between a 2018 cohort-month ratio and a window-period composition resample is an assumption rather than a measurement. The hull's lower edge is thus narrower than the layer the paper calls binding; the two must be quoted together, not the hull alone.

\subsubsection{The floor's provenance and the off-window re-anchor}

The floor's provenance now has an out-of-agency read as well as an out-of-window one, and it does not narrow the range. The committed off-window rule, applied unchanged to the Fannie Mae cohort-month panel, returns 5.52\% on the headline cell against Freddie's 5.19\%, on 421 cohort-months against 438 (run \texttt{fannie\_floor\_read}). That panel carries the same fourteen columns and the same 2017 start, and it was staged for the external replication of Section~\ref{sec:pathb-fannie}; the Freddie panel replays every committed cell of that grid through the identical code path first. That read clears the clean band's 5.334\% top, but it is no second read of the clean leg: off window the rule covers pooled 2017--2019, whose Freddie counterpart is that 5.19\%, not the 2018 leg's 4.991\%, so the agency difference is 0.337 points of CPR (5.522\% against 5.185\%) and the rest is period. The age-standardized 5.51\% lands 0.01 points away, but it perturbs something else: an age correction to the 2018 leg itself, on 84.0\% imputed weight. The near-coincidence of the two is therefore not a second measurement of either. What they share is direction---both above the clean top, both saying the clean read understates the mature book's turnover---and that is what the clean leg cannot test on its own. The 2018 leg holds no out-of-the-money cohort-month past the twenty-four-month age cut, so its mature-seasoning test is not computable (Appendix~\ref{app:floormech}), and a test that cannot be run has not been passed. The one age contrast the leg does support is confounded by the PSA ramp. Where a mature contrast is measurable at all (two agencies, two rate regimes, three depth cuts), it runs one way in twelve cells of twelve, the age-$\geq$-24 read the higher one in each. In window that read is higher by 0.08 and 0.19 points of CPR on Freddie and Fannie, the rate configuration suppressing refinancing as on the 2018 leg; on the refi-contaminated pooled 2017--2019 period it is higher by 3.74 and 4.61 points, which bounds the correction from above without measuring it. A floor read confined to cohort-months aged twelve to twenty-four therefore understates mature baseline turnover and overstates the marginal: the error is signed against the headline, its size not priced anywhere in this paper. On the committed floor-to-marginal grid a 5.52\% floor brackets the marginal between the 5.334\% row's $+4.27$ points and the 6.0\% row's $+2.28$, below the $+4.3$ lower edge this paper quotes. I do not restate the headline on the Fannie floor: the production calibration is the Freddie one, and the two books differ in composition as well as in this read. So the out-of-agency read establishes only that the band's \emph{lower} edge is the soft one, the direction the age-standardized caveat already pointed.

That sensitivity makes the floor's provenance decisive, so Section~\ref{sec:robustness-floor} re-measures the floor outside the evaluation window and re-runs both legs there. The marginal lands at $+4.3$ to $+6.8$ points across the defensible off-window range (inside the box, well clear of either edge) and stays positive at all twenty-seven floor-by-band cells run. The covariate-prior swing of Section~\ref{sec:pathb} is the design's largest single level sensitivity at both calibrations, and at the headline floor it exceeds the marginal: re-run at the 4.991\% floor, the two covariate endpoints span 8.2 points of recovery against the $+5.6$-point marginal. The marginal is the counter-fact. Switching the covariate block between its endpoints moves the off-window marginal by only $+0.1$ points, against $+1.9$ in-window, because both legs shift together where the floor censors more.

\subsubsection{The sign check}

The marginal's positivity in all twenty-seven cells is not the robust identified content. Over the forty window months the cohort panel covers (June 2022 through September 2025; the panel ends two months before the window does), the exposure-weighted mean coupon on the surviving book is 3.134\%, with 97.40\% of exposure at or below 4.79\% and 99.53\% at or below 5.09\%. Set those figures against a window-minimum 30-year rate of 5.2311\%. That minimum is the August 2022 monthly mean --- the window's rate dip --- not a June 2022 figure (\texttt{sign\_forcing\_stats}). The two uncovered months cannot overturn it: the book is closed at 2017--2021 originations, so no higher-coupon loan can enter, and the out-of-the-money share can only hold or rise as low-coupon cohorts season. Essentially the entire book is out of the money for essentially the entire window, so the rate gap is negative in every cell of every sweep reported here. A specification in which a more negative gap suppresses prepayment cannot then return a negative marginal at any floor or band in the box. The sign is fixed by the rate configuration, not by the estimated elasticity. What the twenty-seven-cell check verifies is accordingly the specification: it would catch a wiring error that flipped a coefficient or crossed the central and null legs, and I retain it on that footing. This is why the reported range, and not the sign, carries the paper's identified content. The band's monotone response is forced the same way: \eqref{eq:beta1} is monotone in $\delta$ and the hazard monotone in $\beta_1$ at that gap, so the ordering composes monotone maps. That ordering I retain as a wiring check on the same footing, not as identified content; the sweep's endpoints stay the convention range they are.

\subsubsection{What the design does not claim}

A second qualification: timing carries no credential. The $\beta_1 = 0$ null also peaks at lag $-3$ (Section~\ref{sec:hazard-interp}), so the lead cannot discriminate the elasticity from the rate-inelastic model. A third qualification: the marginal's sampling precision is not the operative uncertainty. Its premise is now the cluster scheme's, not the retracted one's. At this calibration the stratum-cluster bootstrap puts the paired marginal at $[+8.27, +10.19]$ points, while the within-stratum $[+9.17, +9.23]$ understates by 30.8$\times$ (Section~\ref{sec:pathb}). The calibration box, 11.1 points wide against the cluster read's 1.9, still governs the interval.

A fourth qualification concerns composition. The decomposition exhibits in this subsection, computed at the in-sample calibration point, partition the $+\$70.3$ billion figure and report shares and intensities --- quantities the floor's level shifts far less than it shifts the total. Group ablation (zeroing $\beta_1$ cell by cell, everything else at production) shows the $+\$70.3$ billion broad-based, not cohort-specific. All fifteen vintage $\times$ coupon cells are positive, they sum to within 1.03\% of the paired-run total, and no cell exceeds 21\%. The 2020--2021 vintages carry 72\% of the marginal on 73\% of balance, and per-balance intensity runs from $0.82\times$ to $1.75\times$. On the dimensions that grid misses --- FICO ($<$680 / 680--740 / 740$+$), LTV ($\leq$80 / $>$80), Census region --- the same engine returns broad-based on all three. Every cell in those partitions is positive as well, and the sums reproduce the $+\$70.3$ billion total within interaction residuals of 0.8, 0.4, and 0.5 percent. Per-balance intensity across the three lies in 0.88--1.26, tighter than the vintage $\times$ coupon span. That verdict does not, on its own, discriminate, and I say so: twelve random three-way partitions of the same 75,000 loans, run through the identical engine and pre-committed rule, return \emph{broad-based} in all twelve, on intensities of 0.995--1.007 and additivity residuals of 0.49 percent (\texttt{subgroup\_marginals\_placebo}). A partition with no economic content earns the label, so the label alone carries no evidential weight. On the benchmark basis, the 0.2-point proximity of the projection's uniform-spread 88.5\% upper end to the null's in-sample 88.7\% recovery is a doubly conditional coincidence, not corroboration. Section~\ref{sec:method-benchmark} withdraws that presentation, states both conditions with their magnitudes, and reports projection range and null separately; I claim no agreement from their proximity. Every percentage-point magnitude quoted for this decomposition is cap-basis, while dollars and rankings are basis-invariant.

\paragraph{Distributional incidence.} The placebo strips the homogeneity verdict of its evidential content, but a distributional finding survives it: who bears the mobility cost, not whether a label discriminates. The comparison establishes the reverse: the random partitions' intensities concentrate on unity, the widest of the twelve spans 0.011, while the economic groupings span 0.381, more than thirty times wider. The informative statement is therefore not homogeneity, which chance alone delivers, but dispersion well beyond chance: the tilts are structure, not sampling noise, and the generality claim rests on sign homogeneity and the bounded magnitude of that dispersion, not on the verdict label. Those tilts run mildly toward below-740-FICO ($1.23\times$--$1.26\times$) and above-80-LTV ($1.23\times$) borrowers, regional intensities spanning only 0.93--1.06 (Figure~\ref{fig:marginalcells}). The marginal is a book-wide property of the rate-gap distribution interacting with seasoning and the floor, so the lock-in cost falls mildly more than proportionally on weaker-credit and higher-leverage borrowers, close to evenly across regions --- a statement about where a book-wide effect concentrates, measured against a placebo that fixes what chance alone would deliver; it is not evidence of a subgroup-specific mechanism, and this design does not carry it to household welfare.

The monthly-path and coupon-bucket timing exhibits are collected in Appendix~\ref{app:patha-detail}. Neither simulated leg carries calendar-month seasonality, so seasonal structure enters these exhibits only through the empirical comparison. Per the pre-committed posture of Section~\ref{sec:hazard-interp}, they are fit diagnostics, not timing credentials.

A fifth qualification states what kind of claim ``identifies'' makes. The marginal is a within-model counterfactual decomposition: the difference between two runs of one fully disclosed hazard structure, differing only in the elasticity term. It is not causal identification in the econometric sense. This paper exploits no quasi-experimental variation in rate gaps across otherwise-comparable loans, and it claims none. The causal content of $\beta_1$ is inherited, not asserted: I import the elasticity unmodified from \citepos{liebersohn2024} estimated band, with \citet{fonseca2024} and the transaction-record band corroboration of \citet{graybill2026} as independent support. No Path B parameter is estimated on the benchmark (Section~\ref{sec:pathb}). The switch-off adds accounting at a stated price: the marginal inherits any misspecification of the held-fixed structure. So I sweep the channels I can name --- floor level and functional form, the rate-gap transform, the covariate priors --- each directly (Sections~\ref{sec:pathb} and \ref{sec:robustness-floor}), and Table~\ref{tab:uncertainty} reports calibration, not sampling, as the operative uncertainty.

A sixth qualification concerns three readings I decline in three separate places. Path A's fitted rate-gap coefficient is $+0.29$ per 100 bp against the imported $\beta_1$'s $0.069$, on different scales. I therefore compare the two on sign not magnitude (Section~\ref{sec:patha}). The realized cross-cohort gap gradient is roughly $4.5\times$ what the central $\beta_1$ implies. The real-covariate cross-design variant recovers 59.3\% of the benchmark, 76.3\% reweighted to the book's composition, against the production ABM's 13.6\% seed mean (Section~\ref{sec:robustness-crossdesign}). I stand by each reason. The first is not unit-comparable and does not survive a bias-respecting sign test. By contrast, the second survives composition standardization ($+3.44$ against the implied $+0.94$, run \texttt{episode\_confrontation\_within}). I decline it on the narrower ground that remains. Within-window refinancing on shallow buckets is inseparable from moving, and a level gradient is not a dollar marginal. Finally, the third is a recovery level, which this design does not identify, and the same reweight moves the frozen cross-design variant down to 12.6\%. Stated in one place: none of the three points toward a smaller lock-in channel. I keep the imported value for Section~\ref{sec:pathb}'s reason. Path B is a structural verification exercise whose elasticity is estimated in a different setting and never touches this benchmark, so re-fitting it in-sample would end the exercise. Among available literature estimates the adopted band is already the conservative one. \citepos{fonseca2024} anchor puts the marginal at $+11.5$ points against the central calibration's $+9.2$ at the production floor. Other disclosed variants move it the other way, the Ginnie-speed overlay to $+4.4$ points and the age-standardized floor read to a measured $+3.7$ (Table~\ref{tab:uncertainty}). So I do not read the three as a case for a larger headline.

\subsubsection{The assembled corrections}

A seventh qualification assembles those readings against the ones that run the other way, because this paper states each correction where it arises and never lines them up. Lined up they show what the sections do not. Holding the production hard-maximum form and the imported elasticity fixed, every correction to the floor or to the accounting basis moves the headline down; re-solved under the additive form, the calibrated companion below moves it up instead. The in-sample calibration returns $+9.2$ points. Re-anchoring the floor outside the window returns the headline $+5.6$. The Ginnie-speed overlay at that same floor returns $+4.4$. The age-standardized floor read of 5.51\% returns a measured marginal of $+3.7$, against the $+3.8$ its position on the frozen grid implies. The Fannie Mae read of the same rule on the pooled off-window period, 5.52\%, brackets the marginal below the $+4.3$ edge of the clean band. Its like-for-like Freddie counterpart is 5.19\%, not the 2018 leg's 4.991\%. The anchor convention does not flatter the headline. The middle of the five reads the 10\% support rule admits is 4.869\%, below the adopted 4.991\%, and a lower floor implies a marginal above the headline. Carrying the overlay's $0.797$ scaling up to the age-standardized floor projected a marginal near $+3.0$. The joint cell has now been run under a pre-committed landing rule (run \texttt{b5\_joint\_cell}) and measures $+2.9$ points, an interaction within $0.1$ of proportional. By contrast, this paper's first joint cell returned an interaction of $-1.47$ points (Section~\ref{sec:robustness-floor}). Corrections do not compose additively, and a composed figure would overstate precision. Nor are they independent: floor corrections and the censoring share are one mechanism read twice, since raising the floor crowds the elasticity out. At the 4.991\% headline floor the central leg is already pinned to the floor in 68.8\% of evaluated loan-months, against 36.3\% at the in-sample calibration. Further increases have progressively less hazard left to act on. That share is no claim that the estimate rests on a third of the data. A pinned month's contribution is truncated at the floor rather than erased. The $\beta_1 = 0$ leg is pinned in only 35.8\% of loan-months there, so the legs separate wherever the null clears it. These are unweighted loan-month counts rather than balance-weighted shares, so nothing follows about the uncensored months' share of the dollar marginal. What the assembly settles is where in the interval the mass sits, not whether the finding holds. Within the production form the widest disclosed layer is the baseline level. Sweeping the seasoning ramp over 75 to 150 PSA on both legs carries the marginal from $+0.9$ to $+15.6$ points at the off-window floor and $+3.3$ to $+16.5$ at the in-sample calibration. That span has no coverage property: it is a range over a convention I did not estimate, its endpoints not draws from anything. Aggregate fit is no defense of the 100~PSA convention. At the off-window floor the $\beta_1 = 0$ null recovers 101.9\% of the benchmark under the standalone scorer at 75~PSA against 94.8\% at the production 100, while the central leg is closer at 100 (100.4\%) than at 75 --- the legs rank the ramps opposite ways. A conditional estimate of the ramp's \emph{shape} now exists and does not favour the headline. Transporting the age profile of Path~A's estimated seasoning spline onto $h_0$ returns a median marginal of $+9.44$ points, 5th-to-95th-percentile spread $+8.99$ to $+9.81$, across the 107 of 200 bootstrap replicates whose normalised profile stays inside the usable range (run \texttt{h0\_reanchor}). The level is held fixed by construction, absorbed by the fitted constant and stratum fixed effects and not separately identified. Every replicate lies above the $+5.6$ this paper headlines and all but one above the $+8.7$ upper edge of the binding interval quoted below, though all lie inside the $75$--$150$~PSA span. With 46.5\% discarded it is a conditional estimate, not the ramp's sampling interval, so the convention span stays labelled unestimated. I had pre-committed to the opposite sign: that re-anchoring would push more loan-months onto the floor and so \emph{lower} the marginal. The bind share did rise as predicted, from 35.8\% to between 50.4\% and 58.2\% on the null leg, and the marginal rose anyway: a failed prediction, reported as one. The floor read itself transports in two directions, both fixed before either ran (run \texttt{floor\_transport\_two\_sided}). Standardizing its calendar raises it: the 2018 leg is read in August--December---two July cohort-months aside---on the declining side of the committed seasonal profile. Re-weighting to the window's calendar mix carries 4.991\% to 5.2156\% and the marginal to $+4.70$ points. Matching its housing-activity level lowers it: active listings averaged 1.27 million over those 2018 months against 0.80 million over the window. The same multiplicative transport carries the floor to 3.1533\% and the marginal to $+10.69$. I report the two at their own floors and do not net them: two floor-side adjustments composed with an interaction of $-1.47$ points the one time this paper ran them jointly. A cancellation would be a finding requiring its own cell, not arithmetic. The activity leg is a listings-based proxy for an argument about existing-home \emph{sales}, not constructible here (the sales series carries no observations over the 2018 leg months), and its magnitude inherits a one-for-one activity-to-turnover mapping nothing here estimates. The sign is what it establishes. Post-stratification onto the book's composition moves the marginal less than the paper's own share arithmetic implies. Raking the pool's balance weights to the SOMA coupon marginal and the book's vintage mix within the draw's 2017--2021 support (both margins hit to better than $10^{-10}$) returns $+5.11$ points against the unweighted $+5.57$, a move of $-0.47$. The vintage margin alone returns $+4.67$, so the coupon margin offsets about half the vintage correction. The vintage overlay's $+3.7$, pure share scaling, overstates the composition correction roughly fourfold. Two limits travel with it. The book's coupon-by-vintage \emph{joint} distribution is not observable (the SOMA cohort tabulation collapses vintage to one value-weighted origination date per coupon bucket), so this matches two margins, not a cell grid. The pre-2017 and 2022 vintages, 33.7\% of book face, have no draw support, are not re-weighted here, and stay with the vintage overlay. I pre-committed both a direction and a band: the direction held and the marginal fell, while $+5.11$ lies just outside the $[+2.5, +5.0]$ band. I report that miss. I do not re-draw the band. The widest layer that does have a coverage property is the floor reads' own sampling error, $+2.3$ to $+9.1$ points after wild-cluster correction of the percentile read's under-coverage (run \texttt{floor\_inference\_correction}; the percentile read is $+3.0$ to $+8.0$). That is the interval this paper quotes as binding. The headline sits just below its midpoint, every correction listed above falls in its lower half, and the uncomposed pair near its bottom. Resolution should be read off those magnitudes, not the printed decimals. Against a binding interval roughly seven points wide the weighting scheme moves an endpoint about $0.06$ of a point, so the interval is credible to something like half a point and one-decimal endpoint comparisons fall below it. The two are nested: the ramp sweep is the outer statement, $+2.3$ to $+9.1$ a sampling interval on the floor read alone, nested inside it at the production 100~PSA ramp. Which rung of that inner layer to quote is settled by the pre-committed coverage rule of run \texttt{fewcluster\_coverage}, not by width. The restricted inversion is the narrowest construction clearing the $93.0\%/91.0\%$ bar, and CR2 at Bell--McCaffrey degrees of freedom ($+2.4$ to $+9.1$ points), the other qualifying rung, stands beside it, not behind. A designed coverage simulation at this read's committed leverage profile (run \texttt{fewcluster\_coverage}: 31 clusters, 5.9 effective, maximum leverage 0.33; 5{,}000 draws $\times$ 999 replications) covers 94.3\% under Gaussian and 95.1\% under $t_5$ cluster disturbances for the restricted wild-cluster inversion. The comparators are 90.8\%/91.4\% for the Webb wild-$t$, below the spec's pre-committed $93.0\%/91.0\%$ qualification bar, 84.0\%/85.7\% for CR1, 89.5\%/89.9\% for the demoted percentile read, and 99.5\%/99.6\% (over-covering, and wider) for CR2 at Bell--McCaffrey. The frozen selection rule therefore lands the binding layer on the restricted inversion, whose $[+2.3, +9.1]$ this paper quotes. The Webb wild-$t$'s narrower $[+2.9, +8.7]$ is reported beside it, no longer as the binding layer. So the binding-layer choice is measured, not asserted, and the rule was applied after seeing which way it fell. What that layer prices is itself form-conditional: the same $0.99$-point rise in the floor that carries the read off window moves the marginal by $-3.63$ points under the production form and by $-0.04$ under the additive one. The interval's width is as much a property of the max form's censoring as of the floor read. Two further v20 companions price adjacent layers. Re-run on each leg's own counterfactual balance path, not renormalized monthly to realized holdings (run \texttt{compounding\_consistent\_null}), the marginal is $+4.84$ points (\$37.0 billion; central 92.3\%, null 87.4\% shared). That prices the renormalization's disclosed upper-bound bias at $0.74$ points within the quoted interval. The floor read's sampling layer propagated to the null's level puts a wild-cluster interval of $[80.6, 88.4]$\% on the 85.7\%, beside the floor-range band of Table~\ref{tab:headline}. I attach no posture to where $+5.6$ sits among these readings. With the downward variants lined up beside the upward counterweights, no interior member is privileged, and the range rather than any point is what the design delivers. The form-conditional hull of $+3.5$ to $+13.1$ points remains the wider statement, since the form dimension is not a floor correction and is not disciplined by this list. Three things stop the list from becoming a retraction. The sign is fixed by the window's rate configuration, not by these calibrations, as the sign-forcing argument above establishes, so nothing here bears on the direction. Zero-exclusion is deliberately not among the three, a forced sign making it a restatement of that configuration. Specification choices that are not floor corrections run both ways, and this assembly lines both up. Upward: the additive form returns $+11.2$ points at the same off-window anchors, nearly floor-invariantly, and sets the hull's upper end. \citepos{fonseca2024} anchor would put the marginal at $+11.5$ points at the production floor against the adopted band's $+9.2$, and the housing-activity calibration below, re-solved under that form, joins them. Downward: the concave transform alone returns $+5.06$ at the off-window point floor against the production $+5.57$, and the sub-band elasticity at $\delta = 3.25\%$ returns $+3.47$ (Table~\ref{tab:lowband}). The moving-share bracket at the headline floor returns $+3.36$ points at $\mu = 0.5$ and $+1.82$ at $\mu = 0.25$ against $+5.57$ at $\mu = 1$. It is a convention bracket, not an estimate, mildly super-proportional because censored contributions do not scale with $\mu$. The survival-selection attenuation of the same imported coefficient is a bracket of the same kind and belongs here: $\eta = S_{\mathrm{pre}}^{\xi}$ at the measured pre-window survival share carries the marginal from $+5.1$ at $\xi = 0.25$ to $+3.7$ at $\xi = 1$, $\xi$ free (Section~\ref{sec:pathb}) --- a different axis from the elasticity band, not a widening of it. The assembly is not exhaustive over runs scoring a marginal at a labelled floor. The PSA sweep, the covariate ablation and the floor-form mixture do so too, reported at their own sites. I state no rule admitting these rows while excluding those. A calibrated companion prices the omitted housing-activity level term (Table~\ref{tab:specbox}) from the paper's own catalogued anchor. \citet{aladangady2024} attribute 44\% of the 2021--2022 mobility decline to rate-gap lock-in. Scaling the seasoning baseline in both legs until the non-rate channel traps $0.56/0.44$ times the rate channel's trapped liquidity puts the headline floor's root at $\phi^{*} = 0.754$ (effectively 75.4 PSA) and the marginal at $+0.9$ points at the headline floor and $+3.5$ at the in-sample calibration (run \texttt{scaled\_null\_housing\_activity}). The $\phi = 1$ cells reproduce the committed values bit-exactly. The variant documents an upward bias under both forms; under the production form the corrected member lies below the headline, at the bottom of the ramp span just swept, and at the headline floor falls outside the quoted $+2.3$ to $+9.1$ interval. That is an externally anchored calibration the quoted interval does not cover, not a challenge to the sign. At the in-sample calibration it falls inside. Re-solving the same condition under the additive form barely moves the root ($\phi^{*} = 0.756$ against $0.754$ at the headline floor, a direction I declined to predict). It lands the surviving marginal at $+8.5$ points at both anchors, clearing the $+4.0$-point bar fixed before the run (run \texttt{scaled\_null\_housing\_activity\_additive}; the $\phi = 1$ cells reproduce the committed additive legs bit-exactly). The calibration cuts the additive form's own $+11.2$ to $+8.5$ as it cuts the production form's $+5.6$ to $+0.9$. What the form decides is where the corrected member lands: above the headline at both anchors and near the top of the quoted interval under the additive form, below the headline at both under the production one. Censoring separates them: at the production form's root the central leg is pinned to the floor in 94.4\% of evaluated loan-months and the maximum truncates the elasticity there. The additive form combines the two hazards on the survival scale, so the floor, which dominates no less often at the additive root, truncates nothing. On that side of the fork this reading belongs with the counterweights above rather than with the ladder. That the floor side runs one way is also an absence. The floor is a level of total turnover, not of a strictly involuntary component (Section~\ref{sec:pathb}), so it carries its period's activity. Re-reading it at the window's own level enters at the floor, not the ramp just scaled, and so cuts the other way --- a lower anchor raises the marginal. It is the one floor-side correction I can name that would move the headline up, turning on whether the window's turnover level is the lower of the two, which I do not measure here. And the three readings of the preceding paragraph point to a larger lock-in channel, not a smaller one. I decline those three for the reasons given there and decline to compose these for the reason given here: the identified content is the range. Floor and basis corrections concentrate in its lower half, and the form-side counterweights sit at its top or above it. Table~\ref{tab:assembly} tabulates this assembly.

\begin{table}[!t]
\centering
\caption{The assembly at a glance. Every value repeats a committed figure from the text of this subsection; the binding interval and the form-conditional hull are deliberately not restated here (see the text and Table~\ref{tab:uncertainty}).}
\label{tab:assembly}
\footnotesize
\begin{tabular}{@{}p{3.1cm}p{2.6cm}p{2.4cm}p{7.7cm}@{}}
\toprule
Reading & Marginal (pp) & vs.\ $+5.6$ & Status \\
\midrule
In-sample calibration point & $+9.2$ & above & in-window point, not the headline \\
Off-window re-anchor & $+5.6$ & headline & mid-grid anchor (the middle of the three depth cuts fixed before the reads) inside the binding interval \\
Ginnie overlay at the corrected floor & $+4.4$ & below & change of estimand (re-scoping onto the conventional sub-book), not a correction \\
Ginnie share's gap response priced & $+4.4$ to $+5.8$ & spans it & bracket over the mapping, not an estimate: a level differential does not identify a slope (run \texttt{ginnie\_overlay\_attenuated}) \\ Vintage overlay at the corrected floor & $+3.7$ & below & change of estimand (re-scoping onto the in-window-vintage share), not a correction (run \texttt{vintage\_overlay}) \\
Age-standardized floor (5.51\%) & $+3.7$ & below & indicative bound, 84\% imputed weight; measured directly (run \texttt{b5\_joint\_cell}) against the grid's $\sim +3.8$, a $-0.13$-point wedge; the printed tie with the vintage-overlay row above is a rounding accident ($3.67$ against $3.70$) \\
Fannie floor read (5.52\%) & below $+4.3$ & below & bracketing; second agency, pooled period; like-for-like Freddie 5.19\% \\
Overlay $\times$ age-standardized & $+2.9$ & below & measured composed lower member (run \texttt{b5\_joint\_cell}); within $0.1$ of proportional \\ Moving-share bracket, $\mu = 0.5$, off-window & $+3.4$ & below & convention bracket, not an estimate (run \texttt{moving\_share\_bracket\_offwindow}) \\ Moving-share bracket, $\mu = 0.25$, off-window & $+1.8$ & below the interval & the bracket's deep stress; edge by construction \\ Book-composition post-stratification & $+5.11$ & below & two margins raked (SOMA coupon $\times$ book vintage) on the 2017--2021 support; the joint is unobserved and 33.7\% of face is out of support (run \texttt{book\_composition\_marginal}) \\ Survival-selection attenuation, $\xi = 0.25$ to $1$ & $+5.1$ to $+3.7$ & below & bracket in $\xi$ on the imported coefficient ($\eta = S_{\mathrm{pre}}^{\xi}$, $\xi$ not estimable here); a different axis from the elasticity band (run \texttt{attenuation\_sensitivity}) \\ Housing-activity-scaled baseline (Aladangady-anchored) & $+0.9$ (max) / $+8.5$ (additive) & below the interval / above the headline & form-conditional scaled-null variant; upward-bias entry under both forms, roots $\phi^{*} = 0.754$ and $0.756$ at the headline floor \\ Sub-band elasticity, $\delta = 3.25\%$, off-window & $+3.5$ & below & band-extension member (Table~\ref{tab:lowband}) \\ Concave transform alone, off-window & $+5.1$ & below & transform member (run \texttt{concave\_marginal}) \\
Concave $\times$ additive joint cell & interaction $-1.47$ & --- & the one joint cell run \\
Additive form & $+11.2$ & above & counterweight, functional-form choice \\
Fonseca--Liu anchor & $+11.5$ & above & counterweight, conservative-band evidence \\ Cross-sectional gradient, composition-standardized & gradient $+3.4$ vs implied $+0.9$ & above (directional) & directional; not a marginal re-estimate (run \texttt{episode\_confrontation\_within}); holding loan age too leaves $+3.4$, but its interval $[+0.6, +4.4]$ covers the implied value (run \texttt{episode\_gradient\_age\_bands}) \\
\bottomrule
\end{tabular}
\end{table}

\begin{table}[!t]
\centering
\footnotesize
\begin{threeparttable}
\caption{Out-of-sample floor anchors and the lock-in marginal each implies, in percentage points of the \$764.7 billion benchmark. Floors are exposure-weighted annualized turnover CPR on out-of-the-money cohort-months measured \emph{outside} the June 2022--November 2025 evaluation window, except the early-window calibration row, whose floor is calibrated on the window's own first nineteen months. The pooled and 2019-leg rows read 2017--2019 performance; the three 2018-leg rows that set the headline read 2018 reporting months on 2017-vintage cohorts, and every cohort-month they contain is aged twelve to twenty-four months, so no mature cell enters them. Every re-run row, the early-window one included, reports its floor's \emph{full-window} marginal on a common basis, and the held-out-months evaluation is given in the note; the age-standardized and Fannie rows are grid reads rather than paired re-runs, so they carry the central elasticity alone. Contaminated anchors are reported for completeness and are not this paper's off-window anchor; the pooled anchor's 6.07\% measurement is evaluated at the box's adjacent 6.0\% grid row. Every marginal is basis-invariant and reproduces the committed production artifacts exactly at the 4\% row. Pre-committed run \texttt{oos\_identification}; artifact catalogued in the replication package's ledger.}
\label{tab:oosfloor}
\begin{tabular}{llcccc}
\toprule
Floor anchor & Status & Floor & $\delta = 5.5\%$ & $\delta = 6.5\%$ & $\delta = 7.7\%$ \\
\midrule
2018 leg, gap $\leq -0.005$ & defensible & 4.70\% & $+6.1$ & $+6.8$ & $+7.5$ \\
2018 leg, gap $\leq -0.0025$ & defensible & 4.99\% & $+5.0$ & $+5.6$ & $+6.1$ \\
2018 leg, gap $\leq 0$ & defensible & 5.33\% & $+3.9$ & $+4.3$ & $+4.6$ \\
\midrule
Pooled 2017--2019 & refi-contaminated & 6.07\% & $+2.1$ & $+2.3$ & $+2.4$ \\
2019 leg, gap $\leq 0$ & refi-contaminated & 6.91\% & $+0.7$ & $+0.7$ & $+0.8$ \\
\midrule
2018 leg, age-standardized & indicative bound\tnote{b} & 5.51\% & --- & $+3.8$ & --- \\
Fannie Mae, pooled off-window & cross-agency read\tnote{c} & 5.52\% & --- & $+2.3$ to $+4.3$ & --- \\
\midrule
In-window (production) & in-sample point & 4.00\% & $+8.1$ & $+9.2$ & $+10.4$ \\
Early-window calibration & full window\tnote{a} & 3.97\% & $+8.1$ & $+9.3$ & $+10.5$ \\
\bottomrule
\end{tabular}
\end{threeparttable}
\end{table}

\noindent{\footnotesize \emph{Notes to Table~\ref{tab:oosfloor}.} a Temporal holdout. The 3.97\% floor is calibrated on deeply out-of-the-money turnover in the window's first nineteen months (June 2022--December 2023). The remaining twenty-three months (January 2024--November 2025) enter that calibration nowhere; summed over those strictly held-out months, the central-minus-null marginal is $+\$45.1$ billion at the central elasticity, against $+\$44.8$ billion at the production floor over the identical months. The table row above reports the same floor's full-window marginal, so that every row is comparable; the two quantities are different objects and should not be read off one another. \quad b The 2018 leg's age-bucket reads reweighted to the in-window deep-discount population's age mix, its empty mature buckets imputed from the in-window age gradient: 84.0\% of the weight is imputed, so this is an indicative bound and not a measurement. \quad c The committed off-window selection rule applied to the Fannie Mae panel at gap $\leq -0.0025$ over the pooled 2017--2019 period, not to the 2018 leg. The like-for-like Freddie read of that cell is 5.185\%, so the agency difference is 0.337 points of CPR and the rest of the distance from the 2018 leg's 4.991\% is period. The Fannie row is read off the committed floor-to-marginal grid rather than re-run. The age-standardized row no longer is: it has since been re-run directly (run \texttt{b5\_joint\_cell}), and returns $+3.7$ against that grid's $+3.8$ --- a $-0.13$-point wedge. The table above keeps the printed $+3.8$, the grid read its neighbouring rows share.\par}

\begin{table}[!t]
\centering
\caption{The lock-in marginal as a curve in the imported elasticity. Paired central-minus-null runs at each quarterly mobility decline $\delta$ (the elasticity import of Section~\ref{sec:pathb}, mapped to $\beta_1$ via \eqref{eq:beta1}), with the floor and every other production convention held fixed, at the in-sample 4.0\% floor and the off-window 4.991\% headline floor (run \texttt{band\_low\_extension}; the band-cell rows reproduce the committed production and off-window cells within the sweep tolerance of $\pm\$0.01$ billion). The $\delta = 0$ rows are the $\beta_1 = 0$ nulls, so their marginal is zero by construction; $\delta = 3.25\%$ is half the adopted central; $\{5.5, 6.5, 7.7\}\%$ is the adopted Liebersohn--Rothstein band.}
\label{tab:lowband}
\footnotesize
\begin{threeparttable}
\begin{tabular}{@{}rrrrrr@{}}
\toprule
& & \multicolumn{2}{c}{In-sample floor (4.0\%)} & \multicolumn{2}{c}{Off-window floor (4.991\%)} \\
\cmidrule(lr){3-4}\cmidrule(lr){5-6}
$\delta$ (\%) & $\beta_1$ & marginal (\$B) & (pp) & marginal (\$B) & (pp) \\
\midrule
0.00 & $0.0000$ & $0.0$ & $0.00$ & $0.0$ & $0.00$ \\
1.00 & $0.0103$ & $+12.8$ & $+1.67$ & $+9.4$ & $+1.23$ \\
2.00 & $0.0206$ & $+24.9$ & $+3.26$ & $+17.7$ & $+2.31$ \\
3.00 & $0.0311$ & $+36.4$ & $+4.76$ & $+24.9$ & $+3.25$ \\
3.25 & $0.0337$ & $+39.1$ & $+5.12$ & $+26.5$ & $+3.47$ \\
4.00 & $0.0417$ & $+47.1$ & $+6.16$ & $+31.0$ & $+4.06$ \\
5.00 & $0.0523$ & $+57.1$ & $+7.46$ & $+36.3$ & $+4.74$ \\
5.50 & $0.0577$ & $+61.7$ & $+8.07$ & $+38.6$ & $+5.04$ \\
6.50 & $0.0686$ & $+70.3$ & $+9.20$ & $+42.6$ & $+5.57$ \\
7.70 & $0.0817$ & $+79.5$ & $+10.39$ & $+46.6$ & $+6.09$ \\
\bottomrule
\end{tabular}
\begin{tablenotes}[flushleft]\footnotesize
\item The $\beta_1$ column follows \eqref{eq:beta1}'s sign convention, under which $\beta_1$ is positive and suppression enters \eqref{eq:pathB} through the signed gap $g_{i,t}$, negative for a locked-in borrower. The production helper \texttt{rothstein\_beta1} returns the same quantity without that leading minus, so a replicator reading the code will see these nine values negated; the marginal columns are unaffected either way.
\end{tablenotes}
\end{threeparttable}
\end{table}

\begin{figure}[tbp]
\centering
\includegraphics[width=0.78\textwidth]{fig11_marginal_decomposition}
\caption{Group-ablation decomposition of the lock-in marginal over the vintage $\times$ coupon grid (production configuration; each cell's $\beta_1$ zeroed with everything else held fixed). All fifteen cells contribute positively; the 2020--2021 rows carry 72\% of the cell sum on 73\% of sample balance, and no cell exceeds 21\%. Cell values and shares are read from the committed decomposition artifact, whose additivity gate bounds the cells' sum within 1.03\% of the paired-run \$70.3 billion in-sample total.}
\label{fig:marginalcells}
\end{figure}

Table~\ref{tab:uncertainty} consolidates the three uncertainty layers --- sampling, simulation seed, and calibration --- that attach to the headline level and marginal on both accounting bases. The hierarchy it displays is the one argued above: sampling error under the operative cluster schemes sits an order below the calibration layers, the draw-seed contribution is smaller still (under \$0.01 billion across ten paired reruns), and the operative uncertainty on every headline quantity is calibration. The headline marginal is the exception. There the binding layer is the floor read's own propagated sampling error (Section~\ref{sec:robustness-floor}) --- the layer Table~\ref{tab:uncertainty}'s headline row quotes, binding among the layers with a coverage property. The unestimated baseline-level convention spans wider (run \texttt{psa\_level\_sweep}, Section~\ref{sec:identification}).

{\footnotesize
\begin{longtable}{@{}p{2.4cm}p{3.3cm}p{4.0cm}p{4.8cm}@{}}
\caption{Consolidated uncertainty for the headline Path B quantities. Sampling intervals are the stratified loan-level bootstrap's 95\% percentile intervals (200 replicates, Section~\ref{sec:pathb}); that scheme resamples within strata and holds the 130 three-way draw strata fixed, so the headline row also carries the stratum-cluster interval, which is the wider of the two loan-sampling reads; shared-basis intervals follow from the standalone intervals by the exact common netting of Section~\ref{sec:robustness-hybrid} (9.10 points of benchmark), which cancels from the marginal entirely. Calibration entries restate Sections~\ref{sec:pathb} and~\ref{sec:robustness-floor}.}
\label{tab:uncertainty}\\
\toprule
Quantity & Point (shared; standalone) & Sampling 95\% CI & Calibration range \\
\midrule
\endfirsthead
\toprule
Quantity & Point (shared; standalone) & Sampling 95\% CI & Calibration range \\
\midrule
\endhead
Path B central level (in-sample calibration) & 97.9\%; 107.0\% & $[97.9, 98.0]$\%; $[107.0, 107.1]$\% & covariate priors 92.2--99.2\%; elasticity band 96.8--99.1\% (shared basis); off-window floor 91.3\% shared (88.2--93.7\% across its range); covariate endpoints at the off-window floor 94.3--102.5\% standalone, an 8.2-point swing exceeding the marginal there (run \texttt{covariate\_ablation\_offwindow}) \\
$\beta_1 = 0$ null level (in-sample calibration) & 88.7\%; 97.8\% & $[88.7, 88.8]$\%; $[97.8, 97.9]$\% & moves one-for-one with the floor level: 85.7\% shared at the off-window floor, 83.9--86.9\% across its range; floor-read wild-cluster $[80.6, 88.4]\%$ on that level (Section~\ref{sec:robustness-floor}) \\
Lock-in marginal, off-window floor (headline) & $+5.6$pp (basis-invariant) & floor-read wild-cluster inversion $[+2.3, +9.1]$pp at the central elasticity $\delta = 6.5\%$, the binding layer (it prices the floor read's own sampling error alone --- the notes below state what it does not price; percentile read $[+3.0, +8.0]$pp; the elasticity dimension is Table~\ref{tab:lowband}); loan/stratum cluster bootstrap $[+4.63, +6.92]$pp (width $0.34\times$ the corrected floor read, meaning the restricted inversion, against $0.455\times$ the demoted percentile read; 25.8 effective clusters); the two layers convolved as independent, $[+2.80, +8.99]$pp (run \texttt{layer\_convolution}, computed on the Webb rung and not re-derived under the adjudicated binding layer; independence is assumed, not measured --- under maximal positive dependence the width is $8.12$pp) & $+4.3$ to $+6.8$ across the defensible off-window floor range (Table~\ref{tab:oosfloor}); age-standardized floor 5.51\% $\to$ $+3.7$ measured, against $+3.8$ read off the grid (84\% imputed weight; band open below $+4.3$); the out-of-agency read, the form and transform variants, the overlays and the convention sweeps are catalogued in Appendix~\ref{app:floormech} \\
\addlinespace
Lock-in marginal, in-sample point & $+9.2$pp (basis-invariant) & within-stratum $[+9.17, +9.23]$pp (sd $0.016$); stratum-cluster $[+8.27, +10.19]$pp (sd $0.52$; run \texttt{bootstrap\_pathb\_cluster}) & box $+2.1$ to $+13.2$; floor form $+11.2$; within-run floor dispersion $+7.1$ to $+9.4$; regime-split baseline break $+5.0$ to $+12.1$; PSA level sweep $+3.3$ to $+16.5$ (run \texttt{psa\_level\_sweep}); covariate endpoints move this marginal $+1.9$ \\
\bottomrule
\end{longtable}
}

\noindent{\footnotesize \emph{Notes to Table~\ref{tab:uncertainty}.} Seed noise: the marginal varies by less than \$0.01 billion across ten paired draw seeds (Section~\ref{sec:hazard-interp}). The Fannie Mae external replication reproduces the marginal at $+8.68$ points against Freddie's $+9.20$. It clears the pre-committed envelope, but that envelope is 11.06 points wide and the agreement, not the gate, is what carries the information (Section~\ref{sec:pathb-fannie}). \quad Both cluster bootstraps above rest on few clusters: 31 carry the floor read, and the loan/stratum scheme's 130 strata are unequal enough to be worth 25.8 effective ones. Percentile cluster bootstraps are known to under-cover at cluster counts in that range. This rung's printed lower edge additionally turns on an unlabelled convention: 26 of its 1{,}000 draws fall outside the committed floor grid, all 26 above the 6.0\% edge and therefore all on the low-marginal side, and the row excludes them. Truncating them to the edge instead would print $+2.28$ to $+8.01$ in place of $+3.0$ to $+8.0$. Runs \texttt{floor\_inference\_correction} and its Webb-weight successor \texttt{floor\_inference\_correction\_v2} correct the floor read. Table~\ref{tab:ladder} sets the corrected ladder out layer by layer with its degrees of freedom, and the headline row quotes the wild-$t$ interval as the binding layer among those with a coverage property. That interval prices one layer: the floor read's own sampling error around the 4.991\% point read, propagated through the frozen floor-to-marginal grid at the production hard-maximum form, the 100~PSA ramp and the central elasticity $\delta = 6.5\%$. None of its width comes from the elasticity, which is swept separately in the calibration column and in Table~\ref{tab:lowband} --- swept over the published 5.5\%--7.7\% range as a convention, not estimated, so no layer of this paper's accounting prices sampling error in the elasticity itself. Nor does any of its width come from the read-to-read spread across the floor reads in the calibration column. No such correction is run for the loan/stratum scheme, whose interval is still better read as a lower bound on sampling uncertainty than as calibrated 95\% coverage. One input to the convolved line being a lower bound makes the convolved line a lower bound as well. The ladder's reads agree in location and none crosses zero. The closest agreement among them carries no information: the Rademacher and Webb constructions land with both endpoints within a tenth of a point, a re-printing rather than a substantive move. Both are sign-symmetric reweightings of the same residuals, so each bootstrap distribution turns on the draw's sign for the one cluster carrying a third of the total leverage---nearly twice the next largest (Table~\ref{tab:ladder}). Their coincidence re-prints that cluster's influence instead of testing it. The small-sample degrees of freedom and the restricted inversion are wider. The bootstrap imposes the point estimate as the data-generating truth over unrestricted residuals (9{,}999 replications), studentizes on the floor's own scale, and maps endpoints through the frozen floor-to-marginal grid. \quad The within-run floor dispersion range is the $\kappa$ grid of Section~\ref{sec:pathb}. It bounds the marginal's sensitivity to floor variation within the window; it does not bracket floor cyclicality, and should not be read as doing so. Scrambling the floor path's time order leaves 85.0\% of its span standing, and switching the combination rule from the hard maximum to the additive form eliminates 98.0\% of it (the $\kappa$ grid then spans $+11.22$ to $+11.27$ points). Appendix~\ref{sec:robustness-seasonalfloor} diagnoses the Jensen mechanism responsible. The co-movement component that does survive is monotone in $|\kappa|$ and sign-consistent, but reaches at most \$5.29 billion in any cell. \quad The calibration catalogue for the headline marginal row --- the out-of-agency read, the form and transform variants, the overlays, and the convention sweeps --- is in Appendix~\ref{app:floormech}.\par}

\clearpage
\subsection{Interpretation}\label{sec:hazard-interp}

Table~\ref{tab:estimators} summarizes the three estimators evaluated in this paper against the common \$764.7 billion empirical benchmark.

\begin{table}[!tp]
\centering
\footnotesize
\begin{threeparttable}
\caption{Estimator comparison against the \$764.7 billion empirical trapped-liquidity benchmark.}
\label{tab:estimators}
\renewcommand{\arraystretch}{0.92}
\setlength{\itemsep}{0pt}
\begin{tabular}{@{}>{\raggedright\arraybackslash}p{3.2cm}>{\raggedleft\arraybackslash}p{2.4cm}>{\raggedleft\arraybackslash}p{2.4cm}>{\raggedright\arraybackslash}p{2.3cm}>{\raggedright\arraybackslash}p{4.0cm}@{}}
\toprule
Estimator & Trapped liquidity & \% of benchmark & CPR path corr. (lag) & Path diagnostics \\
\midrule
Empirical benchmark (SOMA, phased cap, 30+15yr) & \$764.7B & 100.0\% & -- & mean CPR 5.14\% (ABM-basis back-out; hazard-basis 5.52\%, 5.79\% at the note-rate WAC) \\
Production ABM (Section~\ref{sec:abm-results}) & \$103.7B seed mean (frozen draw \$91.0B) & 13.6\% (seed band 10--17\%; frozen draw 11.9\%) & $r = -0.318$ (0) & mean CPR 11.68\%; raw $R^2$ = $-6.984$ \\
Hazard Path B: literature microsimulation & \$818.5B & 107.0\% (band 105.9--108.2\%) & $r = +0.404$ ($-3$) & mean CPR 4.76\%; $r = +0.190$ at lag 0; 3-month offset (empirical leads; reference only, Section~\ref{sec:pathb}) \\
Hazard Path B: concave-gap variant (Section~\ref{sec:pathb}) & \$808.4B & 105.7\% & $r = +0.281$ (0) & piecewise transform, kink 200 bp; robustness on the log-linear extrapolation \\
Hazard Path A: stratum-month Poisson PML (spec v4, calendar-month) & \$928.9B & 121.5\% & $r = -0.438$ ($-2$) & mean CPR 3.34\%; holdout diagnostics under spec v3 \\
Cross-design ABM, real covariates, recalibrated (Section~\ref{sec:robustness-crossdesign}) & \$453.5B & 59.3\% (fifty-seed mean 60.2\%) & $r = -0.336$ (0) & mean CPR 8.14\%; no positive alignment at any lag \\
Cross-design ABM, real covariates, frozen calibration (Section~\ref{sec:robustness-crossdesign}) & \$159.5B & 20.9\% & $r = -0.311$ (0) & mean CPR 11.63\%; no positive alignment at any lag \\
Hazard Path B: Fannie Mae external replication (Section~\ref{sec:pathb}) & \$821.0B & 107.4\% & -- & null 98.7\%; marginal $+8.68$pp \\
\bottomrule
\end{tabular}
\end{threeparttable}
\end{table}

\noindent{\footnotesize \emph{Notes to Table~\ref{tab:estimators}.} The table reports the three primary estimators, the two cross-design variants of Section~\ref{sec:robustness-crossdesign}, and the Fannie Mae external replication of Path B, whose path-correlation cell is dashed because the committed replication artifact carries headline statistics only (Section~\ref{sec:pathb}). The Danish counterfactual specifications are collected in Table~\ref{tab:danish}. Dashes denote quantities that are not applicable to the row. The benchmark row's ``30+15yr'' scope is not symmetric across estimators: 15-year holdings enter the empirical back-out and the ABM through real weights and real scheduled amortization, with voluntary CPR for those cohorts drawn from the same-coupon 30-year surface (the structural-only fold-in of Appendix~\ref{sec:robustness-pipeline}). Every production ABM figure in this row --- dollars, share, path correlation, $R^2$ and mean CPR --- is that fold-in run, while Section~\ref{sec:robustness-scale}'s scale test sits on the native-gate side of the offset between the two and is reconciled against this row there. The hazard rows are 30-year-collateral estimates from which 15-year face is excluded (Table~\ref{tab:composition}, Term row). Peak-lag entries select the maximum-magnitude cross-correlation (max $|r|$) within a $\pm$3-month window. Under a most-positive-$r$ convention the production ABM's best alignment is $+0.19$ at lag $-3$ (Section~\ref{sec:abm-results}); the two conventions summarize the same lag profile, and Path B's lag $-3$ peak satisfies both. \quad Path B is the only estimator whose aggregate recovery is close to the benchmark-consistent basis (97.9\%, and 100.0\% composed as committed --- 98.1\% converted to a single coupon convention --- at the in-sample calibration point; 91.3\% at the off-window floor the paper headlines). Its raw lag-0 correlation ($+0.190$) is carried by the window's shared trend: in first differences ($-0.23$) and after linear detrending ($-0.20$) the simulated--empirical co-movement is indistinguishable from zero ($\pm 0.31$ at $n = 41$), as it is for every estimator. The two cross-design rows substitute real Freddie Mac structural covariates into the ABM under two calibration choices. Neither reproduces the empirical path's timing --- no estimator does, as above --- but the recalibrated variant's dollar recovery exceeds my own ex-ante threshold for undercutting the paradigm interpretation (Sections~\ref{sec:robustness-crossdesign} and~\ref{sec:conclusion}). Moving-block bootstrap 95\% intervals (block length 6): ABM lag-0 $[-0.60, +0.04]$; Path A lag-0 $[-0.68, -0.02]$; Path B lag-0 $[-0.13, +0.44]$, peak lag $[-0.21, +0.49]$ (no positive lag-0 coefficient is distinguishable from zero; block-length sensitivity at 4/6/8 changes no zero-inclusion verdict in this table except Path A's marginally negative levels endpoint, detailed in Section~\ref{sec:pathb}). Path A's contemporaneous levels correlation is negative on the levels construction, which strengthens the offsetting-errors reading of its aggregate fit (detrended, all three estimators cluster at $-0.20$ to $-0.30$; Section~\ref{sec:pathb}). No timing evidence is claimed for any estimator. Table~\ref{tab:theil} decomposes each estimator's monthly error into bias, variance, and covariance shares and reports cumulative dollar-error profiles; it is a fit diagnostic and changes no verdict in this table. \quad Correlation-column convention: the ABM row's correlations are computed against the ABM pipeline's own empirical back-out (cohort-weighted, term-aware scheduled amortization), while the hazard rows use the hazard-side back-out. Against the common hazard-side series the ABM lag-0 value is $-0.316$, a third-decimal, no-verdict difference. Section~\ref{sec:method-benchmark} notes the two constructions. \quad Shared-accounting-basis equivalents, together with the full-book and composed bases, are collected in Table~\ref{tab:bases} (basis defined in Section~\ref{sec:method-benchmark}, derived in Section~\ref{sec:robustness-hybrid}; the common curtailment netting is a flat 9.1pp for every U.S. hazard leg). Path A's point estimate carries no useful forward precision and is excluded from headline figures. Section~\ref{sec:patha} reports its refit-and-resimulate interval --- since rerun jointly at both specifications, with the joint interval floor-censored and reported beside, not in place of, the committed one --- and the spec v3 provenance of every Path A inferential object.\par}

\begin{table}[H]
\centering
\footnotesize
\begin{threeparttable}
\caption{Recovery of the \$764.748 billion benchmark by accounting basis. Every U.S. hazard leg is quoted on one of these bases: the shared basis nets the common \$69.56 billion curtailment flow from the standalone-scorer figure (a flat 9.1 percentage points), and the composed basis applies the same netting to the book-coupon-reweighted (full-book) figure. The lock-in marginal (central minus $\beta_1 = 0$ null) is basis-invariant because the netting cancels.}
\label{tab:bases}
\begin{tabular}{@{}lrrrrr@{}}
\toprule
Hazard leg & Standalone & Shared & Full-book & Composed & Cum.\ runoff \\
 & (scorer) & ($-9.1$pp) & (book WAC) & ($-9.1$pp) & error (\$B) \\
\midrule
Path B, central (in-sample calibration) & 107.0\% & 97.9\% & 109.1\% & 100.0\% & $+53.8$ \\
Path B, $\beta_1 = 0$ null (in-sample calibration) & 97.8\% & 88.7\% & -- & -- & $-16.6$ \\
\addlinespace
Path B, central (off-window floor) & 100.4\% & 91.3\% & -- & -- & $+2.8$ \\
Path B, $\beta_1 = 0$ null (off-window floor) & 94.8\% & 85.7\% & -- & -- & $-39.8$ \\
\addlinespace
Path A (spec v4) & 121.5\% & 112.4\% & 127.5\% & 118.4\% & $+164.1$ \\
Path B, concave-gap & 105.7\% & 96.6\% & -- & -- & $+43.7$ \\
\bottomrule
\end{tabular}
\end{threeparttable}
\end{table}

\noindent{\footnotesize \emph{Notes to Table~\ref{tab:bases}.} The book-WAC and composed columns are the committed mixed-basis measurements (note rates bucketed against pass-through coupons). Converted to a single pass-through convention they read 107.2\% and 98.1\% for Path B central, with the identified marginal bit-invariant (run \texttt{coupon\_convention\_reweight}). \quad A prior fact frames the whole table. The engine renormalises the simulated pool to realized Fed holdings every month, so a leg's dollar roll-off is a \emph{rate} applied to an exogenous level, and its balance path cannot drift from the realized one. No recovery percentage here is a balance-path fit, and none can compound error across months: each is close to a restatement of that leg's window-mean CPR against the empirical 5.14\% (ABM basis). That is why the sampling intervals are as tight as they are, and it is a limit on what any of these levels can corroborate. The final column is each leg's terminal cumulative runoff error, realized minus simulated roll-off (equivalently, simulated minus realized trapped liquidity) --- the sign-transparent statement of the same information. Recovery is an affine transform of it, $\text{recovery} = 1 + \text{error}/764.7$, so a large additive constant sits inside every percentage in the first four columns and makes all of them look better. Two things the percentages hide. Path~B's in-sample 107.0\% is an \emph{over}-prediction of trapped liquidity of \$53.8 billion, which is an 8.2\% \emph{under}-prediction of the \$652.8 billion of runoff actually realized---the referent flips even though the map does not. And the headline off-window central leg, at 91.3\% shared, carries the smallest error of any leg here on the standalone scorer: $+\$2.8$ billion, 0.4\% of realized runoff. Errors as a share of realized runoff: $+8.2\%$, $-2.5\%$, $+0.4\%$, $-6.1\%$, $+25.1\%$, $+6.7\%$ down the rows. The lock-in marginal is unaffected by any of this: it is a difference of two trapped aggregates from which the constant cancels, which is the same fact as its basis-invariance. Standalone and shared are related by the flat 9.1-percentage-point curtailment wedge of Section~\ref{sec:method-benchmark} (\$69.56 billion over the window), derived in Section~\ref{sec:robustness-hybrid}. Path B's central elasticity band is 105.9--108.2\% standalone (96.8--99.1\% shared); the covariate-prior band is 92.2\%--99.2\% shared (approximately 94.3\%--101.3\% composed). The lock-in marginal is $+9.2$ points ($+\$70.3$ billion) on the U.S. leg at the in-sample calibration point, and $+5.6$ points ($+\$42.6$ billion) at the off-window floor (Section~\ref{sec:robustness-floor}). Both are basis-invariant because the netting is common to the central and null legs. The off-window rows carry the same flat netting: the layer is a constant curtailment flow over a constant cap benchmark, identical across all five committed estimators spanning 29.6 points of standalone recovery, so it is path-invariant and therefore floor-invariant (\texttt{calibration\_reconciliation}). Full-book and composed columns were run only at the in-sample calibration and are not extrapolated to the off-window floor. Path A, the concave-gap variant, the elasticity band, and the covariate-prior band are all in-sample-calibration quantities. The Danish leg's counterfactual image, run at the production floor and so comparable to the in-sample point, is $+\$61.2$ billion, or 8.0 points, the difference reflecting the Danish leg's dynamic-balance compounding (Section~\ref{sec:abm-danish}). The ABM rows of Table~\ref{tab:estimators} run through the accounting layer natively (Figure~\ref{fig:recovery}) and are not restated here.\par}

The two hazard paths share no household-level decision machinery with the ABM. They land closer to the empirical benchmark than any synthetic-population ABM specification tested here (the falsified burnout attempt included) at 97.9\% and 112.4\% in-sample on Section~\ref{sec:robustness-hybrid}'s benchmark-consistent shared accounting basis (107.0\% and 121.5\% under their standalone scorers). At the off-window floor the paper headlines, Path B is 91.3\% shared; Path A was not re-run there. Overshooting less is not validation. Path A overshoots by \$164.1 billion (its terminal cumulative runoff error, a 25.1\% under-prediction of the \$652.8 billion realized) with a monthly path fit no better than the ABM's. Path B's recovery across its full elasticity sensitivity band (\$809.9--\$827.7 billion, 105.9\%--108.2\%; Section~\ref{sec:pathb}) rests on an elasticity imported unmodified from \citet{liebersohn2024}, its verification content limited by Section~\ref{sec:pathb} to the marginal's bounded interval. My claim is therefore narrower than ``the hazard framework fits'' and narrower again than ``lands within 2.1\% of the benchmark'': that 2.1-point miss belongs to the in-sample calibration point, so quoting it beside a marginal measured at the off-window floor mixes two calibrations. The 2.1-point miss is 97.9\%, its composed companion basis-dependent (100.0\% as committed on the mixed coupon basis, 98.1\% on a single pass-through convention, Section~\ref{sec:robustness-hybrid}), retained as the model-fit diagnostic it is, ex-ante covariate priors spanning 92.2\%--99.2\% and approximately 94.3--101.3\% composed, none re-run off-window. At the headline calibration Path B (loan-level survival, no household-utility content) recovers 91.3\% of the aggregate benchmark on the benchmark-consistent accounting basis, a miss of 8.7 points, and 88.2--93.7\% across that floor's defensible range, a miss of 6.3 to 11.8 points. What survives is order-of-magnitude comparison, not precision. Path A lands at 112.4\% sharing data but no estimated parameters, which I read as order-of-magnitude agreement and nothing stronger. It has no useful propagated precision in forward simulation, and a rate-gap coefficient not statistically distinguishable from zero under any bias-respecting construction (Section~\ref{sec:patha}). The same resampling pass puts its mean at 78.3\% of the benchmark with a percentile interval spanning zero, corroborating no sign. No synthetic-population household-choice specification I tested lands within 45\% of it. I attach no timing credential. Path B's raw $+0.190$ lag-0 correlation is carried by the window's shared trend, the co-movement indistinguishable from zero in first differences and after detrending (Section~\ref{sec:pathb}). Table~\ref{tab:estimators}'s moving-block intervals caution against weighing any single timing coefficient, Path B's lag $-3$ peak (interval spanning zero) included. The benchmark side is itself disqualified, its four exact-zero months clip-induced boundary artifacts of the reporting series (Appendix~\ref{sec:robustness-seasonalfloor}, run \texttt{h1\_zero\_months\_diagnosis}). Timing alone does not identify the lock-in channel. The $\beta_1 = 0$ null also peaks at lag $-3$, recovering 88.7\% of the benchmark on the shared basis in-sample and 85.7\% at the off-window floor (Table~\ref{tab:bases}), both under the production floor form. It recovers 35.6\% at that floor under the additive form (Section~\ref{sec:robustness-floor}). The elasticity's identified contribution is its marginal recovery over that null, defined and bounded once in Section~\ref{sec:identification}. The design identifies the marginal, not the lag structure. Across 10 paired in-sample reruns with different draw seeds (central and null sharing each seed) that marginal is \$70.3 billion, varying by less than \$0.01 billion: the engine drains prepayments fractionally and the seed touches only the small delinquency channel. The marginal's uncertainty is therefore the calibration band and structural assumptions, not simulation noise. I restrict this to synthetic-population specifications: Section~\ref{sec:robustness-crossdesign}'s real-covariate cross-design ABM variant does land within 45\% (59.3\%, and 76.3\% once that covariate population is reweighted to the book's own coupon $\times$ vintage composition, within 24\%), complicating the paradigm interpretation this paragraph would otherwise support---and the reweighted reading makes the complication larger, not smaller.

The hazard framework's only behavioral content is loan-level covariates, a cohort-level burnout term, and literature-calibrated elasticities---no analogue to mobility desire, loss aversion, or utility maximization. Its comparative success might therefore be read as evidence for the loan-level and institutional mechanisms Section~\ref{sec:abm}'s synthetic-population ABM does not represent: FICO, LTV, and geography heterogeneity; servicer forbearance and pipeline delay; cohort-level burnout. The paper's own measurements cap that reading. The named mechanisms ablate to a point or less each (burnout 0.86 points at the in-sample floor, 0.45 at the off-window floor; FICO/LTV 1.3 points; the delinquency pipeline \$0.39 billion, with no servicer channel beyond it modeled). Section~\ref{sec:robustness-synthesis}'s synthetic companion, meanwhile, recovers a comparable aggregate level with no real loan-level structure at all. And Section~\ref{sec:robustness-scale}'s data-source asymmetry remains unresolved. The level therefore supports the rate-inelastic decomposition, those mechanisms bounded small and candidates only for the residual. Section~\ref{sec:robustness-scale}'s scale-sensitivity replication rules out the ABM/hazard sample-size difference as the driver of this gap. Section~\ref{sec:robustness-crossdesign}'s cross-design test addresses the data-source difference only partially: real structural covariates there change the ABM's explanatory share substantially. That shift is enough, under one calibration choice, to raise the question of whether ``cannot represent'' is the right description at all. The rate-gap coefficient's sign stability across both hazard specifications does not indicate that the underlying par-payoff penalty is doing real work, and the reason bears restating: only Path A estimates that coefficient. Path B imports it from \citet{liebersohn2024} as a signed elasticity band that cannot come out the other way. Their agreement, having no freedom to disagree, is therefore an orientation check on unit and sign conventions, not corroboration. It was worth running---it would catch a transposed sign between the paths, and it passes. But the evidence that the penalty has a real effect is the external literature, Path A's estimated coefficient being directionally consistent without being statistically distinguishable from zero under bias-respecting constructions (Section~\ref{sec:patha}). Household rationality is not irrelevant, but the aggregation structure of loan-level survival is the minimum a model of this shortfall must get right, and Section~\ref{sec:robustness-crossdesign} leaves open how much more. That sharpens the two-sided exchange this paper describes, with the qualification Section~\ref{sec:robustness-crossdesign} adds: the price of the TBA market's homogeneity appears to be paid at least in part through loan-level and institutional frictions that the synthetic-population household model does not represent---frictions a market-value repurchase design would substantially relieve, as Section~\ref{sec:abm-danish}'s counterfactual suggests.
\clearpage
\section{Discussion}\label{sec:discussion}

Sections~\ref{sec:abm} and~\ref{sec:hazard} presented the paper's central contrast. Household simulation, rational or behaviorally extended, does not close the gap to the empirical benchmark. A structurally distinct hazard framework---no household utility, a substantially larger non-synthetic sample---recovers 91.3\% of it in microsimulation form on the benchmark-consistent accounting basis at the headline off-window floor, and 97.9\% at the in-sample calibration point. Across that floor's range the recovery runs 88.2--93.7\%, a miss of 6.3 to 11.8 points. But no estimator's monthly path co-movement, Path B's included, survives removing the window's shared trend (Section~\ref{sec:pathb}). The contrast is easy to over-read. The $\beta_1 = 0$ null recovers 85.7\% on the same basis at that headline floor (88.7\% in-sample), so components common to any prepayment model---scheduled amortization, the seasoning baseline, the turnover floor---deliver the large majority of the level, and the caps sat far above what any plausible prepayment environment would have produced. Named loan-level mechanisms are measured small wherever I can ablate them. Switching burnout off moves Path B by \$6.58 billion ($0.86$ points) at the in-sample floor and \$3.48 billion ($0.45$ points) at the headline off-window floor. Zeroing the FICO/LTV covariates moves it $1.3$ points, and the delinquency transition matrix is bounded at \$0.39 billion (Appendix~\ref{app:params}). No servicer channel beyond that pipeline is modeled at all. Three mechanisms therefore remain \emph{candidates} for the residual and for the timing failure every estimator shares, not demonstrated drivers of the level. That residual is modest at the in-sample calibration (2.1 points) but 8.7 points, roughly 1.6$\times$ the headline marginal, at the headline off-window calibration. Servicer behavior beyond the delinquency pipeline, updated-equity dynamics, and home-equity extraction carry no bound of any kind here, each plausibly rate-cycle-correlated. First, then, loan-level heterogeneity in credit score, geography, and life circumstance makes mobility no smooth function of coupon and market rate alone; the hazard framework's 295 stratum fixed effects are one mechanism letting Path A track pool-level heterogeneity the ABM's coupon-cohort buckets cannot represent. Second, servicer behavior---loss mitigation, forbearance, and modification practices that vary by servicer and were unusually active in the higher-rate, higher-delinquency-risk QT window---reaches prepayment and default outcomes outside any borrower-side decision model, which Path B's competing-risks structure and Markov delinquency pipeline represent directly if partially. Third, prepayment burnout is representable at the cohort level (as in both hazard paths) but not inside a closed, per-agent household population, as Section~\ref{sec:abm-results}'s falsified survivor-selection experiment showed.

\subsection{The Dual Systemic Fallout}\label{sec:discussion-fallout}

The macroeconomic literature often focuses on the housing mobility implications of the lock-in effect, but its transmission into the fixed-income sector creates a two-pronged systemic risk.

The first channel runs through the central bank itself. Extending MBS durations deprives the SOMA of its scheduled cash flows and breaks the predictability effective Quantitative Tightening requires; Appendix~\ref{sec:robustness-wal} states the extension in approximate weighted-average-life units. The \$764.7 billion shortfall-versus-cap is, per the decomposition above, mostly rate-inelastic: a gap to a cap schedule no plausible prepayment environment would have met. The rate-driven object is the duration extension itself---a 9.4-year portfolio WAL under the empirical path against 3.4 years at 2021 realized speeds (Appendix~\ref{sec:robustness-wal}), that 3.4-year comparator a refinancing-boom baseline. The extension, not the cap arithmetic, feeds the fragility channels below.

The second channel runs through private banks, which, unlike the Fed, must mark to market: the structural duration extension of these low-coupon MBS is one source of the unrecognized mark-to-market losses that accumulated across the banking sector. \citet{jiang2023} show duration extension alone was not sufficient to precipitate failure---it was mark-to-market losses interacting with a high share of uninsured deposits that made specific regional institutions vulnerable to self-fulfilling runs, uninsured depositors withdrawing once reported asset values no longer covered liabilities. The duration extension that is a necessary precondition for this fragility channel comes from the below-market coupon stock meeting rate normalization under any payoff rule: Section~\ref{sec:abm-danish}'s rule-only counterfactual shortens the portfolio by only 0.6 of the roughly six extension-years (Appendix~\ref{sec:robustness-wal}), and those six years are measured against 2021's refinancing-boom speeds, not normal turnover. Read instead against the baseline turnover floor, the empirical path's approximate weighted-average life is 9.4 years against 9.5 at the window open and 8.5 against 8.6 at its close, so the extension measures refinancing-versus-turnover, not lock-in-versus-normalcy. The elasticity channel I identify is therefore a second-order contributor to that precondition, and neither it nor the extension is, on its own, a sufficient condition for bank failure. On the institutional side I price cash-flow timing alone (principal arriving later than the cap schedule assumed), not a market-value loss. Appendix~\ref{sec:robustness-wal} declines option-adjusted duration and convexity as out of scope, so no mark-to-market or DV01 image is put on that 0.6 of an extension-year. The bank-side losses \citet{jiang2023} price have no counterpart here. The fragility channel is invoked as context, not as something this design measures.

Beyond fixed-income transmission, the lock-in effect imposes a real economic cost by keeping households from relocating toward higher-productivity labor markets. \citet{hsieh2019} estimate substantial aggregate output losses from such spatial misallocation of labor under housing supply constraints in high-productivity U.S. cities. Their headline magnitudes are cumulative 1964--2009 level effects revised in a published correction, so I cite the mechanism, not a point estimate. Theirs is a supply-side housing-availability constraint; the demand-side mobility constraint I quantify runs through the same channel, since households who cannot economically justify relocating cannot respond to geographic labor-demand differences, whatever the housing supply elsewhere. No per-household or per-move welfare estimate enters the paper, so the labor-reallocation leg of the mobility argument is directional only. The per-move surplus that would turn this leg into a magnitude is nowhere quantified. Nor do I translate the \$764.7 billion shortfall into welfare terms: it is an accounting gap against a cap schedule, mostly rate-inelastic by the decomposition above, and neither a lower nor an upper bound on the par-payoff rule's welfare cost. The welfare-relevant objects are the constrained moves, whose fiscal image is the lock-in marginal ($+\$42.6$ billion of trapped roll-off at the headline off-window floor, $+\$70.3$ billion at the in-sample calibration point), and the forgone labor-reallocation gains this literature identifies, which my accounting framework does not measure.

A back-of-envelope prices the extension in fiscal units. At the window's policy-rate levels the funding spread over the book's 2.49\% weighted-average coupon ran roughly 250--300 basis points, so the \$42.6 billion of trapped roll-off carries on the order of \$1 billion per year of negative carry on the consolidated public balance sheet for as long as the gap persists. That figure restates in remittances terms the maturity extension Table~\ref{tab:wal} prices in years. It derives from committed quantities---the book coupon, the marginal---and the public policy-rate path, and it is subject to the same face-versus-cash incidence reading as the buyback bracket (Appendix~\ref{app:params}).

\subsection{Redemption-Cap Design: The Ex-Ante Arithmetic}\label{sec:discussion-capdesign}

That the caps would sit above achievable runoff is not my observation. New York Fed staff projected \$20--\$30 billion in monthly principal payments when the \$35 billion ceiling took effect \citep{nyfed2022}, and their May 2022 baseline path implied agency MBS runoff of \$740.6 billion over the window against the caps' \$1{,}417.5 billion. What this paper adds is the quantification, constructed ex ante. The \$35 billion cap was set roughly 1.7 to 1.9 times the Federal Reserve's own contemporaneous projection of achievable runoff, the multiple depending on the same disclosed intra-2022 allocation choice as the anticipated share (Section~\ref{sec:method-benchmark}). From the other side, both components of the $\beta_1 = 0$ null recover 88.7\% of the realized shortfall at the in-sample calibration point and 85.7\% at the off-window floor under the production floor form (35.6\% under the additive form). That is how far the fixed \$35 billion cap sat above the book's rate-inelastic principal path. Those two components are scheduled amortization from the book's coupon, age, and term composition, and baseline turnover at any assumed floor. Table~\ref{tab:wal} prices that path in maturity units. Under the empirical CPR the book's approximate weighted-average life is 9.4 years against 3.4 years at 2021 realized speeds, and that 3.4-year comparator is a refinancing-boom baseline, not a turnover norm. So the speed that would have filled the cap was the one the window's rate configuration had already shut off.

\begin{table}[H]
\centering
\small
\begin{threeparttable}
\caption{Approximate SOMA portfolio WAL by prepayment scenario (years).}
\label{tab:wal}
\begin{tabular}{@{}lrr@{}}
\toprule
Scenario (mean CPR) & WAL, June 2022 & WAL, Nov.\ 2025 \\
\midrule
Scheduled amortization only (0\%) & 14.7 & 12.6 \\
Empirical path, ABM-basis back-out (5.14\%) & 9.4 & 8.5 \\
Empirical path, note-rate WAC basis (5.79\%) & 8.9 & 8.1 \\
Production ABM (11.68\%) & 6.0 & 5.6 \\
Hazard Path A (3.34\%) & 10.9 & 9.7 \\
Hazard Path B (4.76\%) & 9.7 & 8.7 \\
Danish rule-only leg, production anchor (5.61\%) & 9.1 & 8.2 \\
Danish-level bracketing leg (3.39\%) & 10.8 & 9.6 \\
Baseline turnover floor, off-window headline read (4.99\%) & 9.5 & 8.6 \\
No-shock counterfactual, 2021 realized speeds (22.81\%) & 3.4 & 3.3 \\
\bottomrule
\end{tabular}
\begin{tablenotes}[flushleft]\footnotesize
\item Single-pool-per-cell approximation: vintage face shares from the SOMA CUSIP tabulation (2021: 43.9\%; 2020: 16.3\%; 2017--19: 6.0\%; 2022: 23.1\%; pre-2017: 10.6\%), a 90.7\%/9.1\% 30-year/15-year term split, coupon at the book's 2.49\% WAC, and each scenario's window-mean CPR held constant. Level-payment amortization with monthly SMM; November 2025 values age the same cells 42 months. Cell WAL formula: for remaining term $n$ months, $\mathrm{WAL} = \sum_{t=1}^{n} t\, p_t / (12 \sum_{t} p_t)$, with $p_t$ month-$t$ principal (scheduled level-payment principal plus $\mathrm{SMM} = 1-(1-\mathrm{CPR})^{1/12}$ of the surviving balance). Worked anchors for the scheduled-only row: a 2.49\% 30-year pool has origination WAL 16.9 years (the 2021-vintage cell, aged 12 months at June 2022, is at 16.3); the 9.1\% 15-year sleeve (origination WAL 8.0 years) and the older vintages (32.9\% of face is 2020 or earlier) pull the book blend to 14.7. Representative cell ages ($\{2022{:}\,2;\ 2021{:}\,12;\ 2020{:}\,32;\ 2017\text{--}19{:}\,60;\ \text{pre-}2017{:}\,120\}$ months at June 2022) are single-pool stand-ins recovered by matching the printed table (a reconstruction of the original computation's assumptions, disclosed as such), and anchored externally: ages derived from the SOMA CUSIP tabulation's own back-derived origination dates (maturity minus term, face-weighted within cells, computed with no reference to this table) run $\{2021{:}\,12;\ 2020{:}\,28;\ 2017\text{--}19{:}\,50\}$ months and give a scheduled-only blend of 14.9 years; naive vintage midpoints likewise give 14.9. The printed 14.7 sits within 0.3 years of both alternatives. Coverage reconciliation: the 90.7\%/9.1\% split (which sums to the 99.8\% covered) is the audit parse of the \$1{,}940.8 billion book (the 2026-07-15 as-of of the same object whose committed June-2022 as-of total is \$2{,}700.6 billion; the difference is the interval's runoff, not a discrepancy), with 0.2\% of face in other terms never modeled; the 90.6\% of Appendix~\ref{sec:robustness-pipeline} is the production run's own as-of-date 30-year-only coverage share; the two parses differ by as-of date and their components round independently.
\item Accounting basis of each row. The two empirical rows are the same SOMA series backed out under two coupon conventions, and the basis is named because it moves the row: 5.14\% is the ABM-basis back-out, 5.79\% the note-rate WAC --- the corrected amortization basis, since the note rate is what amortizes the loan --- and the hazard-basis back-out at 5.52\% sits between them at 9.1/8.3 years. Moving from the ABM basis to the note-rate basis shortens the empirical path by 0.5 years at the window's open and 0.4 at its close (run \texttt{wal\_note\_rate\_basis}, through the identical calculator), which is the same order as the 0.6-year Danish rule-only effect this table also reports. Every other row is a \emph{simulated} mean CPR from its own leg and carries no back-out convention. The 6.0-year extension is deliberately \emph{not} restated on the note-rate basis: the no-shock row's 22.81\% is a 2021 back-out on the ABM basis with no committed note-rate counterpart, so differencing across the two would mix conventions inside one statistic.
\item Turnover benchmark: the final row is a refinancing-boom baseline rather than a normal-turnover one --- 22.81\% is the 2021 realized speed, most of which is refinancing, and refinancing is what the QT window's rate configuration shut off. The turnover-floor row is the normal-turnover counterpart. The paper's own out-of-window turnover measurement is the floor read (Table~\ref{tab:oosfloor}): defensible 2018 reads of 4.70\%, 4.99\% and 5.33\% by depth of out-of-the-moneyness, against refi-contaminated pooled 2017--2019 and 2019 anchors at 6.07\% and 6.91\%. The turnover-floor row prices the middle read --- the paper's headline floor --- through the identical calculator (run \texttt{wal\_normal\_turnover}): it falls between Path B's 4.76\% and the empirical 5.14\% rows, and monotonicity in CPR puts the band ends on either side of it.
\end{tablenotes}
\end{threeparttable}
\end{table}

The design question that follows is speculative: design space, not a finding these estimates support. Whether this episode generalizes is a question this design cannot answer but can pose precisely. \citet{du2024}'s cross-country record covers seven central banks running active and passive runoff, and the mechanism here is not specific to the Federal Reserve: a fixed nominal cap above the rate-inelastic principal path of a book whose coupons sit below market. What would settle it is narrow and public: each bank's announced cap or redemption schedule against its realized redemptions, month by month, over its own tightening window. I do not assert the other six were non-binding; I record that the comparison is constructible from published data and is not made here. A substitution instrument is a standing policy of substituting Treasury runoff or outright MBS sales when passive runoff undershoots. Pre-committing to one would convert extension risk from a silent shortfall into an explicit, priced policy choice, and leaves the cap free to communicate the path, not force it. Indexing the cap to the observable state of the outstanding stock (the share of the book sitting deeply below market coupon) is the weaker alternative: it still works through the instrument the shortfall defeats. What the decomposition supplies is computable in advance of the cycle, as a band: scheduled amortization is fixed by the book's composition; the turnover floor is measured, its defensible out-of-window reads running from 4.70\% to 5.33\% with the depth at which out-of-the-moneyness is cut, carrying the central-elasticity marginal from $+6.8$ points at the deepest cut to $+4.3$ at the shallowest (Table~\ref{tab:oosfloor}). The floor read's own sampling error widens the identified margin further: in the units a cap designer would have to plug in, the mid-cut read's binding wild-cluster inversion runs from 4.033\% to 6.181\% (the Webb rung's narrower read, 4.177\% to 5.800\%, is the one the cap grid below was run at) (Section~\ref{sec:robustness-floor}). On my committed projections, over the 42 window months the phased ceilings averaged \$33.75 billion per month of allowable MBS redemption against the passive principal the book could deliver, \$17.6 billion per month under the uniform-spread allocation and \$20.0 billion per month under the settlement-aware one. The two comparisons give ratios of 1.91 and 1.69, so the ceiling was near twice a rate the book could not passively reach under either allocation, which is what it means to say these caps never bound. Propagated through the committed floor-to-marginal grid, that sampling error prices the lock-in marginal at \$0.42 to \$1.66 billion per month of foregone passive principal. That band attaches to the identified marginal, not to the two projection levels, which are accounting constructions with no sampling interval of their own. The indexing proposal can now be priced: against the window-minimum market rate, the share of book face sitting below market is 99.9\% at a zero out-of-the-moneyness cut and 99.6\% at cuts of 25 and 50 basis points, the last two coinciding because the book's coupons sit on a half-point grid. Carrying each cut's committed floor read through the identical amortization gives achievable passive principal, hence an implied cap, of \$17.07, \$16.47 and \$15.95 billion per month (run \texttt{state\_contingent\_cap\_grid}). The grid supplies the missing function and shows it is worth little: the whole depth range moves the implied cap by about \$1.1 billion per month, while the mid-cut floor read's own wild-cluster interval moves it by about \$2.9 billion, from \$15.02 to \$17.88 billion per month. The observable is near-degenerate, and the little it does carry is dominated by the measurement error in the floor it would be indexed to. One scope limit travels with it: the observable is the SOMA book's bucketed coupon distribution while the floor is estimated on the Freddie panel's loan-level gaps---two populations at two grains. What the cap is for bears on whether binding is desirable at all: a redemption cap paces and communicates the portfolio's composition path, and Section~\ref{sec:method-benchmark} concedes these caps never bound over this window, so a binding level trades the communication objective for a speed the schedule was not built to deliver. A cap intended to bind would sit at or below the lower edge of the book's projected scheduled-plus-turnover principal path, with the identified elasticity margin ($+2.1$ to $+13.2$ points of the shortfall across the calibration box) as the behavioral band around it. That is the form the Federal Reserve's own contemporaneous projection took: a range. Nothing in my estimates identifies the welfare or market effects of such designs; active sales in particular would realize the mark-to-market losses passive runoff defers. I record the option to mark where the design space lies, not to recommend a point in it.

Two boundaries belong beside this design lesson. First, non-bindingness may have been a choice, not an error: a Treasury-symmetric cap is simple to communicate and forecloses forced sales, and nothing in the public record establishes that the Committee intended the MBS cap to bind. Still, the FOMC's stated longer-run intention to hold primarily Treasury securities gives the MBS-specific shortfall a standing the cap alone does not: against that composition objective, slower runoff is delayed progress whatever the cap's intent. Second, the design space is priced in partial equilibrium: a cap set to bind, or active sales, would transmit to primary-market spreads and through them to the lock-in margin itself. Pricing that feedback is outside this paper's accounting.

\paragraph{What a designer of a future runoff program should and should not take.} Three quantities transfer: the baseline turnover floor in CPR units (the off-window read's binding sampling span is 4.033--6.181\%, Section~\ref{sec:robustness-floor}), the scheduled-amortization path (computable from the book at purchase time), and the elasticity band as the behavioral budget ($+2.1$ to $+13.2$ points of a cap-relative shortfall across the calibration box). Two numbers should not travel alone: the $+5.6$-point anchor, a convention inside a range rather than a central tendency, and the \$764.7 billion cap-shortfall, a statement about a never-binding schedule's composition rather than a missed objective. One accounting rule: face value answers the benchmark question, cash the reserve-drain question, and the two differ by the par windfall---a transfer, not a resource cost (Appendix~\ref{app:params}).

\subsection{Evaluating Alternative Securitization Policies}\label{sec:discussion-policy}
In transaction units the marginal is small against the market it operates in. Dividing the central-minus-null differential by the mean balance of the sample's surviving loans (\$237{,}317 across 40{,}077 loans) gives about 179{,}500 foregone payoffs on the SOMA book over the 42 window months, or 51{,}300 a year. Both figures are computed at the production $\mu = 1$ convention, which treats the whole voluntary gap response as the moving channel and therefore returns an \emph{upper bound}. The all-loan mean averages in prepaid loans at zero balance and would roughly double the count. Confining the response to a half or a quarter of the voluntary hazard gives 108{,}300 and 58{,}700 (run \texttt{marginal\_transaction\_counts}). Grossing to the whole market on a SOMA share of 20--22\% of 1--4 family mortgage debt puts the upper-bound count at 233{,}000--256{,}000 payoffs a year. That is 5.7--6.3\% of the 4.08 million annual existing-home-sale run rate. This is the first outcome-side check the design admits, and it does not flatter the estimate. Even at its upper bound the identified channel is a small minority of transaction volume. One comparison is deliberately not made. \citet{fonseca2026} document a 40\% fall in existing home sales between 2022 and 2024, but that series has no pre-2025 history on this project's data access. So the marginal's share of \emph{that decline} cannot be formed without assuming the 2024 level equals the 2025 one. I decline the assumption and report the run-rate share instead.

In response to the mobility collapse, U.S. policymakers revisited assumable and portable mortgages. In late 2025 the Federal Housing Finance Agency stated that it was actively evaluating portability, and that Fannie Mae and Freddie Mac---which permit fixed-rate assumption only in limited circumstances such as death or divorce---were examining wider assumption. No rulemaking had issued as of writing, so I cite this as the policy setting, not an enacted rule, and what follows as a taxonomy of that design space, not a finding. Section~\ref{sec:intro}'s concession applies to any evaluation of it. The statutory right already exists for the 20.4\% of SOMA face that is Ginnie Mae collateral. Realized take-up, though, is measured nowhere in this design and in no public assumption-volume series I could find, so the 20.4\% share is used throughout as a ceiling on the carve-out, not an estimate of it. Fixed-income analysis suggests both address only part of the problem: each eases the consumer mobility constraint while explicitly preventing prepayment upon relocation. Under assumption the below-market coupon transfers to the buyer and the extension stays on the investor's book, while the buyer must separately finance the gap between purchase price and assumed balance, typically at prevailing rates, which limits take-up. Both instruments therefore relocate duration extension risk, solving the microeconomic mobility constraint while cementing the macroeconomic extension risk onto the investor's balance sheet.

\subsection{The Danish Market-Value-Repurchase Alternative}\label{sec:abm-danish}

\subsubsection{The transplant}

A third alternative acts on the payoff rule itself, not on portability or the redemption cap. Section~\ref{sec:abm-results} shows that no synthetic-population household-choice specification tested here explains the majority of the cap shortfall. Section~\ref{sec:robustness-crossdesign} confines that finding to the synthetic population: its recalibrated real-covariate variant crossed the pre-committed threshold past which the paradigm reading cannot be cleanly maintained. The question is therefore how much of that trap the par-payoff rule itself causes, rather than household behavior under any institutional regime. The Danish market-value repurchase system severs the link between market rates and payoff cost, which lets me ask that quantitatively. The operative specification calibrates each framework's Danish response against \citet{berger2026}'s separately estimated moving and refinance-in-place channels, then transplants them onto the U.S.\ loan population so that only the payoff rule changes. \citet{berger2026} is a July 2026 SSRN working draft, and the single unrefereed source this channel rests on. The paper's finding on this question is an institutional gap of $+\$61.2$ billion (Table~\ref{tab:danish}).\footnote{\label{fn:manifest}CPR means below aggregate the 42 active QT months (June 2022--November 2025), and the institutional wedge is those months' mean cohort-weighted Danish-minus-U.S.\ simulated CPR difference. This operative specification supersedes an initial mechanism-substitution that had implied a much larger gap; that superseded figure and the frozen run manifests underlying every ABM headline are catalogued in the replication package's ledger.}

\subsubsection{The two estimated channels}

The recalibration rests on two behaviorally distinct channels that \citet{berger2026}---my source for the Danish institutional mechanism itself---estimate separately and a single extrapolated mechanism would conflate. The moving channel is whether removing the par-payoff penalty changes how often households relocate. It is well-identified and close to flat in the coupon gap: Danish households move at an unconditional 3.2\% per year regardless of how far their coupon sits below market \citep[\S4.6, structural parameters in tab.~3]{berger2026}. Their Danish administrative microdata put that slope near zero and not consistently signed: a 100 bp \emph{decrease} in the yield gap moves the annual moving rate by about $+14$ basis points under their specification and about $-19$ under \citeauthor{fonseca2024}'s control set \citep[\S3.3.2]{berger2026}. The U.S. moving-hazard slope, by contrast, is 0.57--1.20 percentage points per year per 100 bp of coupon gap \citep{fonseca2024}. The refinance-in-place channel---whether a market-discount buyback induces discount-buying independent of moving---is substantially a Danish tax artifact: buybacks are attractive there because mortgage interest is deductible at a flat 33\% rate and realized capital gains on the buyback are tax-exempt. Under U.S. tax treatment---a 22\% mortgage-interest deduction and a 15\% capital-gains tax on the discount, per Berger et al.'s own counterfactual---the channel loses that shield. Their general-equilibrium counterfactual \citep[\S4.10.2]{berger2026} finds that introducing the option raises the equilibrium mortgage rate by about 20 basis points on average, 16--23 across their robustness variants. They read that as leaving equilibrium rates virtually unchanged. That estimate is itself draft-sensitive: the January 2026 version of the same paper reported about one basis point. That is one more reason this channel is carried as an unrefereed import rather than a settled magnitude. I read that characterization differently, though I do not adjudicate it. On the ordinary reading a discount extinguishment is cancellation-of-debt income at ordinary rates, IRC \S 61(a)(11) and Treas.\ Reg.\ \S 1.61-12(c), while \S 108(a)(1)(B) and \S 108(a)(1)(E) exclude it for insolvent borrowers and for qualified principal-residence indebtedness respectively, so the treatment is borrower-state-dependent, not uniform in either direction. Post-TCJA, the assumed interest deduction also overstates most households' marginal benefit. Read as a directional bracket, both considerations push the U.S.-transplant channel \emph{below} Berger et al.'s already-small estimate, and neither touches the realized Danish evidence, which is measured under Danish tax.

Recalibrating both frameworks against these two estimated channels first requires a choice about the moving channel's \emph{level}. The choice arises because Berger et al.'s identified fact about the moving margin is one of slope: Danish moving is flat in the coupon gap under market-value payoff. The 3.2\% per year, by contrast, is Denmark's descriptive unconditional moving rate, not an estimate of how often U.S. households would move under a Danish payoff rule. Importing that level onto the U.S. loan population---the \emph{Danish-level} anchor---therefore bundles two changes, the payoff rule and Danish baseline mobility itself. It produces Danish legs whose mean CPR (3.36--3.39\%, Table~\ref{tab:danish}) sits \emph{below} the 4\% turnover floor both frameworks impose on the U.S. side (Appendix~\ref{sec:method-frictions} and Section~\ref{sec:pathb}). Death, divorce, and forced relocation do not become rarer because the payoff rule changed, so a rule-only counterfactual on the same U.S. households cannot prepay below that floor. The construction that isolates the rule is the \emph{U.S.-intercept} anchor: keep the identified flatness, but anchor the level at the production U.S. hazard evaluated at zero rate gap---turnover floor, seasoning baseline, burnout, and covariate terms all retained, only the rate-gap suppression switched off---plus the same U.S.-transplant refinance-in-place channel. I adopt it as the production reading, the only one consistent with the floor's stated semantics. I retain the Danish-level anchor as a labeled bracketing case in which U.S. households are assumed to adopt Danish baseline mobility along with the Danish rule. Holding the floor fixed while importing only the Danish moving slope treats baseline turnover and rate-elastic moving as separable channels. That separability is motivated by the floor's semantics. It is not identified in the data available here.

\subsubsection{The anchor choice and the sign of the gap}

The anchor choice determines the sign of the institutional gap. Under the shared-accounting hybrid pipeline (Section~\ref{sec:robustness-hybrid}), the U.S.-intercept anchor yields a Danish mean CPR of 5.61\% against the locked-in U.S. leg's 4.76\% and a gap of $+\$61.2$ billion (8.0\% of the benchmark) at Berger et al.'s general-equilibrium best estimate of negligible refinance-in-place. The U.S. leg reproduces the committed production artifacts exactly. That estimate is a price statement (an equilibrium-rate effect of about 20 basis points), not a near-zero buyback \emph{volume}. Realized Danish data contradict the near-zero volume: over 2022M07--2023M12 deep-discount ($\leq 2\%$ coupon) Danish callable bonds redeemed at 26.4\% per year against a 4--5\% baseline. The extraordinary component is at least 22\% per year after a generous scheduled-amortization allowance (run \texttt{danish\_redemption\_validation}; Danmarks Nationalbank securities statistics, exact stock--transaction closure on the coupon-blocked universe). The transplant imports the payoff rule alone: not the Balance Principle's series-level match funding, the large tap-issued series carrying Denmark's liquidity, the advisory distribution prompting conversions, or the Danish product mix. Each is a held-fixed feature of the U.S. leg, not a measured invariance, and one --- the large interest-only share --- now carries a magnitude. I therefore report the rule-only gap as a band over the refinance-in-place contribution, not a point at zero. Swept in 0.25\%-CPR steps under the production anchor, the gap runs from $+\$61.2$ billion at 0\% to $+\$256.8$ billion at 3\% (run \texttt{danish\_refi\_finegrid}; slope $\approx \$65$ billion per CPR point). The large interest-only share cuts scheduled amortization, the null's largest component, and so moves the shortfall against the payoff rule. What run \texttt{danish\_interest\_only\_share} adds is the crossing, since the direction was already conceded here. Zeroing scheduled amortization on a share $\chi$ of the Danish leg alone cuts the face gap by $\chi$ times \$198.67 billion. That figure is the U.S. leg's scheduled path over the window and this paper's own proxy for the scheduled component. I zero it on the Danish leg alone because scoring it in both legs would neutralize it exactly, and the U.S. leg's effective zero stays held fixed. The gap crosses zero at an interest-only share of 30.8\%. Deferred principal survives to prepay. Carrying it forward at the Danish leg's own 5.61\% speed returns \$9.1 billion of extra early roll-off at a 45\% share and lifts the crossing to 34.3\%; the offset is reported beside the first-order effect rather than folded into it, and it was the one limb I could not sign in advance. Denmark's share sits above both crossings. Danmarks Nationalbank puts interest-only loans at 45\% of outstanding mortgage volumes \citep{nationalbanken2020}, at which the face gap is $-\$19.2$ billion after that feedback. Three limits travel with the 45\% wherever it is read. It is pre-window (February 2020 against a June 2022--November 2025 window). Danish ``interest-only'' is deferred amortisation (\emph{afdragsfrihed}), conventionally capped at ten years, so it is the share of volume currently \emph{in} an interest-only period and not a permanent product feature. Its population is whole-market outstanding volumes rather than the owner-occupied thirty-year callable segment this transplant models. The crossing carries none of them: it is arithmetic on the committed chain and needs no Danish input, which is why I lead with it. This is an accounting-layer calculation: the hazard is untouched and only the amortization path moves. It establishes that the rule-only gap survives the product mix only below that crossing, not what a re-simulated interest-only book would return. The same leg reads in duration units. At each leg's mean CPR through the identical calculator, the rule-only transplant shortens the book from Path B's 9.7 years to 9.1 at the window's open and from 8.7 to 8.2 at its close --- roughly 0.6 years. The Danish-level bracketing leg runs the other way, to 10.8 and 9.6 (Table~\ref{tab:wal}). Duration separates the two anchors in the same direction the dollars do, without resting on the incidence choice the dollar metric cannot settle. The realized-implied anchor prices at $+\$1{,}212$ billion, beyond the rejected partial-equilibrium ceiling. Four confounds make it an upper anchor on the transplant, not an estimate of it. The realized rate is gross, so conversions bundle in. It includes non-household collateral. It is matched to the environment, not to the book. And it is a Danish-tax realization of an option whose U.S.-transplant value the tax attenuation is meant to shrink. Both tax regimes read off this one axis. That is worth naming because the sweep parameterizes refinance-in-place CPR, not a tax code: a regime enters only through the refinancing it induces. The U.S.-tax regime sits at \citeauthor{berger2026}'s general-equilibrium best estimate of $\approx$0\% refinance-in-place and gives $+\$61.2$ billion. The Danish-tax realization sits at the realized-implied 22.4\% per year and gives $+\$1{,}212$ billion, above the swept range entirely. They are not equally credentialed (the second carries those four confounds), but quoting only the first would present a tax-attenuated anchor as though it were the institutional gap itself. This returns to a construction the pipeline already contained. Leg for leg, the U.S.-intercept hybrid coincides numerically with the NPV-identity hybrid the replication package's technical documentation records from before the Berger recalibration. That is the buyback-at-market-value identity zeroing the effective rate gap --- computationally the zero-gap U.S. hazard. What it lacked is empirical evidence that Danish-rule moving is flat in the gap. Berger et al.'s moving-flatness estimate supplies it. The sign is invariant under this anchor, and it could not have come out otherwise. The specification defines the Danish moving hazard as the production U.S. hazard at a zero rate gap. The U.S. book's gap is negative in essentially every cell over essentially the whole window (99.53\% of exposure below the window-minimum market rate, Section~\ref{sec:identification}). The hazard is monotone in it. The zero-gap leg must therefore prepay faster, the Danish trapped balance come in lower (\$687.8 billion against \$749.0 billion), the gap be positive. Sweeping refinance-in-place from 0\% to the rejected partial-equilibrium ceiling ($\approx$18\% CPR) only adds Danish prepayment, so the sweep's sign-robustness is forced by construction too. The gap is positive at every point, rising to $+\$1{,}038.9$ billion at the ceiling (Figure~\ref{fig:gapsweep}). I retain both as wiring checks and read neither as evidence. That premise is an inequality about this episode, not a property of the design. The book is closed at 2017--2021 originations, and the window's market rate stood above essentially all of its coupons in essentially every month. Where a rate gap is mixed in sign (another cycle, a book open to new originations at the prevailing rate), it fails. The sign then becomes a composition question this design does not answer, the elasticity being imported, not estimated. Neither this sign nor the marginal transports without re-checking it. What is not forced is the magnitude. \$61.2 billion, 8.0\% of the benchmark, is what the recalibration actually delivers, and the Danish-level anchor below shows that a different and defensible transplant choice reverses the sign outright. This positive gap is the lock-in marginal of Section~\ref{sec:hazard-interp} from the counterfactual side. The $\beta_1 = 0$ null switches off the same rate-gap suppression, so the rule-only Danish counterfactual and the paper's lock-in contribution at the in-sample calibration point are one object on two accounting legs. The legs read $+\$61.2$ billion here, where the Danish leg compounds its own counterfactual balances, and $+\$70.3$ billion as the U.S.-leg marginal. At the 0\% edge of the swept band they nearly coincide, a property of that edge and of the in-sample calibration, not of the counterfactual. On that second leg the Danish magnitude inherits every conditionality the marginal carries --- the baseline seasoning convention and the floor's age transport among them. The transplant therefore supplies no independent information about the payoff rule: it re-denominates a quantity this paper has already conditioned, conditional in exactly the same ways. The near-equality is about the demoted in-sample point, not the paper's headline. Against the $+\$42.6$ billion off-window marginal the Danish rule-only gap is about 44\% larger, not equal, and both legs were run at the production floor. Re-run at the off-window floor itself, it is $+\$28.2$ billion, 3.7\% of the benchmark, 34\% below the off-window marginal (run \texttt{danish\_offwindow\_floor}; the U.S. leg reproduces the committed off-window central exactly). No equality survives at matched floors. Under the Danish-level anchor, by contrast, the gap is $-\$99.9$ billion (U.S. 4.76\%, Danish 3.39\% mean CPR), crossing zero at a 1.4\% CPR refinance-in-place contribution. The fragility reflects the bundled mobility-culture change. The ABM's recalibrated run under that anchor shows $-\$728.4$ billion (crossing at 12.6\%). Its U.S. leg, however, is over-predicted against the empirical benchmark (11.76\% simulated mean CPR under the Berger run's own calibration-stage freeze versus 5.14\% empirical (ABM basis); cf.\ the production tag's 11.68\%, Section~\ref{sec:abm-results}). Under the U.S.-intercept anchor its direction follows a fortiori without a run. Its lock-in penalty only ever suppresses moving, so its zero-gap intercept exceeds its locked-in U.S. leg month by month and its rule-only gap is positive by construction. Table~\ref{tab:danish} collects all four specifications.

\begin{figure}[H]
\centering
\includegraphics[width=\textwidth]{fig4_institutional_gap_sensitivity}
\caption{Sensitivity of the institutional gap to the U.S.-transplant refinance-in-place contribution, swept from 0\% to the rejected partial-equilibrium ceiling ($\approx$18\% CPR), under both transplant anchors for the hybrid pipeline. Under the production U.S.-intercept anchor (rule-only counterfactual) the gap is $+\$61.2$ billion at the Berger et al.\ general-equilibrium best estimate ($\approx$0, marked) --- one anchor among the swept range, not a bound: the fine grid of run \texttt{danish\_refi\_finegrid} spans $+\$61.2$ to $+\$256.8$ billion over 0--3\% refinance-in-place CPR, and the realized-implied Danish anchor (22.4\% per year) lies above the plotted ceiling --- and positive at every sweep point; that positivity is forced by the zero-gap anchor and cannot be overturned by this sweep, since additional refinance-in-place only raises Danish prepayment further. Under the bracketing Danish-level anchor (rule plus Danish baseline mobility) the gap is $-\$99.9$ billion at the best estimate and crosses zero at a contribution of 1.4\% CPR. The anchor, not the refinance channel, decides the sign.}
\label{fig:gapsweep}
{\footnotesize\noindent Notes: The order-of-magnitude reduction against the superseded \$925.5 billion extrapolation is a property of the Berger et al.\ best estimate, not of the sweep. At the best estimate both anchors sit an order of magnitude below that figure. But the gap rises steeply in the refinance contribution, and at the rejected partial-equilibrium ceiling both anchors exceed it ($+\$1{,}038.9$ billion under the U.S.-intercept anchor, $+\$972.0$ billion under the Danish-level anchor). That ceiling is rejected on partial-equilibrium grounds and is not reached by any calibration this paper adopts. The upper half of the sweep should not be read as a quantitative statement at all: at refinance contributions of 12\% and above the Danish leg's trapped balance turns negative ($-\$14.4$, $-\$158.0$ and $-\$290.0$ billion at 12, 15 and 18\%), which is not an economically meaningful quantity. Those points are retained only to show the direction of the sweep.}
\end{figure}

\begin{table}[H]
\centering
\small
\begin{threeparttable}
\caption{The Danish market-value-repurchase counterfactual under four specifications.}
\label{tab:danish}
\begin{tabular}{@{}p{3.4cm}rrr>{\raggedright\arraybackslash}p{4.0cm}@{}}
\toprule
Specification & U.S. leg & Danish leg & Gap (US$-$DK) & Mean CPR US/DK \\
\midrule
Mechanism-extrapolated ABM\tnote{a} & $+\$91.0$B & $-\$834.5$B & $+\$925.5$B & 11.68\% / 47.14\% \\
Berger ABM, Danish-level\tnote{b} & $+\$84.5$B & $+\$812.9$B & $-\$728.4$B & 11.76\% / 3.36\% \\
Berger hybrid, Danish-level\tnote{c} & $+\$749.0$B & $+\$848.9$B & $-\$99.9$B & 4.76\% / 3.39\% \\
Berger hybrid, U.S.-intercept\tnote{d} & $+\$749.0$B & $+\$687.8$B & $+\$61.2$B & 4.76\% / 5.61\% (band $+\$61.2$--$+\$256.8$B over 0--3\% refinance-in-place; run \texttt{danish\_refi\_finegrid}) \\
\bottomrule
\end{tabular}
\end{threeparttable}
\end{table}

\noindent{\footnotesize \emph{Notes to Table~\ref{tab:danish}.} Sign convention: each leg is signed trapped liquidity (the net shortfall of simulated roll-off relative to the phased redemption cap over the QT window), so positive values mean runoff fell short of the cap (liquidity trapped) and negative values mean simulated runoff exceeded the phased cap path. Gap $=$ U.S. leg $-$ Danish leg; a negative gap means the Danish regime traps more than the U.S. regime. \quad a Superseded: applies the U.S.-calibrated mobility function to a Danish-style rate gap it was never estimated to describe (see text). \quad b The Berger recalibration's U.S. leg is re-simulated under that run's own calibration-stage freeze (native 15-year gate), not the \$91.0 billion production headline, so both legs and the gap share one frozen calibration. The Danish leg prepays more slowly under the tax-adjusted elasticities, trapping more than the U.S. leg. Danish-level anchor: imports Denmark's descriptive 3.2\%/yr moving level, a rule-plus-country bundle whose Danish mean CPR sits below the 4\% turnover floor (see text). \quad c Bracketing case under the shared-accounting hybrid pipeline (Section~\ref{sec:robustness-hybrid}). Both legs are printed from the shared-layer scoring artifact at the in-sample calibration point, the floor the Danish legs were run at (U.S. trapped \$749.0 billion = 97.9\% of benchmark; Danish trapped \$848.9 billion); neither leg was re-run at the off-window floor. Its sign is not robust: the sweep of the refinance-in-place input crosses zero at a 1.4\% CPR contribution. \quad d The paper's production estimate (rule-only counterfactual): Berger et al.'s moving-flatness fact anchored at the production U.S. hazard's zero-gap intercept, turnover floor retained, refinance-in-place at the general-equilibrium best estimate ($\approx$0). The U.S. leg reproduces the committed production artifacts exactly. The gap is positive at every point of the refinance-in-place sweep (Figure~\ref{fig:gapsweep}), which under this anchor is forced by construction rather than found: see Section~\ref{sec:abm-danish}. This row is face-denominated and ex-post on the U.S. coupon stack. Under the cash-haircut reading its gap reverses sign, and the wedge that reverses it is a transfer from the SOMA portfolio to moving households rather than a resource cost (Appendix~\ref{app:params}). No ex-ante repricing of the repurchase option into the coupon is netted from either leg, an omission signed against the reported relief. And the row holds the Danish product mix fixed, which is not a neutral choice: above a Danish interest-only share of 34.3\% the face gap reverses too, and Denmark's own whole-market, pre-window read is higher (see text).\par}

The mechanism-extrapolated specification this one corrects had implied a far larger gap; the replication package's ledger carries both the correction's derivation and the superseded figure's entry.

\section{Robustness Checks}\label{sec:robustness}

This section collects the methodological checks referenced throughout the paper. Two of them changed reported results. The recalibrated cross-design variant of Section~\ref{sec:robustness-crossdesign} exceeds my ex-ante threshold and qualifies the household-choice reading at the aggregate level. The off-window floor re-measurement of Section~\ref{sec:robustness-floor} is where the paper's headline marginal and level are actually derived ($+9.2 \to +5.6$ points; 97.9\% $\to$ 91.3\% shared), demoting the in-window calibration to an in-sample point. On the shared basis the paper quotes throughout, the hazard framework undershoots the benchmark at the headline calibration (91.3\%); it overshoots only on the standalone basis (107.0--121.5\%). The remaining checks confirm the central decomposition.

Four of this section's checks are mechanics rather than results, and the online appendix collects them: the behavioral-extension mechanics (Appendix~\ref{sec:robustness-extensions}), the weighted-average-life quantification (Appendix~\ref{sec:robustness-wal}), the fifteen-year fold-in and the rate-gap-units correction (Appendix~\ref{sec:robustness-pipeline}), and the seasonal-floor program (Appendix~\ref{sec:robustness-seasonalfloor}). The calibration box already carries the seasonal-floor verdict: what a seasonally varying floor moves is floor dispersion under the hard maximum, not cyclicality. Monthly timing earns no credential there---the frozen timing rule is cleared both by a wrong-month rotation and by the $\beta_1 = 0$ null itself.

\subsection{Scale Sensitivity}\label{sec:robustness-scale}

The ABM's 10,000-household population and the hazard framework's 75,000-loan sample differ in both scale and data source. To test whether population size alone drives the gap between the two frameworks, I re-ran the ABM's core household mechanic at populations of $N = 10{,}000$ and $75{,}000$, 50 seeds each. That mechanic is the production code's exponential, dollar-scale mobility-desire heterogeneity mapped through a rate-gap-driven CPR function. Mean simulated CPR is 15.616\% at both, a difference of 0.0003 percentage points. The 95\% confidence interval narrows 2.4$\times$ at the larger $N$, about 12\% short of the 2.74$\times$ an idealized $\sqrt{N}$ law predicts ($\sqrt{7.5} \approx 2.74$), consistent with a non-i.i.d.\ component of seed noise; I do not lean on the exact ratio. Population size changes that mechanic's precision, not its central tendency. Two limits, both large, so I state the conclusion conditionally. The check ran the household cost-benefit mechanic \emph{in isolation}, its 15.616\% sitting 3.9 points (about a third) above the production ABM's 11.68\%, so the object shown scale-invariant is not the estimator whose recovery is at issue. And no committed script, artifact or gate stood behind it, the figure surviving from an earlier draft. A pre-committed rerun on the full production behavioral specification addresses both (run \texttt{production\_scale\_test}, fifty seeds per size). It carries the 43\% front-end DTI wall, the discrete twenty-percent wait-and-see freeze trait, income-scaled curtailment, eleven-cohort multi-vintage weighting with its fifteen- and thirty-year term split, and a settlement-lag kernel --- none of them present in the isolated mechanic. It recovers 12.66\% of the benchmark at $N = 10{,}000$ against 12.18\% at $N = 75{,}000$ ($0.48$ points against a fifty-seed standard deviation of $3.25$ points, well inside one standard deviation), at mean simulated CPRs of 11.60\% and 11.64\%, close to the native-gate run's 11.76\%, not the isolated mechanic's 15.6\%. One dimension of the production specification this run does not carry, and it is the one that reconciles it with Table~\ref{tab:estimators}: the rerun evaluates the fifteen-year book on the native behavioral gate, where production folds that book in on structural terms only (Appendix~\ref{sec:robustness-pipeline}). Table~\ref{tab:estimators}'s 13.56\% fifty-seed mean is the production fold-in; this arm's 12.66\% is the native gate; the 0.90 points between them are that documented level offset, which the run was specified to expose, and not an effect of population size. Both arms here sit on the same gate, so the scale comparison itself is like-for-like. Seed precision improves with population, the standard deviation falling from $3.25$ to $1.05$ points (a ratio of 3.1 against the same $\sqrt{7.5} \approx 2.74$), and again I do not lean on the exact ratio. Population size therefore does not drive the gap between this estimator and the hazard framework, and that now rests on the production specification, not a fragment of it. It does not follow that the gap is modeling paradigm alone: resampling cannot close that distinction, and Section~\ref{sec:robustness-crossdesign} takes it up with the completed cross-design pair.

\subsection{Benchmark Sensitivity to Data Source and Cap-Schedule Assumptions}\label{sec:robustness-benchmark}

The \$764.7 billion empirical benchmark (Section~\ref{sec:method-benchmark}) is constructed from the New York Fed's SOMA current-face-value series under the actual phased redemption cap schedule, netting realized roll-off against the cap within the active QT window. I treat net accumulation, not one-sided accumulation of shortfall months only, as the correct construction because it is the only version comparable to the Danish counterfactual of Section~\ref{sec:abm-danish}, where simulated roll-off routinely exceeds the cap; a one-sided construction would overstate the empirical benchmark relative to that counterfactual and bias the recovered percentages reported for both estimators. Researchers extending this work with alternative SOMA vintages or a different cap-schedule assumption should expect the benchmark level, and therefore the recovered percentages in Sections~\ref{sec:abm} and~\ref{sec:hazard}, to shift. The qualitative ranking of the three estimators reported in Table~\ref{tab:estimators} is not sensitive to this choice, since the ABM's shortfall against any plausible benchmark construction remains an order of magnitude smaller than either hazard path's. Section~\ref{sec:method-benchmark} separately addresses the conceptual status of the cap-relative construction itself. One such alternative is now measured. The benchmark nets cash arriving in month $t$ against the cap in force in month $t$, which is what the cap constrains: the Committee set a ceiling on principal not reinvested in a given month, and that is a statement about cash. Agency remittance, however, means the cash arriving in month $t$ was generated by prepayment decisions in months $t$, $t-1$ and $t-2$, on the settlement kernel this paper already applies to simulated roll-off (0.10, 0.60, 0.30 at lags zero through two). Agency remittance is itself a deterministic lag of roughly one month, so that profile is not a description of the remittance cycle: it is an empirical smearing kernel over the months in which a given month's prepayment decisions arrive as cash, and the alternate kernels below price the choice. Convolving the cap schedule forward with that same kernel and re-netting the unchanged realized series gives a settlement-aligned benchmark of \$722.7 billion against the calendar \$764.7 billion, a shift of $-5.5$\% (run \texttt{settlement\_months\_benchmark}). The shift is mechanical, and its size is exactly accounted for. Because the QT window is half-open, cap convolved past its final month is dropped: October 2025 loses its two-month tap and November 2025 its one- and two-month taps, \$42.0 billion of the \$1{,}417.5 billion schedule. The realized series is untouched and both legs sum over the same months, so the benchmark moves by exactly that \$42.0 billion. Two alternate kernels bracketing that profile put the shift at $-\$28.0$ and $-\$52.5$ billion. Aligning the realized leg as well moves the benchmark less, not more. The pre-QT cap is zero, so the cap leg loses mass only at the window's end, whereas realized roll-off is not zero before the window and the two months the kernel reaches back into settle into its start: the realized leg gains \$6.8 billion there against \$23.4 billion leaving at the end, giving back \$16.6 billion of the cap-side \$42.0 billion. Convolving both legs therefore lands the benchmark at \$739.4 billion, a shift of $-\$25.4$ billion or $-3.3$\%, which the same two alternate kernels bracket at $-\$19.0$ and $-\$29.2$ billion (run \texttt{realized\_boundary\_allocation}). That the two edge effects would partially offset was fixed before the run; the sign of the residual deliberately was not, and at 3.3\% of the committed benchmark the boundary treatment stays a disclosed sensitivity rather than a re-basing: the production convention remains calendar on both legs. I retain the calendar convention, for the reason Section~\ref{sec:method-benchmark} gives for retaining the cap at all: it is the only policy-committed quantity available, and it is committed in cash-arrival months. What the cap-only variant establishes is that re-basing moves no dollar quantity in this paper (the identified marginal is \$42.6 billion on either denominator) and moves the percentages only by the common factor $1.058$: the central Path~B level reads 96.6\%, not 91.3\%, the $\beta_1 = 0$ null 90.7\%, not 85.7\%, and the marginal $+5.9$ points, not $+5.6$. The decomposition is unchanged in substance, and the null still recovers the large majority of the benchmark under either convention. The cap-only variant also bears on the path-level timing question every estimator here leaves open. By a cash-arrival object I mean reported roll-off smeared over the months in which a month's prepayment decisions arrive as cash and, in its CPR form, clipped by a back-out whose non-negativity rule pins four months to exactly zero and posts the displaced mass into the following month's reading (Section~\ref{sec:method-benchmark}). If the monthly series each estimator's path is scored against is such an object, then a monthly path is not a behavioral object, and the timing question may be ill-posed. I state that as a candidate account, not a finding: reporting mechanics and behavior arrive in the same series, and nothing here separates them. A related window-construction pitfall merits a reproducibility note: accumulating from the QT start date without an upper bound admits eight post-QT months (December 2025 onward). Where the dollar aggregation already applies the bounded active-QT mask, as it does throughout this paper, only the path statistics drift (the headline CPR means by roughly 0.3--0.5 percentage points); where it does not, the benchmark level itself moves. That asymmetry (path statistics silently sensitive to a window filter the masked dollar aggregates are robust to) is why every headline figure in this paper cites a single tagged run manifest (footnote~\ref{fn:manifest}).

\subsection{Cross-Design Test: Real Covariates in the ABM}\label{sec:robustness-crossdesign}

Section~\ref{sec:method-estimators} flags an unresolved confound: the two estimators differ simultaneously in modeling paradigm (household choice versus loan-level survival) and in data source (synthetic population versus real Freddie Mac loan records), so a low ABM share against a high hazard-framework share is equally consistent with either explanation. To make partial progress on separating them, I re-ran the ABM's decision rule on structural covariates (coupon, loan age, original LTV) from the same 75,000-loan Freddie Mac sample as Section~\ref{sec:hazard}, replacing the log-normal, FRED-median-calibrated draws used elsewhere. Behavioral draws (mobility desire, loss aversion, patience, transaction cost) stay synthetic, since the Freddie Mac data contain no analogue for them. The test therefore varies real structural heterogeneity (and, until the post-stratification below, book composition jointly), not a fully real population.

The ABM's mobility-scale parameter is calibrated to hold a real-world baseline turnover floor fixed (Appendix~\ref{sec:method-abm}), so swapping in a population with different LTV and coupon moments raises a calibration question that a naive covariate swap does not resolve on its own: re-derive the turnover-floor anchor on the new moments, or hold it fixed? I report both variants. The recalibrated variant applies Appendix~\ref{sec:method-frictions}'s recalibration rule, re-anchoring the mobility scale to reproduce the same 4--5\% turnover floor under the real population's moments. The frozen variant suspends that rule and leaves the original synthetic-population calibration unchanged.

Before running either variant I fixed the interpretive thresholds ex ante: a recovery above 50\% of the empirical benchmark would undercut the paradigm interpretation advanced in Sections~\ref{sec:abm}--\ref{sec:conclusion}; a recovery below 35\% would corroborate it, consistent with the synthetic-population ABM's own specifications; and a recovery between 35\% and 50\% would be ambiguous.

The recalibrated variant recovers \$453.5 billion, or 59.3\% of the benchmark, above my ex-ante 50\% threshold; the frozen variant \$159.5 billion, or 20.9\%, within the corroborating band. Both are single draws. A pre-committed fifty-seed sweep over both randomization layers (behavioral draws and covariate subsample; run \texttt{cross\_design\_seeds}) shows the threshold crossing is the estimator's, not the seed's. The recalibrated variant averages 60.2\%, standard deviation 3.3 points, 95\% band 55.1--66.0\%; the committed draw sits at the 48th percentile of its own distribution, and all fifty seeds classify as undercutting. The frozen variant averages 22.0\%. The band is wide enough to report (10.9 points against the 7.7-point disclosure threshold fixed before the run) but too narrow to reach the calibration swing. The operative uncertainty here is therefore calibration choice, not seed noise, and it is now measured. The same sweep records each variant's deep-gap floor, and the two separate cleanly on it: the recalibrated CPR at an 8\% market rate averages 4.59\%, inside the observed 4--5\% turnover band on all fifty seeds; the frozen averages 7.05\%, outside it on all fifty (range 6.31--7.63\%). Applied consistently, the admissibility test I use on the externally anchored variant earlier in this subsection counts against the frozen leg, not the recalibrated one: the leg reported as within the corroborating band is the one that misses the turnover moment. I therefore read the 20.9\% as the less admissible of the two calibrations, not as corroboration, which strengthens the undercutting reading. One caveat on ``frozen'': re-deriving that variant's mobility-scale anchor per seed gives mean 43{,}493, standard deviation 1{,}181, minimum 39{,}984 against the committed 43{,}883, the very top of its own range. The anchor is frozen by convention, not because the data pin it there. Table~\ref{tab:crossdesign} summarizes both against the synthetic-population control, and Figure~\ref{fig:crossdesign} plots the paired comparison.

\begin{figure}[H]
\centering
\includegraphics[width=\textwidth]{fig5_cross_design}
\caption{The cross-design test as paired panels. Panel (a): dollar recovery for the two real-covariate variants against the pre-committed thresholds (above 50\% undercuts the paradigm interpretation; below 35\% corroborates it) and the synthetic-population control band; the recalibrated variant's 59.3\% crosses the undercutting threshold while the frozen variant's 20.9\% stays within the corroborating band. Panel (b): lag-0 path correlations, all negative, so neither variant reproduces the empirical path's timing regardless of calibration choice, and after Section~\ref{sec:pathb}'s detrending result no estimator in this paper does.}
\label{fig:crossdesign}
\end{figure}

\begin{table}[H]
\centering
\small
\begin{threeparttable}
\caption{Cross-design test results: real Freddie Mac structural covariates substituted into the ABM under two calibration choices, against the synthetic-population production ABM.}
\label{tab:crossdesign}
\begin{tabular}{@{}p{4.0cm}lll@{}}
\toprule
 & Recalibrated (primary) & Frozen (robustness) & Synthetic control \\
\midrule
Mobility-scale parameter & 36,086 & 43,883 & 43,883 \\
Trapped liquidity & \$453.5B (59.3\%; fifty-seed mean 60.2\%) & \$159.5B (20.9\%; 22.0\%) & \$91.0B (11.9\%) \\
Mean CPR & 8.14\% & 11.63\% & 11.68\% \\
CPR path corr. (lag 0) & $-0.336$ & $-0.311$ & $-0.318$ \\
Peak lag & 0 (no positive alignment) & 0 (no positive alignment) & 0 \\
\bottomrule
\end{tabular}
\begin{tablenotes}[flushleft]\footnotesize
\item Peak lag follows the max-$|r|$ convention of Table~\ref{tab:estimators}. Under the alternative most-positive-$r$ convention, the synthetic control's best alignment is $+0.195$ at lag $-3$, the ``$+0.19$ at the three-month lead'' of Section~\ref{sec:abm-results}; the two conventions summarize the same cross-correlation profile, and the substantive conclusion (no positive contemporaneous alignment, which, after Section~\ref{sec:pathb}'s detrending result, no estimator including Path B shows) holds under either.
\end{tablenotes}
\end{threeparttable}
\end{table}

By my own ex-ante criterion the primary recalibrated variant's 59.3\% recovery, 60.2\% averaged over fifty seeds, moves against the paradigm interpretation: with real structural covariates and a like-for-like turnover anchor, the household-choice ABM explains far more of the benchmark than the synthetic-population specifications of Sections~\ref{sec:abm} and~\ref{sec:robustness-scale} suggested. But the \$294 billion gap between the recalibrated and frozen variants (produced solely by how the turnover-floor anchor is set, covariates fixed) indicates that the ABM's aggregate share is highly sensitive to a calibration choice orthogonal to data source. The test therefore cannot, on its own, cleanly attribute Sections~\ref{sec:abm} and~\ref{sec:hazard}'s contrast to modeling paradigm rather than to data source: recalibration choice does comparable work to whichever the test was designed to isolate.

I therefore withdraw the paradigm reading as a claim of this paper. The reading was pre-committed to a threshold before the run, and the threshold was crossed: on the committed draw, on all fifty seeds, and by a wider margin once the sample carries the book's own composition. A reading that survives only by narrowing, after the fact, what the threshold was taken to test is not one I will defend, and that narrowing is what the pre-commitment architecture exists to prevent. Two results the ABM sections do establish stand without it, and neither is a claim about modeling paradigm: no synthetic-population household-choice specification tested here explains the Federal Reserve's cap shortfall, and no estimator in this paper reproduces the monthly path once the window's shared trend is removed.

A pre-committed post-stratification separates composition from heterogeneity. I reweight the real-covariate population to the SOMA coupon-by-vintage-group composition, the book the benchmark is measured on, and re-derive the mobility scale on it under the same floor rule. That moves the recalibrated variant from 59.3\% to 76.3\% and the frozen variant from 20.9\% to 12.6\%. The reweight uses importance weights, effective sample size 1{,}332 of 10{,}000, disclosed. A seed-42 re-draw within SOMA-weighted cells gives 70.6\% for the recalibrated variant; the frozen variant's re-draw gives 36.0\%. Composition mismatch therefore understated rather than manufactured the recalibrated recovery. At book composition, real-covariate household choice with a like-for-like floor anchor explains roughly seventy to seventy-six percent of the benchmark. The ex-ante undercutting verdict survives composition adjustment, and the calibration-choice spread widens to roughly \$488 billion. The reweight is a level exercise; the path-level failure of Table~\ref{tab:crossdesign} stands.

A third variant drops the turnover-floor anchor entirely and takes both author-calibrated behavioral elements from Path B's external elasticity source. The mobility-desire scale comes from \citepos{liebersohn2024} zero-gap moving level: a 1.5\% quarterly ZIP-code moving hazard at an essentially zero average rate gap, 5.9\% annualized. The wait-and-see freeze share comes from what their 6.5\%-per-quarter-per-100 bp mobility decline implies at the production 150 bp trigger: 18.3\% against the author-calibrated 20\%. The DTI wall and loss-aversion coefficient are already externally sourced. Recovery does rise without the floor recalibration, to 150.6\% of benchmark, but only by predicting essentially zero prepayment at QT-window gaps: CPR at an 8\% market rate is 0.0\%, against observed 4--5\% discount-cohort turnover and the roughly 1.44--1.45\% quarterly moving hazard \citet{liebersohn2024} themselves observe at 2022--2023 gaps. The externally anchored gate structure can match the zero-gap moving level or the observed deep-gap turnover, not both: the exponential mobility-desire distribution that fits the former (scale 4{,}472 against the recalibrated 36{,}086) shuts moving down once the par-payoff penalty binds. I read that as clarifying the floor anchor's role: it keeps the ABM on observed turnover data rather than suppressing its rate response. The anchored dimension is never claimed as fit, since matching the floor is tautological by construction (Appendix~\ref{sec:method-frictions}). What the variant measures is what that tautology purchases: empirical admissibility at deep gaps, another instance of this subsection's deflate-as-well-as-inflate indeterminacy (same harness). Scored against the pre-committed bands, as its own specification required, its 150.6\% classifies as undercutting the paradigm reading, and the frozen artifact records that classification. I do not read it as informative about the paradigm question, and I label that reading post hoc: the pre-run specification committed only to reporting the deep-gap floor diagnostic ``alongside whatever the recovery number is, whichever direction it moves'', not to removing the recovery from the band when the diagnostic fails. The admissibility precondition (an empirically admissible deep-gap turnover level before a recovery is interpretable) was formulated after the result was seen, and I report it here as a labelled post-run qualification, not a pre-registered exclusion. A reader who declines it should score the variant at 150.6\% against the 50\% band; that strengthens the undercutting reading this subsection concedes, and no headline quantity turns on the choice.

I ran a minimal simulated-minimum-distance calibration to reconcile the two external moments: a joint two-parameter fit of the mobility-desire scale plus a forced-move point mass added to the structural family. Moments, admissibility bands, distance, and verdict vocabulary were fixed before execution, and a parity gate replayed the committed 59.303\% first. The verdict is joint infeasibility, now with two degrees of freedom, not one: at the distance-minimizing point (scale 1{,}000, mass 3.9\%) the zero-gap moving level overshoots its band (0.0644 against $[0.055, 0.063]$) just as the deep-gap turnover reaches its own (0.0394 against $[0.040, 0.050]$), the forced mass lifting both together at the scale's floor. The near-fit recovers 107.1\% of benchmark through the same harness, far above the pre-committed cross-design bands: external anchoring pushes this family past the benchmark, not toward it. I deliberately do not fit QT-window CPR moments. The recovery statistic is a functional of that path, so fitting it would make the test a fitted quantity---the overfitting the pre-commitment architecture exists to prevent. Hazard-side headline quantities are not functions of ABM calibration and cannot move under any outcome of this run.

Both variants retain the ABM's central path-level failure: the monthly CPR path correlation stays negative at lag 0 ($-0.336$ and $-0.311$), as every estimator's does once the window's shared trend is removed (Section~\ref{sec:pathb}). Path B's distinction is level accuracy, not timing. The timing critique of Section~\ref{sec:abm} therefore survives regardless of which calibration variant is preferred, even as the dollar-recovery critique does not straightforwardly survive under the recalibrated variant.

\subsection{Cross-Design Synthesis: The Symmetric Companion Test}\label{sec:robustness-synthesis}

The symmetric companion test estimates the hazard framework's survival structure on a synthetic population. It materially updates the cross-design picture. Path B recovers 106.0\% of the benchmark with no real Freddie Mac data at all (96.9\% on the shared accounting basis, whose common netting applies equally here), close to its 107.0\% real-data recovery at the same in-sample calibration point. The raw levels correlation, however, drops from $+0.190$ to $-0.06$ at lag 0 (a levels comparison; neither value survives detrending, Section~\ref{sec:pathb}). That 106.0\% is a recovery level, and this paper demotes recovery levels everywhere else for a reason. The mechanism-removed leg on the same synthetic population ($\beta_1 = 0$, nothing else changed) recovers 95.6\%, so the test identifies a differential of 10.4 points, or \$79.3 billion, not the 106.0\% level (\texttt{synthetic\_companion\_null}). That differential is of the same order as the 9.2-point real-book gap between the central leg's 107.0\% and the $\beta_1 = 0$ null's 97.8\%. Section~\ref{sec:identification} treats that gap as too small to license a level-based reading. A superseded cross-basis figure put that gap at 18.3 points, which was an error of basis rather than of arithmetic. It differenced a standalone-scorer central leg against a shared-basis null. On any single basis the real-book gap is 9.2 points at the in-sample calibration point and 5.6 points at the off-window floor, so the synthetic differential slightly exceeds its real-book counterpart. The synthetic recovery is floor-dominated like the real one. This decomposes the tension. The household-choice paradigm needed real covariates to reach 59.3\%, while the loan-level paradigm reaches comparable aggregate recovery with none. At the path level both paradigms still depend on real structure for correct timing. I treat the aggregate-level contrast of Sections~\ref{sec:hazard} and~\ref{sec:conclusion} as corroborated by this result on its differential rather than substantially strengthened by its level, and the path-level claim as still open. The paradigm reading itself is withdrawn (Section~\ref{sec:robustness-crossdesign}); nothing here is offered in support of it.

\subsection{Cross-Foundation Checks: Hybrid Pipeline and Full-Book Weighting}\label{sec:robustness-hybrid}

The hybrid pipeline of Section~\ref{sec:abm-danish} holds the paper's accounting fixed and swaps the micro-foundation for loan behavior. The ABM's macro-accounting layer runs unchanged (SOMA balances, phased-cap netting, income-scaled curtailment, term-aware amortization, the Danish dynamic-balance loop) while Path B's microsimulated CPR path replaces its CPR surface. No settlement kernel applies, since Path B settles prepayments in the month they occur. The U.S.-side recovery is 97.9\%, and the shared layer adds the curtailment flow Path B's standalone scorer omits: \$69.56 billion over the window. Under the production U.S.-intercept anchor, the Berger-recalibrated Danish re-run yields the U.S. 4.76\% / Danish 5.61\% mean-CPR pair and the $+\$61.2$ billion institutional gap Section~\ref{sec:abm-danish} adopts as its production estimate. The Danish-level bracketing anchor gives 3.39\% and $-\$99.9$ billion. One accounting layer with one empirically grounded micro-foundation carries both regimes, so accounting differences cannot confound the institutional comparison.

The same layer scores every estimator, which makes the recoveries commensurable with the benchmark. The actual SOMA roll-off behind the empirical shortfall contains borrowers' curtailment dollars, which no hazard simulation models, so the shared layer adds them to each simulated roll-off. The netting is identical across estimators (\$69.56 billion for every run), and identical \emph{by construction}, not by coincidence. Hazard dynamics evolve on each estimator's own balances (Path A cohort balances, Path B loan balances, the Danish legs a deliberately compounded dynamic counterfactual, Appendix~\ref{sec:method-abm}). But every \emph{U.S.-leg} dollar conversion is common --- roll-off rescaled month-by-month to the actual holdings path (WSHOMCB) by the book-to-sample ratio. So the curtailment flow, computed on those holdings, is common, not per-run. It is a flat 9.1-percentage-point wedge between the two bases. The invariance covers U.S. legs only: the Danish legs curtail on their own counterfactual balances, whose slower drain retains more curtailment dollars. I quantify and print that differential: \$70.33 billion against the U.S. legs' common \$69.56 billion for the central hybrid (+\$0.77 billion). It is already embedded in every printed gap and immaterial to its sign or size.

Table~\ref{tab:bases} collects every hazard leg's recovery on all four accounting bases; the lock-in margin is basis-invariant at either floor and at every cell of Section~\ref{sec:robustness-floor}'s calibration box. The netting is bit-identical across legs and cancels exactly. The fully corrected figure comes from a joint run, the book-coupon-reweighted simulation scored through the shared layer in one pass, recovering \$765.1 billion (unrounded: \$765.077 billion): 100.0\% of the \$764.7 billion benchmark at one decimal, exactly 100.04\% as the unrounded ratio $765.077/764.748$. That is a mixed-basis measurement. The reweight matches note rates to pass-through buckets, and converting the sample side to the pass-through basis first (note rate less the vintage guarantee fee and base servicing, $\Delta = 0.80$ points, legs at $0.60$--$0.90$) puts the like-for-like composed figure at 98.1\% (97.6--99.1 across the legs; run \texttt{coupon\_convention\_reweight}) --- so 100.0\% is the committed mixed-basis reading, not the corrected number. Path A's joint run reads \$905.1 billion, 118.4\%. The composed figure is conditional on the ex-ante covariate priors (92.2\%--99.2\% on the shared basis, approximately 94.3\%--101.3\% composed, Section~\ref{sec:pathb}), so its precision is a property of that calibration, not a tuned result. Because the netting is common, the joint results equal the composed one-at-a-time deltas ($107.0 + 2.1 - 9.1$) exactly. That equality is a cross-check on the joint run, not its source. The abstract and Sections~\ref{sec:hazard-interp} and~\ref{sec:conclusion} state recoveries on this basis.

The curtailment layer's assumed rate is an external input, not a fitted quantity, so the layer invites one further stress. The monthly flow is an income-driven fraction, independent of that rate, times the rate on the actual holdings path (WSHOMCB), so the netted total is exactly proportional to the rate and interacts with no estimator's dynamics. A $\pm 25\%$ stress (realized window average approximately 0.84\% annualized, Section~\ref{sec:method-benchmark}) scales the \$69.56 billion netting to \$52.17 and \$86.95 billion and the 9.1-point wedge to 6.8 and 11.4 points. The stress shifts every U.S. leg's shared-basis recovery uniformly, leaving the lock-in marginal and every cross-estimator gap unchanged to the decimal, since both legs of each comparison net the identical flow. Because the identity holds month by month, any time-varying shape subtracts out of the marginal term by term; only the common level wedge and the Danish leg's sub-\$1 billion differential can move. Demonstration runs confirm it: seasonal, regime-split, and servicer-mix profiles all leave the marginal unchanged to below $10^{-13}$ points (Appendix~\ref{app:params}). The Danish legs, curtailing on their own counterfactual balances, scale only to first order. An exact accounting-layer rerun (CPR paths bit-identical across scales, curtailment entering only that layer) puts their differential at $+\$0.76$, $+\$0.77$, and $+\$0.66$ billion at $0.75\times$, $1.0\times$, and $1.25\times$. Those figures sit on the \emph{Danish-level bracketing} leg, whose anchor built the committed microsim cache the scaling run consumes, and they do not carry over: re-scored on the U.S.-intercept cache (run \texttt{curtailment\_danish\_scaling}, production anchor), the differential reverses sign and crosses the materiality threshold --- $-\$1.68$ billion at the unit point, $-\$2.39$ billion at $1.25\times$. Curtailment accrues on retained balance, so the sign follows which leg drains faster: the bracketing Danish leg is \emph{slower} than the U.S.\ leg (3.39\% against 4.76\% mean CPR) and accrues more, the production leg \emph{faster} (5.61\%) and less. I therefore withdraw the claim that the curtailment differential is immaterial to the Danish comparison: at $-\$1.68$ billion it is 2.7\% of the $+\$61.2$ billion in-sample gap and 6.0\% of the $+\$28.2$ billion off-window gap, small enough not to threaten either sign but large enough to belong in the figure.

The second check reweights the sample to the book. Both headline recoveries are balance-weighted on the composition of the Freddie object each path uses: Path A the full 2017--2021 origination universe, Path B the 75,000-loan draw, whose unweighted mean note rate of 3.86\% sits above the book's 2.49\% face-weighted pass-through coupon. That draw is where the Path B reweight starts; the unindented bucket shares Table~\ref{tab:composition} prints beside the book are the universe's, exposure-weighted at the window open --- different populations. Reweighting to the book's actual coupon mix shifts Path B from 107.0\% to 109.1\% and Path A from 121.5\% to 127.5\% on the committed mixed coupon basis. Those roughly two percentage points for Path B bound the sample-to-book extrapolation error along the coupon dimension. On a single pass-through convention the same reweight moves Path B only from 107.0\% to 107.2\%, so the like-for-like composition effect is $+0.2$ points and the bound tightens by an order of magnitude (run \texttt{coupon\_convention\_reweight}; full sample-versus-book comparison in Appendix~\ref{app:composition}). The agency dimension (Ginnie Mae collateral) has no internal bound; it is externally bounded and now scored by a published-series overlay in Section~\ref{sec:pathb}, whose Ginnie-specific contribution ($+\$38.5$ billion) lands inside the static bound.

\subsection{Floor and Elasticity-Band Sensitivity: The Calibration Box}\label{sec:robustness-floor}

\subsubsection{The calibration box}

Path B's dominant level-setter is the baseline turnover floor, whose 4\% calibration Section~\ref{sec:pathb} concedes is measured inside the evaluation window. I swept it across seven values from 2\% to 6\% annual CPR jointly with the full Liebersohn--Rothstein elasticity band (5.5\%, 6.5\%, 7.7\%), rerunning the central and $\beta_1 = 0$ null legs at every floor under the production convention: same committed 75,000-loan sample, same seed and regime tuple, one shared macro frame. Both sweeps were pre-committed --- grid, parity gates, and interpretive thresholds entered the repository before any run executed (commit ordering in the replication package's ledger). The fresh runs at the production point reproduce the committed artifacts exactly (central \$818.530 billion, null \$748.185 billion, marginal $+\$70.345$ billion), and the instrumented floor-bind share reproduces the 36\% of 1,683,124 evaluated loan-months quoted in Section~\ref{sec:pathb}.

The sweep mechanics and the contamination separation behind the off-window anchors are collected in Appendix~\ref{app:floormech}. The exhibits they produce --- the calibration box (Table~\ref{tab:floorband}), the marginal as a curve in the elasticity (Table~\ref{tab:lowband}), and the out-of-sample anchors (Table~\ref{tab:oosfloor}) --- sit in Section~\ref{sec:hazard}; what follows are the readings the headline interval rests on.

\subsubsection{Level versus form}

The box sweeps the floor's level, never its functional form: the choice that censors the elasticity. A form test specified ex ante reran the central and null legs with \eqref{eq:pathB}'s hard maximum replaced by a competing-risks combination on the survival scale, $h^{\mathrm{prep}} = 1 - (1 - \underline{h})(1 - h^{\mathrm{vol}})$. Involuntary and voluntary hazards combine on the survival scale, and the elasticity is never inert. Spec and interpretive threshold were fixed in the script's header before its runs, the convention the replication package's ledger records. Script and artifact were committed in one revision. The max-form pair reproducing the committed artifacts exactly at the production floor served as the parity gate. The marginal moves to $+11.25$ points at the production floor, a $+2.05$-point shift past the $\pm 1$-point immateriality threshold fixed ex ante. The form-dependence is real. Under the additive form, though, the marginal is nearly floor-invariant, spanning only $+11.29$ to $+11.21$ points across the 3--5\% floors whose max-form span ($+10.8$ to $+5.5$) failed the stability threshold of Appendix~\ref{app:floormech}. That gate was measuring the max form's censoring mechanics, not instability in the elasticity's contribution. The aggregate level is strongly form-dependent. The additive central leg recovers 67.2\% at the production floor, versus 107.0\% under the max form. The reason is that a 4\% involuntary hazard added to moderate voluntary hazards raises book-wide prepayment. The max form remains production for its semantics: a minimum on \emph{total} turnover. That minimum matches the observed 4--5\% on cohorts aged 12--24 months sitting in the depth plateau that carries 61.5\% of gap-zero exposure. That is the reading the deep-discount calibration anchors, and the anchor at which the two forms nearly coincide. The aggregate-level comparison is context for the level, not the selection rule for the marginal's form, since the design does not identify levels. Run at the off-window anchors themselves (pre-committed run \texttt{floor\_form\_offwindow}), the additive form returns $+11.22$/$+11.21$/$+11.19$ points at the 4.695/4.991/5.334\% floors, within 0.06 points of its production-floor value. The max-form marginal's entire in-window-to-off-window demotion operates through this censoring channel. The near-coincidence defense at the deep-discount anchor does not reach the shallow-gap off-window ones. The form dimension is a measured curve, not a two-point hull (run \texttt{floor\_form\_mixture}; both endpoints reproduce the committed forms bit-exactly). Let $\omega$ be the share of the measured floor entering as an additive competing involuntary cause (a different object from Path A's stratum index $s$) and $1-\omega$ a floor on total turnover. So $\omega = 0$ is the production hard maximum, $\omega = 1$ the additive form. The curve rises steeply and plateaus, while the central leg's standalone-scorer recovery falls from 100.4\% to 55.9\%, the level cost the additive end pays. At the off-window anchors the marginal is $+7.5$ points by $\omega = 0.1$, $+9.7$ at $\omega = 0.25$, $+10.7$ at $\omega = 0.4$, and flat near $+11.2$ from $\omega = 0.6$ up (peak $+11.24$ at $\omega = 0.8$, a 0.03-point plateau dip to the additive endpoint). So Section~\ref{sec:pathb}'s reading that the strictly-involuntary share is plausibly well under half does \emph{not} pull the marginal toward the max end. Modest additive shares already close most of the gap, an additive component acting wherever the max form censors. The headline stays at the $\omega = 0$ endpoint on the max form's semantics, now the conservative endpoint of a measured curve. It is not a position defended by fit. That level cost carries to the off-window anchors, large enough to report beside the hull. At the 4.991\% headline anchor the additive central leg recovers 55.9\% of the benchmark against a 44.7\% null, and 59.3\% and 52.0\% at the 4.695\% and 5.334\% anchors at the central band point. So the additive form undershoots the max form's aggregate level by 43.5 to 45.3 points at every off-window anchor. The hull's upper endpoint is its least accurate cell: $+13.10$ points, additive at the 4.695\% floor and the band's 7.7\% high edge. Its central leg recovers 61.2\% of the benchmark it is scored against. The form-conditional range therefore runs over specifications that differ sharply in aggregate fit, and its upper end should not be read as an equally credentialed member. Reported side by side and never crossed, form and transform were presumed to add. The pre-committed joint cell (\texttt{concave\_additive\_marginal}) shows they do not. The concave transform under the additive form puts the marginal at $+9.25$ points at the production floor and $+9.22$ at the 4.991\% anchor, inside the hull at both. The range itself stands. But the interaction is $-1.47$ points at the headline anchor, past the same $\pm 1$-point convention. The reason is censoring. The transform costs only $-0.52$ points under the max form, which absorbs most of it into the floor, against $-1.99$ under the additive form, which never censors and passes the whole effect through. A successor run closes that grid (run \texttt{concave\_hull\_band\_ends}). The concave transform's remaining sixteen off-window cells put the form $\times$ transform product at thirty-six of thirty-six, and the hull does not survive intact. Seventeen of the eighteen concave cells fall inside it. But concave under the max form at the highest defensible floor and the band's low edge returns $+3.53$ points, below the $+3.89$ lower edge the log-linear cells set. So the form-conditional hull widens at the bottom to $+3.5$ to $+13.1$ points. Its upper edge is unchanged: the concave transform is signed downward, its highest cell reaching only $+10.83$. Any reading that composes a form adjustment with a transform adjustment additively is unsupported at the calibration this paper headlines. The form-robust statement of the paper's identified quantity is a lock-in marginal of roughly $+9$ to $+11$ points at the production floor, positive under every level, band, and form combination tested.

\paragraph{A transportable lesson: the floor's functional form is a first-order choice.} The involuntary-turnover floor is calibrated, not estimated, and exists to keep ordinary turnover from being attributed to the rate gap --- a nuisance parameter by construction. Its combination rule is nevertheless first-order for the reported marginal, and its leverage is largest exactly where this paper takes its headline: at the off-window anchor, switching from a hard maximum on total turnover to an additive competing cause carries the marginal from $+5.6$ to $+11.2$ points, a doubling, against a rise from $+9.2$ to $+11.2$ points --- a little over a fifth --- at the in-window calibration. The portable device is the mixture curve above: indexing the two forms by the share $\omega$ entering additively replaces a two-point fork with a measured curve whose endpoints reproduce the committed forms bit-exactly. Its shape is what transports: it rises steeply and plateaus, so modest additive shares buy most of the distance between the forms, and reporting the two endpoints alone understates the exposure of any design that calibrates a floor, a censoring threshold, or a competing risk and then differences across it. The 2018 ladder carries own-data evidence on the same fork, out-of-window and out-of-episode: re-binning its nested reads by exposure-weighted differencing gives per-bin turnover of 7.09\% and 6.99\% in the two shallowest bins, then 4.67\% and 4.65\% in the two that follow --- a step, then a plateau across bins holding 61.5\% of gap-zero exposure (run \texttt{depth\_ladder\_shape}). The forms predict different shapes: a floor on total turnover censors the voluntary hazard once the gap is deep enough and per-bin turnover should stop falling; an additive competing-risks form never censors and predicts continued decline. The plateau is what the max form predicts. Two qualifications travel with it, the first the one an additive reading would press: the ladder does eventually fall, to 4.35\% and 3.61\% in the deepest bins --- but those hold 3.85\% and 0.02\% of exposure and the source artifact's own rule flags them as not well supported, as it does the bin between them, which \emph{rises} to 5.14\% rather than falling and so is not evidence of continued decline either. The second is that this is a shape comparison on one leg, not a test: I read it as own-data evidence consistent with the production form, not as a resolution of the fork.

\subsubsection{The off-window re-anchor}

Re-running the central and $\beta_1 = 0$ legs with the floor at each defensible off-window value, production convention otherwise, gives Table~\ref{tab:oosfloor}. The fresh runs reproduce the committed production artifacts exactly at the 4\% floor: central \$818.530 billion, null \$748.185 billion, marginal $+\$70.345$ billion, floor-bind share 36\% of 1,683,124 loan-months. Across the 2018 leg's defensible band the marginal runs $+6.8$ to $+4.3$ points ($+\$51.8$ to $+\$32.6$ billion), centered at $+5.6$ points ($+\$42.6$ billion) at that leg's mid-grid 4.99\% reading, and $+3.9$ to $+7.5$ points once the full Liebersohn--Rothstein band is swept jointly with it. The small numbers come from the contaminated anchors: the pooled 6.07\% read, evaluated at the box's adjacent 6.0\% row, gives $+2.3$ points, the 2019 leg's 6.91\% read $+0.7$. An off-window floor is higher than the in-window 4\%, and the marginal decreases in the floor, so a higher anchor shrinks the identified contribution. It does not reverse it. The direction of the in-window concession therefore stands, but not the conclusion that the plausible contribution sits at the calibration box's lower edge. That edge was an artifact of the refinancing-contaminated pooled read; the defensible off-window floor places the contribution in the box's interior, at roughly $+5.6$ points, not at $+2.1$ to $+2.4$. Every read in that band sits on cohort-months aged twelve to twenty-four, the only ages the 2018 leg contains, so the leg's mature-seasoning test is not computable (Appendix~\ref{app:floormech}); uncomputable is not passed. Where the mature contrast is measurable it runs one way in every cell (Section~\ref{sec:identification}), putting the band below the mature book's baseline turnover and the marginal it implies above what that book would give.

The move from the in-window anchor to the off-window one changes two things at once; until now the paper charged both to the window. The in-window anchor reads 3.972\% at gap $\leq -2$ points, the off-window anchor 4.991\% at gap $\leq -0.25$ points. Mapping both through the committed floor-to-marginal grid makes the move $-\$28.312$ billion, or $-3.702$ points. That grid is PCHIP on the frozen sweep, not re-derived, and it reproduces the committed grid to within \$0.69 billion at every read inside it. Decomposed at matched depth, the move splits. The window component holds depth at gap $\leq -0.25$ points on both sides: $-\$20.874$ billion ($-2.729$ points, 73.7\% of the move). The depth component, taken within the in-window panel from gap $\leq -2$ points to gap $\leq -0.25$ points, is $-\$7.438$ billion ($-0.973$ points, 26.3\%). Re-anchoring the same decomposition at the deepest depth both panels support (gap $\leq -1$ point) gives 66.9\% window and 33.1\% depth. The window dominates either way, and the demotion is not primarily a depth artifact. But depth owns a real and previously unattributed \$7.4 billion, so the description must change. At matched depth the in-sample comparison point is $+\$63.5$ billion, not the $+\$70.3$ billion this paper narrates the demotion against. Roughly \$20.9 billion of the \$27.7 billion demotion is the window, the rest a change in the discount depth at which the anchor is defined, with nothing to do with out-of-sampleness. The reverse ordering (window first at matched deep depth, then depth) is not computable: it would read the off-window leg at gap $\leq -2$ points, where that leg retains 0.02\% of its exposure. So no average of the two orderings is formed, and the depth-first path is reported alone. Extending the 2018 depth ladder from three well-supported depths to five leaves the committed clean band exactly [4.695\%, 5.334\%]. The two added reads, 4.722\% at gap $\leq -0.75$ points, 4.869\% at gap $\leq -1$ point, fall strictly inside. The marginal band therefore stays \$33.31 to \$51.75 billion around the \$42.58 billion point. The \$33.31 billion grid-read and \$32.6 billion direct-run endpoints quoted earlier differ by the PCHIP grid-read wedge, whose committed maximum is \$0.69 billion. Two further pre-committed measurements discipline this band. First, the floor reads carry sampling error the depth cuts do not express. A 1{,}000-replicate stratum-cluster bootstrap puts the clean-band reads' standard errors at 0.39--0.46 points of CPR. The 2018 leg rests on 31 clusters---31 of the panel's 296 four-way strata are occupied by that selection. The in-window reads, on 226 clusters, come in at 0.16--0.19 points. Propagating the point anchor's replicates through the committed grid gives a percentile 95\% interval of $+3.0$ to $+8.0$ points on the off-window marginal. That is double the depth-cut band's width, so the binding uncertainty on the headline marginal is the floor's sampling error, not the depth convention. That percentile read under-covers at that cluster count and is demoted. The pre-committed corrections widen it mostly at the top. They are run \texttt{floor\_inference\_correction}, whose percentile layer is reproduced bit-exactly as its parity gate, and its Webb-weight successor \texttt{floor\_inference\_correction\_v2}. The corrected ladder, with its degrees of freedom, is Table~\ref{tab:ladder}. The restricted wild-cluster inversion's $+2.3$ to $+9.1$ points is the binding layer, its lower edge above zero on every corrected read. The pre-committed coverage rule of run \texttt{fewcluster\_coverage} selects it, with the Webb wild-$t$'s $+2.9$ to $+8.7$ reported beside it. Second, age-standardizing the off-window floor reweights the 2018 leg's age-bucket reads to the in-window deep-discount population's age mix. It also imputes the 2018 leg's empty mature buckets from the in-window age gradient (84.0\% of weight imputed, and labeled as such). The standardized floor reads 5.51\%, above the clean band's top, for a measured marginal of $+3.7$ points against the $+3.8$ its grid position implies. That is an indicative bound, not a measurement, but signed one way: the clean band is open below $+4.3$, and that endpoint is not a conservative floor. This read holds the conventional accounting fixed. The Ginnie overlay of Section~\ref{sec:pathb} holds the floor at 4.991\%. Both correct the same $+5.6$-point cell, both cut the same way, and the joint cell is now run (run \texttt{b5\_joint\_cell}). They do not look like the same adjustment counted twice. One moves the floor's level; the other confines the marginal to the conventional face share. The overlay's share scaling is identical across the three speed series it was tested on. So carrying that $0.797$ scaling up to the higher floor---a carry the overlay run does not itself establish---would put the marginal near $+3.0$ points, at the bottom of that sampling interval and below the form-conditional hull's $+3.5$ lower end. The joint cell returns $+2.9$ points, an interaction of $+0.0$ against the proportional projection. The first joint cell this paper ran for two adjustments to the marginal was the concave transform under the additive form above; it returned an interaction of $-1.47$ points at the headline anchor. That is why a composed point here would not be credible. The same matched-depth machinery undercuts the argument by which the 2019 leg was rejected in Appendix~\ref{app:floormech}, though not the rejection itself. That argument compares the 2019 leg's gap-$\leq$-0 read of 6.910\% against the retained 2018 leg's gap-$\leq$-0.25-point read of 4.991\%. The 1.919-point spread is taken across two different depths, which is not a valid comparison as written. At matched depth the two legs do not converge either. Over the four depths at which both clear the 10\% support rule, the 2018-versus-2019 spread runs 1.576, 0.491, 2.112, and 0.857 points, non-monotone in depth. At half a point it exceeds the mismatched 1.919 points the earlier argument rested on. So neither ``the legs disagree by 1.9 points'' nor ``the legs agree once depth is matched'' is a defensible summary; the second would be a cherry-pick of one depth in four. The pooled 2017--2019 read is not a third independent leg: it contains both 2018 and 2019 (and 2017). Any appeal to three legs agreeing double-counts; there are two. And the 2019 leg's own exposure collapses from \$1{,}085 billion at gap $\leq$ 0 to \$199 billion by half a point and \$0.29 billion by one and a half. Its deep reads are noise on the same grounds that disqualify the deep off-window cell. The honest verdict is therefore that the refinancing-contamination story is unadjudicable below gap $\leq$ 0 rather than refuted there. The only depth carrying solid exposure for both legs is gap $\leq$ 0, where 2019 sits 1.576 points above 2018. That is directionally consistent with refinancing contamination and with the seasoning diagnostic that independently disqualifies the 2019 leg. The conclusion stands; the argument as earlier stated does not, and I state the weaker one. What it leaves standing is a definitional alternative, not a measurement dispute. Suppose the floor's level is itself an equilibrium outcome of the episode (turnover depressed by the same gaps the elasticity prices, as the refi-era reads' higher levels would suggest under that reading). Then part of what the rate-inelastic baseline absorbs belongs to the lock-in channel, and the rate-inelastic share overstates what is invariant to household behavior. Pricing that channel requires a cross-regime turnover-versus-gap relationship estimated outside this window's rate configuration, which this design does not contain. I name it as future work. I do not fold an unestimated correction into the headline.

\subsubsection{Floor stability over months}

A second instrument addresses the months dimension directly, not the population dimension. It is a floor-stability check, not an outcome holdout: no realized prepayment from the held-out months enters either simulated leg, so the construction can falsify only the floor calibration's temporal stability, never the marginal against realized data. Calibrating the floor on deeply out-of-the-money turnover confined to the window's first nineteen months (June 2022 through December 2023) returns 3.97\%. The central-minus-null marginal then sums over the remaining twenty-three months (January 2024 through November 2025) alone, which enter that calibration nowhere. That held-out marginal is $+\$45.1$ billion, against $+\$44.8$ billion at the production floor over the identical months: very nearly the magnitude, and with it the sign, survives evaluation months the calibration never touches. The decomposition is exact, because trapped liquidity is a monthly sum against the redemption target: the target cancels from the central-minus-null difference and the monthly marginals add back to the aggregate. Across every floor and elasticity-band combination run here---twenty-seven cells spanning floors from 3.97\% to 6.91\%---the marginal is positive without exception, a consistency check the window's rate configuration cannot let fail (Section~\ref{sec:identification}). The identified content of the lock-in channel is therefore a magnitude that survives both instruments as a bounded interval. Under the production max form it is $+4.3$ to $+6.8$ points at the defensible off-window floor, and $+3.9$ to $+7.5$ jointly with the elasticity band. Under the additive form at the same anchors it is $+9.6$ to $+13.1$, giving a form-conditional hull of $+3.5$ to $+13.1$ points. Against that, the in-window production calibration returns $+9.2$ points, now reported as this paper's in-sample calibration point, not its headline.


\section{Conclusion}\label{sec:conclusion}

The shortfall of the Federal Reserve's mortgage runoff beneath its Quantitative Tightening caps was not solely a product of high interest rates. It followed from the structural design of the U.S. mortgage contract interacting with a prepayment environment mostly rate-inelastic relative to the caps. Under the production floor form the $\beta_1 = 0$ null recovers 85.7\% of the benchmark at the headline calibration (35.6\% under the additive form), so the caps could not have bound under any rate response above the baseline turnover the null embeds. How much of the remainder sits outside household-level decision-making is a question the cross-design test of Section~\ref{sec:robustness-crossdesign} leaves open. This paper's central empirical finding is the rate-inelastic decomposition, established in Section~\ref{sec:hazard} on real Freddie Mac loan-level data. The agent-based model of Section~\ref{sec:abm}, run on the synthetic household population specified in Section~\ref{sec:method}, stands beside it as a stress-test companion, since Section~\ref{sec:robustness-crossdesign} confines what its verdict can mean. The two are complementary and neither is derived from the other, though they are not independent in calibration: both are anchored to a baseline turnover floor read off discount-cohort turnover, and at the in-sample calibration point to the same 2023--24 read (Section~\ref{sec:pathb}).

Rational, behaviorally extended, or subjected to a direct burnout test, the ABM cannot explain the majority of the Federal Reserve's cap shortfall. A structurally distinct hazard framework, built entirely from loan-level survival dynamics with no household utility function at all, recovers 91.3\% of the full empirical benchmark in its microsimulation form on the benchmark-consistent accounting basis, at the off-window floor this paper headlines. It recovers 88.2--93.7\% across that floor's defensible range and 97.9\% at the in-sample calibration point, so the miss is 8.7 points on the headline calibration and 2.1 points on the in-sample one. That is a level-accuracy distinction, not a timing one: no estimator's monthly co-movement survives detrending (Section~\ref{sec:pathb}). The estimated cohort-hazard form lands at a comparable aggregate level but does not corroborate it: its own coefficient's sign is not statistically distinguishable from zero. It carries no valid production-spec uncertainty and is excluded from headline figures (Section~\ref{sec:patha}).

The composition of that recovery disciplines what the contrast can mean. Under the production floor form the $\beta_1 = 0$ null recovers 85.7\% on the same basis at the headline floor (35.6\% under the additive form) and 88.7\% in-sample. The paper's own ablations bound the named loan-level mechanisms: burnout $0.45$ points at the headline floor and $0.86$ in-sample, FICO/LTV covariates $1.3$ points, the delinquency pipeline \$0.39 billion. No servicer channel beyond that pipeline is modeled at all. What delivers the level is scheduled amortization, the seasoning baseline, and the turnover floor. The shortfall against the phased caps is mostly rate-inelastic, because the caps sat far above what any plausible prepayment environment would have delivered. But against the Federal Reserve's own ex-ante projection rather than the non-binding cap, the lock-in channel accounts for roughly half the genuine surprise under the central allocation. That share runs 23--49\% at the off-window floor and 38--80\% at the in-sample calibration point across the two disclosed intra-2022 allocations (Section~\ref{sec:method-benchmark}). The elasticity's contribution is a bounded margin over that rate-inelastic base: $+2.3$ to $+9.1$ points under the production floor form ($+9.2$ at the in-sample calibration point; $+11.2$ under the competing-risks form). Inside that interval, $+5.6$ is the value at the mid-grid anchor. Measured corrections to the floor and the accounting basis move that value in both directions within the interval (Section~\ref{sec:identification}). The cross-design answer is asymmetric. The hazard framework's aggregate recovery survives a fully synthetic population. The ABM's share moves with calibration choice and data source: 59.3\% recalibrated, 76.3\% at the book's composition. That range sits above the pre-committed threshold at which Section~\ref{sec:robustness-crossdesign} concludes the contrast cannot, on that test alone, be cleanly attributed to modeling paradigm. The path-level timing question remains open for every estimator.

Section~\ref{sec:abm-danish}'s Danish counterfactual retires a superseded figure. Recalibrated against \citepos{berger2026} separately-estimated Danish elasticities (unrefereed working paper, draft July 2026), the \$925.5 billion mechanism-extrapolated gap collapses by an order of magnitude. Anchoring the transplant so that it changes only the payoff rule---the anchor consistent with the baseline turnover floor's semantics---signs it at $+\$61.2$ billion at the band's zero-refinance edge, 8\% of the benchmark and the image of the lock-in marginal on the in-sample calibration both legs share. At the off-window floor the paper headlines, the rule-only gap is $+\$28.2$ billion, 34\% below that marginal (Section~\ref{sec:abm-danish}). The sign is anchored, not estimated: pinning the Danish leg at the U.S. zero-gap hazard guarantees it against a book that is out of the money almost everywhere. The bracketing anchor that also imports Danish baseline mobility reverses it to $-\$99.9$ billion, on a Danish leg whose 3.39\% mean prepayment sits below the turnover floor the U.S. legs impose; that inconsistency demoted it to a bracket. The payoff-rule question therefore ends at a modest magnitude and a stated anchor. One qualification is priced as a bracket of its own: the Danish leg credits retired balance at par, and the pre-committed bracket (runs \texttt{buyback\_credit\_bracket} and \texttt{buyback\_discount\_rederived}) prices both incidence readings---the gap stands at $+\$61.2$ billion under the balance-adjustment reading and reverses to $-\$51.0$ billion under the cash-haircut reading at the re-derived discount. The relief therefore stands on the benchmark's own face denomination, while as reserve drain the rule costs the Federal Reserve exactly the par windfall it would otherwise have extracted from moving households. On that reading, legs (2) and (3) of the trade-off are the same dollars seen from two sides. The mobility relief such a rule would provide remains the better-supported and larger effect.

The remaining leg of the trade-off is secondary-market liquidity. The Danish system relies on strict one-to-one match funding (the Balance Principle); the United States relies on the To-Be-Announced (TBA) forward market, which requires highly homogenized, pooled mortgages. \citet{vickery2013} show that this homogeneity produces a measurable liquidity premium: TBA-eligible pools trade at significantly lower cost than comparable ineligible specified pools. \citet{gao2017} measure the same differential directly in transaction data, \citet{song2019} document the dollar-roll market that finances TBA positions, and \citet{campbell2013} frames the U.S.--Danish contrast as a mortgage-design problem in exactly these terms. A market-value repurchase feature prices each payoff at the market value of the funding bond series (a series-level price common to every borrower funded by that series on a given day). What it strains is therefore not homogeneity itself but its \emph{locus}: Denmark standardizes at the level of large, tap-issued bond series, the TBA market at the level of the deliverable pool. The Danish callable market's own depth shows that market-value redemption and liquid secondary trading coexist under the series convention \citep{berg2018}. The binding cost for the United States is migrating its issuance architecture from pool-level to series-level standardization, not a liquidity impossibility. Adopting the rule without that redesign would nonetheless risk a documented source of liquidity in the transition. A hedging-side literature sharpens what the two payoff rules allocate. Prepayment risk is borne and priced by specialized MBS investors, not by diversified households \citep{gabaix2007}; its aggregate hedging feeds back into the level and volatility of long-term rates \citep{malkhozov2016}; and the behavioral component is the part that cannot be hedged away, only bounded \citep{perotti2024}. The Danish rule does not remove that behavioral residual --- the borrower retains the par call on the rate-decline side --- but it removes the rate-rise extension leg on which this cycle's exposure accrued. The choice between payoff rules is thus also a choice about which side of the convexity profile the imperfectly hedgeable risk sits on.

U.S. policymakers weigh three considerations: (1) deep, homogenized secondary market liquidity (the TBA market); (2) continuous household mobility (avoiding the lock-in effect); and (3) convexity-neutral duration profiles (predictable institutional cash flows). This paper's evidence is that they do not form the strict trilemma a three-way ``trade-off'' would imply. That reading holds under its production floor form; the additive form roughly doubles the margin at issue. The evidence of Sections~\ref{sec:abm} and~\ref{sec:hazard} is that the U.S. has prioritized (1) at real cost to (2): household mobility is genuinely constrained, while the cash-flow shortfall against the caps is mostly rate-inelastic. The Federal Reserve staff's own record and the ex-ante projection corroborate that composition as composition, not as numerical agreement, and not independently of it, since the projection may already have embedded part of the channel (Section~\ref{sec:lit}). Within that composition, the lock-in channel contributes an identified margin of $+2.3$ to $+9.1$ points at the off-window floor, $+5.6$ at the mid-grid anchor inside it ($+9.2$ at the in-sample calibration point). That constraint is not evenly borne. The marginal's per-balance intensity tilts toward below-740-FICO ($1.23\times$--$1.26\times$) and above-80-LTV ($1.23\times$) borrowers (Section~\ref{sec:identification}), and the statutory assumability carve-out that would relieve it is bounded above by the 20.4\% Ginnie share, with take-up unmeasured; the mobility cost therefore falls disproportionately on the households with the fewest refinancing alternatives. What I no longer find support for is the assumption that (2) and (3) trade off sharply under face accounting, in which fixing mobility would cost face-value cash flow. Under the rule-only transplant, Section~\ref{sec:abm-danish}'s recalibrated Danish counterfactual finds that a market-value repurchase rule would have relieved (2), while its effect on (3) is incidence-conditional: an $\approx\$61$ billion improvement under face accounting, a $-\$51.0$ billion cost under cash accounting at the re-derived discount (the foregone household-to-bondholder transfer). That reading inherits the marginal's floor-and-band conditionality and its statement at the in-sample calibration point (Section~\ref{sec:identification}). The dissolution rests on thinner ground than the framing suggests, so I state it plainly. The rule-only Danish gap is the lock-in marginal viewed from the counterfactual side, so ``(2) and (3) do not trade off sharply'' reduces to ``the marginal is small''. Under the production floor form the marginal is bounded small by construction. By contrast, under the disclosed additive form, which never censors the elasticity, the same counterfactual is about $+11$ points and nearly floor-invariant (Section~\ref{sec:robustness-floor}). Because \eqref{eq:pathB} combines the turnover floor with the voluntary hazard by a hard maximum, switching the elasticity off can raise the hazard only by the excess of the $\beta_1 = 0$ hazard over the floor---and not at all where the floor binds on that leg itself, 14.3\% of loan-months at the in-sample calibration point (Section~\ref{sec:identification}). The dissolution is therefore a statement about a calibrated ceiling, not an independent finding, and it is form-conditional. What survives across forms is that the trade-off is modest, not that it is absent. On either form, the binding constraint against adopting such a rule is not a mobility--cash-flow trade-off but the TBA liquidity premium described above. That is an argument from the cited literature, not a result of this design: nothing here prices a series-level call, so the investor-side cost of the transplant is named and left unpriced (\citet{diep2021} show that prepayment risk is compensated in the cross section of MBS returns, which is the price this design does not measure). This is the two-sided exchange the paper opened with: (1) purchased in the currency of (2), with a small side payment in (3) whose sign depends on the buyback's incidence. As long as the U.S. maintains the par-payoff MBS structure within the TBA framework, severe duration extension should be expected to recur in future tightening cycles that begin, as this one did, with a large outstanding stock of deeply in-the-money, below-market-coupon mortgages. The state dependence documented by \citet{berger2021} and \citet{eichenbaum2022} makes the extension conditional on that stock, not on tightening per se. The recurrence claim is itself an extrapolation beyond what this design tests. The accounting machinery is also episode-shaped: with essentially the whole book out of the money, the marginal's sign is forced and the refinancing margin is inert, so a mixed-moneyness window would require re-deriving both. This design tests one episode. The duration extension itself, however, is mostly the arithmetic of scheduled amortization and baseline turnover against caps set above any plausible prepayment path; an alternative payoff rule would have trimmed it by the lock-in margin, not eliminated it, and would carry the liquidity cost above.

\subsection{Limitations and Avenues for Future Research}\label{sec:limitations}

Several limitations remain on both sides of the framework.

\textbf{Housing supply as a competing channel}: The mobility margin this paper imports is estimated in an environment of historically tight for-sale inventory, and housing supply enters here only as a friction scaler on the moving decision, never as a channel in its own right. A household that does not move because nothing is listed is observationally identical, in these data, to one that does not move because its coupon is below market, and this design does not separate them. The direction is argued rather than measured. A binding supply constraint would absorb some of what is attributed to lock-in, which is why the scaled-null exercise that removes the non-rate share of the moving decision is reported alongside the headline rather than as a robustness footnote (Section~\ref{sec:identification}; \citealp{aladangady2024,hsieh2019}). No inventory series is used to price it.

\textbf{Distributional silence}: This paper measures a book-wide institutional cost and never asks who bears it. Ginnie Mae borrowers---20.4\% of SOMA face, and the population whose structurally faster speeds Section~\ref{sec:pathb} documents---sit outside the Freddie Mac sample entirely, so nothing here establishes whether the mobility cost falls differently on FHA and VA households than on the conventional borrowers the sample contains. I record that as a distributional limitation of the design, not merely as a coverage bound on the estimate.

\textbf{Sample scale and comparability}: The ABM's 10,000-household synthetic population and the hazard framework's 75,000-loan real Freddie Mac sample differ in both size and data source. Section~\ref{sec:robustness-scale} settles the size question on the production specification: population size changes precision, not central tendency. The completed cross-design pair of Section~\ref{sec:robustness-crossdesign} addresses the data-source question asymmetrically --- the hazard framework's aggregate recovery is largely data-source-independent, while the ABM's is not.

\textbf{Data:} A dedicated ingest pipeline integrates the Freddie Mac 2017--2021 loan-level data (20 origination quarters) that both hazard paths use; where real data files are unavailable, results are validated against a synthetic fixture before real-data estimation. Fifteen-year MBS, approximately 9\% of SOMA face value, remain excluded from cohort modeling on the hazard side only---the ABM side folds them in (Appendix~\ref{sec:robustness-pipeline}).

\textbf{Model:} On the ABM side, the monthly CPR path fit remains weak (raw $R^2$ of $-6.984$) despite a directionally sensible level response to rates, and the production Danish counterfactual is a rule-only transplant onto U.S. households (Section~\ref{sec:abm-danish}): it does not estimate how Danish institutional frictions would themselves migrate, and the Danish-level bracketing anchor shows the gap's sign is not robust to bundling baseline-mobility differences into the transplant. The debt-to-income check remains front-end only, without a measure of total household debt. On the hazard side, Path A's Poisson pseudo-maximum-likelihood estimator requires a mild ridge penalty to remain well-conditioned at 295 stratum fixed effects (a stabilization device: Section~\ref{sec:patha} shows the prior-spec (v3) coefficients equal the unpenalized fit on the estimable strata), and its holdout RMSE (spec v3: approximately 38 percentage points unweighted, 2.5 points exposure-weighted; the frozen spec v4 fit re-scored on the same holdout: 37.9 and 3.0) together with a lag-0 CPR correlation of $-0.378$, more negative even than the ABM's $-0.318$, confines its evidential contribution to the aggregate dollar level, as Section~\ref{sec:patha} details; Path B's literature-calibrated hazard magnitudes, while directionally and quantitatively closer to the empirical benchmark than any ABM specification, are simulated on the 75,000-loan real Freddie Mac sample with elasticities calibrated from secondary literature rather than estimated jointly with the empirical benchmark itself, and its recovery percentage is reported across the calibration band (105.9\%--108.2\%) and now also with loan-level sampling uncertainty from an ex-ante-specified stratified bootstrap over the 75,000-loan draw (Section~\ref{sec:pathb}). An external second-agency replication (Fannie Mae loan-level data, same pipeline end-to-end) reproduces the lock-in marginal at $+8.68$ points against Freddie's $+9.20$, inside its pre-committed envelope: a sign-and-envelope replication, not a hazard-level match, and one whose envelope gate is 11.06 points wide, so the 0.52-point agreement, rather than the gate, is what should be read from it (Section~\ref{sec:pathb}). One global fact about the design, each of whose components is disclosed only locally: no outcome-holdout months exist anywhere in this paper. The nearest construction is the marginal's temporal holdout (Section~\ref{sec:robustness-floor}), in which the turnover floor is calibrated on the window's first nineteen months and the central-minus-null difference is then summed over the remaining twenty-three alone, months that enter that calibration nowhere; it returns $+\$45.1$ billion. That construction holds for the marginal, not for the level, and it is an input-stability check rather than an outcome validation: realized prepayment from the held-out months enters neither simulated leg, so it cannot falsify the marginal against data, and the statement above stands without exception. The production baseline turnover floor is still calibrated on 2023--24 turnover inside the evaluation window (the pre-committed off-window re-measurement of Section~\ref{sec:robustness-floor} finds the defensible off-window level higher than 4\%, so the direction of any anchoring error shrinks rather than inflates the identified marginal, and the headline marginal is now reported at that off-window level rather than the in-window one); Path A trains on 19 of the 42 QT months and its ridge-selection holdout doubles as its out-of-sample diagnostic; and Path B's \emph{level}, as distinct from its marginal, is still set by a floor measured in the window on which it is evaluated. Every recovery percentage in the paper is accordingly an in-sample or calibrated quantity in the months dimension, disciplined by ex-ante-specified sweeps, nulls, parity gates, and bootstraps rather than by clean holdout months. All headline figures are frozen in tagged run manifests, with reproduction details recorded in the replication package's run ledger. The marginal itself, at either floor, is a model-based decomposition conditional on the disclosed hazard structure (Section~\ref{sec:identification}): causal identification in the econometric sense would require exogenous variation in rate gaps across otherwise-comparable loans, which this design does not exploit.

The Danish-regime recalibration (Section~\ref{sec:abm-danish}) calibrates both frameworks' Danish CPR response against \citepos{berger2026} separately-estimated moving and refinancing elasticities rather than an extrapolated U.S. mechanism, and it anchors the rule-only transplant at the U.S. zero-gap intercept. That collapses the institutional gap by an order of magnitude, to a signed $+\$61.2$ billion. That source is an unrefereed working paper, draft dated July 2026, and it is moving: its general-equilibrium refinance-in-place estimate rose roughly twentyfold between its January and July 2026 drafts (Section~\ref{sec:abm-danish}). Path B's full 5.5\%--7.7\% elasticity band is propagated through the microsimulation (Section~\ref{sec:pathb}). That exercise also surfaced and corrected a units error in the rate-gap coefficient's application (Appendix~\ref{sec:robustness-pipeline}). The symmetric companion cross-design test is complete (Section~\ref{sec:robustness-synthesis}). The hazard framework's survival structure recovers 106.0\% of the benchmark on a fully synthetic population, closely matching its real-data recovery, while its path-level timing does not transfer to synthetic data. Three questions remain open. The first is whether to add a seasonal multiplier to Path B's baseline hazard (Section~\ref{sec:pathb}). The second is a total-debt measure, or threshold sensitivity, for the ABM's front-end DTI gate (Appendix~\ref{sec:method-frictions}). The third is the monthly-frequency question the timing diagnosis leaves: the empirical path's variation at that frequency is seasonal, not rate-driven, and no estimator here models it. The second-agency replication is complete (Section~\ref{sec:pathb}; the envelope caveat is in the Model paragraph above). The cyclical-floor variant is likewise complete. What its grid measures is floor dispersion under the hard maximum rather than cyclicality, with the surviving co-movement component small and bounded (Section~\ref{sec:pathb}; Table~\ref{tab:uncertainty} notes; Appendix~\ref{sec:robustness-seasonalfloor}). The Ginnie composition bound has been upgraded from two published snapshots to a time-varying overlay built from the published monthly CPR series. That overlay leaves the conventional-share marginal at $+7.3$ points at the in-sample calibration point ($+4.4$ composed with the off-window headline floor), positive under every variant, with the Ginnie-specific level contribution inside the static bound (Section~\ref{sec:pathb}). Replication alone could not retire the vintage residual: the 2022 and pre-2017 vintages (23.1\% and 10.6\% of book face) lie outside the 2017--2021 universe of both agencies' samples used here. The observed-speed bound of Section~\ref{sec:pathb} now disciplines that extrapolation directly (\$11.7 billion, 1.5\% of benchmark, signed toward overstating trapped liquidity). What remains open on this dimension is therefore no longer the absence of a bound but the bound's own proxy assumption: observed Fannie segment speeds standing in for the SOMA book's, the same single-agency posture the production calibration itself carries. A modeled overlay in place of that proxy is not constructible under the committed sampler: the staged replication sample enforces the 2017--2021 vintage universe, the pre-QT snapshot rule leaves second-half-2022 originations with no admissible observation month, and the acquisition window closes at 2022Q4, so 2022 originations acquired in 2023 were never staged. The 2022 leg of the observed-speed bound (\$3.3 billion of the \$11.7 billion) therefore stands as that vintage's discipline. Appendix~\ref{app:composition} collects the formal composition comparison behind these bounds: sample versus book by agency, vintage, coupon, and term, with a cross-agency stability check on the dimensions the book's CUSIP-level disclosure cannot supply.

Sections~\ref{sec:abm} and~\ref{sec:hazard} establish that contrast under the synthetic household population this paper specifies. Section~\ref{sec:robustness}'s cross-design pair leaves the attribution asymmetric, and it inherits Section~\ref{sec:robustness-crossdesign}'s own hedge. The hazard framework's recovery survives a fully synthetic population (106.0\% versus 107.0\%), evidence that its survival structure, not its data source, carries the aggregate level. But the recalibrated real-covariate ABM's 59.3\%, and its 60.2\% fifty-seed mean, cross the pre-committed threshold above which the contrast cannot, on that test alone, be cleanly attributed to modeling paradigm rather than to calibration and data source. The path-level timing question, meanwhile, remains open for every estimator. That result cuts both ways. Recovering 106\% on a fully synthetic population shows that the calibration constants, not real loan-level structure, deliver the level, which is consistent with the rate-inelastic decomposition above. On the paradigm question itself I make no claim: the reading was pre-committed to a threshold, the threshold was crossed, and Section~\ref{sec:robustness-crossdesign} withdraws it. The decomposition this paper does deliver never rested on it.

\clearpage
\bibliographystyle{plainnat}
\bibliography{references}

\clearpage
\appendix

\noindent\textbf{Online Appendix.} The appendices that follow are for online publication. The run ledger of superseded figures and the pre-committed-rule adjudication record form a standalone document in the replication package (\texttt{replication\_appendices}); its cross-references resolve against this manuscript.

\section{Agent-Based Model: Specification}\label{sec:method-abm}

I adopt an agent-based approach because the phenomenon under study is heterogeneous. Households differ in coupon, loan age, and exposure to search and sentiment frictions, and these differences interact non-linearly with the level of interest rates. \citet{geanakoplos2012} argue for, and demonstrate the feasibility of, agent-based models of mortgage prepayment in environments where behavioral and search frictions vary across the population and over time. That property is directly relevant here, since the frictions I model (housing search costs and consumer sentiment) are themselves time-varying. The broader methodological case for agent-based computational economics in settings with meaningful heterogeneity is set out in \citet{tesfatsion2006}.

The ABM simulates a population of N = 10,000 heterogeneous households. Initial agent endowments (income, home value) are drawn from log-normal distributions calibrated to Federal Reserve Economic Data (FRED) medians. In the production multi-vintage specification of Appendix~\ref{sec:method-frictions}, each household's mortgage coupon and origination date are drawn from the Federal Reserve's actual holdings composition: the live CUSIP-level SOMA parse of Section~\ref{sec:abm-results}, across seven coupon cohorts. A single-cohort variant, used for exposition and mechanism analysis, instead endows every agent with a 30-year fixed-rate mortgage at a 3.0\% coupon aged 60 months at the simulation start. Where the text below discusses a single representative coupon, it refers to this pedagogical case, not to the production population.

Section~\ref{sec:robustness} reports a scale-sensitivity replication showing that the gap in explanatory power is not an artifact of sample size. Whether it reflects modeling paradigm or data source is taken up in Section~\ref{sec:robustness-crossdesign}.

The decision to relocate weighs the utility of moving against the aggregate financial penalty of doing so. The monthly payment follows the standard annuity formula
\begin{equation}\label{eq:annuity}
M = P\,\frac{r\,(1+r)^{n}}{(1+r)^{n}-1},
\end{equation}
where $P$ is the outstanding principal, $r$ is the monthly interest rate, and $n$ is the remaining term in months. The mobility penalty evaluates \eqref{eq:annuity} twice. The first evaluation applies the borrower's locked-in note rate to the original principal and term, which gives the current payment; by the amortization identity, that payment equals the payment on the outstanding balance at the note rate over the remaining term. The second applies the prevailing market rate to the outstanding balance over a fresh full original term, which gives the payment the borrower would take on by refinancing. The replacement loan re-amortizes---a term-preference persistence the code implements deliberately. The penalty scales with the difference between the two. A native short-term variant of the penalty, active only for terms at or below 240 months, floors payment decreases at zero and is specified with the full decision rule in Appendix~\ref{app:abmspec}. Here $r$ is expressed in decimal monthly terms: the annual note rate divided by twelve, so that a 3.0\% coupon gives $r = 0.0025$. It is a rate \emph{level}, distinct in both role and units from the annualized rate \emph{gaps} $g_{i,t}$ (in 100 bp units) and $\tilde g_{i,t} = g_{i,t}/100$ (decimal) that drive the hazard framework of Section~\ref{sec:pathb}. No rate symbol is reused across the two. I evaluate the population's move/stay decisions household-by-household at each point of a three-dimensional grid of market rate, dynamic friction, and rate velocity (and, in the multi-vintage specification, across seven coupon cohorts), then aggregate them to a pool-level annualized Conditional Prepayment Rate (CPR) surface. CPR is a pool-level convention; the household-level object is an annualized move intensity, consumed by the accounting layer at one-twelfth per month (Appendix~\ref{app:abmspec}). That surface is precomputed once and then interpolated month-by-month against the live population of simulated market conditions, which cuts the cost of the full grid by roughly two orders of magnitude and makes fifty-seed Monte Carlo estimation feasible. The surface stores each node's direct household-by-household evaluation, so interpolation reproduces the micro-simulation exactly at grid points. The operative approximation is the trilinear segment between 50 bp-spaced nodes. A spot-check validates that approximation directly: after an on-grid parity gate reproduced the committed surface values, direct micro-simulation and interpolation at worst-case off-grid midpoints differ by at most 0.23 points of CPR away from the freeze threshold, with the Danish column exact. The one systematic difference is the wait-and-see freeze itself. The engine applies it as a step at the 150 bp trigger, while the surface renders it as a ramp across the adjacent velocity nodes. That gap binds in four realized months (August--November 2022; SOMA-weighted mean absolute difference 0.22 points of CPR over the window). These are surface semantics that the production figures share, and the binary-step leg of Appendix~\ref{sec:robustness-extensions}'s freeze sweep prices their trapped-liquidity consequence directly. Figure~\ref{fig:cprsurface} renders one slice of that surface (the 3.0\% cohort of the expository variant) with the realized QT window located on it.

\begin{figure}[H]
\centering
\includegraphics[width=0.85\textwidth]{fig8_cpr_surface}
\caption{The precomputed CPR surface: the production surface's slice at the 3.0\% thirty-year SOMA cohort, at its parsed seasoning. Left: the U.S. par-payoff surface over the market-rate $\times$ dynamic-friction grid (no-freeze regime), with the 42 QT-window months overlaid at their realized (rate, friction) coordinates; the window sits entirely in the deep out-of-the-money corner, near the baseline turnover floor. Right: rate slices at the friction grid point nearest the window mean. Velocity enters the grid only through the wait-and-see freeze of Appendix~\ref{sec:method-frictions} (20\% of otherwise-mobile households held in place above a 150 bp trailing six-month rise), so the two regimes shown are its entire velocity structure; the Berger-recalibrated Danish counterfactual (Section~\ref{sec:abm-danish}) is near-flat, moving roughly a quarter point across the grid. The U.S. slice is identical across the frozen production surface and the current code head; the Danish curve is the Berger production specification.}
\label{fig:cprsurface}
\end{figure}

\begin{figure}[H]
\centering
\includegraphics[width=0.9\textwidth]{fig10_abm_cpr_paths}
\caption{Production ABM simulated CPR against the empirical SOMA back-out over the QT window (frozen fold-in run, seed 42). The simulated path runs at roughly twice the empirical level and does not track its month-to-month shape; the annotated means, lag-0 $r$, and $R^2$ are the manifest values quoted in the text. This is the level-and-shape mismatch behind the negative path statistics; the cross-correlation structure across lags is in Figure~\ref{fig:ccf}.}
\label{fig:abmpaths}
\end{figure}

\begin{figure}[H]
\centering
\includegraphics[width=0.82\textwidth]{monte_carlo_trapped_liquidity}
\caption{Monte Carlo population-draw distribution of ABM trapped liquidity (50 seeds, fold-in specification; CPR surface rebuilt per draw). The seed-level dispersion (SD \$24.5 billion) dwarfs the distance between the frozen seed-42 production draw (\$91.0 billion, 34th percentile) and the seed mean (\$103.7 billion); the empirical benchmark of \$764.7 billion lies far outside the distribution's support, which is the visual form of the ABM's 10--17\% recovery band.}
\label{fig:mc}
\end{figure}

The simulation contrasts two institutional payoff rules. Under U.S. par-payoff, the cost to retire the mortgage is strictly equal to the outstanding principal. This is the negatively convex case: the investor's upside is capped, while extension risk is bounded only by the scheduled-amortization profile---a 14.7-year weighted-average life at zero CPR against the no-shock 3.4-year life (Table~\ref{tab:wal}), an order-of-magnitude extension ceiling. Under Danish market-value repurchase, the payoff cost is the present value of the remaining scheduled payments discounted at the current 30-year primary mortgage rate (FRED MORTGAGE30US), capped at par. Two qualifications attach to this proxy. First, it is an option-free present value discounted at the primary rate, so it ignores the embedded par call and any option-adjusted spread; the primary rate also embeds the primary--secondary spread (guarantee fee, servicing, originator margin; \citealp{fuster2013}), whereas the bond-market discounting a true buyback would use sits below it. Both choices overstate the borrower's buyback discount as rates rise. These biases inflate the mechanism-substitution Danish CPR of Section~\ref{sec:abm-danish} (the since-superseded 47.1\%), but they are exactly immaterial to the Berger-recalibrated production finding. Under neither transplant anchor does that leg's prepayment hazard consume this pricing proxy: the Danish-level anchor imports the near-flat moving curve, and the production U.S.-intercept anchor evaluates the U.S. hazard at zero rate gap. A replication-package diagnostic that re-runs the Danish leg with the discount rate shifted down by 25--150 bp moves trapped liquidity by \$0.00 billion. It flips 2,470--92,495 of the 1,683,082 loan-month present-value classifications\footnote{This Danish-leg classification universe of 1{,}683{,}082 loan-months is 42 fewer than the 1{,}683{,}124 evaluated in the U.S.\ production hazard panel (Section~\ref{sec:pathb}): the two figures count active loan-months in the two counterfactual legs, whose prepayment paths differ marginally even under the shared imported elasticity,} that no downstream quantity in that leg consumes. Second, what the Danish rule neutralizes is extension risk for the bondholder \citep{berger2026}: the borrower retains the par call, so the rate-decline side keeps its price cap, and I use ``convexity-neutral'' below strictly in this rate-rise sense. Under the Danish rule, as market rates rise above 3.0\%, the bond's present value falls, which lets the borrower extinguish the debt at a discount and guarantees continuous cash flow generation. I simulate the Danish counterfactual balance forward on a dynamic, declining-balance basis: each month's elevated Danish roll-off applies to the already-shrunk balance carried over from the prior month, so the simulated payoff path compounds correctly over the multi-year QT window.

\section{Macroeconomic Friction and Structural Extensions}\label{sec:method-frictions}

To quantify the behavioral friction gap in standard models, I scale base transaction friction using two FRED series. (1) I measure search friction as the shortfall in Active Listing Count (ACTLISCOUUS) against a 2017--2019 baseline, the last stable pre-pandemic period before the 2020/21 rate-driven frenzy and the 2022 inventory collapse. (2) I index psychological friction from University of Michigan Consumer Sentiment (UMCSENT) against a baseline of 85.0, scoring readings below that line as incremental friction at +5 bp per point short, capped at +150 bp during acute sentiment stress. The result is a dynamic friction parameter---a one-time moving transaction cost expressed as a share of the household's current home value---that scales between 8.23\% and 9.96\% of home value over the 42 active QT months.

I add three behavioral extensions to the household decision function. No parameter in any of them is tuned to match the empirical benchmark. Their grounding is nonetheless uneven: the first two are parameterized from external sources, the third is an author calibration. The first is a debt-to-income (DTI) constraint that vetoes any move in which the new housing payment would exceed 43\% of gross monthly income. That threshold is a front-end heuristic motivated by the pre-2021 General QM debt-to-income limit (12 C.F.R. \S{} 1026.43(e)(2)(vi)). The regulatory threshold was a back-end total-debt ratio and was superseded by the price-based General QM definition in 2021, so I use it as a behavioral anchor. The model carries no total-debt data against which to test the threshold (Section~\ref{sec:limitations}). The second is asymmetric loss aversion that scales perceived payment increases by 2.25$\times$ while leaving decreases unscaled (loss aversion: \citealp{kahneman1979}; the 2.25 coefficient: \citealp{tversky1992}). \citepos{andersen2022} housing-market estimate of roughly 2.5 is close. Any coefficient in this range triggers the same turnover-floor recalibration described below, so the artifact interpretation of Section~\ref{sec:abm-results} does not hinge on the choice. The third is a wait-and-see freeze that holds 20\% of otherwise-mobile households in place whenever the trailing six-month change in mortgage rates exceeds 150 bp. It is modeled as a fixed per-agent trait, so the CPR surface remains deterministic. The 20\% share and 150 bp trigger are author calibrations with no external source, as are the sentiment-to-friction mapping and the base transaction-friction level of the preceding paragraph: the UMCSENT and ACTLISCOUUS inputs are data, but the $+5$ bp-per-point mapping, its $+150$ bp cap, and the base friction level are mine. Section~\ref{sec:robustness-crossdesign} tests an externally parameterized variant of the mobility-desire scale and the freeze share, drawn from the same source as Path B's elasticity.

Loss aversion steepens the penalty function facing every household. The model's mobility-desire parameter is independently calibrated to reproduce a real-world baseline turnover floor of 4--5\% CPR at 8\% market rates. That floor is the turnover speed observed on deeply out-of-the-money discount cohorts during 2023--2024, when refinancing incentives were absent and remaining prepayment reflected death, divorce, and forced relocation. I therefore hold that empirical floor fixed and re-run the calibration under the steeper penalty function. The re-run requires nearly tripling the mobility-desire scale parameter, and this mechanical recalibration is the principal driver of the results in Section~\ref{sec:abm-results}.

Whenever a specification change alters the penalty or covariate structure facing every household, the mobility-desire scale, and only that parameter, is re-derived so that the simulation continues to reproduce the empirically observed 4--5\% baseline turnover floor. The same rule generates the recalibrated variant of the cross-design test (Section~\ref{sec:robustness-crossdesign}), whose frozen variant is defined by suspending it; Appendix~\ref{sec:robustness-extensions} reports how much work this single choice does. Because the 4--5\% floor is imposed by construction, the ABM's agreement with observed discount-cohort turnover at deep rate gaps is an input, not a finding. The model earns no evidential credit on that dimension; its evidential exposure lies entirely on the dimensions the anchor leaves free, the dollar recovery and the timing of the CPR path. The anchor is also insulated from the behavioral gates by construction: the calibration point is evaluated at zero rate velocity, where the wait-and-see freeze cannot bind, so no freeze parameterization can move the floor (Appendix~\ref{sec:robustness-extensions} verifies the calibrated scale and floor CPR are bit-identical at every freeze share swept). What the anchor purchases, and what dropping it costs, is measured directly in Section~\ref{sec:robustness-crossdesign}: the externally parameterized variant that removes the floor overshoots the benchmark by half while predicting essentially zero turnover at QT-window rate gaps.

I added an income-scaled curtailment channel as an extension specified ex ante. The ABM's monthly CPR over-prediction correlates with year-over-year growth in real disposable income (FRED series DSPIC96) at $r = -0.487$, and the correlation concentrates in the worst over-prediction months (June--September 2022). That diagnostic led me to predict ex ante that curtailment, a strictly additive channel, would worsen the aggregate match. I implemented curtailment as an income-scaled addition to simulated roll-off, mapping real income growth linearly from zero curtailment at $-2$\% growth to a full-scale annualized rate of 1.2\% CPR at +4\% growth, consistent with GSE daily prepayment report estimates.

Finally, as a diagnostic, I replaced the single representative 3.0\% coupon assumed elsewhere with a multi-vintage composition. Bucketing SOMA holdings into seven coupon cohorts yields a weighted-average coupon of approximately 2.55\% on the 30-year book, with the 2.0\% and 2.5\% buckets accounting for roughly 69\% of 30-year face value. A lower coupon implies a larger gap to prevailing market rates, so my ex ante hypothesis was that this composition would increase predicted trapped liquidity toward the empirical benchmark. This pass excludes fifteen-year MBS (approximately 9\% of SOMA face value), and cohort origination dates are approximated from parsed maturity strings under a 360-month term assumption. Section~\ref{sec:abm-results} describes the production replacement of this diagnostic with a live, CUSIP-level pull of SOMA holdings.

\section{ABM Decision Rule and Parameter Provenance}\label{app:abmspec}
% Sources: abm/abm_lockin_simulation.py ("S:"), abm/fed_mbs_extension_risk.py ("F:"),
% abm/cross_design_test.py, abm/abm_external_gates.py. Line pins refer to the committed code.

This appendix specifies the ABM's move decision and every behavioral parameter in the detail needed to reimplement the model without consulting the code. Each equation and value is pinned to its source line in the replication package.

\paragraph{The move decision.}
Let $A(P,r,n) = P\,(r/12)\big/\big(1-(1+r/12)^{-n}\big)$ be the monthly payment on principal $P$ at annual rate $r$ over $n$ months. % S:130-136
Household $i$ carries gross income $y_i$ (\$/year), home value $H_i$ (\$), original principal $P_i = 0.8\,H_i$ (a fixed 80\% origination LTV in the synthetic population), note rate $c_i$, original term $n_i$ in months, and outstanding balance $K_i$ (closed form at the cohort's seasoning; the U.S. par payoff). % S:141-155, S:314
The replacement-loan payment is $M^{\mathrm{new}}_{i,t} = A(K_i, m_t, n_i)$: the par payoff financed at the market rate $m_t$ over a fresh \emph{full original term}, not the remaining term (term-preference persistence). % S:222-225, S:366-381
With $\Delta_{i,t} = M^{\mathrm{new}}_{i,t} - A(P_i, c_i, n_i)$, whose second term is the current payment (by the amortization identity, equal to the payment on $K_i$ at $c_i$ over the remaining term), the loss-aversion-scaled penalty, in \$/month, is term-dependent:
\begin{equation}\label{eq:abmpenalty}
\pi_{i,t} =
\begin{cases}
\lambda\,[\Delta_{i,t}]_{+} + [\Delta_{i,t}]_{-} & n_i > 240 \text{ (30-year book)},\\
\lambda\,\big[M^{\mathrm{new}}_{i,t} - A(K_i, c_i, n_i)\big]_{+} & n_i \le 240 \text{ (native short-term gate)},
\end{cases}
\end{equation}
% S:383-409; the 240-month switch is S:408
with $\lambda = 2.25$, $[x]_{+} = \max(x,0)$, and $[x]_{-} = \min(x,0)$. Payment increases are scaled by $\lambda$. On the 30-year book, decreases pass through unscaled as gains, while the short-term gate floors them at zero to isolate the pure rate-lock cost. The frozen production headline draws 15-year voluntary CPR from the same-coupon 30-year surface, so the short-term branch is operative only in the Berger-run calibration stage (Section~\ref{sec:abm-danish}).
The effective transaction rate is $\tau_{i,t} = \mathrm{clip}\big(\phi_t + (\tau_i - 0.07),\,0.02,\,0.12\big)$: the agent's idiosyncratic draw enters as an offset around the 7\% base mean, then is re-clipped after the macro friction $\phi_t$ shifts it. % S:241-243, S:459-461
Household $i$ moves in month $t$ iff
\begin{equation}\label{eq:abmmove}
\underbrace{\frac{M^{\mathrm{new}}_{i,t}}{y_i/12} \le 0.43}_{\text{DTI veto}}
\;\wedge\;
\underbrace{D_i > S\,\pi_{i,t} + \tau_{i,t} H_i}_{\text{cost--benefit}}
\;\wedge\;
\underbrace{\neg\big(v_t > 0.015 \,\wedge\, u_i < 0.20\big)}_{\text{wait-and-see freeze}},
\end{equation}
% Gates: DTI S:455-457, cost-benefit S:458-463, freeze S:464-467, conjunction S:468
where $D_i \sim \mathrm{Exponential}(\theta)$ is the mobility-desire draw in dollars, $S = 60$ months is the expected-stay horizon converting the monthly penalty into a dollar stock, and $v_t$ is the trailing six-month mortgage-rate change (decimal). The remaining term, $u_i \sim U[0,1]$, is a fixed per-agent trait --- the code's ``patience'' draw, consumed only by the freeze gate. It is not a time-preference parameter. % S:48, S:192, S:302-309
All inequalities are as coded: the DTI gate passes at exactly 43\%, the cost--benefit gate requires strict exceedance, and the freeze binds only strictly above 150 bp. At an evaluation point, the share of a pool passing \eqref{eq:abmmove} is read as that point's annualized CPR --- the calibration anchor is an annual floor --- and the macro accounting layer consumes it at $\mathrm{CPR}/12$ per month. % F:957

\paragraph{Parameter provenance.}
Table~\ref{tab:abmparams} collects every behavioral parameter. None is tuned to the \$764.7 billion benchmark. The mobility-desire scale $\theta$ is the one parameter re-derived under the floor rule, and only ever that one.

\begin{table}[H]
\centering
\footnotesize
\begin{threeparttable}
\caption{ABM behavioral parameters: value, units, distribution, and provenance.}
\label{tab:abmparams}
\begin{tabular}{@{}p{3.3cm}p{5.9cm}p{1.6cm}p{4.0cm}@{}}
\toprule
Parameter & Value / distribution & Units & Provenance \\
\midrule
Population & $N = 10{,}000$; engine seed 42 & households & author choice \\
Income $y_i$ & log-normal, $\sigma = 0.45$, median from FRED MEHOINUSA672N & \$/year & FRED-calibrated draw \\
Home value $H_i$ & log-normal, $\sigma = 0.35$, median from FRED MSPUS & \$ & FRED-calibrated draw \\
Origination LTV & 0.80, fixed & --- & author choice \\
Mobility desire $D_i$ & exponential, scale $\theta$ & \$ & floor recalibration (last row) \\
\quad $\theta$, production (frozen) & 43{,}883 & \$ & floor anchor, synthetic population \\
\quad $\theta$, cross-design recalibrated & 36{,}086 & \$ & floor anchor, real Freddie covariates (Section~\ref{sec:robustness-crossdesign}) \\
\quad $\theta$, external variant & 4{,}472 & \$ & \citepos{liebersohn2024} zero-gap 5.9\%/yr moving level, no floor (Section~\ref{sec:robustness-crossdesign}) \\
Expected-stay horizon $S$ & 60 & months & author calibration \\
Loss aversion $\lambda$ & 2.25 & --- & \citet{tversky1992} \\
DTI wall & 0.43, front-end & --- & 12 C.F.R. \S{} 1026.43(e)(2)(vi), behavioral anchor \\
Freeze share; trigger & 0.20 of agents; $v_t > 0.015$ over six months & ---; decimal & author calibration \\
Patience trait $u_i$ & uniform $[0,1]$, fixed per agent & --- & implementation device (freeze membership) \\
Transaction-cost draw $\tau_i$ & normal$(0.07,\,0.015)$, clipped $[0.02, 0.12]$ & share of $H_i$ & author calibration \\
Dynamic friction $\phi_t$ & $0.07$ base $+$ search penalty (listing shortfall vs.\ 2017--19 baseline, proportional, cap $+0.02$) $+$ sentiment penalty (5 bp per UMCSENT point below 85, cap $+0.015$); realized 8.23--9.96\% over the 42 QT months & share of $H_i$ & FRED ACTLISCOUUS, UMCSENT; mappings and base level are author calibrations \\
Curtailment & linear in real-income YoY growth: 0 at $-2\%$ to 1.2\% CPR at $+4\%$; $/12$ monthly, additive, macro layer only & annualized CPR & GSE daily-prepayment report estimates; FRED DSPIC96 (Appendix~\ref{sec:method-frictions}) \\
Turnover floor & binary search $\theta \in [10^{3}, 5\times10^{5}]$, at most 30 iterations, until CPR at an 8\% market rate under the U.S. rule, base friction, zero velocity lands in $[4\%, 5\%]$, evaluated on the dominant-weight 30-year SOMA cohort & annual CPR & 2023--24 discount-cohort turnover (Appendix~\ref{sec:method-frictions}) \\
\bottomrule
\end{tabular}
\begin{tablenotes}
\footnotesize
\item Committed artifacts carry the unrounded binary-search scales (43{,}882.8125; 36{,}085.9375; 4{,}472.05). The module default $\theta = 12{,}500$ is a placeholder overwritten by calibration at every run. Code pins: draws \texttt{abm\_lockin\_simulation.py} 296--309; constants 38--69; floor search 649--685; friction \texttt{fed\_mbs\_extension\_risk.py} 66--70, 294--334; curtailment 85--87, 141--155; variant scales \texttt{cross\_design\_test.py} 81--111 and \texttt{abm\_external\_gates.py} 88--89.
\end{tablenotes}
\end{threeparttable}
\end{table}

\paragraph{Order of operations.}
Within a simulated month, the three gates of \eqref{eq:abmmove} are evaluated on the surviving population and combined conjunctively. They are side-effect-free, so their order cannot change the result; the code evaluates them in vectorized form in the order DTI, cost--benefit, freeze. % S:449-468
In cohort-path (burnout) mode, movers are removed permanently before the next month's evaluation, so the turnover tail sets a natural floor on the depleted pool. % S:470-518
Two named components are not in-month gates. The 4--5\% floor enters once, before simulation, through the desire-scale search of Table~\ref{tab:abmparams}. That search is re-run, for that parameter only, whenever a specification change alters the penalty or covariate structure facing every household. Because the anchor is evaluated at zero velocity, no freeze parameterization can move it (Appendix~\ref{sec:method-frictions}). % S:649-685, S:520-524
Curtailment never touches the household decision: it is an additive term in the macro layer's simulated roll-off. Nor does production re-solve the decision rule monthly. The CPR surface is precomputed on the rate $\times$ friction $\times$ velocity grid, with all gates applied at each node, and then interpolated month-by-month; Appendix~\ref{sec:method-abm} discloses the freeze's step-versus-ramp semantics. % F:873-895, S:526-548

This paper earlier described the mobility-desire distribution as ``log-normal''; it is exponential, on a dollar scale, and the correction is recorded in the replication package's run ledger.

\section{Path B Parameterization and Delinquency Transitions}\label{app:params}

\begin{table}[H]
\centering
\footnotesize
\begin{threeparttable}
\caption{Objects and notation for the two hazard specifications. Path A's macro coefficients are reported in Section~\ref{sec:patha}, and every Path B parameter with its provenance is collected in Table~\ref{tab:params}.}
\label{tab:hazard-notation}
\begin{tabular}{@{}lllc@{}}
\toprule
Symbol & Object & Units & Path \\
\midrule
\multicolumn{4}{@{}l}{\emph{Sample and exposure}}\\
$s$ & stratum (296 four-way cells) & -- & A \\
$(s,t)$ & unit of observation (stratum, month) & -- & A \\
$i$ & loan (75,000 sampled; 40,234 active) & -- & B \\
$(i,t)$ & unit of observation (loan, month) & -- & B \\
$y_{s,t}$ & prepaid unpaid principal balance & dollars & A \\
$n_{s,t}$ & UPB exposure (log offset) & dollars & A \\
\midrule
\multicolumn{4}{@{}l}{\emph{Link and baseline}}\\
$\alpha_{s}$ & stratum fixed effect (295 enter) & log-rate & A \\
$f(a_{s,t})$ & loan-age seasoning spline & log-rate & A \\
$h_{0}(a_{i,t})$ & PSA baseline hazard (Section~\ref{sec:pathb}) & monthly & B \\
$\underline{h}$ & baseline turnover floor & monthly SMM & B \\
$h^{\mathrm{vol}}$ & voluntary prepayment hazard, \eqref{eq:pathB} before the floor & monthly & B \\
$h^{\mathrm{prep}}_{i,t}$ & prepayment hazard, \eqref{eq:pathB} after the floor & monthly & B \\
$h^{\mathrm{def}}_{i,t}$ & competing default hazard, \eqref{eq:default} & monthly & B \\
\midrule
\multicolumn{4}{@{}l}{\emph{Covariates and coefficients}}\\
$g_{s,t}$ & rate gap & 100 bp units & A \\
$b_{s,t}$ & stratum burnout state & demeaned share & A \\
$\phi_t$ & dynamic friction index & index & A \\
$\beta_g,\ \beta_b,\ \beta_\phi$ & rate-gap, burnout, friction coeff. & log-rate & A \\
$g_{i,t}$ & rate gap (negative when locked in) & 100 bp units & B \\
$\tilde g_{i,t}=g_{i,t}/100$ & rate gap in decimal units & decimal & B \\
$\sigma_{i,t}=\max\{0,-\tilde g_{i,t}\}$ & rate stress & decimal & B \\
$z^{\mathrm{FICO}}_{i},\ z^{\mathrm{LTV}}_{i}$ & standardized static covariates & SD & B \\
$B_{s(i),t}\in[0,1]$ & cohort-level burnout state & cumulative share & B \\
$\beta_1,\ \beta_F,\ \beta_L,\ \beta_B$ & gap, FICO, LTV, burnout coeff. & log-hazard & B \\
$\gamma_r,\ \gamma_F,\ \gamma_L$ & default-hazard coefficients & log-hazard & B \\
\midrule
\multicolumn{4}{@{}l}{\emph{Convention and transport parameters}}\\
$\mu$ & moving-share bracket (production $\mu = 1$) & share & B \\
$\omega$ & floor-form mixture weight (production $\omega = 0$) & share & B \\
$\underline{h}_t = 4\%\,(1+\kappa z_t)$ & within-run cyclical floor variant & monthly SMM & B \\
$\eta = S_{\mathrm{pre}}^{\xi}$ & survival-selection attenuation & factor & B \\
$S_{\mathrm{pre}}$ & pre-window survival share & share & B \\
$\xi$ & gamma-frailty parameter (not estimable here) & -- & B \\
\bottomrule
\end{tabular}
\begin{tablenotes}[flushleft]\scriptsize
\item Case and subscript separate the two paths throughout, and no coefficient symbol is reused across them. That guarantee is about coefficients: the convention and transport parameters in the last block bracket reported results rather than entering either fit, and none of them shares a glyph with the stratum index $s$ or with each other. Path A's burnout coefficient $\beta_b$ (on $b_{s,t}$, estimated on the panel) and Path B's $\beta_B$ (on the cohort state $B_{s(i),t}$, a calibrated prior) are distinct objects carrying distinct symbols, as are Path A's friction coefficient $\beta_\phi$ (on the friction index $\phi_t$) and Path B's FICO coefficient $\beta_F$. The rate-gap elasticities $\beta_g$ (Path A, log-rate, standardized) and $\beta_1$ (Path B, log-hazard, per 100 bp) are on different scales; the paper compares them on sign, not magnitude (Section~\ref{sec:hazard-interp}).
\end{tablenotes}
\end{threeparttable}
\end{table}

{\footnotesize
\begin{longtable}{@{}p{1.5cm}p{6.2cm}p{6.2cm}@{}}
\caption{Reading key for words this paper reuses across designs. Each row gives the sense an unqualified occurrence carries and the qualifiers that mark every other sense in use. Symbols are in Table~\ref{tab:hazard-notation}; no row states a result, and every distinction below is made where it is used.}
\label{tab:terminology}\\
\toprule
Word & Unqualified, it denotes & Marked senses it must not be read as \\
\midrule
\endfirsthead
\toprule
Word & Unqualified, it denotes & Marked senses it must not be read as \\
\midrule
\endhead
stratum & Path A's four-way panel cell, vintage $\times$ coupon $\times$ FICO $\times$ LTV: 296 observed, 295 fixed effects, and the cluster unit that carries the floor reads (31 occupied off window, 226 in window) & the draw stratum: the same key with the LTV axis dropped, 130 of them and 25.8 effective, the unit of Path B's loan bootstrap. Both are resampled under the name ``stratum-cluster''; the path decides which \\
cell & nothing on its own --- the partition is always named, and the counts do not interchange & a panel stratum (296 four-way; 130 three-way for the draw); a stratum-month observation; the standardization cell of Section~\ref{sec:patha}, the panel key with the coupon axis removed (22 cells, 29 with 12-month age bands, 18 in common); a point of the 21-cell calibration box (Table~\ref{tab:floorband}: seven floors $\times$ three elasticities); a point of the twenty-seven-cell off-window floor $\times$ elasticity-band grid (run \texttt{oos\_identification}); a joint agency $\times$ vintage cell of the SOMA book; a printed table cell. ``Every cell positive'' is claimed once over book subgroups and once over the off-window grid, and those are two different claims \\
band & the imported rate-response range, 5.5\%--7.7\% with 6.5\% adopted --- a calibration interval, not sampling error & the 4--5\% turnover band, which is a floor \emph{level} range; the clean band, the admissible off-window floor reads; the seed band, the $\pm 1$~SD band and the moving-block bands, which are sampling dispersion; the corroborating band, a pre-committed acceptance range; age bands and single-band vintages, which are bucket edges \\
anchor & in the hazard sections, the floor level a read is taken at: in-window versus off-window, the mid-grid anchor being the middle of the three depth cuts fixed before the reads & the transplant construction of Section~\ref{sec:abm-danish}, U.S.-intercept versus Danish-level; an outside estimate quoted beside the band; the ABM's mobility-scale anchor, set by the floor search; anchor-gated, meaning checked against the committed SOMA marginals \\
basis & the accounting basis: shared, standalone, full-book, composed (Table~\ref{tab:bases}). ``Basis-invariant'' means invariant across those four and across nothing else & the coupon convention behind a CPR back-out --- ABM basis, hazard basis, note-rate basis --- which does move the printed levels; the weighting behind a share (count, origination balance, surviving balance); basis points; the seasoning spline's power basis \\
floor & the baseline turnover floor $\underline{h}$, a level in annual CPR & the floor \emph{form}, the rule combining it with the voluntary hazard, hard maximum versus additive competing risks --- an axis, not a level; the floor \emph{read}, the measurement of that level off discount cohorts; the scheduled-amortization floor the back-out's non-negativity clip pins against \\
form & the floor form above; ``form-conditional'' is always this axis & the rendering of an estimator, microsimulation form versus cohort-hazard form of the same framework; functional form of a hazard in general; decimal form and CPR form, where only a restatement of one object is meant \\
draw & the stratified 75{,}000-loan sample Path B simulates & a bootstrap or permutation replicate, always plural and counted; a random variate inside a run, such as the patience draw or the default draw; a seed-level realization of the ABM, the frozen draw; draw seeds, which are simulation seeds and not resamples of that sample \\
null & the $\beta_1 = 0$ rate-inelastic leg: a paired simulation run scored in dollars & a null hypothesis --- the $H_0$ tests of Appendix~\ref{app:patha-detail}, and the mean-only null of this paper's $R^2$ convention. A leg that clears a threshold is not a non-rejection \\
\bottomrule
\end{longtable}
}

\noindent{\footnotesize \emph{Notes to Table~\ref{tab:terminology}.} Rows are ordered by how badly a conflation would mislead, most damaging first. ``Cluster'' is deliberately absent: it carries one sense here, the resampling unit of a cluster-robust bootstrap, and what varies is the unit, not the word --- the four-way panel stratum, the three-way draw stratum, the six calendar months of the off-window leg, and the 137 cells where those two intersect. The stratum and cell rows fix each.\par}

Table~\ref{tab:params} collects every Path B parameter. All values are fixed a priori in the configuration, and none is estimated on, or tuned to, the \$764.7 billion benchmark. One parameter is nonetheless measured inside the evaluation window, and benchmark-independence should not be read as window-independence. The 4\% baseline turnover floor is anchored to turnover observed on discount cohorts during 2023--2024 (Appendix~\ref{sec:method-frictions} and Section~\ref{sec:pathb}), so it is in-sample with respect to the QT window even though it never touches the benchmark dollar figure. That floor sets much of the aggregate level, binding in 36\% of loan-months. Section~\ref{sec:robustness-floor} sweeps it jointly with the elasticity band: the aggregate level and the lock-in marginal both move materially with the floor, within the bounded interval that section reports. The marginal's sign does not move, but under this window's rate configuration it could not (Section~\ref{sec:identification}). The covariate and default-hazard coefficients are round calibration priors. Both prepayment covariates carry signs opposite the canonical credit-risk direction (higher FICO $\Rightarrow$ better refinancing access $\Rightarrow$ faster prepayment; higher LTV $\Rightarrow$ less equity $\Rightarrow$ slower prepayment). In the deeply out-of-the-money 2022--2025 window the refinancing channel is shut and prepayment is turnover-dominated, so credit covariates proxy mobility, and high-LTV loans are disproportionately recent purchases that churn. A Path-A-style estimation on the same stratum-month panel yields point estimates consistent with the priors' direction; taken at face value, they are roughly double the priors' magnitudes. That estimation is a Poisson PML with vintage-by-coupon fixed effects, so the credit buckets are identified from cross-bucket variation, and the bucket-mean covariates are z-scored to the microsim pool's own moments. But the intervals reject neither zero nor the opposite signs: $\hat{\beta}_F = -0.39$ (four-way stratum-cluster bootstrap 95\% CI $[-1.51, +0.65]$) and $\hat{\beta}_L = +0.21$ ($[-0.59, +1.00]$). The estimation does not settle the sign question; it fails to contradict the priors while quantifying their sensitivity. The sensitivity brackets are correspondingly modest: a $\beta_F = \beta_L = 0$ rerun moves the recovery from 107.0\% to 108.3\%, and the panel-estimated pair moves it to 101.3\%. Those are in-window figures. At the off-window floor the same endpoints span 8.2 points and `modest' does not carry over to the level, while the marginal's covariate sensitivity shrinks to $+0.1$ points. Benchmark-independence is unaffected, since the estimation touches only the Freddie panel. The $\beta_B = 0$ rerun of Section~\ref{sec:pathb} quantifies the burnout channel the same way.


Delinquent loans transition monthly through the matrix of Table~\ref{tab:delinq}. I estimate it from consecutive stratum-month \emph{modal} servicer states in the Freddie cohort panel, exposure-UPB weighted, and commit it as an artifact in the replication package. The default-hazard draw and the transition matrix divide the labor. The draw moves loans \emph{into} the pipeline (Current $\to$ D30) and terminates already-delinquent loans it hits. The matrix governs month-to-month movement \emph{within} the pipeline. The matrix runs after the draw each month, so a loan entering D30 via the draw can also take a matrix transition in its entry month. The \$0.39 billion extreme-matrix test below bounds that overlap along with everything else. Only the delinquency rows are used in simulation. Terminal states are absorbing by construction. The forbearance state is unpopulated in the estimation panel and never entered in production runs.
I publish the full crosstab. D30 row (21 transitions, \$58.3M row exposure): 13 cures carrying 95.6\% of the exposure weight, two D30 continuations (0.8\%), five D60 progressions (3.1\%), and one Defaulted transition (0.4\%); the printed 0.004 is the exposure share of that single event, whose count share would be $1/21 = 0.048$. D60 row (9 transitions, \$3.1M): six cures carrying 56.8\% and three D90+ progressions carrying 43.2\%, against count shares of $6/9$ and $3/9$. That divergence between count and exposure shares means a few larger-balance stratum-months drive the matrix, which is what an exposure-weighted estimator is designed to do. D90+ row (38 transitions, \$76.2M): seven cures (30.3\%), one D30 (0.8\%), thirty self-continuations (68.9\%), and zero Defaulted events; the ``--'' cells are observed zeros in a thin panel, not structural zeros, which is why serious delinquency shows no direct default exit while D30 shows one. Probabilities are printed to three decimals from these exposure shares and rows may not sum to exactly one. The matrix is accordingly reported for completeness, not precision: with cells this thin it describes the estimation panel and does not estimate a transition law. The paper's estimates rely on the materiality bound below, not on the matrix's point values.

The thinness is immaterial for the paper's estimates: replacing the entire matrix with either extreme --- all delinquent loans cure next month, or all default next month --- moves trapped liquidity by \$0.39 billion (0.05\% of the estimate). The matrix routes the 482 loans delinquent at the window start from the first month, mostly curing them, given the cure-heavy rows. Thereafter the month-end delinquency stock never exceeds 154 of the 40,234 active loans, and delinquent UPB never exceeds 1.1\% of the pool. Both extreme-matrix runs include the initial 482, so the bound covers them.

\begin{table}[H]
\centering
\small
\begin{threeparttable}
\caption{Exposure-weighted delinquency transition matrix (modal stratum-month states).}
\label{tab:delinq}
\begin{tabular}{@{}lrrrrr@{}}
\toprule
From $\backslash$ To & Current & D30 & D60 & D90+ & Defaulted \\
\midrule
D30 & 0.956 & 0.008 & 0.031 & -- & 0.004 \\
D60 & 0.568 & -- & -- & 0.432 & -- \\
D90+ & 0.303 & 0.008 & -- & 0.689 & -- \\
\bottomrule
\end{tabular}
\begin{tablenotes}[flushleft]\footnotesize
\item Probabilities are exposure-UPB shares of thin modal-state transition counts (D30 row: 21 transitions; D60: 9; D90+: 38; full crosstab in the text and replication artifact). ``--'' cells are observed zeros in a thin panel, not structural zeros; rows may not sum to exactly one at printed precision. Reported for completeness, not precision: the paper's estimates rely on the \$0.39 billion extreme-matrix materiality bound in the text, not on these point values.
\end{tablenotes}
\end{threeparttable}
\end{table}
Sampling design: the draw is not a probability sample of the SOMA book, and the allocation rule is why. A pool built at four times the target size takes an equal allocation from each quarterly Freddie origination file. The pool is then drawn down to 75,000 loans proportionally within stratum, and every retained loan carries weight one (\texttt{hazard/loan\_sample.py}). The draw's stratum is the three-way vintage $\times$ coupon $\times$ FICO cell --- the Path A estimation stratum with the LTV axis dropped --- which is why the draw carries 130 of them against the panel's 296. Nothing in that chain reweights toward the book. Each origination year from 2017 through 2021 therefore takes 20.0\% of the draw by count, so 2017--19 is 60.0\% of it against 6.0\% of book face. But that count basis is not the weighting the aggregation applies. The aggregation weights by surviving balance at the window open, and on that basis the picture is materially different: 2017--19 is 27.0\%, 2020 30.6\% and 2021 42.3\%, against book face shares of 6.0\%, 16.3\% and 43.9\%. So the 2021 cohort is close to matched, and the 2017--19 overweight is a factor of about four and a half, not the factor of ten the count basis suggests; origination balance, which compares originations to current face and is the wrong basis for this comparison, gives 55.5\%, 22.2\% and 22.3\%. What no weighting repairs is out-of-support mass. The draw spans 2017 through 2021 only, so pre-2017 (10.6\% of book face) and the 2022 cohort (23.1\%) carry no support at all: 33.7\% of face, which must be netted against the Ginnie share on the book's own joint cells, not added to it. The draw's 3.86\% unweighted mean note rate also sits above the book's 2.49\% face-weighted pass-through coupon. Both tilts move the rate-gap distribution the floor takes its maximum against under the production floor form, and this paper signs neither: Section~\ref{sec:robustness-hybrid}'s coupon reweight rescales the aggregate level and leaves the marginal unchanged, so it does not test a composition effect running through the floor's bind share.

Attrition accounting: of the 75,000 sampled loans, 34,734 prepaid and 32 defaulted before the QT window opens, so 40,234 are active at the June 2022 start (39,752 current, 482 delinquent). The production U.S. run evaluates 1,683,124 loan-months over the 42-month window, a mean of 40,074 active loans per month. The baseline turnover floor $\underline{h}$ binds in 36\% of those loan-months (Section~\ref{sec:pathb}).

The $\delta$-to-$\beta_1$ conversion of \eqref{eq:beta1} is insensitive to the turnover proxy behind it. In the small-$P_q$ limit the transform reduces to $\beta_1 = -\ln(1-\delta) = 0.0672$ at the central $\delta$. Over $P_q \in (0, 0.12]$---from arbitrarily low turnover to double the adopted proxy---the central $\beta_1$ ranges only from $0.0672$ to $0.0701$. Even a large error in $P_q$ therefore moves the conversion by a few percent.

\paragraph{Production step order, clamp, and parameters.}
Each month, both hazards are evaluated for every active loan. In the production configuration prepayment is fractional: each active loan's balance is drained deterministically by $h^{\mathrm{prep}}$ (a per-loan uniform draw for prepayment is a non-production discrete branch). Default instead turns on a uniform draw, which sends hits into the delinquency pipeline (current loans to D30, already-delinquent loans to default), where they transition through the Markov delinquency-resolution matrix of Appendix~\ref{app:params}. A proportional rescaling guards the case $h^{\mathrm{prep}} + h^{\mathrm{def}} > 1$. It is a safety clamp, not a competing-risks construction. In production it never binds: across all 1,683,124 evaluated loan-months per regime the maximum hazard sum is 0.016, two orders of magnitude below the threshold. At that level the exact exponential cause-allocation construction $1-\exp(-(\lambda_p+\lambda_d))$ with shares $\lambda_p/(\lambda_p+\lambda_d)$ differs from mine only at second order in the hazards. In the taxonomy of standard competing-risks formulations \citep{putter2007}, Path B is a discrete-time multi-state model parameterized entirely through cause-specific hazards. Equations~\eqref{eq:pathB} and \eqref{eq:default} give monthly event probabilities conditional on the loan remaining active---the estimand of discrete-time cause-specific regression, as implemented in open-source packages such as PyDTS \citep{meir2025}---and the delinquency pipeline of Appendix~\ref{app:params} extends the two-cause structure to a multi-state process with cures. No subdistribution (Fine--Gray) hazard \citep{fine1999} appears in the specification, and none is needed: that device places covariate effects directly on cumulative incidence functions, whereas the simulation obtains every cumulative-incidence functional (prepaid UPB, the delinquency stock, trapped liquidity) by forward composition of the cause-specific hazards. One caveat qualifies the mapping. In production, prepayment is a fractional balance-flow intensity, not a per-loan absorbing transition, so the cause-specific reading is exact for default and holds at the balance-aggregate level for prepayment. Prepaid balances settle to simulated SOMA cash flow in the month they occur; no settlement delay is modeled on this path.

\begin{table}[H]
\centering
\footnotesize
\setlength{\tabcolsep}{4pt}
\begin{threeparttable}
\caption{Path B parameters.}
\label{tab:params}
\begin{tabular}{@{}llll@{}}
\toprule
Parameter & Value & Units & Provenance \\
\midrule
$\beta_1$ (central) & $0.069$ & per 100 bp gap & \eqref{eq:beta1} at $\delta = 0.065$ \citep{liebersohn2024} \\
$\delta$ band & $0.055$--$0.077$ & quarterly decline & \citet{liebersohn2024} \\
$P_q$ & $0.06$ & quarterly slot & turnover proxy (conversion auxiliary; not a measured quarterly mobility rate) \\
$\beta_F$ (FICO, z) & $-0.15$ & per SD & calibration prior; panel estimate $-0.39$ (see text) \\
$\beta_L$ (LTV, z) & $+0.10$ & per SD & calibration prior; panel estimate $+0.21$ (see text) \\
$\beta_B$ (burnout) & $-0.5$ & per unit of $B \in [0,1]$ & calibration prior \\
$\underline{h}$ (floor) & $0.0034$ & monthly SMM & 4\% annual CPR (baseline turnover) \\
PSA speed & 100 & -- & industry seasoning convention \\
$h^{\mathrm{def}}_0$ & $3\times10^{-4}$ & monthly & calibration prior \\
$\gamma_r$ (rate stress) & $2.0$ & per decimal unit & calibration prior ($\approx$2\% per 100 bp) \\
$\gamma_F$ (FICO, z) & $-0.20$ & per SD & calibration prior \\
$\gamma_L$ (LTV, z) & $+0.25$ & per SD & calibration prior \\
\bottomrule
\end{tabular}
\end{threeparttable}
\end{table}

\paragraph{Cohort-level burnout.}
Burnout is deliberately cohort-level, not per-loan: a fractional burnout measure is not well-defined for a one-time, absorbing transition (Section~\ref{sec:abm-results}). Each stratum instead tracks its cumulative prepaid UPB as a share of its original UPB---the $B \in [0,1]$ state defined under \eqref{eq:pathB}---after all prepayment transitions in a given month. That state is broadcast to all of that stratum's remaining active loans before the next month's hazard is evaluated, and enters \eqref{eq:pathB} with the calibrated prior coefficient $\beta_B = -0.5$ (Appendix~\ref{app:params}). Path A's separately estimated burnout coefficient, statistically indistinguishable from zero (Section~\ref{sec:patha}), plays no role here. Because the simulated burnout state is endogenous, I test whether the channel is load-bearing with a $\beta_B = 0$ rerun on the production seed: trapped liquidity moves by \$6.58 billion ($-0.86$pp of benchmark, to 106.17\%), with the peak lag unchanged at $-3$. Two qualifications travel with that figure. It is measured at the in-window 4.0\% floor; at the 4.991\% off-window floor that the paper headlines, the same ablation is worth only \$3.48 billion ($-0.45$pp), because the floor's bind share rises from 36.3\% to 68.8\% of evaluated loan-months and the hard maximum makes burnout inert wherever it binds. And because the burnout state is endogenous to realized prepayment, the ablation moves the $\beta_1 = 0$ null as well as the central leg, so the identified marginal is not burnout-invariant: switching the coefficient off raises it from $+9.20$ to $+9.94$ points at the production floor and from $+5.57$ to $+6.51$ at the off-window floor. Both moves sit inside the $\pm 1$-point materiality convention, the second only barely, so I report them without resting anything on them. The FICO and LTV covariates receive the same treatment. A $\beta_F = \beta_L = 0$ rerun moves the recovery to 108.3\% ($+1.3$pp), and a rerun under coefficients estimated Path-A-style from the panel moves it to 101.3\% ($-5.7$pp), both at the in-window 4.0\% floor. At the 4.991\% off-window floor the same endpoints span more, not less: 102.5\% against 94.3\% standalone, an 8.2-point swing. The marginal they bracket nevertheless moves by only $+0.1$ points, against $+1.9$ in-window, because both legs shift together where the floor censors more. Appendix~\ref{app:params} reports the estimation, whose imprecise point estimates are consistent with the priors' turnover-regime signs without rejecting alternatives. Production retains the ex-ante priors---the estimates are too imprecise to license replacing a pre-specified calibration mid-stream---and I flag the covariate block as the specification's largest single sensitivity: the two endpoints correspond to 99.2\% and 92.2\% on the shared accounting basis (the common netting makes the conversion exact), a swing comparable in size to the lock-in marginal itself.

\paragraph{The time-varying Ginnie overlay.}
The static Ginnie bound is now extended to a time-varying overlay built from the published series the bound's two snapshots came from. I extract the full monthly agency CPR series (June 2017 through November 2025) from the December 2025 Global Markets Analysis Report and score the 20.4\% Ginnie face share by the observed Ginnie speed in \emph{both} Path B legs. That report's aggregation was revised in that issue to weighted-average unpaid principal balance, one consistent dollar-weighted methodology across the whole window. The overlay moves the shared-basis central from 97.9\% to 89.3\% and the null from 88.7\% to 81.9\% in the raw construction, both legs at the in-sample calibration point. It confines the identified marginal to the conventional share: $+7.3$ points, positive under every variant tested (total-CPR, voluntary-only, and placebo), scaling with the conventional face share almost exactly ($0.797$ against $0.796$). Composed with the off-window headline calibration, the overlay pair is 84.0\%/79.5\% central/null on the shared basis with a $+4.4$-point marginal, and the overlay marginal is pure conventional-share scaling ($0.797\times$) --- a change of estimand, a re-scoring onto the conventional sub-book, identical across the primary, voluntary-only, and placebo series. It is not a correction to the headline: the headline marginal therefore carries the conventional-extrapolation convention wherever quoted, with $+4.4$ points its overlay-scored member. The raw level shift overstates the Ginnie-specific effect, though. A placebo that scores the same 20.4\% share by the observed GSE-mean series moves the level by $+\$27.8$ billion on its own: the model's level calibration plus the full-universe-versus-SOMA composition difference, common to any observed-series overlay. The Ginnie-specific differential contribution is therefore $+\$38.5$ billion, 5.0\% of benchmark. That sits inside the static \$20--\$47 billion bound above, and is consistent with the window-mean Ginnie--Freddie differential of 2.1 points the series itself shows against the bound's 1.2--2.8-point snapshots. The placebo and Ginnie-specific components, $\$27.8 + \$38.5 = \$66.3$ billion, reconcile with the raw 8.6-point shared-basis shift on the \$764.7 billion benchmark, $\approx\$65.8$ billion, to within rounding of the 0.1-point recovery shares. The full-universe series if anything overstates the correction for SOMA's seasoned, low-coupon Ginnie holdings, the same signed direction as the static bound; and structural Ginnie inclusion remains out of scope for the modeling reason above: the overlay treats the buyout-inclusive Ginnie speed as data, which is what runoff accounting requires of a passthrough holder. This overlay and the age-standardized floor read of Section~\ref{sec:robustness-floor} cut the same way, and the joint cell has now been run under a pre-committed landing rule (run \texttt{b5\_joint\_cell}): composed, they return $+2.9$ points ($+\$22.4$ billion), within $0.1$ points of the proportional projection, and that measured point is the assembly's lower member. Pure conventional-share scaling is what a gap-inert Ginnie share gives. Scoring that share at a response instead (run \texttt{ginnie\_overlay\_attenuated}) prices it at a fixed proportion of the conventional share's per-unit-face wedge (transplanted, not estimated on Ginnie collateral), with the wedge fixed by requiring that a full-strength Ginnie response return the un-overlaid marginal. What the series measures is a level, and a voluntary-only one: across the window's 42 months Ginnie's voluntary speed runs $1.35$ points above Freddie's, $1.228$ times as fast, against the 2.1-point total-CPR differential above, and read off the same full-universe series with the same overstatement. An elasticity is a slope, and nothing here maps a level onto one, so I bracket the mapping rather than choose it. The four are: zero (the committed overlay); $81.4$\% of the conventional response, the reciprocal of that speed ratio, on which a book running $1.228$ times faster carries a proportionally smaller semi-elasticity; $100$\% (no differential); and $122.8$\% (proportional-hazard preservation, which is amplification, not attenuation). In-sample the marginal runs $+7.3$, $+8.9$, $+9.2$ and $+9.6$ points across those four; composed with the off-window calibration, $+4.4$, $+5.4$, $+5.6$ and $+5.8$. That the scaling rises above the committed $0.797\times$ at every positive response is arithmetic, not evidence---a share contributing exactly zero must contribute more at any positive response---and I rest nothing on it. The bracket, not any cell in it, is the result: none of the four is a measured Ginnie marginal, and the overlay-scored $+4.4$ stands unrevised, scoped by the bracket rather than corrected by it.

\paragraph{Danish pipeline mechanisms.}
Two Danish counterfactual mechanisms exist in this framework's pipeline, and only one is live in production. The first is a \emph{rate-gap swap}: holding the baseline hazard, rate-gap coefficient, covariates, and random draw logic fixed, the Danish leg expresses the market-value buyback in rate-equivalent units, so that the buyback discount offsets the locked-in spread. The swap sets $g_{\mathrm{DK}} = g_{\mathrm{US}} + \Delta y_{\mathrm{buyback}} \approx 0$ whenever the loan is out of the money, and $g_{\mathrm{DK}} = g_{\mathrm{US}}$ when rates have fallen and the buyback price is capped at par. At the representative 3.0\% coupon / 6.8\% market state the U.S. gap is $-380$ bp while the Danish effective gap is zero. The U.S. hazard is therefore suppressed by the full locked-in spread while the Danish hazard resets toward its baseline seasoning curve. This swap defines the \emph{mechanism-substitution} variant (the specification behind the superseded 47.1\% Danish CPR of Section~\ref{sec:abm-danish}), and it is \emph{not} the production mechanism. The production Danish leg is the Berger-recalibrated specification adopted in Section~\ref{sec:abm-danish}: it replaces the rate-gap swap with the imported flat moving elasticity, so the production Danish hazard does not respond to the buyback discount at all. The swap machinery remains in the codebase and its PV-versus-par classification is still computed each month, but no production quantity consumes its output, which is why the discount-rate diagnostic of Appendix~\ref{sec:method-abm} moves the Danish leg by exactly \$0.00 billion. Table~\ref{tab:discount} reports that inertness: shifting the buyback discount rate down by 25 to 150 bp (through the primary--secondary spread range) flips thousands of monthly present-value-versus-par classifications while moving the Danish CPR path, roll-off path, and trapped-liquidity dollar by exactly zero on every scoring basis. Two limits apply to that result. I state both because together they retract the immunity I previously claimed from it. First, the diagnostic executed on the Danish-level bracketing leg, not on the production U.S.-intercept leg: its own artifact records a mean Danish CPR of 3.39\%, the bracketing anchor's speed. Under the production anchor the same independence does hold, but by construction, not by measurement, because that branch replaces the rate gap with zeros before the hazard is evaluated. Second, and more important, the result is about the hazard, not about the accounting. The Danish leg credits every retired balance at par: its dynamic-balance loop books roll-off as balance times drain, with no price term anywhere in it. Netting the roughly \$199 billion scheduled component from the leg's \$669.3 billion of Danish roll-off leaves on the order of \$470 billion of buyback-triggered retirement credited at par, against a committed proxy discount grid of 32 to 38\%. A market-value credit would reduce the Danish leg's cash receipts by an amount larger than the $+\$61.2$ billion gap itself, and in the direction that shrinks it. The pipeline does not distinguish a cash haircut from a balance adjustment, and the incidence depends on which the buyback is. So a pre-committed bracket prices both readings from the committed pair (run \texttt{buyback\_credit\_bracket}; gates and the interpretive rule fixed in the script header before execution). Under the balance-adjustment reading face retirement is unchanged and the gap stands at $+\$61.2$ billion. Under the cash-haircut reading receipts fall by the discount on the leg's \$470.7 billion of early face and the gap reverses. The committed grid was asserted, not derived: run \texttt{buyback\_discount\_rederived} prices each month's retired face at its own prepayment-consistent present value---the SOMA book aged to that month at the 2.49\% weighted-average coupon, prepaying at that month's Danish CPR, discounted at that month's market rate---for a face-weighted mean of 23.8\%, a \$112.2 billion haircut, and a cash gap of $-\$51.0$ billion. The price and the face come from two populations (the SOMA book's own vintage tabulation and the microsimulated leg's roll-off): the seam the committed early-face netting already carries, inherited here rather than repaired. That is a point, not a range: the committed 32--38\% span was a range only by virtue of a four-point asserted grid. The monthly discounts do range, from 18.4\% to 28.1\%, and none of the forty-two months reaches the committed grid's floor. The committed grid is consistent with a no-prepayment annuity priced at a representative state: the same primitive returns 35.3\% at zero prepayment on a 3.0\% coupon against a 6.8\% market rate, inside the committed span, where the leg itself prepays at 5.61\% against a window averaging 6.61\%. That identification is plausible rather than exactly reproduced: a band check on the primitive, not a target. The re-derivation shrinks the reversal without retiring it; the threshold it clears is arithmetic on the committed chain, fixed before the run rather than produced by it: the verdict turns at a 13.0\% discount, below every month's priced value. The two readings are not symmetric answers to one question. The benchmark is a face-value object: Section~\ref{sec:method-benchmark} constructs it from the SOMA current face value series against the par-denominated caps. So the balance-adjustment reading is the benchmark-consistent one, and on the benchmark's own denomination the gap stands. The cash-haircut reading answers a different question, reserves actually drained, and its reversal has a name: the $1-P$ wedge on bought-back face is exactly the par windfall a moving household pays the bondholder under the U.S. rule and keeps under the Danish one, so the cash-flow cost is the disappearance of a household-to-bondholder transfer, not an added resource cost. The gap's sign is therefore incidence-conditional, the channel is priced, and the discount-rate diagnostic still should not be read as immunity to the pricing critique.

\paragraph{The transform's superseded variant and the concave-transform test.}
Two records travel with the transform of \eqref{eq:beta1}. The superseded $0.0693$ is a distinct object: the continuous-hazard variant $-\ln\!\big[\ln(1-(1-\delta)P_q)/\ln(1-P_q)\big] = 0.0693$ (unrounded $0.069258$), which the simulation does not use. The two forms differ by $0.0007$ and both round to $0.069$, but only \eqref{eq:beta1} reproduces the reported results.

The direct test of the log-linear extrapolation is a piecewise-concave transform---full marginal effect over the first 200 bp of lock-in, half beyond. It moves the recovery from 107.0\% to 105.7\% (\$808.4 billion) and slightly improves the contemporaneous path correlation ($+0.281$ versus $+0.190$, peak lag unchanged at $-3$). The log-linear extrapolation is not load-bearing for the level, and is for the marginal---but the asymmetry is one of proportion, not of magnitude, and the same absolute $\pm 1$-point convention cannot support both readings. A pre-committed rerun of the paired legs under the concave transform puts the concave marginal at $+7.87$ points against the production $+9.20$ at the in-sample floor, and at $+5.06$ against $+5.57$ at the off-window point floor. At the in-sample floor the concave null is bit-identical to the production null---\$748.185 billion on both legs, because $\beta_1 = 0$ leaves nothing for the transform to act on---so the level move and the marginal move are numerically the \emph{same} 1.32 points. That 1.32 points is 1.2\% of a 107.0\% recovery and 14\% of a $+9.20$-point marginal. It is that relative reading, not the absolute one, on which the transform is immaterial to the level and material to the marginal. On the absolute $\pm 1$-point materiality threshold the floor-form test of Section~\ref{sec:robustness-floor} applies, and both moves clear it. The transform dimension travels with the marginal's operative uncertainty in Table~\ref{tab:uncertainty}.

\paragraph{Sampling schemes for Path B's level and marginal.}
The Liebersohn--Rothstein band of Section~\ref{sec:pathb} is a calibration interval, not sampling error. The sampling dimension is quantified separately, by a stratified loan-level bootstrap specified ex ante in its script header: 200 replicates resampling the 75,000-loan draw with replacement within strata, central and null legs paired per replicate, engine seed held at production. That bootstrap puts the central recovery's 95\% interval at $[107.0, 107.1]\%$ under the standalone scorer (replicate standard deviation 0.03 points), the null's at $[97.8, 97.9]\%$, and the paired lock-in marginal's at $[+9.17, +9.23]$ points (standard deviation 0.016). Re-run under the cluster scheme at this same floor, the paired marginal reads $[+8.27, +10.19]$ points (sd 0.52; run \texttt{bootstrap\_pathb\_cluster}), an understatement of 30.8$\times$ in-window against 37$\times$ off-window. Those intervals are narrow because the scheme holds the design fixed. It resamples loans within each stratum and preserves every stratum's count, so between-stratum composition variance is never drawn. Drawing the 130 three-way draw strata themselves with replacement---at the off-window floor the paper headlines---widens the marginal's interval to $[+4.63, +6.92]$ points with a standard deviation of $0.60$ points (run \texttt{bootstrap\_pathb\_cluster}, 200 replicates). That is a width of $2.29$ points against the within-stratum scheme's $0.062$, so the committed scheme understated loan-sample uncertainty by roughly a factor of $37$. Two properties of that interval are not corrected for and should be stated. Drawing strata with replacement makes the loan count itself random: it runs from 54,134 to 101,404 across replicates against a fixed-design 75,000. So every composition weight is a ratio of random totals, not of fixed ones, and the percentile interval carries no correction for the resulting ratio bias. A fixed-$n$ variant that rescales weights to the design total would settle how much of the width is that bias. It is not run here. The committed point marginal falls inside that interval. Two limits apply to the wider figure. The first is that the strata are very unequal, so the effective cluster count is $25.8$, not $130$ (the largest single stratum is $9.77\%$ of sample balance), and a percentile interval on $26$ effective clusters is itself noisy. The second is that the resampled loan count varies from roughly $54{,}000$ to $101{,}000$. That variation is harmless here only because the engine renormalises the pool to Fed holdings each month, so the level is set exogenously and only composition varies.

\paragraph{Curtailment-layer demonstration runs.} On the accounting identity of Section~\ref{sec:robustness-hybrid}: a $\pm 25\%$-amplitude seasonal profile and an income-elasticity regime split across 2022--23 and 2024--25 move the netted total to \$70.06 and \$68.09 billion and the wedge to 9.16 and 8.90 points, the marginal unchanged to below $10^{-13}$ points on both legs and the CPR paths bit-identical; a deterministic two-servicer mix (fixed 0.60/0.40 shares, front-loaded and flat group factors) moves the wedge to 9.30 points (\$71.12 billion netted) with the marginal again unchanged, since any servicer-heterogeneous pattern aggregating to one monthly rate is a special case of the same identity.

\section{Path A Estimation Detail and Timing Diagnostics}\label{app:patha-detail}

This appendix collects the estimation detail of Path A (Section~\ref{sec:patha}), the direct input-timing diagnosis, the Danish discount-rate diagnostic, and the monthly-timing exhibits of Section~\ref{sec:identification}.

The second hazard path aggregates Freddie Mac loan-level performance data to the stratum-month level. Strata $s$ are the 296 observed four-way cells of vintage $\times$ coupon $\times$ FICO bucket $\times$ LTV bucket, and the unit of observation is the stratum-month pair $(s,t)$. The estimation configuration sets the bucket edges: FICO below 680, 680--740, and 740 and above; LTV at or below versus above 80; coupons in 0.5-percentage-point steps; vintages 2017--2021. Geography is not a stratum axis, and no geographic term enters \eqref{eq:pathA}. The model is a discrete-time hazard with a log-exposure offset,
\begin{equation}\label{eq:pathA}
\log \mathbb{E}\!\left[y_{s,t} \mid x_{s,t}\right] = \log n_{s,t} + \alpha_{s} + f\!\left(a_{s,t}\right) + \beta_g\, g_{s,t} + \beta_b\, b_{s,t} + \beta_\phi\, \phi_t ,
\end{equation}
where $y_{s,t}$ is stratum $s$'s prepaid unpaid principal balance (UPB) in month $t$, in dollars, $n_{s,t}$ its UPB exposure, $x_{s,t}$ the right-hand-side covariates, and $\alpha_{s}$ the stratum fixed effect (295 enter, one absorbed as reference). The term $f(\cdot)$ is a loan-age seasoning spline in mean loan age $a_{s,t}$: a truncated quadratic power basis, a linear term plus $(a - k)_+^2$ at knots $k \in \{12, 24, 36, 60, 84, 120\}$ months, seven basis columns. The remaining terms are $g_{s,t}$, the rate gap in 100 bp units; $b_{s,t}$, the stratum burnout state; and $\phi_t$, the dynamic friction index. The outcome is dollars, not a count. With the log-exposure offset the modeled object is the prepayment rate $y_{s,t}/n_{s,t}$, a nonnegative fraction, so the model is estimated by Poisson pseudo-maximum likelihood, whose score equations remain consistent for the conditional mean without any distributional assumption on $y_{s,t}$; no dispersion statistic is meaningful for this pseudo-likelihood. The 295 stratum fixed effects raise no incidental-parameters problem for the coefficients of interest. Because the model is Poisson, concentrating the fixed effects out of the log-likelihood returns exactly the conditional-likelihood score for the slope vector: the unconditional dummy-variable fit computed here yields the same slope estimates as the conditional estimator, and stays consistent as the number of strata grows with the 36 training months held fixed. The fixed-effect logit or probit does not---its slope estimates are inconsistent at fixed panel length \citep{wooldridge1999}. All that consistency requires is the multiplicative conditional mean $\mathbb{E}[y_{s,t}\mid x_{s,t},\alpha_{s}] = n_{s,t}\exp(\alpha_{s} + x_{s,t}^{\top}\beta)$, with no restriction on the distribution, dispersion, or serial dependence of $y_{s,t}$. All uncertainty statements below come from the two bootstrap schemes of Table~\ref{tab:bootstrap}. The first is a stratum-level cluster bootstrap, robust to arbitrary within-stratum dependence and overdispersion but structurally unable to measure uncertainty in the friction coefficient: $\phi_t$ varies only over time, so resampling strata reproduces the identical friction series in every replication. That burden (and, partially, burnout's) falls to the second, a temporal moving-block bootstrap that resamples the 36 training months in blocks of six with each month's full cross-section intact. Table~\ref{tab:panel} reports the estimation panel's accounting. Both it and Table~\ref{tab:bootstrap} are collected in Appendix~\ref{app:ridge}, because Path A is excluded from every headline range by pre-committed specification and neither table enters a headline quantity. The displayed form \eqref{eq:pathA} is the no-seasonal spec v3 design. The production specification (spec v4) adds eleven calendar-month dummies and is the design under which the 121.5\% recovery reported below is computed, though that level is not calendar-attributable (Section~\ref{sec:patha}); every Path A inferential object stays on spec v3. Table~\ref{tab:hazard-notation} collects the notation shared by the two hazard paths.

The 295 stratum fixed effects absorb each cohort's baseline prepayment speed so that the burnout term captures within-pool adverse selection over time. The burnout regressor is a within-stratum demeaned measure of the surviving loan stock, orthogonalized against the loan-age spline so that it is not mechanically collinear with seasoning. This design is the estimation counterpart of \citepos{lesniewski2026} selection identity. The exposure-offset outcome is itself a pool-level, survival-weighted (here UPB-exposure-weighted) hazard. The 295 strata cut the cross-sectional dispersion that drives the identity's negative selection term down to within-cell variation. The demeaned surviving-stock regressor proxies the residual composition drift that the time-invariant fixed effects cannot absorb. An earlier specification using four vintage-year dummies in place of the 295 stratum fixed effects returned a burnout coefficient of the wrong sign (+0.76). Replacing the coarse vintage fixed effects with the full 295-stratum specification and demeaning burnout within stratum removes that wrong-signed artifact. The corrected estimate ($-0.13$) carries the theoretically expected sign: a more heavily burned-out pool should prepay more slowly, holding age and rate gap fixed. But the sign statement is a point-estimate statement only. The temporal bootstrap scheme of Table~\ref{tab:bootstrap} centers this coefficient on the opposite sign (mean $+1.23$, median $+0.52$). Under both schemes it is imprecise and not statistically distinguishable from zero. I therefore treat it as a control, not a finding. The fit carries a nominal ridge penalty $\lambda_{\mathrm{ridge}}$ on the full coefficient vector, because 295 fixed effects make the model ill-conditioned under plain iteratively reweighted least squares. In 20 strata no training prepayment event is recorded at all, so the likelihood is monotone in their fixed effects and the maximum-likelihood estimates diverge to $-\infty$. I choose $\lambda_{\mathrm{ridge}}$ over the grid $\{10^{-5}, 10^{-4}\}$ on a temporally blocked holdout comprising all stratum-months from January 2024 onward, so that no within-stratum information crosses the train/holdout split. The penalty is not what stabilizes the fit. At these magnitudes it is a numerical no-op. The two grid points return coefficient vectors identical to machine precision from the IRLS warm start, so the production coefficients are the IRLS solution and the grid selection is inconsequential. The 20 zero-event strata are held at the finite value the warm-started optimizer leaves them at, and the penalty shrinks them only negligibly. Because the split is temporal over a fixed set of strata, that same fixed effect scores each stratum's post-January-2024 holdout months, so no fixed effect is estimated on, or transferred to, a stratum absent from the training window. The penalty's operative role is as the fallback for resampled refits in the bootstrap machinery. Nor is the production point an artifact of the device. Dropping the 20 zero-event strata and fitting unpenalized Poisson PML reproduces the spec v3 macro coefficients to within $0.002$ in standardized units (rate gap $+0.6727$ against the spec v3 $+0.6734$; under spec v4 the same check holds to within $0.0022$, Appendix~\ref{app:ridge}). Appendix~\ref{app:ridge} records the verification details under both specifications, including a reference-stratum sensitivity of penalized cold-start refits that motivates the warm-started bootstrap below. The battery has been recomputed under spec v4 and the device is unchanged. Two caveats still attach. The same holdout later serves as the out-of-sample diagnostic, so the RMSE reported below is at least weakly favored by selection. Months never touched by any selection step do not exist in this design, one reason Path A is held to a robustness role.
Estimated on this panel, the production specification --- spec v4, the calendar-month seasonal design adopted at the freeze --- yields a rate-gap coefficient of $+0.65$, a burnout coefficient of $-0.17$, and a friction coefficient of $-0.007$, each with the theoretically expected sign. The rate-gap coefficient is standardized: the training standard deviation of the rate gap is 221 bp, so the equivalent per-100 bp coefficient is $+0.29$. The friction coefficient's shrinkage relative to the prior spec ($-0.037$) is expected under the seasonal design, whose eleven month dummies absorb most within-year variation in the time-only friction index. Point estimates alone overstate what is identified here. Table~\ref{tab:bootstrap} therefore reports cluster-bootstrap uncertainty for all three, including the bias diagnostics. That uncertainty is estimated under the prior spec v3 (see the table note). The rate-gap point estimate sits near the 3rd percentile of its own bootstrap distribution (bootstrap mean $+1.49$, median $+0.76$, estimated bias $+0.82$). That is a signature of resampling-induced upward skew, under which percentile and basic intervals disagree. The percentile interval excludes zero; the bias-reflecting basic interval does not. I therefore rest no claim on interval-based significance for any Path A coefficient. The defensible statements are about direction and imprecision. The rate-gap coefficient's point estimate is positive, and positive in 197 of 198 converged replications (99.5\%). That share is a property of a right-skewed resampling distribution, not a test. The pre-committed bias-respecting test of $H_0\!: \beta_g \le 0$ does not reject under any construction (BCa $p = 0.093$ stratum / $0.412$ temporal; permutation $p = 0.241$; run \texttt{patha\_sign\_test}). The coefficient is therefore directionally consistent with the imported elasticity without being statistically distinguishable from zero. Its magnitude is identified only up to a wide, construction-dependent interval. Burnout and friction are not distinguishable from zero on any construction. All three enter the production simulation as fitted inputs, none as findings, and no result in this paper depends on their point values. Path A's aggregate recovery now carries the interval that pass produces. I resimulate the QT window under each of the 198 converged replications, swapping in the three macro coefficients per draw and holding intercept, age spline, and fixed effects at the production fit. That yields a median of \$922.6 billion and an interquartile range of \$900.1--\$1{,}130.6 billion, but a 95\% percentile interval of $[-\$4{,}459.8, +\$1{,}190.1]$ billion. A joint rerun swaps in each replication's full coefficient vector: intercept, age spline, and at spec v4 the month effects, drawn together with the macros. That rerun returns $[+\$689.0, +\$1{,}217.6]$ billion, far narrower, and it is not adopted. Of its 198 draws, 106 sit at exactly the zero-voluntary-prepayment floor of \$1{,}217.6 billion, among them all 79 rogue-mode draws --- rogue mode meaning a drawn rate-gap coefficient nearer the reference-swap attractor $+1.96$ than the production $+0.67$. That floor is the trapped figure a pool returns when voluntary prepayment is switched off and only scheduled amortization rolls off, measured by a probe through the same coefficient handoff. The committed interval remains the quoted one, because its $-\$4{,}459.8$ billion tail and the joint rerun's floor pile are the same rogue-mode replicates rendered without and with their compensating intercepts. That is one optimizer pathology of penalized refits on resampled panels, not two sampling statements. Measured at spec v4, the same joint rerun fails 1 of 200 replications against spec v3's 2, and its draws avoid the floor entirely (0 of 199). The seasonal design's instability therefore manifests as rogue-mode attraction (73 of 199 draws), not as fit failure. Fifteen replications with extreme burnout or friction draws (the two coefficients the bootstrap classifies as statistically zero) destabilize the 42-month forward compounding entirely, and 83.8\% of draws land above the benchmark. Those two facts pull in opposite directions, and the unfavorable one belongs in the text as much as the favorable one. The same 198 draws have a \emph{mean} of \$598.9 billion, or 78.3\% of the benchmark --- below the rate-inelastic $\beta_1 = 0$ null's 85.7\% and below Path~B's headline 91.3\% --- against a median of 120.6\%. On the share basis the percentile interval is $[-583.2\%, +155.6\%]$ with a resampling standard error of 201.9 points. I report the mean alongside the median for the same reason the coefficient table does. These distributions are severely right-skewed, so no single summary is sufficient; quoting only the median and the fraction above the benchmark would flatter the estimator. Propagated coefficient uncertainty therefore leaves the forward simulation with no useful precision. That is an additional reason the aggregate dollar figure is held to the aggregate-level role assigned below and is not carried in the abstract. Every Path A inferential object remains estimated under the prior spec v3: Table~\ref{tab:bootstrap}, this interval, and the train/holdout decomposition below. That is because bootstrap replications under the spec v4 seasonal design are branch-unstable. The resampled penalized likelihood has a second optimum that attracts replications away from the production branch. On the production panel itself 1 of 50 dispersed starts lands on the production branch, and no start reaches a better penalized objective (run \texttt{pathA\_cold\_starts}). The objective evaluated at the frozen coefficient vector is the best value any of 150 starts attains, and every non-production terminus sits at least seven orders of magnitude above it. That is an optimizer stall of a strictly convex objective. Until the resampled instability is stabilized, the production specification carries no valid uncertainty of its own. Path A is excluded from the paper's headline figures for this reason, retained as an aggregate-context, non-corroborating exhibit. Simulating forward, this model recovers \$928.9 billion, or 121.5\%, of the \$764.7 billion empirical benchmark under the production calendar-month specification (spec v4). The forward simulation drains each cohort's balance deterministically each month by balance times hazard, with fractional prepayment settling directly to simulated SOMA roll-off and no settlement-lag convolution. That specification was adopted at the pre-submission freeze under pre-committed parity gates that reproduce the committed seasonal robustness artifact exactly. The prior specification is preserved as a spec-v3 exhibit.

Three diagnostics temper this aggregate result, because they bear directly on how much weight the \$928.9 billion figure can carry. First, the holdout fit, evaluated under spec v3, is weak in both weightings. Re-scoring the frozen spec v4 fit on the same January-2024-onward holdout, without refitting, gives 37.9 unweighted and 3.0 exposure-weighted percentage points against spec v3's 37.8 and 2.5, so the seasonal terms do not improve the out-of-sample fit. The units matter here. The unweighted RMSE of approximately 38 percentage points of annualized stratum-month CPR is dominated by small-exposure strata with unstable realized rates, while the exposure-weighted RMSE, the weighting the fit itself uses, is 2.5 percentage points against an empirical mean CPR of 5.14\% (ABM basis), half the mean. Neither weighting supports out-of-sample predictive power at the stratum-month level, and the ridge grid's two candidate fits are numerically identical on this holdout, so no selection hinges on the weighting. A flexible-learner ceiling test bounds how much of this weakness is functional form. A gradient-boosted Poisson learner given the same covariate information on the identical pre-committed holdout improves the exposure-weighted RMSE materially, to 0.9 points against the frozen spec v4's 3.0. It leaves the unweighted RMSE essentially unmoved (37.6 against 37.9) and out-of-sample $R^2$ at zero ($+0.003$). A penalized Poisson GLM on the exact v4 design does no better than the production fit (2.6 and 38.1 points). Read against the run's pre-committed thresholds, flexibility buys recalibration of the high-exposure cells but no cross-cell predictive power in the 2024-onward regime. The landmarking-with-recalibration strand receives the same test. An exposure-weighted isotonic map fit on the 2023 inner-validation year only, applied to the frozen spec v4 holdout predictions, improves the exposure-weighted RMSE to 2.53 points against the frozen 3.04, short of the pre-committed 0.80$\times$ threshold. It leaves the unweighted RMSE at 37.8 and out-of-sample $R^2$ at zero ($-0.007$). The pre-committed verdict is no material change: a monotone level recalibration fit in the training regime does not transfer predictive power across the 2024 boundary. Second, the recovery is concentrated where the model is unsupervised. The decomposition is measured under spec v3 and was not recomputed at adoption. Of the spec v3 fit's \$915.1 billion, \$331.0 billion accrues in the 19 QT months overlapping the training window (June 2022--December 2023), against \$337.5 billion empirical (98\%). A further \$584.1 billion accrues in the 23 post-training simulation months from January 2024 onward, against \$427.3 billion empirical (137\%); 21 of those months have performance data, and the empirical comparison uses those. Components are independently rounded, and the overshoot lives almost entirely out of sample. Third, the raw CPR correlation against the empirical path is $r = -0.444$ at zero lag, more negative even than the ABM's $-0.318$, and \eqref{eq:pathA} omits calendar-month seasonality, one of the four factors of \citepos{richard1989} canonical prepayment function and a natural first suspect for a monthly path fit this poor. Re-estimating with eleven calendar-month dummies confirms strong, correctly signed seasonality (spring and early-fall log-effects up to $+0.43$ against a January base) and improves the simulated path's contemporaneous correlation from $-0.444$ to $-0.378$, with the peak moving from lag $-3$ to $-2$. That timing improvement is calendar-attributable, and I test it. Permuting the assignment of the twelve estimated month log-effects to calendar months preserves their multiset exactly, and therefore every level channel they drive, while destroying only their alignment with calendar time. I score the permutation on the conjunctive criterion the claim actually states: a lag-0 correlation better than $-0.378$ \emph{and} a peak at lag $-2$. Scrambled assignments win 1 of 64 draws, a one-sided $p$ of 0.031. An independent sixteen-draw family at a different seed also returns a single win (\texttt{pathA\_seasonal\_permutation}). The two components scored separately, 11 and 8 of 64, are diagnostics only: each necessarily exceeds the joint count, so reading either alone understates a conjunctive claim. The recovery figure is a different matter, and I no longer credit it to the calendar terms. Across those same 64 scrambles the recovery spans 121.14--122.29\%, a range of 1.15 points that contains the adopted 121.5\%, so the level is calendar-invariant: permutation holds the multiset of month effects fixed, and it is the multiset, not its alignment, that sets the level. Nor is the spec-v4 total decomposable as production plus seasonality in the first place. The seasonal refit with all twelve month multipliers zeroed moves the recovery by $+\$80.27$ billion on its own, and restoring the multipliers claws back $-\$66.44$ billion of that, against a net reported seasonal effect of $+\$13.83$ billion. The timing improvement is real but far from sufficient: the named first suspect accounts for only a modest share of the negative path fit. A model that recovers 121.5\% of the dollar benchmark while mismatching the monthly path this badly may be recovering the aggregate partly through offsetting errors, over-predicting prepayment in some cohort-months and under-predicting in others in a way that nets out near the benchmark without the underlying mechanism being correctly timed. I therefore treat Path A's recovery as evidence about the aggregate level only. No timing claim rests on Path B either: its apparent levels advantage does not survive detrending (Section~\ref{sec:pathb}), and every claim this paper makes for Path B is a level-accuracy claim. Table~\ref{tab:pathadiag} collects Path A's diagnostics in one place.

\begin{table}[H]
\centering
\caption{Path A diagnostic summary. An aggregate-context, non-corroborating exhibit: Path A is excluded from the paper's headline figures by pre-committed specification, and no cell of this table enters a headline quantity. Values repeat committed figures derived in this section and Appendix~\ref{app:theil}.}
\label{tab:pathadiag}
\footnotesize
\begin{threeparttable}
\begin{tabular}{@{}p{4.6cm}p{9.6cm}@{}}
\toprule
Diagnostic & Value \\
\midrule
Aggregate recovery \emph{(spec v4)} & \$928.9 billion $=$ 121.5\% standalone; 112.4\% shared basis. Order of magnitude coincides with Path B; the sign does not corroborate it, the coefficient being statistically indistinguishable from zero (row below and Section~\ref{sec:patha}) \\
Fit diagnostics \emph{(spec v4)} & exposure-weighted holdout RMSE 3.0 points; Theil $U_1$ 0.431, $U_2$ 1.010; terminal cumulative error $+\$164.1$ billion \\
\bottomrule
\end{tabular}
\begin{tablenotes}
\footnotesize
\item \emph{Notes:} Every row above is measured under the production calendar-month specification (spec v4), so this table is single-specification and no cell need be read against a neighbour measured under a different design. The spec v3 objects the table previously carried in its body are recorded here instead, demoted because they were not recomputed at adoption: the \emph{refit-and-resimulate distribution} --- point \$915.1 billion; median \$922.6 billion; \emph{mean \$598.9 billion $=$ 78.3\% of benchmark}; IQR \$900.1--\$1{,}130.6 billion; 95\% percentile interval $[-\$4{,}459.8, +\$1{,}190.1]$ billion ($[-583.2\%, +155.6\%]$ on the share basis); 83.8\% of draws above benchmark --- the \emph{train vs post-training split}, \$331.0 billion (98\% of empirical) vs \$584.1 billion (137\%); and the spec v3 counterpart of the holdout RMSE, 2.5 points against the spec v4 row's 3.0. The resampling distribution has since been rerun jointly at both specifications (Section~\ref{sec:patha}), and the committed spec v3 interval remains the quoted one --- the joint rerun is narrower for the mechanical reason stated there. The spec v3 point (\$915.1 billion) is recorded here, rather than beside the spec v4 aggregate, so that the median is not read as centred on a level measured under a different design. The 95\% interval spanning zero is the stated reason for the pre-committed headline exclusion; the flexible-learner and isotonic comparators of this section bound what refitting the same information could recover.
\end{tablenotes}
\end{threeparttable}
\end{table}

\begin{figure}[H]
\centering
\includegraphics[width=0.9\textwidth]{fig2_ccf}
\caption{Cross-correlation functions between each estimator's simulated CPR path and the empirical SOMA-derived path, lags $-6$ to $+6$. Under the implemented convention (Section~\ref{sec:pathb}), a peak at negative lag means the empirical series leads and the simulated path trails. Path B's $-3$/$-4$ peak and Path A's and the ABM's negative contemporaneous correlations summarize the same fact from Table~\ref{tab:estimators}: no estimator shows positive contemporaneous alignment, and once the window's shared trend is removed no estimator's co-movement is distinguishable from zero under the moving-block bands.}
\label{fig:ccf}
\end{figure}

A moving-block bootstrap (block length 6, 4,000 replications) puts the 95\% interval on the peak-lag correlation at $[-0.21, +0.49]$ and on the lag-0 correlation at $[-0.13, +0.44]$. With seven lags searched, individual band-point differences in peak $r$ carry no evidential weight. These intervals are not sensitive to the block-length choice in any verdict-relevant way: at block lengths 4 and 8 the peak-lag endpoints move to $[-0.27, +0.39]$ and $[-0.22, +0.51]$, and the lag-0 endpoints to $[-0.11, +0.45]$ and $[-0.17, +0.42]$, every interval still containing zero.

A Theil decomposition formalizes where each estimator's monthly error lives. It is a fit diagnostic, not a timing credential. On levels, the ABM's mean-squared error is 68\% bias share (the level failure already visible in its aggregate). Path B's is 11.3\% bias and 82.5\% variance share on the corrected amortization basis (run \texttt{coupon\_convention\_amortization}). Its residual monthly error is under-dispersion against a seasonal empirical path, not level bias. Path A is intermediate (26.5\% bias, 63.0\% covariance share). After linear detrending the variance share still dominates for Path B (93.7\%). On first differences, Theil's $U_2$ is 1.23 for the ABM, 1.01 for Path A, and 1.00 for Path B. No estimator beats a naive no-change forecast of month-over-month CPR changes, which restates the timing concession above in forecast-accuracy units. Table~\ref{tab:theil} collects the full decomposition by transform, with cumulative dollar-error profiles whose terminal values reproduce Table~\ref{tab:estimators}'s aggregates by construction.

A direct input-timing diagnosis rules out the dating of the rate input and narrows, but does not isolate, the mechanism. In monthly \emph{changes}, the simulated CPR responds to the contemporaneous rate change with the correct sign and near-mechanical strength ($\Delta$-correlation $-0.94$ at lag 0). The empirical CPR's month-to-month changes show no detectable rate response ($+0.25$, inside the $\pm 0.31$ zero band at $n = 41$). At monthly frequency the empirical path's variation is dominated by non-rate factors the simulation does not model, calendar seasonality above all. One construction asymmetry belongs in this reading as well. The empirical CPR back-out retains the income-scaled curtailment component (Section~\ref{sec:method-benchmark}), while the simulated CPR paths exclude it, since curtailment enters only as a common dollar flow in the accounting layer. The two series being correlated therefore differ by a small, slowly-varying component as well as by seasonality.

\begin{table}[H]
\centering
\caption{Danish-leg sensitivity to the buyback discount rate. The diagnostic executed on the Danish-level Berger-recalibrated leg (mean Danish CPR 3.39\%), which is the bracketing anchor of Section~\ref{sec:abm-danish} rather than the production U.S.-intercept leg. Neither leg consumes the present-value classifications the discount rate enters, so the zero deltas hold on the standalone, shared, and printed bases alike; under the production anchor that independence is structural rather than run-verified. The baseline CPR path reproduces the frozen parquet for the leg that ran exactly. The zeros bound the hazard channel only: the leg's dollar conversion credits retired balance at par, an un-priced channel sized in Section~\ref{sec:pathb}.}
\label{tab:discount}
\begin{tabular}{lccc}
\toprule
Discount-rate shift & PV classifications flipped & $\Delta$ CPR path (pp) & $\Delta$ trapped (\$B) \\
\midrule
0 bp (production) & -- & 0.00 & 0.00 \\
$-25$ bp & 2{,}470 of 1{,}683{,}082 & 0.00 & 0.00 \\
$-50$ bp & 7{,}805 of 1{,}683{,}082 & 0.00 & 0.00 \\
$-75$ bp & 17{,}493 of 1{,}683{,}082 & 0.00 & 0.00 \\
$-100$ bp & 32{,}488 of 1{,}683{,}082 & 0.00 & 0.00 \\
$-125$ bp & 55{,}210 of 1{,}683{,}082 & 0.00 & 0.00 \\
$-150$ bp & 92{,}495 of 1{,}683{,}082 & 0.00 & 0.00 \\
\bottomrule
\end{tabular}
\end{table}

The decomposition also carries its time axis. The monthly marginal path is the central-minus-null difference of the two committed simulated legs, and its sum reproduces the $+\$70.3$ billion bit-exactly. That path is positive in all 42 window months, splits 21.3\%/41.9\%/36.8\% across the window's thirds, and peaks at \$2.29 billion in November 2023, the window's eighteenth month. By coupon bucket the terminal marginal divides 38.0\%/43.7\%/18.3\% across the $<3.0\%$/3.0--4.0\%/$\geq 4.0\%$ groups on 40.3\%/40.4\%/19.3\% of sample balance. Figure~\ref{fig:marginaltiming} draws the stack.

\begin{figure}[H]
\centering
\includegraphics[width=0.94\textwidth]{fig12_marginal_timing.png}
\caption{Monthly lock-in marginal by coupon bucket. Stacked areas are the per-bucket monthly marginals from the month-by-cell capture of the group-ablation engine; the dark line is the aggregate central-minus-null path, whose sum reproduces the $+\$70.3$ billion in-sample marginal bit-exactly. Dotted rules mark the window's thirds; the marker is the November 2023 peak. The stack differs from the aggregate by the decomposition's 1.03\% terminal additivity residual, asserted in the generator.}
\label{fig:marginaltiming}
\end{figure}

\section{Path A Estimation Panel, Coefficient Uncertainty, and Ridge Device}\label{app:ridge}

The replication package verifies three facts about the ridge device of Section~\ref{sec:patha}, under both specifications: the spec v3 records stand, each fact has been recomputed under spec v4, and the device is unchanged. First, at production magnitudes the penalty is a numerical no-op on the full panel. The two grid points ($\lambda_{\mathrm{ridge}} \in \{10^{-5}, 10^{-4}\}$) return coefficient vectors identical to machine precision from the IRLS warm start (maximum parameter difference $4.5\times10^{-17}$ in the replication artifact; $8.4\times10^{-17}$ under spec v4). Grid selection is therefore inconsequential, and the production coefficients are the IRLS solution to machine precision; ``mildly shrunk'' would overstate an effect that is nil. Second, ridge on the full coefficient vector with one stratum absorbed as reference is not invariant to the reference choice. In cold-start refits on resampled panels, swapping the reference stratum can move the standardized rate-gap coefficient by more than its own magnitude --- an optimizer-path artifact of penalized refits with the 20 inestimable fixed effects. The bootstrap machinery accordingly warm-starts every replication and discards non-converged draws. The warm-started spec v4 swap refit is reference-invariant to $7.3\times10^{-14}$, consistent with the warm-start policy that concession motivated. That invariance is not evidence that the cold-start sensitivity is gone. Third, the production point itself is not an artifact of these choices. A refit that drops the 20 zero-event strata (423 of 10,176 cells) and fits unpenalized Poisson PML converges, is exactly invariant to the reference choice, and reproduces the spec v3 macro coefficients to within $0.002$ in standardized units: rate gap is $+0.6727$ against the spec v3 $+0.6734$, and friction $-0.0373$ against $-0.0371$, third-decimal agreement. Burnout is $-0.1317$ against $-0.1301$; the larger burnout gap reflects that dropping the 423 cells changes the burnout regressor's within-stratum demeaning and age-orthogonalization sample. At spec v4 the identical 20 strata and 423 cells drop out, since the month dummies do not change the zero-event set. The same refit is again exactly reference-invariant ($1.1\times10^{-14}$) and reproduces the spec v4 macro coefficients to within $0.0022$: rate gap $+0.6463$ against $+0.6469$, friction $-0.00715$ against $-0.00711$, burnout $-0.1734$ against $-0.1713$, against a tolerance pre-committed at $0.005$, loosened from v3's $0.002$ because the month dummies absorb time variation the friction coefficient formerly carried.

Table~\ref{tab:panel} reports the estimation panel this fit is built on, and Table~\ref{tab:bootstrap} the coefficient uncertainty the two bootstrap schemes of Section~\ref{sec:patha} return. Both sit here because Path A is an aggregate-context, non-corroborating exhibit that pre-committed specification excludes from every headline range; no cell of either table enters a headline quantity.

\begin{table}[H]
\centering
\small
\begin{threeparttable}
\caption{Path A estimation panel.}
\label{tab:panel}
\begin{tabular}{@{}lrr@{}}
\toprule
 & Training & Holdout \\
\midrule
Performance months & Jan.\ 2021--Dec.\ 2023 & Jan.\ 2024--Sep.\ 2025 \\
Stratum-months & 10,176 & 6,077 \\
Prepaid UPB (PML outcome) & \$650.0B & \$131.3B \\
UPB exposure (month-sums) & \$72.5T & \$37.4T \\
\bottomrule
\end{tabular}
\begin{tablenotes}[flushleft]\footnotesize
\item The panel aggregates the full Freddie Mac 2017--2021 origination universe (8.84 million active loans at the window's peak month, not the 75,000-loan sample that Section~\ref{sec:pathb} simulates) to 296 strata, and carries no sampling weights: exposure and the outcome are raw dollars, so the pseudo-maximum-likelihood fit is UPB-weighted by construction (a dollar of prepayment counts the same in any stratum-month). The temporal split at January 2024 is the ridge-selection holdout of the text; the performance window is January 2021 through September 2025.
\end{tablenotes}
\end{threeparttable}
\end{table}

\begin{table}[H]
\centering
\footnotesize
\setlength{\tabcolsep}{4pt}
\begin{threeparttable}
\caption{Path A coefficient uncertainty under two bootstrap schemes (estimated under the prior specification, spec v3; see note).}
\label{tab:bootstrap}
\begin{tabular}{@{}lrrrrccr@{}}
\toprule
Coefficient & Point & Mean & Median & SE & Percentile CI & Basic CI & Draws $\le 0$ \\
\midrule
\multicolumn{8}{@{}l}{\textit{Stratum-cluster scheme (198 of 200 replications converged)}} \\
Rate gap (standardized) & $+0.673$ & $+1.491$ & $+0.763$ & $1.152$ & $[+0.60,\,+4.11]$ & $[-2.76,\,+0.74]$ & 0.5\% \\
Burnout (orthogonalized) & $-0.130$ & $-0.056$ & $-0.130$ & $1.363$ & $[-1.76,\,+1.43]$ & $[-1.69,\,+1.50]$ & 75.3\% \\
Friction & $-0.037$ & $+0.352$ & $+0.015$ & $0.506$ & $[-0.25,\,+1.38]$ & $[-1.46,\,+0.18]$ & 48.0\% \\
\midrule
\multicolumn{8}{@{}l}{\textit{Temporal moving-block scheme, block length 6 months (196 of 200 converged)}} \\
Rate gap (standardized) & $+0.673$ & $+3.313$ & $+2.556$ & $3.845$ & $[+0.44,\,+12.18]$ & $[-10.84,\,+0.91]$ & 0.5\% \\
Burnout (orthogonalized) & $-0.130$ & $+1.229$ & $+0.523$ & $2.586$ & $[-3.16,\,+6.16]$ & $[-6.42,\,+2.90]$ & 40.8\% \\
Friction & $-0.037$ & $+1.050$ & $+0.567$ & $2.802$ & $[-3.71,\,+7.69]$ & $[-7.76,\,+3.63]$ & 37.2\% \\
\bottomrule
\end{tabular}
\begin{tablenotes}[flushleft]\footnotesize
\item Stratum-cluster scheme: the observed vintage--coupon--FICO--LTV strata are resampled with replacement, each resampled stratum receives a fresh fixed-effect label so duplicated strata contribute independent fixed effects, and per-replication standardization (including the burnout orthogonalization, which is refit inside every replication) is rescaled to production standardized units. This scheme is structurally silent about the friction coefficient: the friction index varies only over time, so every replication sees the identical series, and the stratum-scheme friction row understates that coefficient's uncertainty by construction. Temporal scheme: the 36 training months are resampled in moving blocks of six with each month's cross-section kept intact, same refit and rescaling; it is the informative scheme for friction and, partially, burnout; its friction SE (2.80) is roughly $5.5\times$ the stratum scheme's (0.51). The ridge penalty is held at the production $\lambda_{\mathrm{ridge}} = 10^{-4}$ rather than re-selected per replication (re-selection would bootstrap a different, coefficient-plus-model-selection estimator). Replications returning non-finite or implausible ($|\beta| > 20$) coefficients are counted as failures, not draws; the point refit reproduces the production coefficients to four decimals on the same 10,176 training cells. Coefficients are in production standardized units (training SD of the rate gap: 221 bp). The stratum-scheme burnout median coincides with the point estimate at reported precision only (unrounded: median $-0.1297$, point $-0.1301$). Mean, median, and the implied bias (mean $-$ point) are reported because both schemes' rate-gap distributions are heavily right-skewed; under that skew the percentile and basic 95\% intervals disagree about zero, so interval-based significance is construction-dependent and the text relies on sign stability only; but the bootstrap sign-share is itself percentile logic, and the pre-committed bias-respecting test of $H_0\!: \beta_g \le 0$ (run \texttt{patha\_sign\_test}; BCa with the 296-refit jackknife acceleration this note previously deferred, plus a month-permutation leg) does not reject under any construction ($p = 0.093$/$0.412$/$0.241$), so the 99.5\% positive share is a description of a right-skewed resampling distribution, not a sign test, and no sign claim for this coefficient survives.
\end{tablenotes}
\end{threeparttable}
\end{table}

\section{Dynamic-Fit Decomposition}\label{app:theil}

Table~\ref{tab:theil} collects the Theil fit diagnostics summarized in Section~\ref{sec:pathb}: $U$ statistics, the bias/variance/covariance decomposition of each estimator's mean-squared error on three transforms of the monthly CPR paths, and the terminal cumulative dollar errors. The $\beta_1 = 0$ null's cumulative-error profile follows from the committed legs by arithmetic. Its terminal cumulative error against the empirical path is $-\$16.6$ billion, versus Path B's $+\$53.8$ billion: the null under-runs the empirical roll-off by approximately the lock-in marginal, by construction. The full monthly profile is committed with the marginal-timing data of Section~\ref{sec:identification}.

\begin{table}[H]
\centering
\caption{Theil dynamic-fit decomposition for the three estimators' monthly CPR paths over the 42-month QT window. Fit diagnostics, not timing credentials: no estimator's co-movement survives detrending (Section~\ref{sec:pathb}), and $U_2 > 1$ means worse than a naive no-change forecast of month-over-month changes.}
\label{tab:theil}
\footnotesize
\begin{threeparttable}
\begin{tabular}{@{}lrrr@{}}
\toprule
 & ABM & Path A & Path B \\
\midrule
Theil $U_1$, levels & 0.412 & 0.438 & 0.271 \\
Theil $U_1$, detrended & 0.807 & 0.839 & 0.980 \\
Theil $U_2$, first differences & 1.226 & 1.009 & 1.003 \\
\midrule
\multicolumn{4}{@{}l}{\textit{MSE shares, levels (\%)}} \\
\quad Bias & 68.1 & 30.6 & 11.3 \\
\quad Variance & 0.4 & 10.5 & 82.5 \\
\quad Covariance & 31.5 & 58.8 & 6.3 \\
\multicolumn{4}{@{}l}{\textit{MSE shares, first differences (\%)}} \\
\quad Bias & 0.1 & 0.3 & 0.1 \\
\quad Variance & 14.0 & 71.4 & 96.8 \\
\quad Covariance & 85.8 & 28.3 & 3.1 \\
\multicolumn{4}{@{}l}{\textit{MSE shares, linearly detrended (bias $\equiv 0$, \%)}} \\
\quad Variance & 0.5 & 43.6 & 93.8 \\
\quad Covariance & 99.5 & 56.4 & 6.2 \\
\midrule
Terminal cumulative error (\$B) & $-673.8$ & $+164.1$ & $+53.8$ \\
\bottomrule
\end{tabular}
\begin{tablenotes}[flushleft]\footnotesize
\item $U_1 \in [0,1]$, where zero is a perfect fit; $U_2$ is relative to the naive no-change forecast. Path A and Path B columns are computed against the hazard-side empirical CPR back-out at the pool's note-rate WAC (the corrected amortization basis; run \texttt{coupon\_convention\_amortization}, committed 2.5\%-convention values in \texttt{theil\_data.json}); the levels-bias shares are the cells the correction moves ($+4.6$pp for Path B), the timing rows move at or below a thousandth. The ABM column retains its committed basis --- the converted ABM reconstruction failed its own replay gate and is not quoted (the ABM's cumulative-error leg uses its own committed back-out). Terminal cumulative errors equal each estimator's trapped-liquidity aggregate minus the empirical \$764.7 billion on its own accounting basis, reproducing Table~\ref{tab:estimators} by construction. Artifact \texttt{theil\_data.json}; inputs and correlations are gated to the committed run manifests.
\end{tablenotes}
\end{threeparttable}
\end{table}

Grouped calibration localizes the same verdict (run \texttt{grouped\_calibration}). I bin the frozen Path A cell predictions into five exposure-weighted prediction-CPR quintiles, with edges fit on the training window only, and cross them with calendar half-years. The training halves calibrate at a pooled predicted-to-realized ratio of 1.000, exact by construction of the exposure-offset fit, with per-half ratios spanning 0.80--1.17. The post-2024 holdout instead drifts into monotonically worsening under-prediction: 0.64, 0.57, 0.44, 0.43 across 2024H1--2025H2, pooled predicted 2.24 against realized 4.21 points. The departure is temporal drift, not a reliability-slope failure: pooled across periods the four upper prediction bins calibrate at ratios 0.94--1.11, and the shortfall concentrates in the lowest-prediction bin (43.5\% of exposure), where the 2024--25 cells accumulate. That is the finer-resolution form of the $U_2 > 1$ diagnosis above. This is a cell-panel CPR object, distinct from the cumulative dollar-runoff split of Section~\ref{sec:patha}, and like every Path A diagnostic it moves no headline figure.

\section{Composition Shift: Estimation Sample versus SOMA Book}\label{app:composition}

Table~\ref{tab:composition} reports a formal covariate-shift analysis between the Freddie Mac data behind the hazard paths and the SOMA book the benchmark spans, on the dimensions both sides carry. The book side is the New York Fed's CUSIP-level SOMA holdings, parsed by the same production code that builds the ABM cohorts and anchor-gated to the committed marginals (book WAC 2.49\%, Ginnie share 20.4\%, term split, vintage-group shares). The API also served the historical June 29, 2022 as-of, whose composition is consistent with those anchors; the artifact reports that as-of as a non-gated secondary block. The shifts on coupon, vintage, and agency are large --- the expected result, disclosed in prose throughout the paper --- and each carries a committed, signed bound as its consequence, listed in the table's final column. The joint agency$\times$vintage cells settle the one netting question the marginals leave open: of the 2022 vintage's 23.1\% of book face, 20.7\% is conventional and 2.4\% is the Ginnie intersection, so vintage and agency marginals cannot be added without double-counting that cell. Geography, loan size, FICO, and LTV do not exist at the CUSIP level of the Fed's published holdings, which carry par value and security description only. Panel (b) supplies the deliverable substitute: the Freddie and Fannie loan-level samples agree to negligible shift on exactly those dimensions. The hazard specification's exposure to the loan-level covariates is bracketed directly by the printed covariate-block endpoints (99.2\%/92.2\% shared-basis, Section~\ref{sec:pathb}).

\begin{table}[H]
\centering
\caption{Composition comparison. Panel (a): the two Freddie objects the hazard paths use --- the estimation universe Path A is fitted on (exposure-weighted at the June 2022 window open) and, indented beneath it, the 75,000-loan draw Path B simulates (by count and on its surviving balance at the window open) --- against the SOMA agency-MBS book (CUSIP-level holdings, anchor-gated). The three are distinct populations and none is mixed within a row; only the unindented rows carry a PSI against the book. Panel (b): cross-agency stability of the loan-level samples on the dimensions the book cannot supply. PSI is the population stability index $\sum_i (p_i - q_i)\ln(p_i/q_i)$; values above 0.25 are conventionally read as large shift, below 0.10 as negligible.}
\label{tab:composition}
\footnotesize
\begin{threeparttable}
\begin{tabular}{@{}p{2.1cm}p{4.7cm}p{4.7cm}p{3.6cm}@{}}
\toprule
Dimension & Freddie side (\%) & Book (\%) & PSI; committed consequence \\
\midrule
\multicolumn{4}{@{}l}{\emph{(a) Freddie side vs SOMA book: estimation-universe rows, with Path B's draw indented beneath}} \\
Agency & UMBS 100.0; Ginnie 0.0 & UMBS 79.6; Ginnie 20.4 & 2.54; static bound \$20--47B, overlay differential $+\$38.5$B \\
Vintage & pre-2017: 0.0; 2017--19: 12.5; 2020: 37.3; 2021: 50.2; 2022: 0.0 & 10.6; 6.0; 16.3; 43.9; 23.1 & 4.33; vintage residual bound \$11.7 billion \\
\quad Path B draw & \emph{count}: 0.0; 60.0; 20.0; 20.0; 0.0 \newline \emph{surviving balance}: 0.0; 27.0; 30.6; 42.3; 0.0 & & no committed draw-vs-book vintage PSI; equal allocation per origination quarter, never reweighted (Appendix~\ref{app:params}) \\
Coupon & $<3.0\%$: 18.0; 3.0--4.0\%: 68.2; $\geq 4.0\%$: 13.8 & 73.9; 18.7; 7.3 & 2.29; full-book reweighting 107.0\% $\to$ 109.1\% \\
\quad Path B draw & \emph{count}: 7.1; 34.7; 58.2 \newline \emph{surviving balance}: 14.0; 58.8; 27.2 & & 3.38 on the draw's origination-balance shares; draw mean note rate 3.86\%; the reweight moves the level, not the marginal \\
Term & 30-year collateral (15-year face excluded from hazard cohorts, disclosed in Section~\ref{sec:limitations}) & 30yr 90.7; 15yr 9.1; other 0.2 & book-only dimension \\
\midrule
\multicolumn{4}{@{}l}{\emph{(b) Freddie vs Fannie loan-level samples (75{,}000 loans each, origination-balance weighted)}} \\
State & \multicolumn{2}{l}{PSI 0.018} & negligible \\
Loan size & \multicolumn{2}{l}{PSI 0.000} & negligible \\
FICO & \multicolumn{2}{l}{PSI 0.005} & negligible \\
LTV & \multicolumn{2}{l}{PSI 0.000} & negligible \\
\bottomrule
\end{tabular}

\end{threeparttable}
\end{table}

\noindent{\footnotesize \emph{Notes to Table~\ref{tab:composition}.} Three populations appear in panel (a), and no row mixes them: the unindented rows are the hazard estimation universe, exposure-weighted at the window open; the indented rows are the 75,000-loan Path B draw, by count and on the surviving balance the aggregation applies; the Book column is SOMA face. All three sides use rounded half-point coupon buckets, but the two Freddie sides bucket borrower note rates while the Book column buckets security pass-through coupons --- a wedge of $+0.80$ points --- so the buckets are aligned in width, not in object. Within the draw the $\geq 4.0\%$ share moves with both conventions, and the two must be read separately: on count weights, rounding to half points raises it from 47.0\% to 58.2\%; holding the rounding fixed, the weighting lowers it from 58.2\% by count to 53.7\% on origination balance and 27.2\% on surviving balance. The WAC pair is mixed-basis, as Section~\ref{sec:pathb} records: the draw's 3.86\% is the unweighted mean of borrower note rates, the book's 2.49\% the face-weighted mean of security pass-through coupons. Every shift is signed in the same conservative direction: the out-of-sample segments prepaid faster, so trapped liquidity is if anything overstated (Section~\ref{sec:pathb}). That signing covers observed speeds in the segments outside the draw's support; it does not cover what the draw's composition within 2017--2021 does to the floor's bind share, which no committed run scores (Appendix~\ref{app:params}). Panel (b) uses the committed Fannie replication sample under its aggregates-only license posture. Artifact \texttt{composition\_shift\_results.json}.\par}

\nocite{cfpb_atr,frbny_soma,fred_actliscou,fred_wshomcb,fred_dspic96,fred_umcsent}

\section{Floor-Sweep and Contamination Mechanics}\label{app:floormech}

This appendix collects the mechanics behind the calibration box of Section~\ref{sec:robustness-floor}: the floor-level sweep, the marginal as a curve in the imported elasticity, the separation of the two refinance contaminations behind the off-window floor anchors, the out-of-sample anchor construction, the calibration catalogue for the headline marginal row, and the inference ladder on the floor read (Table~\ref{tab:ladder}). The three exhibits themselves are read where the headline is set, in Section~\ref{sec:hazard} (Tables~\ref{tab:floorband}, \ref{tab:lowband}, and~\ref{tab:oosfloor}).

Table~\ref{tab:floorband} reports the lock-in marginal at every cell of the resulting calibration box. Three facts shape the reading. First, the pre-committed stability threshold ($\pm 2$ points across the empirically plausible 3--5\% floor range) fails: at the central elasticity the marginal spans $+10.8$ to $+5.5$ points across that range. The mechanism is structural. The hazard specification takes a maximum against the floor, so wherever the floor binds the elasticity is inert, and raising the floor mechanically crowds out the lock-in channel. Second, both pre-committed monotonicities hold: the marginal decreases in the floor within every band row and increases in the band at every floor. The box maximum lands at the predicted corner, $+13.2$ points ($+\$100.7$ billion) at a 2\% floor on the 7.7\% band edge. Third, every cell is positive, the minimum $+2.1$ points at a 6\% floor on the 5.5\% band edge, and every marginal in the box is basis-invariant, because the shared-accounting netting is common to the central and null legs of each cell (Section~\ref{sec:robustness-hybrid}). That positivity is a consistency check on the specification, not an identifying result, since the window's rate configuration fixes the sign before any floor or band is chosen (Section~\ref{sec:identification}). The verification reading of Section~\ref{sec:pathb} accordingly weakens from a point to a bounded interval. The lock-in elasticity's identified contribution is $+9.2$ points at the production calibration, $+5.0$ to $+12.6$ points across the empirically plausible floor range, and $+2.1$ to $+13.2$ points across the full box. The off-window re-measurement below places the headline reading at $+4.3$ to $+6.8$ points. The aggregate level remains floor-dominated throughout.

The box sweeps the elasticity only over its published band, and that band's lower edge is a literature boundary, not an evidential one. The in-sample estimates cannot reject zero for the elasticity (Section~\ref{sec:pathb}), so a reader who weighs that non-rejection over the imported band needs the marginal at values the box never visits. A pre-committed extension (run \texttt{band\_low\_extension}) re-runs the paired legs at quarterly-decline values $\delta \in \{0, 1.0, 2.0, 3.0, 3.25, 4.0, 5.0\}\%$ below the band at both anchor floors; its grid, its parity gates against the committed production and off-window cells, and its monotonicity check were fixed in the script header before any run. Table~\ref{tab:lowband} reports the resulting curve beside the band's own cells. The curve is monotone at both floors, passes through zero at $\delta = 0$ by construction, and is concave at the off-window floor through the censoring mechanism of Section~\ref{sec:identification}: raising the response crowds the hazard into the floor (the bind share climbs from 35.8\% to 73.5\% along that curve), so each increment of elasticity buys less marginal. At half the adopted central, $\delta = 3.25\%$, the marginal is $+5.1$ points ($+\$39.1$ billion) at the in-sample floor and $+3.5$ points ($+\$26.5$ billion) at the headline off-window floor. This curve prices the dependence on the imported elasticity rather than identifying it: how small the marginal could be, if the moving response is weaker than the literature band, is read off the curve, and the identified content remains the bounded range stated in Section~\ref{sec:identification}.


Those two checks vary the floor's level and its functional form. A third pre-committed check interrogates the level's provenance: whether the 4\% figure is an artifact of being measured in-window on the locked-in cohorts the marginal later decomposes. Two runs carry it, \texttt{out\_of\_window\_floor} and \texttt{oos\_identification}, and spec, anchor definitions, acceptance thresholds, and parity gate were fixed in each script header before its runs, the ex-ante convention the replication package's ledger records. Run in-window on the deeply out-of-the-money 2023--24 discount cohorts, the identical exposure-weighted-CPR procedure reproduces the production floor at 3.8--3.9\% annual CPR, flat across seasoning cuts and validating the measurement where it is used. Applied out-of-window to pooled 2017--2019 discount-cohort turnover, before both the QT window and the 2020--21 refinancing wave, it recovers a materially higher 6.1\%. That pooled reading, however, does not survive its own diagnostics, and I do not take it as the operative off-window anchor.

Two contaminations are separable, and separating them is what changes the reading. The first is seasoning. A floor meant to capture strictly involuntary turnover---death, divorce, relocation---should be approximately flat in loan age, but the pooled 2017--2019 out-of-the-money turnover \emph{rises} monotonically with it: 3.68\% below twelve months, 5.78\% at twelve to twenty-four, and 10.99\% at twenty-four to thirty-six. The second is rate regime. Splitting the pooled window at the 2018/2019 turn separates a rising-rate leg (the 30-year rate moved 4.03\% to 4.64\% over 2018, suppressing refinancing incentive) from a falling-rate leg (4.46\% to 3.72\% over 2019, a refinancing wave). The legs read very differently: the 2018 leg returns 4.70--5.33\% across the out-of-the-money gap grid at the age-$\geq$-12 cut, while the 2019 leg returns 6.91\% and carries the entire 10.99\% age-$\geq$-24 cell. The pooled 6.1\% is therefore a refinancing-inflated composite, not a baseline turnover read, and the 2019 leg is disqualified on both diagnostics at once. The rate-regime half of that disqualification, as stated, compares the two legs at different discount depths; the matched-depth reconciliation below repairs the comparison and reports what survives it, which is less than the sentence just written implies.

What remains is the 2018 rising-rate leg, and its status must be stated precisely, because the paper's off-window anchor now rests on it. Its mature-age seasoning test is \emph{not computable}: the leg contains no out-of-the-money cohort-months at all past the twenty-four-month age cut, because a loan that is out of the money while rates are rising is necessarily a recent, low-coupon origination and therefore young. The 2018 leg is thus clean by rate-regime construction, not by test, and its one computable age contrast (3.78\% to 5.33\% across the twelve-month cut) is confounded by the PSA ramp, on which young loans genuinely turn over less, so it is not evidence about the contrast the leg is missing. The out-of-agency read is that in every cell where a mature contrast exists at all, mature turnover is the higher one, and Section~\ref{sec:identification} reports by how much. A second limitation concerns depth. The gap-$\leq$-two-point cell at which the in-window anchor is measured does exist off-window: the 2018 leg's gap runs to $-2.87$ points, and it holds 35 cohort-months reading 3.440\% annual CPR. Grids that stop at half a point do so because their thresholds were fixed at zero, a quarter point, and half a point in the script header, not because the rate history forbade anything deeper. The true limitation is exposure collapse, and a matched-depth run (\texttt{matched\_depth\_reconciliation}) measures it: down the 2018 leg's depth ladder the share of its gap-$\leq$-0 exposure that survives runs 100\%, 83.7\%, 72.9\%, 35.9\%, 11.3\%, 3.9\%, and 0.02\%, so the two-point cell carries \$0.282 billion of \$1{,}264.5 billion. That cell is present and worthless, and an ex-ante support rule excludes it before any reading is taken: a depth is admissible for a leg only if the leg retains at least 10\% of its gap-$\leq$-0 exposure there. These are the reasons the off-window floor resolves to a \emph{range}, and I report it as one.


\paragraph{Calibration entries for the headline marginal row.} Fannie Mae read of the same selection rule on the pooled off-window period 5.52\%, above the clean band and bracketing the marginal below $+4.3$, its like-for-like Freddie counterpart 5.185\% rather than the 2018 leg's 4.991\%, so 0.337 points of that distance is agency and the rest is period; $+3.9$ to $+7.5$ joint with the band; additive form $+11.2$ (form-conditional hull $+3.5$ to $+13.1$); concave transform $+5.1$; the two \emph{jointly} $+9.2$, an interaction of $-1.47$ points, so these two entries do not add; Ginnie overlay $+4.4$ (cuts the same way as the age-standardized read; now run jointly with it in run \texttt{b5\_joint\_cell}: composed $+2.9$, within $0.1$ of proportional); floor-stability check $+\$45.1$B; baseline-level sweep 75--150 PSA $+0.9$ to $+15.6$ (a convention range, no coverage property; run \texttt{psa\_level\_sweep}); covariate endpoints move this marginal by only $+0.1$ while spanning 8.2 points of level; survival-selection attenuation $+5.1$/$+4.6$/$+3.7$ at $\eta = 0.86/0.73/0.54$, the frailty-identity mapping $\eta = S_{\mathrm{pre}}^{\xi}$ at the measured $S_{\mathrm{pre}} = 40{,}234/75{,}000$ and $\xi = 0.25/0.5/1$ (run \texttt{attenuation\_sensitivity}; a transport sensitivity on the imported coefficient, listed in Table~\ref{tab:assembly} as a bracket in $\xi$ rather than as a correction).

\paragraph{The inference ladder on the floor read.} Section~\ref{sec:identification} quotes the restricted wild-cluster inversion rung of Table~\ref{tab:ladder} as the binding layer; the full ladder, with its degrees of freedom, follows.

\begin{table}[!t]
\centering
\caption{The inference ladder on the binding layer---binding among the layers with a coverage property; the unestimated baseline-level convention spans wider. Every row is one rung on that layer --- the floor read's own sampling error, propagated through the frozen floor-to-marginal grid at the central elasticity $\delta = 6.5\%$ --- so no row's width comes from the elasticity, which is swept separately in Table~\ref{tab:lowband}. All reads are on the 2018 leg's mid-grid anchor, which 31 stratum clusters carry; the last two rungs re-cluster that same read on the six calendar months it populates and on both margins at once. The restricted wild-cluster inversion is the primary construction and the interval this paper quotes, the Webb-weight wild-cluster bootstrap-$t$ being retained as a reported rung; the percentile read is reported demoted rather than as a competing interval. Runs \texttt{floor\_uncertainty} (percentile), \texttt{floor\_inference\_correction} (Rademacher wild-$t$; CR1--CR3 at $t(30)$), \texttt{floor\_inference\_correction\_v2} (Webb primary, Bell--McCaffrey degrees of freedom, restricted inversion), and \texttt{floor\_inference\_correction\_v3} (month and two-way clustering).}
\label{tab:ladder}
\footnotesize
\begin{threeparttable}
\begin{tabular}{@{}llll@{}}
\toprule
Rung & Marginal (pp) & Degrees of freedom & Status \\
\midrule
Percentile cluster bootstrap & $+3.0$ to $+8.0$ & --- (1{,}000-replicate) & demoted; under-covers at 31 clusters \\
CR1 $t$ & $+3.2$ to $+8.3$ & $t(30)$, $G-1$ & no leverage adjustment \\
CR2 $t$ & $+3.0$ to $+8.6$ & $t(30)$, $G-1$ & conventional df \\
CR3 $t$ & $+2.8$ to $+8.8$ & $t(30)$, $G-1$ & conventional df \\
Wild-$t$, Rademacher & $+2.8$ to $+8.7$ & bootstrap, 9{,}999 replications & the committed construction, replayed \\
Wild-$t$, Webb & $+2.9$ to $+8.7$ & bootstrap, 9{,}999 replications & covers 90.8/91.4; below the pre-committed bar \\
CR1 $t$, Bell--McCaffrey & $+2.8$ to $+8.8$ & $t(6.2)$, data-driven & small-sample df \\
CR2 $t$, Bell--McCaffrey & $+2.4$ to $+9.1$ & $t(5.1)$, data-driven & small-sample df \\
CR3 $t$, Bell--McCaffrey & $+2.3$ to $+9.6$ & data-driven, CR3's own & small-sample df \\
Restricted wild-cluster inversion & $+2.3$ to $+9.1$ & 401-point grid inversion & primary; the binding layer (covers 94.3/95.1) \\
CR1 $t$, month clusters & $+3.1$ to $+8.4$ & $t(5)$, $G-1$ & wider on df alone \\
Two-way, stratum $\times$ month & $+2.3$ to $+9.3$ & $t(5)$, $\min(G)-1$ & lower edge censored \\
\bottomrule
\end{tabular}
\end{threeparttable}
\end{table}

\noindent{\footnotesize \emph{Notes to Table~\ref{tab:ladder}.} The Bell--McCaffrey degrees of freedom are data-driven and estimator-specific, so one value cannot serve all three cluster-robust variants: the 6.2 and 5.1 shown are CR1's and CR2's, and CR3's own is smaller still. CR1 applies no leverage adjustment --- it is the narrowest of the three sandwiches --- and is the standard error the wild rows studentize with, so its two rows are the ladder's baseline rather than competing reads. Its Bell--McCaffrey row printing CR3's conventional-df interval is an accident of this leverage profile, not an identity. Unrounded, CR1 at Bell--McCaffrey degrees of freedom runs $[+2.753, +8.796]$ and CR3 at conventional degrees of freedom $[+2.773, +8.774]$ --- two distinct intervals that agree only once each endpoint is taken to one decimal. The effective cluster count is 5.9 by the same Herfindahl-inverse convention as the loan/stratum scheme's 25.8, against 31 nominal clusters. The run's recorded maximum single-cluster leverage is 0.33, in the 2017-vintage, 4.0\% coupon, FICO $740+$, LTV ${\leq}80$ stratum. That is ten times an equal 31-cluster share and 1.9 times the next largest (the same cohort at LTV ${>}80$, 0.17). That concentration is why the wild-cluster correction is applied and the percentile read demoted. It does not rank the wild construction against the Bell--McCaffrey degrees of freedom, and I do not claim it does. The wild rows re-draw cluster-level weights at that same leverage profile; the Bell--McCaffrey rows keep the sandwich and correct the reference distribution's degrees of freedom for it. At 5.9 effective clusters neither correction dominates. So the primary read is now adjudicated, not asserted. Run \texttt{fewcluster\_coverage}'s pre-committed rule --- narrowest construction with Gaussian coverage $\geq 93.0\%$ and heavy-tail coverage $\geq 91.0\%$ --- selects the restricted inversion, the artifact's own \texttt{binding\_construction} field, whose $[+2.3, +9.1]$ this paper quotes as the binding layer. The Webb wild-$t$, at 90.8\%/91.4\%, falls below that bar, and its $[+2.9, +8.7]$ is retained as a reported rung, not the quoted interval. The run's own \texttt{landing\_branch} string records L1 (Webb retains), which contradicts the rule its spec froze. The string is sha-pinned and left unedited, and this note is the adjudication. CR2 at Bell--McCaffrey degrees of freedom ($+2.4$ to $+9.1$ points), the other qualifying rung, stands beside the restricted inversion, over-covering and wider. The three rungs printed wider still understate their own width. CR3 at Bell--McCaffrey degrees of freedom, the restricted inversion and the two-way rung have \emph{lower} endpoints at floors above the committed grid's 6.0\% end and are censored to it, and the marginal falls as the floor rises, so all three run below the $+2.3$ printed. Their upper endpoints are interior. The last two rungs re-cluster the same read rather than re-estimating it: the month rung groups the 137 cohort-months by the six calendar months the 2018 leg populates, July to December, and the two-way rung adds the CR1 variance clustered on the 31 strata to the one clustered on those six months and subtracts the one clustered on their 137 intersection cells, the Cameron--Gelbach--Miller construction, carried at $\min(G)-1 = 5$ degrees of freedom. Both come in wider than the CR1 rung they re-cluster, so the month-level shock is priced here rather than dismissed. They widen for different reasons, and only the two-way rung widens because its variance is larger. Month clustering widens the ladder, as I pre-committed it would and for the reason the specification gave, the drop from 31 clusters to six: CR1 goes from 5.20 to 5.33 points of width, both unrounded. Its variance runs the other way, which that specification allowed. The month-clustered CR1 standard error comes in 18\% \emph{smaller} than the stratum-clustered one, so at this cluster unit the common shock adds no variance, and the widening is the six-cluster reference distribution alone. Clustering on both margins at once does add variance---the Cameron--Gelbach--Miller standard error runs 16\% \emph{above} the stratum-clustered one---and that rung is the wider of the two new rows; whether it would be computable at all I deliberately did not predict. My second pre-commitment missed. I expected the month rung's lower marginal edge to be censored at the grid's 6.0\% end, as the two widest stratum rungs already are, and it is not---both endpoints of the printed rung are interior, and of the two new rows only the two-way one censors. Among the month rung's own variants only the unprinted Rademacher one reaches past that edge. Six clusters is small in more than one direction. The month rung's own Bell--McCaffrey degrees of freedom is 3.9---below the four my pre-committed too-few-clusters rule was written to protect, though that rule's own $G < 5$ trigger did not fire. Six clusters admit only $2^{6} = 64$ distinct Rademacher sign vectors, so a Rademacher wild-$t$ tail probability at this cluster unit cannot resolve below $1/64$, which is why the Webb six-point weights are primary here, as they are above.\par}

\section{Behavioral Extension Mechanics}\label{sec:robustness-extensions}

This appendix gives the mechanism behind each of the three behavioral extensions summarized in Section~\ref{sec:abm-results}. As Appendix~\ref{sec:method-frictions} describes, loss aversion steepens the penalty every household faces. On its own that would suppress simulated mobility below the historically observed 4--5\% baseline turnover floor, so to preserve the floor I re-calibrate the mobility-desire parameter upward, nearly tripling it. That recalibration mechanically inflates the model's overall responsiveness to rate gaps, so the jump to 54.9\% is an artifact of holding the turnover floor fixed, not a validation of loss aversion as a mechanism. Because that floor is anchored at maximum lock-in, the sign of the recalibration's effect on rate-sensitivity is not fixed. A specification change that eases turnover at the anchor lowers the recalibrated mobility scale and deflates rate-sensitivity: in the real-covariate cross-design of Section~\ref{sec:robustness-crossdesign} the scale falls to 36,086 from the synthetic 43,883. The headline recovery, itself resting on a floor-anchored calibration, inherits that indeterminacy. That is why the ABM's negative verdict is read only alongside the cross-design hedge and never on its own. The second extension, income-scaled curtailment, pulls the estimate down to 45.1\%. Curtailment is a strictly additive drag on prepayment, so this decline is what I predicted ex ante from a diagnostic correlation between the ABM's over-prediction and disposable-income growth (Appendix~\ref{sec:method-frictions}). The third extension, multi-vintage coupon composition, pulls the estimate down further, to 33.7\%. A lower weighted-average coupon widens the gap to prevailing market rates, and my hypothesis was that this should increase predicted trapped liquidity, not reduce it. The inversion suggests the ABM's CPR mechanism is not translating a wider coupon gap into slower prepayment the way the household-choice framework predicts it should. I have not diagnosed this inversion, and its implication cuts in both directions. An aggregate response to the rate gap running backward is equally consistent with ``household choice explains little'' and with ``this specific ABM's gap response is mis-specified.'' So the ABM's recovery share is never quoted in this paper without the cross-design qualifier of Section~\ref{sec:robustness-crossdesign}, and the paradigm attribution inherits that section's hedge. It is also why the two-estimator contrast, not the ABM's point estimate, is what the paper asks the reader to carry. I cannot rule out that a differently specified household model would perform better, but that possibility does not bear on the results reported in Section~\ref{sec:abm}.

The wait-and-see freeze now carries a direct sensitivity run, scored on the frozen Danish-anchor production accounting (trapped \$84.5 billion at the production 150 bp trigger and 20\% share; parity gate $+\$0.02$ billion). Removing the freeze entirely drops trapped liquidity to \$44.6 billion, an isolated contribution of $+\$39.9$ billion, or $+5.2$ points of benchmark. Moving the six-month trigger to 100 or 200 bp spans \$105.9 to \$70.5 billion, so the trigger's placement matters at the few-points scale, and its author-calibrated status remains a disclosed limitation. Sweeping the freeze share with the floor anchor retained gives \$65.1 billion at a 10\% share, \$81.5 billion at the externally implied 18.3\%, and \$103.3 billion at 30\%. Re-anchoring the share to the Liebersohn--Rothstein-implied value moves recovery by only $-0.4$ points of benchmark. The calibrated mobility scale and floor CPR are bit-identical at every share, confirming that the freeze cannot compromise the baseline turnover floor. The front-end DTI gate gets the parallel treatment. Shifting the 43\% wall to 36\% or 50\% on the same frozen surface moves ABM trapped liquidity to \$153.3 and \$119.7 billion, from \$84.5 billion at the production 43\%. The turnover floor is re-anchored into the empirical 4--5\% band at every threshold. The gate is therefore material to the ABM's own estimate while never breaching the floor. That sweep is non-monotone in the wall, and the production 43\% sits at its minimum, which invites the reading that the paper's own parameter choice is the one minimizing its own estimate. That reading does not survive decomposition (run \texttt{dti\_nonmonotonicity\_decomposition}). The floor-retention rule re-calibrates the mobility-desire scale only where the fixed-scale floor would leave the 4--5\% band, and across these three thresholds that happens exactly once: at 50\%, where loosening the wall lifts the fixed-scale floor to 5.09\% and the search cuts the scale from 43{,}883 to 41{,}934. Holding the scale at its production value instead, the sweep is monotone (\$153.4, \$84.6 and \$46.4 billion at 36, 43 and 50\%), and the 50\% leg's entire rebound is the re-calibration, worth $-\$73.4$ billion. The frozen and re-calibrated legs are bit-identical at 36 and 43\%, where the rule is a no-op, which is what identifies the channel. The wall is therefore monotone in the direction the mechanism implies, and the published non-monotonicity is the free parameter absorbing the floor constraint. Nothing here moves a headline: the ABM is a stress-test device and the hazard marginal carries no DTI gate.

\section{Duration Quantification: Approximate Portfolio Weighted-Average Life}\label{sec:robustness-wal}

The mechanism the paper quantifies is duration extension, so I express it in maturity units. Table~\ref{tab:wal} is printed with the cap-design arithmetic in Section~\ref{sec:discussion-capdesign}. It reports the SOMA book's approximate weighted-average life (WAL) under scheduled amortization only, under the empirical CPR path, and under each estimator's implied path, at the QT window's start and end. This appendix documents its construction and reads it against each estimator.

Read against the empirical column, each estimator's CPR error maps into effective-maturity error: the ABM's over-predicted prepayment implies a book roughly 3.4 years too short at window start, Path A's under-prediction a book 1.5 years too long, and Path B lands within 0.3 years. The table's final row gives the no-shock world, with prepayments continuing at the 22.81\% mean 2021 realized CPR and implying a 3.4-year WAL; against it, the empirical 9.4-year WAL is a 6.0-year extension. That duration gap is what the paper's dollar benchmark prices month by month. The two Danish rows read the counterfactual in the same maturity units. Feeding the Danish legs' mean CPRs through the identical calculator gives a rule-only WAL of 9.1 years at the window's start against the U.S. leg's 9.7---a 0.6-year rule-only shortening, with the Danish-level bracketing case at 10.8. The rule-only leg's extension is additionally rate-invariant on the rate-rise side by the identified moving-flatness (Section~\ref{sec:abm-danish}). I deliberately price extension, not convexity: option-adjusted duration and convexity require a term-structure and prepayment-option model outside this paper's scope. \citet{hanson2014} shows that the convexity so omitted has aggregate rate-dynamic consequences of its own, which is a reason to flag the omission. The term ``convexity-neutral'' is used strictly in this rate-rise sense.

\section{Fifteen-Year MBS Fold-In and Path B Rate-Gap Units}\label{sec:robustness-pipeline}

Two features of the production pipeline set headline figures cited throughout this paper: how fifteen-year MBS are treated, and the units in which Path B applies its rate-gap coefficient. The replication package's ledger catalogues the figures each supersedes.

The first is the fifteen-year MBS fold-in. The SOMA cohort parser buckets 15-year holdings separately by coupon and never merges 15-year with 30-year buckets, and 15-year scheduled amortization enters both the empirical CPR back-out and the ABM's simulated roll-off. Restricting to 30-year MBS covers 90.6\% of face value; the 15-year book, approximately 9.1\% of holdings, then contributes to total portfolio size but not to cohort-correct amortization or CPR weighting. Folding it in brings coverage to 99.8\% of SOMA face value at a SOMA weighted-average coupon of 2.49\%. Table~\ref{tab:wal}'s note reconciles these coverage components with the audit parse's 90.7\%/9.1\% term split.

I tested, but did not adopt, a fully native treatment in which 15-year loans also draw voluntary CPR from a 15-year-specific surface. It produced out-of-bounds predicted values. The ABM's payment-delta gate---the change in \eqref{eq:annuity}'s monthly payment upon refinancing at the prevailing rate, which drives loss-aversion and mobility-desire responses elsewhere in this paper---breaks down for seasoned 15-year borrowers, whose same-term replacement payment is nearly flat relative to their current payment. Native 15-year surfaces predicted 32\%--61\% CPR against an empirical range of roughly 5\%--8\%, enough to flip the ABM's aggregate trapped-liquidity estimate to a large negative value, far outside the empirically observed range. Production therefore uses a structural-only fold-in: 15-year cohorts contribute real weights and real scheduled amortization, while voluntary CPR for those cohorts is drawn, in the frozen headline specification, from the same-coupon 30-year surface. This is a scope limitation for the production headline. It does not bind the framework: a 15-year-specific behavioral gate---a pure rate-lock penalty for short-amortization loans---is implemented and is the basis of the Berger run's calibration-stage freeze (Table~\ref{tab:danish}, note b). Its adoption for the production headline is deferred to the pre-submission freeze.

The second is the units in which Path B applies its rate-gap coefficient $\beta_1$, calibrated from \citet{liebersohn2024} via the survival-function transform described in Section~\ref{sec:pathb} and expressed as a log-hazard-ratio per 100 bp increase in the rate gap. Applied to the gap in decimal units, the coefficient is numerically inert: at a representative QT-era state (a 3.0\% coupon against a 6.8\% market rate) it yields a multiplier of approximately $\times$1.0026, against the intended suppression of approximately $\times$0.77. Production applies it as $\exp(100\,\beta_1\,\tilde g_{i,t})$, with $\tilde g_{i,t}$ the decimal form of \eqref{eq:pathB}'s rate gap, equivalently $\exp(\beta_1 g_{i,t})$ in \eqref{eq:pathB}'s 100 bp units. The gap is negative when the borrower is locked in, so \eqref{eq:beta1}'s positive $\beta_1$ suppresses prepayment. This application is consistent with Path A's convention of a positive coefficient applied to the gap in comparable units. Production also expresses the Danish counterfactual's rate-gap function in rate-equivalent units, so that the buyback discount properly offsets the locked-in spread. That function is live in the mechanism-substitution variant. The production Danish leg's prepay response, under either Berger-recalibrated anchor of Section~\ref{sec:abm-danish}, does not consume that function's classification, which is computed all the same, as Section~\ref{sec:pathb} discloses. Section~\ref{sec:pathb} reports the results under this convention; the superseded \$747 billion (97.7\%) microsimulation is catalogued in the replication package's ledger.

Both are structural or mechanical, not changes in modeling philosophy, and both are additive. Disabling either one---30-year-only filtering, or the decimal-unit $\beta_1$ convention---exactly reproduces the corresponding superseded figures catalogued in the replication package's ledger, which I treat as evidence that each is isolated to its intended code paths.

\section{Seasonalizing the Turnover Floor: Floor Dispersion, and Why No Timing Claim Follows}\label{sec:robustness-seasonalfloor}

The three checks above vary the floor's level, its functional form, and its provenance, holding it constant within the window in each case. Two runs in this paper let the floor vary within the window instead: the cyclical variant of Section~\ref{sec:pathb}, which indexes it to the rate path, and the pre-committed seasonal check taken up here, which indexes it to the calendar. Both move the marginal; this check is the one built to isolate why. It asks two questions at once---whether letting the turnover floor vary by calendar month moves the lock-in marginal, and whether the resulting monthly path tracks the empirical benchmark well enough to license a timing claim. The run is \texttt{seasonal\_floor\_timing}; its spec, leg definitions, placebo families, and a frozen timing rule were fixed in the script header before any run executed. The answers are asymmetric, and the second is the more consequential: the marginal does move, but not for the reason the design was built to test, and the timing rule is met without being informative.

The calendar profile is a twelve-cell exposure-weighted annualized turnover CPR measured on deeply out-of-the-money cohort-months. I estimate it on the full cohort-month panel (7{,}756 cohort-months, \$74.2 trillion of exposure, pooled annual CPR 4.187\%), not on the QT window. The window-restricted variant differs in four calendar cells and drives nothing here; I state the basis because an earlier module docstring described the profile as in-window and was wrong. The profile peaks in June, with a peak-to-trough ratio of 2.13. I run four legs, each as a paired central and $\beta_1 = 0$ null: the production constant 4.0\% floor, a constant 4.5\% floor, a \emph{shape-only} leg carrying the calendar profile with its exposure-weighted mean pinned at 4.0\%, and the pre-committed 4.5\%-mean seasonal leg. The shape-only leg is the one that isolates the calendar, because the pre-committed leg confounds shape with the level move the committed sweep already prices. Every leg re-runs the production microsimulation with only \texttt{INVOLUNTARY\_CPR\_ANNUAL} varied by calendar month, and the constant-4.0\% control reproduces the committed production artifacts bitwise (central \$818.530 billion, null \$748.185 billion, marginal $+\$70.345$ billion, $+9.20$ points).

Giving the floor a calendar profile at a held-constant 4.0\% exposure-weighted mean moves the marginal from $+\$70.345$ billion to $+\$65.448$ billion, $+9.20$ to $+8.56$ points: a shape effect of $-\$4.897$ billion. A permutation family disposes of the natural reading of that number. Assign the \emph{same} twelve floor values to scrambled calendar months, so that the multiset of values is preserved exactly and only the month labels change, and 89.2\% of the shape effect survives on average across fourteen draws (69.9\% to 102.8\%); the permuted marginals run \$65.311 to \$66.924 billion against the true \$65.448 billion. Calendar order is very nearly irrelevant to the movement, so the movement is not a calendar effect.

It is instead a Jensen effect operating through the floor's combination rule. Production combines the floor and the voluntary hazard by a hard maximum (\texttt{FLOOR\_MODE = 'max'}, the semantics \eqref{eq:pathB} and the form test above discuss), and at the production floor the floor binds in 36.3\% of evaluated loan-months. Under a maximum, dispersion in the floor is one-sided: where the voluntary hazard already exceeds the floor, the low floor cells are truncated away and contribute nothing, while the high cells bind and lift the effective hazard. Any dispersion in the floor therefore raises book-wide prepayment regardless of which month carries which value---exactly what the permutation family measures. A convexity-matched control separates the small remaining piece. The annual-CPR-to-monthly-hazard transform is itself convex, so the twelve-cell profile is equivalent in level to a flat 4.0090\% floor rather than a flat 4.0000\% one; running that flat leg gives $+\$70.156$ billion, a $-\$0.189$ billion move. Level leakage is thus 3.9\% of the shape effect, and floor dispersion under the maximum rule is the remaining 96.1\%. I report this as a floor-dispersion and functional-form result. It is not evidence that seasonality moves the marginal.

Level and shape do not separate cleanly, and I do not present an additive split as though they did. Measured against the constant-4.0\% control, the level move contributes $-\$12.704$ billion ($-1.66$ points) and the shape move $-\$4.897$ billion ($-0.64$ points); the two together come to $-\$16.040$ billion ($-2.10$ points), leaving an interaction residual of $+\$1.560$ billion ($+0.20$ points). On the simulated monthly paths that residual is larger relative to what it qualifies---mean absolute interaction residual 0.164 points of monthly CPR against mean absolute total movement 0.533 points, so the interaction runs 30.9\% of the total movement. Any level-plus-shape decomposition of this floor must carry the residual; quoting either component alone misstates both.

What the experiment does buy is a narrow bound. Holding the turnover floor flat within the window is a modeling convenience, and this run bounds how much of the constant-floor marginal that convenience was responsible for: at most $\$4.897$ billion of the $+\$70.345$ billion in-sample marginal, 7.0\% of it. That same absolute bound is 11.5\% of the $+\$42.6$ billion headline marginal it is measured against. Of even that, 96.1\% is the maximum rule's response to dispersion. The identified lock-in marginal is not an artifact of a flat floor. That is the whole of what the seasonal design establishes, and it establishes it about levels only.

\paragraph{The timing concession.} The design's second purpose was to test whether a seasonalized floor earns a monthly-timing claim, under a rule frozen before the runs: Theil $U_2$ on first differences below one \emph{and} a positive detrended lag-0 correlation against the empirical benchmark. That rule is met as written by both seasonal legs (shape-only $U_2 = 0.995$, $r = +0.016$; pre-committed leg $U_2 = 0.990$, $r = +0.070$), and I report it as met, because it was frozen ex ante and I will not reopen a rule after seeing which way it fell. I nonetheless claim nothing from the pass, because five diagnostics computed afterwards show the rule cannot discriminate. First, the benchmark carries it. The SOMA back-out has four exact-zero months (June 2022, February 2023, April 2024, September 2025), three of them immediately followed by a spike of up to 14.01\% annualized. The zeros are now traced to their mechanism: each is clip-induced, a month whose reported rolloff falls below the scheduled-amortization floor so that the back-out's non-negativity clip pins it at exactly zero, with the displaced mass posting into the following month's reading (unclipped values $-1.59$/$-0.49$/$-0.84$/$-0.84$ points; run \texttt{h1\_zero\_months\_diagnosis}) --- a month-boundary allocation artifact of the reporting series, not a behavioral observation. The obvious alternative --- rebuilding the series from published monthly principal payments, or from per-CUSIP factor changes --- was probed rather than assumed, and neither route is usable. No monthly SOMA principal-payment series is published: the summary endpoint carries holdings levels only, and both monthly endpoints return HTTP 400. The per-CUSIP route does exist but recovers face levels rather than principal payments --- gross face declines run \$2{,}186 billion against a \$652.8 billion net, since face also falls on reinvestment settlement and roll --- so netted over 192 weekly as-of dates it reproduces the committed month-end differencing, \$652.7517 billion against \$652.7517 billion, and returns the same four low months at $-\$2.0$, \$4.5, \$3.7 and \$3.6 billion against a \$16.8 billion window median. Three of the four sit on a week the New York Fed republished unchanged and June 2022 does not, and republication is ordinary rather than diagnostic --- 59 of those 192 weeks and 36 of the 44 months they span carry one --- so what the rebuild establishes is that the alternative source is unavailable and the zeros survive it, not a second mechanism (run \texttt{benchmark\_monthly\_rebuild}). The June 2022 zero opens the window, is followed by an ordinary month, and so contributes one first difference where the others contribute two. Those pairs touch 7 of 41 first differences in total but carry 0.813 of the empirical differenced signal and 0.796 to 0.811 of every leg's $U_2$ numerator. Repairing them by reallocating each spike evenly across its pair, conserving the pair total, flips both seasonal legs from pass to concede and turns the shape-only leg's detrended correlation from $+0.016$ to $-0.134$. Second, a calendar-rotation placebo moves the true profile to the wrong months while holding its amplitude and dispersion exactly; the rule is cleared by 1 of 11 wrong-month rotations for each seasonal leg, and the five-month rotation \emph{beats} the true profile on $U_2$ (0.9934 against 0.9947). Third, 3 of the 14 scrambled calendar orders above clear the rule, and the best of them attains $U_2 = 0.975$, better than either real seasonal leg. Fourth, and decisively, the $\beta_1 = 0$ null clears the rule on both constant-floor legs, the 4.0\% null doing so with $U_2 = 0.9997$ and a detrended correlation of $+0.115$, stronger than either seasonal central leg's. A criterion satisfied by a specification with the lock-in elasticity removed cannot be evidence for the lock-in mechanism. Fifth, the correlations are not distinguishable from zero on 42 months in any case ($p = 0.921$ and $p = 0.661$), and the leg that isolates the mechanism is the fragile one: the shape-only leg's verdict flips on 5 of 42 single-month deletions while the pre-committed leg's flips on none.

Two further diagnostics point the same way without being needed for the verdict. The seasonal legs concentrate 79.7\% (shape-only) and 89.4\% (pre-committed) of their simulated detrended variance in the twelve-month calendar harmonic; the empirical benchmark carries 0.8\% of its own detrended variance there. The model acquires calendar structure the benchmark can neither confirm nor refute. On the committed cross-correlation metric, meanwhile, the seasonal legs peak at lag $+6$ with a negative correlation, on the boundary of a 42-month window, which I mark non-interpretable.

The right reading of this is a finding about the criterion, not a failure of the experiment. The frozen rule was a reasonable ex-ante test, and it turns out to have no power against the alternatives that matter: a lumpy benchmark, wrong-month profiles, scrambled calendars, and an elasticity-free null all satisfy it. I therefore restrict every claim in this paper to levels. No estimator here earns a monthly-timing claim, which extends to the seasonalized floor the levels-only concession the production microsimulation's monthly path already carries, where the raw lag-0 correlation does not survive detrending (Section~\ref{sec:pathb}). This restricts what the design identifies; it retracts no level result, and no aggregate, marginal, or bound reported elsewhere in this paper depends on timing.

\paragraph{Carried caveats.} Four qualifications travel with the seasonal floor wherever it is quoted, and none of them is repaired by the runs above. First, the profile's identification is circular in the same way the in-window floor's is: the deep-out-of-the-money filter that defines it captures 98.2\% of QT exposure, so it is close to the identity map on the panel and the profile is not measured off a lock-in-free population. Second, the clean off-window leg that rescues the floor's \emph{level} cannot rescue its \emph{shape}: the 2017--2019 out-of-the-money cells number 335 cohort-months on \$86.2 billion of exposure and record zero prepaid balance in February and April, so no twelve-cell profile is estimable there at all. Third, the seasonal share is not stable across windows, the calendar terms explaining $R^2 = 0.011$ of the full-panel variation against 0.755 in the QT window, and the profile inherits that instability. Fourth, this floor profile peaks in June while the committed \emph{voluntary} seasonality of the estimated cohort-hazard form peaks in March and October. They are different objects and need not agree, but the tension is real and I state it rather than let the two profiles sit unreconciled in the same paper.

\begin{table}[H]
\centering
\footnotesize
\begin{threeparttable}
\caption{Seasonalizing the turnover floor. Each leg re-runs the production microsimulation as a paired central and $\beta_1 = 0$ null with only the turnover floor's calendar profile varied; the marginal is the central-minus-null difference and is basis-invariant. $U_2$ and $r$ are scored on the central leg against the SOMA back-out over the 42-month window. The constant-4.0\% row reproduces the committed production artifacts bitwise. Pre-committed run \texttt{seasonal\_floor\_timing}; artifacts catalogued in the replication package's ledger.}
\label{tab:seasonalfloor}
\begin{tabular}{@{}lrrrrrr@{}}
\toprule
Floor leg & Null (\$B) & Central (\$B) & Marginal (\$B) & Marginal (pp) & $U_2$ (diffs) & $r$ (detr.) \\
\midrule
Constant 4.0\% (production) & 748.185 & 818.530 & $+70.345$ & $+9.20$ & 1.0029 & $-0.203$ \\
Constant 4.5\% & 738.869 & 796.510 & $+57.642$ & $+7.54$ & 1.0023 & $-0.255$ \\
Shape only, mean pinned 4.0\% & 742.625 & 808.073 & $+65.448$ & $+8.56$ & 0.9947 & $+0.016$ \\
Pre-committed seasonal, mean 4.5\% & 729.298 & 783.602 & $+54.305$ & $+7.10$ & 0.9896 & $+0.070$ \\
\midrule
\multicolumn{7}{@{}l}{\textit{Placebo families, against the shape-only leg}} \\
Permutation, 14 scrambled orders\tnote{a} & --- & --- & $+65.311$--$66.924$ & --- & 0.975 best & --- \\
Rotation, 11 wrong-month shifts\tnote{b} & --- & --- & --- & --- & 0.9934 best & --- \\
Convexity-matched flat 4.0090\%\tnote{c} & --- & --- & $+70.156$ & $+9.17$ & 1.0029 & $-0.204$ \\
\bottomrule
\end{tabular}
\begin{tablenotes}[flushleft]\footnotesize
\item[a] The same twelve floor values assigned to scrambled months, multiset preserved exactly. The family retains 89.2\% of the shape effect on average (69.9\%--102.8\%), so the movement is floor dispersion under the maximum combination rule, not calendar timing. Three of the fourteen clear the frozen timing rule, and the best beats both seasonal legs on $U_2$.
\item[b] The true profile moved to the wrong months, amplitude and dispersion held. One of eleven clears the frozen rule for each seasonal leg, and the five-month shift beats the true profile's $U_2 = 0.9947$.
\item[c] The flat floor whose monthly hazard matches the calendar-frequency-weighted mean of the twelve seasonal hazards, reproducing the transform's convexity with zero dispersion. Its $-\$0.189$ billion move is pure level leakage, 3.9\% of the shape effect; the remaining 96.1\% is floor dispersion.
\end{tablenotes}
\end{threeparttable}
\end{table}

\paragraph{The within-run cyclical floor: dispersion versus co-movement.}
The within-run cyclical floor variant of Section~\ref{sec:pathb}, $\underline{h}_t = 4\% \times (1 + \kappa z_t)$, sweeps $\kappa$ symmetrically about zero. The decomposition of the resulting grid follows.

That range is mostly not a measure of cyclicality, and I no longer read it as one. A separate pre-committed run (\texttt{floor\_cyclical\_permutation}) reproduces all seven committed $\kappa$ cells bit-exactly and then decomposes them. Scrambling the floor path's time order within each $\kappa$ cell, which preserves the multiset of monthly floor values and their mean exactly, leaves 85.0\% of the range span standing: 72 draws, twelve per non-zero $\kappa$, span $+7.20$ to $+9.12$ points against the committed $+7.09$ to $+9.36$. Switching only the combination rule removes it. Under the competing-risks (additive) form the identical $\kappa$ grid spans $+11.22$ to $+11.27$ points, a span of $0.045$ points, so changing the rule alone eliminates 98.0\% of the max-form span; the additive form moves less across the whole $\kappa$ grid than it moves across the 3--5\% constant-floor level sweep. The mechanism is the one-sided Jensen effect that Appendix~\ref{sec:robustness-seasonalfloor} diagnoses at length for the seasonal floor. Under a hard maximum, dispersion in $\underline{h}_t$ lifts the effective hazard whenever the high cells bind and is truncated away where they do not, regardless of \emph{when} it occurs, so $\kappa$ is a dispersion dial before it is a cyclicality dial. Decomposing the low endpoint confirms it. At $\kappa = +0.5$, which departs $-\$16.09$ billion from production, level equivalence accounts for $-\$2.14$ billion (13.3\%), order-free dispersion for $-\$11.42$ billion (71.0\%), and genuine time-order for $-\$2.53$ billion (15.7\%): 84.3\% of that endpoint is order-free. The high endpoint is not an endpoint at all---it departs from production by $+\$1.23$ billion, $+0.16$ points, at $\kappa = -0.1$.

A cyclicality signal does survive, and it is real but small. The true marginal falls outside the entire twelve-draw permutation range at every one of the six non-zero $\kappa$ (0 of 72 draws bracket it). The order components are monotone in $|\kappa|$ and sign-consistent with the co-movement story the concession in Section~\ref{sec:pathb} states: $+\$5.29$, $+\$4.77$, and $+\$2.69$ billion at $\kappa = -0.5$, $-0.25$, $-0.1$, and $-\$1.42$, $-\$1.75$, and $-\$2.53$ billion at $+0.1$, $+0.25$, $+0.5$. A floor rising with rates does lower the marginal, and one falling with rates does raise it. But scrambling the time order strips most of the grid's asymmetry while leaving its level drop intact: the permuted means differ across the sign of $\kappa$ by only $\$0.22$, $\$1.80$, and $\$6.01$ billion at $|\kappa| = 0.1$, $0.25$, and $0.5$, against a true spread of $\$13.82$ billion at $|\kappa| = 0.5$ ($+\$68.08$ billion at $\kappa = -0.5$ against $+\$54.26$ billion at $+0.5$). A two-sided response that survives losing the time order is the signature of dispersion, not of co-movement. The symmetric part of the grid's movement is large and is dispersion; the asymmetric part is small and is cyclicality. What the grid bounds is therefore the marginal's sensitivity to within-run floor variation of this amplitude, of which at most \$5.29 billion in any cell is attributable to the direction of co-movement. That the concession's residual is bounded by the level sweep already reported holds a fortiori. One correction to the specification as earlier described: the window mean is pinned at the production 4\% only where the floor's non-negativity clip does not bind, and it binds in the two $|\kappa| = 0.5$ stress cells, whose realized window means are 4.0005\% and 4.0828\% (the artifact flags the clip per cell). Of those two, $\kappa = +0.5$ sets the bracket's low end at $+7.09$ points, while $\kappa = -0.5$ sits mid-grid at $+8.90$. The invariant therefore fails in the cell that matters most for the reported range, and the convexity-matched flat floor equivalent to the $\kappa = +0.5$ path is 4.0977\%, not 4\%.

\end{document}
